% =====================================================================
%  Harald Høffding, DEN ENGELSKE PHILOSOPHI I VOR TID.  Kjøbenhavn: P. G. Philipsens
%  Forlag, 1874 (Trykt hos J. H. Schultz).  [IV], 190 pp.
%
%  TRANSCRIPTION -- GENERATED; DO NOT EDIT BY HAND.  Made in steps, all rerunnable
%  (ebook-method/engelske-philosophi/, see its scripts' docstrings):
%    python3 ebook-method/engelske-philosophi/draft.py --emit --force
%      this template + the Internet Archive's ABBYY OCR of the University of Toronto
%      copy (~/bibliotek/Høffding, Harald/1874-den-engelske-philosophi-i-vor-tid.pdf,
%      the Archive's own PDF), then patches.py and spaced.tsv;
%    python3 ebook-method/collate/check.py engelske-philosophi --apply
%      the decisions in decisions.tsv, taken against an INDEPENDENT reading of the same
%      images (tesseract, dan, tessdata_best), with the Oxford copy's layer (Google's scan,
%      ABBYY) as a third, voting witness.  SCANS.tsv has both scans' sha256.
%    python3 ebook-method/engelske-philosophi/draft.py --post
%  No e-book exists.
%
%  STATUS (2026-10-15): TRANSCRIPTION COMPLETE, pp. 1-190 and the front matter.  Translated the same day
%  (translation.tex).
%  HOW ACCURACY WAS ESTABLISHED.
%  (a) Witnesses.  The draft is the Internet Archive's ABBYY OCR of the Toronto copy (its layer, word for word);
%      it was collated word by word, with punctuation, against tesseract (dan, tessdata_best) on the same images:
%      1488 disagreements.  A second copy of the same edition (Oxford, Google's scan) was read twice, by its ABBYY
%      layer and by tesseract, and voted where the two Toronto readings disagreed (rules.py umn-*, vote2.py ox2-*).
%  (b) Rules: 920 disagreements settled by rule or by the Oxford vote (decisions.tsv, each with its rule's name);
%      40 of these decisions were mixed unmarked into the blind crops as an audit: 2 were wrong (an Oxford vote that
%      imported a speck; a capital Ø both ABBYY readings set as ø) and led to two mends of their own (no vote may
%      import a mark the reader does not see; nouns set with a capital Ø by a reviewed list, draft.py --post).
%  (c) The image: the other 435 disagreements, 48 places where both Oxford readings agreed against the text
%      (oxsweep.py) and the audit spots were read BLIND in 52 crop sheets (round b1, six readers, two hidden controls
%      a sheet: 103 of 104 right, the one miss the reader's).  OPEN 0, --unapplied mended by hand.
%  (d) The footnotes (66) were read whole at the image (notes-read.txt -> patches.NOTESET): three one-line notes the
%      draft had set as text, two notes run over a page and one page of text the draft had taken into a note (p. 162)
%      were found that way and mended (patches.NOTEFIX, POST2).
%  (e) Letterspacing: candidates by the gaps between letters (spacedsheets.py, 113) and by word WIDTH against widths
%      fitted per character over the book (widthsheets.py, 180), read at the image (rounds l1, l2).  Italics: by stroke
%      slant (check.py --italics): none in the text; the Indhold's names are italic.  Front matter and the chapter
%      heads read whole (structure reading f01-f10).
%  (f) Sampled page check: 8 pages (2434 words) chosen at random and read whole and blind: 0 word errors (95 % upper
%      bound 1.2 per 1000 words), 2 punctuation and 3 letterspacing faults (short words: det, endnu), all mended.
%  (g) The translators' odd spots: 208, read blind at the image (round b3, three readers, controls 52/52): 71
%      transcription faults mended (POST2: specks read as stops, a dropped kj, I read as 1, words joined or split,
%      a note numbered wrongly, paragraph breaks inside a sentence), 53 misprints of the print marked PRINTED AS IS.
%      The rate of what the odd spots found (~0.4 transcription faults per page where the first collation had passed) shows the
%      limit of two OCRs that share a worn sort: misprints and faults where ABBYY, tesseract and the Oxford copy all
%      agree remain possible, most of them in non-words a reader of the text would notice.
%  (h) Unattested-word sweep: words of the text found in no other transcription of the repo and no word list, read
%      at the image where suspect: 2 more misprints marked.  In all, 65 places are marked PRINTED AS IS.
%  EDITORIAL.  The print closes quotations both ways (»...« and »...»); the notes use „...“; both kept as printed.
%  „ɔ:“ (id est) as printed.  The ø rule (CLAUDE.md) applied by oefix's reviewed list and this book's additions
%  (rules.py OE_BOOK).  Misprints are kept and marked PRINTED AS IS at the line.  Printer's signatures („4*“) and
%  folios are not transcribed.  In the Indhold, a damaged m in „William“ and a speck are not transcribed.
%
%  CONVENTIONS.  "% --- p. N ---" + \opage{N} at the first word of printed page N
%  (a word broken over the turn is split with %).  Footnotes as printed, marked
%  *), **), ***) and numbered per page (\footnote[k]).  Italics \textit; letterspacing
%  \emph.  A note running over a page is set whole at its mark.  Folios and signature
%  marks are not transcribed.  The print governs; misprints are kept and marked
%  PRINTED AS IS.
% =====================================================================
\documentclass[12pt,a4paper,openany]{book}
\usepackage[utf8]{inputenc}
\usepackage[T1]{fontenc}
\usepackage{libertinus}
\usepackage{libertinust1math}
\usepackage[danish]{babel}
\usepackage[protrusion=true,expansion=true]{microtype}
\usepackage{setspace}
\usepackage[margin=1.2in]{geometry}
\usepackage{fancyhdr}
\setlength{\headheight}{28pt}
\usepackage{hyperref}
\hypersetup{pdftitle={Den engelske Philosophi i vor Tid (1874)}, pdfauthor={Harald Høffding}, hidelinks}
\pagestyle{fancy}
\fancyhf{}
\fancyhead[LE,RO]{\thepage}
\fancyhead[RE]{\textit{Harald Høffding: Den engelske Philosophi i vor Tid}}
\fancyhead[LO]{\textit{1874}}
\fancypagestyle{plain}{\fancyhf{}\fancyhead[LE,RO]{\thepage}}
\setstretch{1.3}
\usepackage{marginnote}
\usepackage{graphicx}
\DeclareUnicodeCharacter{0254}{\reflectbox{c}}   % ɔ (id est)
\renewcommand{\thefootnote}{%
  \ifcase\value{footnote}\or *)\or **)\or ***)\else \arabic{footnote})\fi}
\newcommand{\tocline}[2]{\noindent #1\dotfill\ #2\par}
\newcommand{\opage}[1]{\marginnote{\footnotesize\textit{[#1]}}}
% a chapter head of the print (structure reading f03-f10): „Første Kapitel.“ large, a short rule, the name in bold
% sans serif; the chapters begin where the print has them (pp. 18, 116 in mid-page), so no page break
\newcommand{\chapterhead}[2]{\par\bigskip\begin{center}{\Large #1}\\[0.5em]\rule{2.5em}{0.4pt}\\[0.6em]{\large\bfseries\sffamily #2}\end{center}\medskip\addcontentsline{toc}{chapter}{#2}}
\newcommand{\closinghead}[1]{\par\bigskip\begin{center}{\Large\bfseries\sffamily #1}\\[0.4em]\rule{2.5em}{0.4pt}\end{center}\medskip\addcontentsline{toc}{chapter}{#1}}
\pdfstringdefDisableCommands{\renewcommand{\footnote}[2][]{}\renewcommand{\opage}[1]{}}


\begin{document}

% --- front matter: read whole at the image (structure reader G, f01-f02, 2026-10-15).  Title PDF 7, verso PDF 8 (blank),
% Indhold PDF 9.  Handwritten shelf and accession marks and two small printer's ornaments on the Toronto copy not
% transcribed; in the Indhold a damaged m in „William“ (Hamilton) and a speck before „om“ are not transcribed. ---
\begin{titlepage}
\begin{center}
\vspace*{3em}
{\LARGE Den engelske Philosophi}\\[0.8em]
{\LARGE i vor Tid.}\\[3em]
{\small Af}\\[1em]
{\large Dr. \textbf{H. Høffding.}}\\[1.5em]
\rule{3em}{0.4pt}\\[-0.8em]\rule{3em}{0.4pt}
\vfill
{\bfseries Kjøbenhavn.}\\[0.4em]
P. G. \emph{Philipsens Forlag.}\\[0.4em]
{\footnotesize Trykt hos J. H. \emph{Schultz.}}\\[0.4em]
1874.
\end{center}
\end{titlepage}
\clearpage
\thispagestyle{empty}
\begin{center}{\large\bfseries \emph{Indhold.}}\end{center}
\begin{center}\rule{3em}{0.4pt}\end{center}
\hfill{\small Side.}\par
\noindent\emph{Første Kapitel:} \textbf{Almindelig Karakteristik}\dotfill\ 1.\par
\begin{quote}\small Den philosophiske Erkjendelses Betingelser. --- Den tidligere engelske Philosophi. --- Vor Tids engelske Philosophi og dens Forhold til andre Retninger.\end{quote}
\noindent\emph{Andet Kapitel:} \textbf{Den rene Empirisme}\dotfill\ 18.\par
\noindent\hspace*{2em}I. \textit{John Stuart Mill}\dotfill\ 18.\par
\begin{quote}\small Historiske Forudsætninger (Hartley og James Mill). --- Logik og Psykologi. --- Det Absolutes Idé. --- Ethik og Sociologi.\end{quote}
\noindent\hspace*{2em}II. \textit{Alexander Bain}\dotfill\ 66.\par
\begin{quote}\small Sandsning og Tænkning. --- Følelsernes Klassifikation. --- Ethik. --- Villie og Tro.\end{quote}
\noindent\emph{Tredie Kapitel:} \textbf{Den kritiske Skole i England}\dotfill\ 88.\par
\noindent\hspace*{2em}I. \textit{William Whewell}\dotfill\ 90.\par
\begin{quote}\small De induktive Videnskabers Philosophi. --- Forhold til Kant og til Theologien.\end{quote}
\noindent\hspace*{2em}II. \textit{William Hamilton}\dotfill\ 99.\par
\begin{quote}\small Polemik mod det Ubetingedes Philosophi. --- Hamiltons Discipel \textit{Mansel} og hans Kritik af Theologien. --- Hamiltons Logik og Psykologi. --- Læren om Kausaliteten.\end{quote}
\noindent\emph{Fjerde Kapitel:} \textbf{Udviklingsphilosophien}\dotfill\ 116.\par
\begin{quote}\small \textit{Herbert Spencer}. --- Almindelig Karakteristik. --- Det Absolutes Idé (Forholdet mellem Religion og Videnskab). --- De første Principer. --- Biologi. --- Psykologi og Erkjendelseslære. --- Ethik og Sociologi.\end{quote}
\noindent\textbf{Slutning}\dotfill\ 184.\par
\begin{center}\rule{3em}{0.4pt}\end{center}
\clearpage


% --- p. 1 ---
\setcounter{footnote}{0}%
\chapterhead{Første Kapitel.}{Almindelig Karakteristik.}
\opage{1}


Den menneskelige Udvikling gaaer ikke for sig i lige
Linie. Paa ethvert Omraade finde vi, at Tider fulde af
lykkelig og storartet Fremadskriden vexle med tilsyneladende
Stillestaaen eller Tilbagegang, hvor Kræfterne røre sig i det
Skjulte eller virke paa andre Omraader. Og dette gjælder
i høiere Grad, jo mere sammensat en Bestræbelse er, jo
flere Betingelser den forudsætter for at kunne lykkes. Om
ingen gjælder det derfor mere end om den philosophiske
Stræben. At philosophere er ikke det Samme som at hengive
sig til umiddelbare Intuitioner og geniale Forsøg paa
at gribe og fremstille det Ideale. Dersom det at philosophere
virkelig skal have videnskabelig Betydning, maa der
til Grund for den hele Bestræbelse lægges en klar og sikker
Bevidsthed om Sandhedens og Virkelighedens Væsen. Det
Sande og det Virkelige maa være Magter, for hvilke Tanken
bøier sig, fordi den i dem erkjender sin egen Lov og i
Underkastelsen under dem sin sande Frihed. Tankens
Stræben maa gaa ud paa Skridt for Skridt at lade sig
binde; først ved helt igjennem at være bunden har Tanken
Gyldighed. Al Philosopheren forudsætter derfor et dobbelt
System af Betingelser. Tanken maa være vaagen, Driften
til en sammenhængende og altomfattende Verdensanskuelse
maa være levende; men paa den anden Side maa ogsaa
det faktiske Grundlag, Enkelthedernes Masse ligge saaledes
beredt, at det virkelig bliver muligt for Tanken at komme
ud over sine subjektive og tilfældige Ideer og iføre sig
% --- p. 2 ---
\setcounter{footnote}{0}%
\opage{2}Sandhedens og Virkelighedens Former. Alt eftersom disse
Betingelser vexle, vexle ogsaa den philosophiske Forskens Epoker.
Det ligger i Sagen selv og bekræftes af Historien. Den
nyere tydske Philosophi, som paa saa energisk og genial en
Maade blev indledet af Kant, og som i denne sin største
Repræsentant endnu stod i nøie og frugtbar Vexelvirkning med
Erfaringen, udtømte alle sine Muligheder i Systemer, hvor
det givne Fond af Realitetsbestemmelser i den Grad absorberedes,
at Tilværelsens Lov tilsidst gjordes til Et med
Tankens egen dialektiske Bevægelse. Saa kom der en
Tid, hvor man var tilbøielig til at kaste al Philosophi overbord;
Interessen vendte sig mod Realiteten; det var en
Glandsperiode for den exakte, experimentale Naturforsken;
man indsamlede Stof paa alle Omraader. Men de gamle
Problemer kunne ikke holdes ude; de trænge sig atter frem
med uimodstaaelig Kraft, og vi finde, tildels endog netop
fra saadanne Mænds Side, som tog Del i den experimentale
Forsken, en Række af nye Forsøg paa at behandle de
philosophiske Opgaver. En lignende Udviklingsgang finde
vi i England. Efter en storartet philosophisk Virksomhed
i det 17de og den første Del af det 18de Aarhundrede
var England fattig paa philosophiske Forsøg, indtil Interessen
igjen vaktes henimod Midten af vort Aarhundrede.
»Der var en Tid», siger Herbert Spencer (Essays. 1868.1.
p. 300),« da vore Landsmænd spekulerede --- og det sikkert
i lige saa stor en Udstrækning som noget andet Folk --- over
alle de høie Spørgsmaal, der fremstille sig for den menneskelige
Tanke; og et Blik paa de philosophiske Systemer,
der ere eller have været herskende paa Fastlandet, er tilstrækkeligt
til at vise, hvor meget andre Nationer skylde
vore Forfædres Opdagelser. Men i en Menneskealder eller
to ere disse mere abstrakte Gjenstande blevne forsømte, ja
endog (af dem, der bryste sig af at være »praktiske») foragtede.
Tildels maaske en Følge af vor hurtige materielle
Væxt, skyldes denne intellektuelle Phase for en stor Del
dog Udtømmelsen af Stof og Nødvendigheden af bedre
Data. Ikke saa meget med bevidst Anerkjendelse af det
% --- p. 3 ---
\setcounter{footnote}{0}%
\opage{3}Maal, der skulde tjenes, som med ubevidst Underordnen
under den Rhythme, der kan spores i sociale Forandringer
saavel som i andre Ting, er en Periode af Theoretiseren
uden Iagttagen bleven efterfulgt af en Periode af Iagttagen
uden Theoretiseren. Under den længe fortsatte Hengivelse
til konkret Videnskab er der blevet ophobet en umaadelig
Mængde Raastof for abstrakt Videnskab, og nu er en Periode
aabenbart i Færd med at begynde, i hvilken dette ophobede
Raastof vil blive organiseret til en sammenhængende Theori.
Paa alle Punkter, i de uorganiske Videnskaber, i Videnskaben
om Livet og i Samfundsvidenskaben, kunne vi bemærke
en Tendents til at gaa fra det Overfladiske og
Empiriske til det Dybere og Rationelle». Naar vi tale om
den engelske Philosophi i vor Tid, er det altsaa Mere end
en tilfældig Række af Phænomener, vi fremdrage; det er
en Række af Bestræbelser, der fremkomme i Medfør af en
naturlig Udviklingslov og derfor med en eiendommelig, ved
deres Plads i Udviklingen betinget Karakter.

Den dobbelte Forudsætning for den philosophiske
Forsken, som bevirker, at dens Udviklingsgang bliver rhythmisk
eller næsten springende, indeholder tillige Muligheden
af en afgjørende Modsætning indenfor Philosophien. To
forskjellige Udgangspunkter blive mulige. Den almindelige
Opgave er, at Tanken skal tilegne sig det Sande og Virkelige.
Men, kan der saa siges, kun ved og for Tanken er
Noget virkeligt og sandt; Tankens Væsen er derfor et væsentligt
og selvgyldigt Element i Virkeligheden; det faktisk
Existerende er kun Middel for Tankens Udvikling, Erkjendelsen
deraf kun Veien, ad hvilken Tanken bliver sig selv
bevidst, og i denne Bevidsthed indeholdes den høieste
Sandhed. Her lægges altsaa Eftertrykket paa det ideelle
Element, paa Tankens oprindelige Aktivitet som sidste Forudsætning
for al Erkjendelse. Men, siges der fra den anden
Side, en Erkjendelse uden virkeligt, objektivt gyldigt Indhold
er kun subjektiv og illusorisk. Naar Tanken vil være
selvvirksom, holder den sig i sine egne selvskabte Ideer,
forsaavidt den ikke følger traditionelle Fordomme. Skal
% --- p. 4 ---
\setcounter{footnote}{0}%
\opage{4}Virkeligheden erkjendes, maa Tanken derfor forholde sig
aldeles passiv og modtagende. Al Erkjendelse maa ske ved
Opsamlen af Indtryk og Erfaringer, og kun det kan gjøre
Krav paa Sandhed, der direkte eller indirekte kan føres
tilbage til denne Kilde. --- Modsætningen mellem de to
Hovedretninger indenfor Philosophien beroer altsaa paa den
forskjellige Besvarelse af det Spørgsmaal: Hvori bestaaer i
sidste Instants Sandhedens og Virkelighedens Kriterium?
Det ene Svar henviser os til Tankens Nødvendighed, det
andet til Erfaringens Virkelighed. Dette er Grundmodsætningen;
men paa Basis af den udfolder der sig en Mængde
afledede og for hele Verdensanskuelsen overordentlig vigtige
Modsætninger. Det Spørgsmaal maa da fremstille sig, om
det er muligt at forsone de to Principer og dermed de to
Verdensanskuelser.

Dog kan dette Spørgsmaal først komme frem, naar
der er indledet en Vexelvirkning, en virkelig Debat mellem
de to philosophiske Skoler. Men dette kan tildels hindres
ved, at de høre hjemme hos forskjellige Nationer. En saa
afgjørende Modsætning som den, der her har udviklet sig
med Hensyn til den philosophiske Forskens hele Retning,
viser tilbage til mere omfattende psykologiske og nationale
Forskjelligheder. Tankegang og Verdensanskuelse ere ikke
noget Isoleret; det er Faktorer, der staa i levende Forbindelse
med Individets og Folkets hele Karakter. Den
idealistiske Skole har i den nyere Tid hovedsagelig havt
sit Hjem i Tydskland, den realistiske i England. Tydskere
og Englændere ere de egentlige philosopherende Nationer
i den nyere Tid, de staa i Philosophiens Historie ved Siden
af Grækerne. Den nationale Begrændsning har her virket
hemmende, men ikke blot hemmende. Thi i den philosophiske
Forsken kommer det væsentlig an paa Klarhed og konsekvent
Gjennemførelse af Tankegangen, og hertil udfordres
en nøiagtig og grundig Bearbeidelse af Grundtankerne, som
kun er mulig under Forudsætning af en vis Ensidighed.
Herpaa beroer de forskjellige philosophiske Skolers Betydning,
og herpaa beroer ogsaa Betydningen af den individuelle
% --- p. 5 ---
\setcounter{footnote}{0}%
\opage{5}og den nationale Karakters Indflydelse paa den philosophiske
Tænkning. Vigtige og afgjørende Hovedbegreber vilde aldrig
være blevne gjennemtænkte uden saadanne Begrændsninger og
Ensidigheder. Saaledes har navnlig den tydske og den engelske
Philosophi ved deres konsekvente, Skridt for Skridt
foregaaende Udvikling kastet et klart Lys over store Tankekredse,
der ellers ikke vilde kunne udmaales. Og kun
derved kan hin Grundmodsætning mellem de philosophiske
Hovedretninger fremstille sig i saa klar og præcis
en Form, at et Arbeide paa dens Overvindelse bliver
muligt. ---

Af det hidtil Udviklede følger, at der gives tre væsentlige
Betingelser for Fremskridt i den philosophiske Erkjendelse.
Philosophien forudsætter den i Erfaringen givne
Virkelighed som sit Stof. Den gaaer ud paa en konsekvent
Gjennemførelse af Principet. Men da dette Princip ifølge
Sagens Natur bliver ensidig udviklet i de forskjellige Skoler,
stræber den at korrigere denne Ensidighed ved en Forening
af de tilsyneladende modsatte Synspunkter. --- Det skal nu
vises, at den engelske Philosophi i vor Tid er af stor Betydning
i alle tre Henseender.

Den engelske Philosophi har til alle Tider havt en
empirisk og realistisk Karakter. I Middelalderen finde vi
Englændere som de mest fremtrædende Repræsentanter for
Nominalismen, den Retning indenfor Scholastiken, som svarer
til den moderne Realisme. Duns Scotus († 1308) polemiserer
mod Thomas Aqvinas, Repræsentanten for den kirkelige
Spekulation, og hævder navnlig Individualitetens Betydning.
Ifølge Thomas var den universelle Ide den sande
Virkelighed, og det Individuelle fremkom kun ved en Begrændsning,
der skyldtes Materien. Men, siger Duns Scotus,
hvad der gjør en Ting til denne bestemte Ting, maa være
noget Positivt; det Individuelle er det Fuldkomne, den
høieste Virkelighed. Roger Bacon († 1292) haanede Albertus
Magnus og Thomas Aqvinas, fordi de vilde philosophere
uden at have lært Sprog, Mathematik og Physik,
og han maatte vansmægte i sit Fængsel, fordi hans ivrige
% --- p. 6 ---
\setcounter{footnote}{0}%
\opage{6}Syslen med og Lovprisen af scientia experimentalis gjorde
ham mistænkt for magiske Kunster. William Occam
(† 1347) førte Nominalismen videre. Kun det Enkelte
er virkeligt; det kan ikke forklares og opløses ved almene
Begrebsbestemmelser; de almene Begreber have ingen Realitet
udenfor Forstanden; de ere Tegn og Symboler, som
udtrykke de enkelte Phænomeners Indbegreb. Al Viden
grunder sig derfor paa Anskuelse, der igjen viser tilbage
til indre eller ydre Erfaring. Da vi ikke kunne erfare eller
anskue Guds Væsen, kunne vi altsaa ikke opnaa nogen
Viden om Gud; Theologien er ikke nogen Videnskab, men
grunder sig ene og alene paa overnaturlig Autoritet. Occam
gjennenfører, hvad allerede Duns Scotus havde udtalt, at % PRINTED AS IS (gjennenfører: O-A03)
Noget kan være sandt i Theologien, som ikke er det i
Philosophien. --- Den nominalistiske eller empiriske Retning
sprænger paa denne Maade Scholastiken. Den Modsætning,
som i nyere Tid fremtræder mellem den engelske Philosophi
og Fastlandets, kunne vi saaledes forfølge tilbage til
Middelalderen.

Baco fra Verulam († 1626) er Stifteren af den nyere
engelske Philosophi. Fra de spekulative og theologiske
Studier kalder han. Menneskeaanden tilbage til Naturen og
dens Erkjendelse. Kun der findes en sikker Grundvold for
Videnskaben. Overladt til sig selv forvilder Tanken sig i
illusoriske og vilkaarlige Forudsætninger (idola mentis), som
dels udspringe af den menneskelige Natur, dels af Samfundsforholdene.
Disse upaalidelige Forudsætninger befrier
man sig fra ved at underkaste sig Erfaringens Lov, ved ad
Induktionens Vei at undersøge det Enkelte for derved
at naa til de sande Grundbegreber (notiones rerum primæ).
Den induktive Naturvidenskab er for Baco alle Videnskabers
Moder, til hvilken ogsaa Logik og Moralphilosophi skulle
føres tilbage. Dens Formaal er ikke blot Erkjendelse, men
Fremskridt i Civilisation, idet Menneskelivet ved dens Opfindelser
skal blive rigere udstyret med alle Slags Goder,
og Menneskets Magt over Naturen udvides. --- Naar Baco
saaledes forlanger, at Mennesket skal afføre sig al Selv%
% --- p. 7 ---
\setcounter{footnote}{0}%
\opage{7}virksomhed i Erkjendelsen og saa meget som muligt forholde
sig passiv, baner han derved Veien for Sensualismen,
der udleder al Erkjendelse fra Sandsningen og betragter alle
Forestillinger og Tanker som omdannede Sandsefornemmelser.
Denne Lære blev udviklet af Thomas Hobbes († 1679) og
John Locke († 1704). De Love, efter hvilke Fornemmelserne
kombineres og sammensættes og derved omdannes
til Tanker, kunde ikke være andre end Ideassociationernes
Love; idet den aandelige Aktivitet forkastes, bliver der kun % PRINTED AS IS (aandelige: the final e printed turned (ə); blind reads b1-03-12, b1-20-04)
de passive Sjælephænomener tilbage. David Hartley
(† 1757) var den Første, der søgte at forklare alle Bevidsthedsphænomenerne
som Ideassociationernes Resultater.
Endelig drog David Hume († 1776) de sidste Konsekventser
af Empirismen og Sensualismen. Saavel den indre
som den ydre Verden, saavel Jeget som Verden, kjende vi
kun gjennem vore Sandsefornemmelser og Forestillinger;
og den indre Sands, Forestillingskredsen, er kun et afbleget
Billede af den ydre Sands; i sidste Instants kjende vi
da altsaa kun Sandsefornemmelsernes Række. Men saa er
det kun grundet i en Illusion, naar vi antage et Jeg, en
indre aandelig Substants som Forestillingskredsens Bærer,
eller en ydre materiel Substants som Sandsefornemmelsens
Aarsag. Det er i det Hele en Illusion, naar vi mene at
erkjende en kausal Sammenhæng, et Forhold mellem Aarsag
og Virkning; vi kjende kun Phænomenernes Rækkefølge i
Tiden, ingen indre og nødvendig Sammenhæng imellem
dem. Det er kun Vanen, som igjen er et Produkt af Ideassociation,
der fører os til saadanne Antagelser; men Vanen
kan ikke begrunde nogen Erkjendelse.

Ved skarp og konsekvent Gjennemførelse var saaledes
Empirismen gaaet over til Skepticisme og havde derved
opløst sig selv. Alligevel var det Humes og Hartleys
Principer, til hvilke den engelske Philosophi støttede sig
da der henimod Midten af det 19de Aarhundrede igjen
begyndte en Blomstringsperiode for den. Vi maa da nu
nærmere angive Forholdet mellem vor Tids engelske Philosophi
og den tidligere. Grundtanken er den samme: den
% --- p. 8 ---
\setcounter{footnote}{0}%
\opage{8}nominalistiske Fastholdelse af det Enkelte og Individuelle
som det eneste Virkelige, den empiristiske Antagelse af
Sandsning og Induktion som eneste Erkjendelseskilder. Men
denne Grundtanke faaer en eiendommelig Gjennemførelse
ved den Vægt, der lægges paa Ideassociationens Love som
afgjørende for Bevidsthedslivets og særlig for Erkjendelsens
Væxt og Udvikling. Hume havde sagt, at vi ikke kunne
opnaa nogen sand Erkjendelse, fordi vi kun have Sandsefornemmelser
og deres Forbindelser at bygge paa. Men,
sige vor Tids engelske Philosopher hertil, selve Erkjendelsen
er ikke andet end en fast og uopløselig Ideassociation.
Tanken er ikke en mystisk Kraft, der atter opdager et
mystisk Baand som nødvendig Forbindelse mellem Tingene.
Naar visse Forestillinger ifølge Erfaringen og hele Udviklingens
Gang have indgaaet en fast Forbindelse, saa maa
vi betragte dem som sammenhørende. Hvad ere vi selv,
vort Jeg Andet end en fast og inderlig Sammensætning,
en Kjede af associerede Forestillinger? Hvad vi kalde nødvendigt,
er kun et Udtryk for Fasthed i Associationen.
Men vi kunne ikke blot psykologisk paavise, hvorledes Bevidstheden
faktisk udvikler sig efter Ideassociationens Love;
vi kunne ogsaa ved logisk Analyse sondre mellem berettigede
og uberettigede Associationer. Baco havde i al Almindelighed
lovprist Induktionen og bekæmpet Syllogismens
Anseelse som høieste Erkjendelsesform. Men nu først
gjennemføres den induktive Logik, idet der paavises, hvorledes
den experimentale Methode fører til at eliminere
de tilfældige Tankeforbindelser og fastslaa og hævde de
gjennemgaaende og derfor berettigede. Først paa Basis
af en saadan forudgaaende, experimentalt fastslaaet Sammenknytten
af Forestillinger blive almene Love mulige, og først
derved vindes Udgangspunkter for Deduktion, Syllogisme
og mathematisk Analyse. Hume polemiserede mod Naturvidenskaben,
men vor Tids engelske Philosophi stiller sig
netop i Kraft af Humes Principer paa Naturvidenskabens
Grund. Den experimentale Naturvidenskab er i sit Væsen
skeptisk. Den vil kun finde Lovene for Phænomenernes,
% --- p. 9 ---
\setcounter{footnote}{0}%
\opage{9}indbyrdes Forbindelser, deres Rækkefølge og deres Sammenhæng.
Kun fordi Hume ikke kjendte den experimentale
Methode kunde han vende sin Skepsis mod selve Naturvidenskaben.\footnote[1]{Stuart Mill (Examination of Sir William Hamiltons philosophy p. 555) mener endog, at Humes Skepticisme kun var en Forklædning: „han vilde hellere kaldes Skeptiker end betegnes ved et mere hadefuldt Navn.“}


Den engelske Philosophi i vor Tid minder her om
den franske Positivisme. Man finder endog meget ofte
Positivismen brugt som Fællesnavn, saa at det ogsaa omfatter
den engelske Skole. Og dog have denne Skoles
første Repræsentanter med Rette fralagt sig Navn af
Positivister, naar dette Ord skal betegne dem som Disciple af
Auguste Comte. Vi have allerede set, at vor Tids engelske
Philosophi er en videre Udvikling af den tidligere engelske
Tænknings Standpunkt og Resultater, men saaledes, at den
ifølge sine logiske og psykologiske Principer kom til at
træde i en nøie Forbindelse med den experimentale
Forsken Denne Forsken og dens Resultater er den fælles % PRINTED AS IS (Forsken Denne: O-A05: no stop after Forsken)
Grund, hvor den mødes med den franske Positivisme; det
Fælles for begge Retninger skriver sig altsaa fra den Kilde,
hvoraf de begge øse. Man maa skjelne mellem Positivismen
som særlig Skole, og hvad Franskmændene kalde l'esprit
positif, en Aandsretning, som er udsprungen af hele den
nyere Kultur og Videnskab, og af hvilken Positivismen
selv kun er en særegen Form. Modsætningen mellem den
franske Positivisme og den engelske Skole kan i sin Almindelighed
angives saaledes: Positivismen gaaer fra Fagvidenskaberne
til Philosophien, den engelske Skole fra
Philosophien til Fagvidenskaberne. Den franske Positivisme
vil egentlig ikke vide af andre Problemer at sige end de
fagvidenskabelige. Den samler i en klassifikatorisk Encyklopædi
de enkelte Videnskabers Principer og Resultater
og lader sig dermed nøie. Kun ét afgjørende Problem
fremdrager den tilsidst og forraader derved uvilkaarlig
% --- p. 10 ---
\setcounter{footnote}{0}%
\opage{10}sit eget inderste Motiv: Problemet om Menneskehedens og
særlig det menneskelige Samfundslivs Udvikling (det sociologiske
Problem). Men den forsømmer de logiske og
psykologiske Undersøgelser, som ene kunne danne en Basis
for en Philosophi, og som ogsaa ere en nødvendig Forudsætning
for Behandlingen af hint store Problem. I ethvert
Tilfælde anenkjender man ikke Psykologiens Selvstændighed % PRINTED AS IS (anenkjender: O-A06)
men lader den ganske gaa ind under Physiologien (saaledes
Littré). Den engelske Philosophi tager derimod sit Udgangspunkt
i Logik og Psykologi og finder herfra Veien
til de speciellere Problemer. ---

Betragtningen af vor Tids engelske Philosophi i dens
Forhold til den tidligere engelske Tænkning førte os tillige
til at berøre dens Forhold til Fagvidenskaberne. Men det
vigtigste Bidrag til denne Philosophis Karakteristik staaer
endnu tilbage. Det bestaaer i en Angivelse af dens Forhold
til den philosophiske Hovedretning, til hvilken den
tidligere engelske Skole havde dannet en afgjørende Modsætning,
den idealistiske eller spekulative.

Hvad vi for Kortheds Skyld have betegnet som den
engelske Skole, har forsaavidt ingen udelukkende Ret til
dette Navn, som der baade tidligere og i vore Dage har
gjort sig philosophiske Retninger gjældende i England, der
stille sig i Modsætning til den empirisk-realistiske Philosophi.
I det 17de Aarhundrede finde vi en saadan Retning
i den mystiske Platonisme, som herskede ved Universitetet
i Cambridge, og hvis mærkeligste Repræsentanter
vare Henry More og Ralph Cudworth; Lockes Polemik mod
de »medfødte Ideer» er maaske snarere rettet mod dem
end mod Cartesianerne. I det 18de Aarhundrede udviklede
der sig ved de skotske Universiteter som Reaktion mod
den engelske Empirisme og dens Resultater en særlig skotsk
Skole, hvis egentlige Stifter var Thomas Reid († 1796).
Denne Skole havde for en Del praktiske Motiver; den vilde
afværge de formentlig farlige Konsekventser af Lockes og
Humes Lære og var, som Cousin udtrykker sig, »en Protest
fra den sunde Menneskeforstand og Samvittigheden mod
% --- p. 11 ---
\setcounter{footnote}{0}%
\opage{11}Lockes Philosophi.« Reid bestrider den empiriske Skoles
Love om Sandsefornemmelserne som de oprindelige Elementer,
ved hvis Association og Kombination hele Sjælelivet
udvikler sig. Han lærer, at vor Bevidsthed umiddelbart
opfatter de ydre Gjenstande. Vi fælde instinktmæssig
Domme om Tingene; og disse uvilkaarlige Antagelser indeholde
de første og naturlige Principer for al vor Erkjendelse.
De ere Bevidsthedens oprindelige Kjendsgjerninger,
hvis Indbegreb danner den saakaldte »almindelige Menneskeforstand«
(common sense). Af saadanne »naturlige Principer«
opregner Reid en stor Mængde, bl. A.: den Sætning,
at enhver Fornemmelse viser hen til en ydre Gjenstand;
Troen paa vor personlige Identitet; fremdeles mathematiske
og metaphysiske Principer (f Ex.: Enhver Ting har sin
Aarsag), og moralske Bud (f. Ex.: Man skal sætte Pligten
over Nytten). Alle disse Principer skulle umiddelbart ligge
i Bevidstheden. --- Det er ikke ved at opstille »den sunde
Menneskeforstands Principer» men ved den Undersøgelse
af Empirismen, de satte i Bevægelse, og ved den Sands for
psykologiske Undersøgelser, de vakte, at Reid og hans
Skole have øvet en ikke ringe Indflydelse. I Frankrig
rykkede Royer Collard og Victor Cousin i Marken mod
Condillac med Reids Vaaben. Men ogsaa i England har
den skotske Philosophi været af stor Betydning, navnlig
siden den i sin største Repræsentant i vor Tid, William
Hamilton, indgik en eiendommelig Kombination med den
kantiske Philosophi. Dette fører os til at undersøge, hvorledes
den tydske Philosophi griber ind i den engelske
Tænknings Udviklingsgang i vor Tid. Kant, som efter sin
egen Erklæring fik Stødet til sin epokegjørende Forsken
fra David Hume, har siden selv virket betydningsfuldt tilbage
paa Humes Fædreland.

Kant træder ligesom Reid op mod Humes Skepticisme.
Men medens det for den ængstelige Reid fremfor
Alt gjaldt om at sikre sig Jegets og den ydre Verdens
Realitet, hvorfor hans Forsøg ogsaa fik en mere overfladisk
Karakter, gik Kant dybere tilbage og fandt i Kausalitetens
% --- p. 12 ---
\setcounter{footnote}{0}%
\opage{12}Problem det egentlige Stridsspørgsmaal mellem begge Retninger.
Han kunde ikke lade sig nøie med den skotske
Skoles Paaberaabelse af den sunde Menneskeforstand. At
den almindelige Menneskeforstand kjendte og anvendte
Aarsagsbegrebet, negtede Hume jo ikke; men Striden dreier
sig netop om dette Begrebs Betydning og Gyldighed. I
ethvert Tilfælde skulde man, mente Kant, hellere bruge
sin sunde Menneskeforstand til at føre Bevis for, hvad
man antager, end paakalde den som et Orakel. Ved
at gjennemtænke Sagen blev Kant nu opmærksom paa, at
der er en hel Kreds af Anskuelsesformer og Begreber, som
ligger til Grund for al Erfaring, fordi de ere de forskjellige
Former og Maader, hvorpaa det i Fornemmelsen givne Stof
sammenfattes til Enhed for Anskuelsen og Tænkningen. Disse
Former maa derfor, sluttede han, være aprioriske, oprindelig
givne i vor Natur. De kunne ikke skrive sig fra Erfaringen,
da al Erfaring netop forudsætter dem, og da de have almindelig
og nødvendig Gyldighed, hvilket blotte Abstraktioner
fra Erfaringen ikke kunne have, da Erfaringen aldrig
viser os det Ubetingede. Netop i Kraft af Erfaringens
Gyldighed hævder Kant saaledes Tankens Oprindelighed.
Han har herved banet Veien for en hel ny Opfattelse af
den aprioriske eller spekulative Erkjendelse. Denne kan
ikke længere, som i Dogmatismen, være et færdigt og
afsluttet System, som skulde kunne udvikles uden Hensyn til
Erfaringen; men den forudsætter stedse Erfaringen, og
optager den i sig i Kraft af sit eget Princip »At al vor
Erkjendelse begynder med Erfaringen«, siger Kant,
»derom er der ingen Tvivl; men derfor udspringer ikke
al vor Erkjendelse af Erfaringen«. Kant selv skjelnede
saa stærkt mellem det Subjektive og det Objektive, at det
ser ud, som om Tanke- og Anskuelsesformerne staa fuldfærdige
paa den ene Side, og Indholdet først kommer til
fra den anden Side; men den sande Konsekvents af den
kritiske Philosophi har Schelling udtalt i den Sætning, at
vor Erkjendelse lige saa lidt er apriorisk som aposteriorisk,
idet denne Modsætning først skyldes Reflektionen. Schelling
% --- p. 13 ---
\setcounter{footnote}{0}%
\opage{13}viser fremdeles, at der vel gives aprioriske men ikke medfødte
Begreber. »Ikke Begreber, men vor egen Natur og
dens hele Mekanisme er os medfødt. Denne Natur er af
en bestemt Beskaffenhed og virker paa en bestemt Maade,
men fuldkomment bevidstløst, thi den er selv ikke Andet
end netop denne Virken.« Sjælen er hverken en beskreven
eller en ubeskreven Tavle.

Det var Samuel Coleridge († 1834), som især blev
Mellemleddet, gjennem hvilket den tydske Philosphi gjorde % PRINTED AS IS (Philosphi: O-A12)
sin Indflydelse gjældende i England. I sin Ungdom havde
han været en ivrig Beundrer af Hartley og Hume, men
frastødtes senere ved disse Tænkeres Benegtelse af Aandens
Aktivitet. Han saa rigtig, at naar Læren om Ideassociationerne
gjøres til den psykologiske Hovedlære, opfattes
alle Bevidsthedsphænomenerne som fuldstændig passiviske.
Fra dette Synspunkt har han allerede (i sin Biographia
litteraria vol. I. pars II. p. 116 ff) givet en træffende
Kritik af den saakaldte »Association psychology.« Studiet
af Kant og hans Efterfølgere førte ham ganske ud over
Empirismen. Han blev saaledes hjemme i den kritiske
Philosophis Principer, at han ad sin egen Vei naaede til
flere af de Ideer, Schelling udtalte i sit Identitetssystem,
saa at dette endog paadrog ham en Beskyldning for Plagiat.
Som Digter virkede Coleridge til at indføre Romantiken
i England; som Tænker og Publicist virkede han til
at udbrede Kundskab til den tydske Philosophis Grundtanker.
Selv en saa afgjort Modstander af disse, som
Stuart Mill, stiller ham ved Siden af Bentham og betragter
disse to Mænd som de to befrugtende Aander for deres
Samtid. Mill anerkjender, at Angrebene paa Lockes Philosophi
saavel fra den tydske som den skotske Skole væsentlig
have ansporet den psykologiske Forsken og gjort
opmærksom paa de Vanskeligheder, der endnu vare uovervundne.
(Dissertations and discussions I.). Buckle omtaler
i sin Civilisationshistorie foruden den skotske Skoles Indflydelse
ogsaa »den pludselige og vel fortjente Beundring for den
tydske Litteratur, hvilken især Coleridge gav Udtryk, og som
% --- p. 14 ---
\setcounter{footnote}{0}%
\opage{14}indgjød den engelske Aand Smag for høiere og dristigere
Anskuelser, end den hidtil havde kjendt«; men han levede
desværre ikke længe nok til at fuldføre en Skildring af
denne Bevægelse.

Den Indvirkning, som den kritiske og spekulative
Philosophi øvede paa Empirismen, bestod for det Første i
at føre denne til en skarpere Bevidsthed om sit eget Princip
og en konsekventere Gjennemførelse deraf, hvorved saa
paa den anden Side den Overbevisning maatte udvikle sig,
at Empirismen nødvendigvis viser udover sig selv. Hume
havde forkastet Kausalitetsbegrebet, men han anerkjendte
Mathematikens Sætninger som nødvendige, fordi de grundede
sig paa Modsigelsens Grundsætning og ikke paa ydre Erfaring.
Kant bemærkede hertil, at Hume konsekvent ogsaa
maatte have negtet Nødvendigheden paa det mathematiske
Omraade. Dette anerkjender Stuart Mill og lærer, at alle
mathematiske Axiomer grunde sig paa Erfaring, paa fast
og stadig Ideassociation, og at Modsigelsens Grundsætning
kun er Udtrykket for det mest abstrakte Resultat af Erfaringen.
Først herved har Empirismen naaet sin Fuldendelse.
Men naar al Nødvendighed saaledes negtes, maa
det Spørgsmaal naturligvis opkastes, hvorledes Empirismen
vil kunne godtgjøre sin egen Sandhed. Herbert Spencer
hævder her overfor Stuart Mill, at hvis hele den empiriske
Erkjendelse ikke skal svæve i Luften, maa der være et Kriterium,
som fra først af og helt igjennem ligger til Grund
for den, og i Kraft af hvilket ethvert Fremskridt i Erkjendelse
gjøres. »At sige, at der ikke gives nødvendige
Sandheder, er det Samme som at sige, at Slutningen ikke
med Nødvendighed indeholdes i sine Præmisser men er en
empirisk Dom, som ikke sikkert er sand. Den rene Empirisme
gaaer ud fra en Uvished og bevæger sig gjennem
en Række af Uvisheder. Man kan standse den ved ethvert
Skridt, saalænge den ikke har gjort Rede for sin oprindelige
Forudsætning.« Før man antager Noget for sandt, maa
man gjøre sig Rede for, hvad Betingelsen i det Hele er
for, at Noget skal være sandt. Bevidstheden maa ved
% --- p. 15 ---
\setcounter{footnote}{0}%
\opage{15}Undersøgelse af sine egne Tilstande komme til en Afgjørelse
af, hvilken Egenskab det egentlig er, der menes, naar Noget
sættes som sandt.

En af Vanskelighederne ved den tidligere empiriske
Skoles Psykologi var den af Coleridge fremhævede, at den
sjælelige Aktivitet forblev uforklaret. Her har den engelske
Psykologi gjort et Skridt, som Stuart Mill træffende sammenligner
med det Skridt, hvorved Laromiguière gik ud
over Condillac, idet han fremhævede Sindets Aktivitet i
Opmærksomheden. Støttet til physiologiske Kjendsgjerninger
har nemlig Alexander Bain, en engelsk Psycholog, der staaer
Mill nær, paavist et oprindeligt aktivt Element som Udgangspunkt
for det sjælelige Liv, nemlig den primitive, automatiske
Bevægelse, som gaaer forud for al Sandsning,
som ved de enkelte Sandsninger ledes i bestemte Retninger,
men uden hvilken disse Sandsninger ikke kunne finde
Sted, --- den første Spire til Tankens og Villiens Aktivitet.


Blev Øiet saaledes aabnet for Sider af Virkeligheden,
man tidligere havde været blind for, saa blev det tillige
ledet til dristigere og mere omfattende Synspunkter. Ved
Siden af Philosopherne optræde her Historikerne og Naturforskerne.
Ideen om Menneskeslægtens og Menneskeaandens
Udvikling paa de forskjellige Omraader baner sig Vei gjennem
den af Kant og Herder grundlagte Historiens Philosophi; den
spekulative Naturphilosophi og den komparative Naturvidenskab
vise, at Udviklingens Ide ikke er indskrænket i sin
Gyldighed til Aandens Verden, men at den omfatter Natur
og Aand som Led af samme Kjede. Den ældre Skoles
ængstelige og snevre psykologiske Empirisme forsvandt for
disse storartede Anskuelser. Der maatte da nødvendigvis
fremstille sig den Opgave at samle alle enkelte Undersøgelser
til en omfattende Verdensanskuelse, behersket af
ét ledende Princip. Det er en Opgave, hvortil der kun
findes spredte Bidrag hos de tidligere engelske Tænkere,
men som de store philosophiske Systemer paa Fastlandet
vare Forsøg paa at løse. Ogsaa vor Tids engelske Philo%
% --- p. 16 ---
\setcounter{footnote}{0}%
\opage{16}sopher fastholde, at vor Erkjendelse kun er relativ og
phænomenal og ikke kan gribe det Absolute; men selve den
relative Viden kan dog organiseres til et sammenhængende
Hele, og da det er umuligt for Tanken at holde det Relative
dualistisk sondret fra det Absolute, som den under en
eller anden Form maa statuere, saa se vi ogsaa den dybsindigste
af disse Philosopher, Herbert Spencer, anerkjende,
at der ved Udmaalingen og Artikulationen af den relative
Videns Omraade vindes indirekte Bestemmelser af den
absolute Kraft, til hvilken Udviklingens Ide, den herskende
Ide i Phænomenernes Rige, viser tilbage.

Idet paa denne Maade den empiriske Philosophi
nærmer sig Spekulationen, dels ved Anerkjendelse af de
oprindelige Forudsætninger, hvorpaa ogsaa Spekulationen
bygger (det aprioriske Sandkedskriterium; Selvets oprindelige % PRINTED AS IS (Sandkedskriterium: O-A14)
Aktivitet), dels ved Bestræbelsen for en altomfattende Verdensanskuelse,
der i det Betingede finder Bestemmelser for
det Ubetingede, i det Relative indirekte Former for det
Absolute, --- nærmer den sig særlig den Retning i vor
Tids tydske Philosophi, som kan kaldes Virkelighedsidealismen,
fordi den gaaer ud fra Betragtningen af Virkeligheden
og kommer til det Resultat, at denne i sidste Instants
er af ideal Natur.\footnote[1]{Se min Bog „Philosophien i Tydskland efter Hegel“ pag. 220---319.} En lignende Gang finde vi i Nutidens
engelske Philosophi. Her er hos Englænderne en
stærkere empirisk og skeptisk Tendents, men den praktiske
Idealisme, der er saa karakteristisk for disse tilsyneladende
saa prosaiske Forskere, maa dog tilsidst ogsaa føre
spekulative Konsekventser med sig, forsaavidt ikke allerede
den theoretiske Interesse leder i denne Retning. Ogsaa
hvad den tydske »Virkelighedsidealisme» angaaer, gjælder
det, at det Absolute væsentlig opfattes som videnskabeligt
og praktisk Ideal. Den Tid synes i det Hele at være forbi,
da man ad den dialektiske Begrebsudviklings Vei troede
% --- p. 17 ---
\setcounter{footnote}{0}%
\opage{17}at kunne gribe det Absolute. Selv det høieste og klarest
tegnede Gudsbegreb er dog kun et Symbol, et Grændsebegreb,
hvortil vi komme, men som vi ikke kunne gjennemføre.
Dog maa det fastholdes, at vi med Nødvendighed
komme til dette Begreb. At det maa være en Hypothese,
er sandt; men det er en Hypothese, vi under en eller anden
Form nødes til at opstille, og det Spørgsmaal kan derfor
ikke skydes tilside i Philosophien, hvilken Hypothese der
kommer Sandheden nærmest. Den spekulative Kritik har
her Arbeide nok, et Arbeide, der aldrig vil ende, saalænge
Menneskeaanden ikke bliver træt af at sysle med sine høieste
Ideer. Ogsaa dette er et af de Punkter, hvori den engelske
Philosophi adskiller sig fra den franske Positivisme. August
Comte vilde vel ophæve al Theologi og Metaphysik, men
statuerede dog et Absolut, et »høieste Væsen,« som faldt
sammen med selve Menneskeheden, betragtet som et Hele
i Fortid, Nutid og Fremtid, i hvis Tjeneste og for hvis
Udvikling alle Enkelte skulle arbeide og fuldstændig opofre
sig. Herimod gjælder den Indvending, som Spencer fremfører,
at det Ubetingede ikke kan søges paa noget enkelt
Omraade, men maa være den universelle Kraft, som rører
sig i Alt og ligger til Grund for hele Tilværelsen. ---

Af denne almindelige Karakteristik kunne vi uddrage
det Resultat, som i den følgende Fremstilling vil finde sin
nærmere Begrundelse, at den engelske Philosophi i vor Tid
viser os, hvorledes Empirismen, netop naar den gjennemføres,
fører ud over sig selv. Som vi have set, tager vor
Tids engelske Philosophi sit Udgangspunkt i den tidligere
engelske Skoles Principer, men ved den Udvikling, den giver
disse Principer, fører den dem ud over Skolens Begrændsning.
Vi maa derfor tilskrive den en ikke ringe Betydning
for Philosophiens videre Udvikling og Bearbeidelsen af dens
universelle Problemer. ---

Som allerede bemærket, forstaa vi ved den engelske
Philosophi særlig den empiriske og realistiske Retning.
Men vi maa dog for at give et klart Billede af den philosophiske
Bevægelse i England ogsaa fremdrage Repræsen%
% --- p. 18 ---
\setcounter{footnote}{0}%
\opage{18}tanter for denne Retnings Modsætning. Efter derfor først
at have beskjæftiget os med Stuart Mill, Empirismens
Fuldender, og Psychologen Bain, som i det Væsentlige
slutter sig til ham, ville vi i Korthed karakterisere Mills
Modstandere Whewell og Hamilton, som især ere paavirkede
af den kritiske Philosophi. Det tredie Afsnit af vor Fremstilling
vil give en Udsigt over Herbert Spencers Lære, som
paa engang har sine Forudsætninger i den empiriske og
den kritiske Skole, som klarest og bestemtest fører ud over
Empirismens Grændser, og som betegner den storartede
Afslutning af de philosophiske Bestræbelser i England i
vor Tid.

\chapterhead{Andet Kapitel.}{Den rene Empirisme.}


I. \emph{John Stuart Mill} (f. 1806, d. 1873) modtog,
som man ser af hans Autobiographi, allerede tidlig
Indtryk, der bleve bestemmende for hans hele theoretiske
og praktiske Retning, Faderen, James Mill, førte ham
tidlig ind i det 18de Aarhundredes Philosophi og lærte ham
at betragte dens Lærdomme ikke blot som Theorier, men
næsten som Udtryk for sin Religion, fordi de alene formaaede
at bekæmpe de herskende Fordomme, som hindre
Menneskehedens Fremskridt til Frihed og Lyksalighed. Det
er igjennem James Mill, at vor Tids engelske Philosophi
staaer i Forbindelse med det 18de Aarhundredes. Han var
væsentlig dannet ved Studiet af Hartley og Hume og betragtede
ligesom disse to Tænkere Forestillingernes Association
som den psykologiske Grundlov, en Opfattelse, som
allerede er forberedt af Thomas Hobbes. Skjøndt Stuart
Mill ikke særlig har behandlet Psykologien, ligger dog denne
psykologiske Theori til Grund for hele hans System; til
den vise, som han stedse selv fremhæver, alle hans logiske,
% --- p. 19 ---
\setcounter{footnote}{0}%
\opage{19}metaphysiske og ethiske Undersøgelser tilbage. Vi maa
derfor, baade for at vise den historiske Sammenhæng
mellem det 19de og det 18de Aarhundredes engelske
Philosophi og for særlig at forstaa Stuart Mills Standpunkt,
i Korthed omtale den nævnte psykologiske Theoris tidligere
Former.

I sin Udvikling fra Locke til Hume førte det 18de
Aarhundredes Philosophi til at eliminere Selvets Begreb og
opfatte Sjælelivet blot som en Række af Forestillinger og
Fornemmelser, en Række, som man saa igjen søgte at
simplificere ved at føre alle Forestillinger tilbage til Sandsefornemmelser.
Lovene for Sjælelivet opfattedes da ikke
som Udtryk og Former for Selvets Energi, men som
Formler for det mekaniske Vexelforhold mellem de enkelte
Led i Rækken af sjælelige Phænomener, i sidste Instants
af Sandsefornemmelser. En fuldstændig Lære om
disse Love, Lovene for Ideassociationerne, søgte først David
Hartley (f. 1705, d. 1757) at give i sit Værk »Observations
on man.» Hartley giver Associationstheorien et
physiologisk Grundlag, idet han, sluttende sig til Newtons
Emissionstheori, som den store Naturforsker ikke blot anvendte
til Forklaring af Lysvirkningerne, men ogsaa til at
forklare Nervevirkningerne, fører alle Sandsninger og Forestillinger
tilbage til Svingninger i Hjernesubstantsen, fremkaldte
ved Indtrykkene udenfra. Naar nu visse bestemte
Svingninger ofte komme igjen, maa der i Hjernesubstantsen
opstaa en Disposition til at gjentage dem under visse Betingelser.
Hvilke disse Betingelser ere, lære vi af Ideassociationernes
Love, som vi finde ved psykologisk Analyse.
Medens Hume endnu skjelner mellem tre forskjellige
Associationslove: Association paa Grund af Lighed, Sammentræf
(Kontiguitet) og Kausalitetsforhold mellem Forestillingerne,
saa fører Hartley disse Love tilbage til en
eneste, Sammentræffets, Berøringens (Kontiguitetens) Lov:
to Sandsefornemmelser forbindes eller associeres, naar de
Indtryk, der fremkalde dem, indtræffe samtidig elter umiddelbart % PRINTED AS IS (elter: O-B01)
efter hinanden. Naar da senere det ene Indtryk ind%
% --- p. 20 ---
\setcounter{footnote}{0}%
\opage{20}træder igjen, vil der være en Disposition til ogsaa at fremkalde
den til det andet Indtryk svarende Forestilling, eftersom
der i Hjernen vil være en Disposition til at gjentage
den til dette andet Indtryk svarende Svingning. En Forestilling
er kun Resten af et tidligere Indtryk, beroende paa
en svagere Svingning, der er ligesom Efterklangen af den
ved selve Indtrykket fremkaldte stærke Svingning. Ved
Sammensmeltning af enkelte Forestillinger formedelst deres
hyppige Gjentagelse i Forening med hverandre opstaaer en
sammensat Ide, som tilsidst synes at danne en uopløselig
Enhed, fordi de enkelte Forestillingselementer, hvoraf den
er sammensat, forsvinde fra Bevidstheden, idet enhver enkelt
af dem overvindes af Summen af alle de andre. Den
sammensatte Ide synes tilsidst slet ikke at have noget
Slægtskab med de Forestillinger og Sandsninger, af hvilke
den er bleven til, ligesom et Lægemiddel i sin Smag kan
være helt forskjelligt fra ethvert enkelt af de Stoffer, hvoraf
det er tilberedt. De sammensatte Ideer kunne nu igjen
sammensættes med hverandre, og saaledes i det Uendelige.
Deraf kommer det, at mange Ideer synes at opstaa ligesom
af sig selv, indenfra, hvilke man da, som Locke, har udledt
af en egen sjælelig Evne, Reflexionen; men en fortsat
Analyse vil kunne føre alle, selv de høieste og mest sammensatte
Ideer tilbage til Enkeltforestillinger, og disse
igjen til Sandseindtryk. Paa Følelsen af Lyst og Ulyst
øve Ideassociationerne en opdragende Magt: de sandselige
Fornemmelser vige efterhaanden for de ideelle Følelser, idet
de umiddelbare Gjenstande for Lyst og Smerte træde tilbage
for Forestillinger, der ikke i sig selv medføre sandselig
Nydelse eller Lidelse. Herpaa beroer Evnen til at uddanne
sin Villie og hæve sig op over det Sandselige. Og ad
denne Vei uddanne den moralske Sands og de moralske
Ideer sig. Men paa den anden Side beroer ogsaa al Fordom
og al Indflydelse af Autoritet paa Ideassociationens
Magt.

Coleridge har med Rette fremhævet den strenge Konsekvents,
hvormed Hartley gjennemfører sin Theori. Naar
% --- p. 21 ---
\setcounter{footnote}{0}%
\opage{21}Sjælelivet opfattes som beroende paa mekaniske Forbindelser
mellem enkelte Elementer, kan det, ligesom ved Forbindelsen
mellem materielle Atomer, kun være den umiddelbare Berøring,
der medfører Forbindelsen. Derfor staaer hos
Hartley den psykologiske Hypothese ogsaa i nøie Forbindelse
med den physiologiske Theori. Thi vel er der, som
Lewes fremhæver i sin Kritik af Hartley, ikke nogen saadan
Lighed mellem physiologiske og psykologiske Phænomener,
at vi af hine kunne lære Noget om disse,
og vort Kjendskab til Fornemmelserne er større end vort
Kjendskab til deres physiologiske Grundlag; men Associationspsykologien
beroer jo netop paa, at de sjælelige
Phænomener opfattes i Analogi med de legemlige. Naar
derfor den physiologiske Hypothese opgives, har Associationspsykologien
allerede mistet sin absolute Karakter.

James Mills (f. 1773, d. 1836) psykologiske Lære
adskiller sig fra Hartleys ved Opgivelse af den physiologiske
Hypothese og ved en bestemtere og omhyggeligere Gjennemførelse
af den psykologiske Analyse. Hans »Analysis
of the phenomena of the human mind» udkom 1829 og
blev 1869 udgivet paa Ny af hans Søn, forsynet med særdeles
interessante Noter af denne, af Psykologen Bain, af Philologen
Findlater og af Historikeren Grote. I disse Noter
gives der Berigtigelser og Forklaringer og føres Undersøgelsen
videre, saa at man her har ligesom et Repertorium
for den empiriske Psykologi i England. --- Det er hos
James Mill ligesom hos Hartley Kontiguitetens (Berøringens)
Lov, der lægges til Grund. Fornemmelser, der pleie at
indtræffe samtidig eller umiddelbart efter hinanden, associere
sig, og Forestillingerne følge efter eller med hverandre
ligesom de Sandsefornemmelser, hvis Billeder de ere. Foruden
Hyppigheden af visse Fornemmelsers Sammentræf, er
det ogsaa den Interesse, dette Sammentræf har for Individet,
der styrker Associationen; Interessen leder os til at
lægge Mærke til de hyppige og konstante Sammenhænge,
thi paa dem grunde vi Planerne for vor Handlen. Snart
danner der sig da for Bevidstheden en Modsætning mellem
% --- p. 22 ---
\setcounter{footnote}{0}%
\opage{22}konstante og variable Forhold. Jeg har ligesaa ofte set
Fuglen flyve fra Vest til Øst som fra Øst til Vest, men
jeg har altid set Stenen tilsidst falde til Jorden, naar den
blev kastet, aldrig vedblive at fare i den modsatte Retning.
To Forestillinger, som paa denne Maade stedse vise sig at
høre sammen, indgaa en uopløselig Forening, en fast
Association, og heri bestaaer, hvad vi kalde Tro eller Antagelse.
Naar vi se det ene Phænomen indtræde (f. Ex.
Kastet af Stenen), vente vi ogsaa, at den andet (Stenens
Fald) vil følge efter. ---

Den psykologiske Mekanisme, som Hartley og James
Mill lærte, ligger nu, som allerede bemærket, til Grund
Stuart Mills hele Lære. Den danner den Forudsætning,
til hvilken han paa ethvert Punkt kommer tilbage. Men
han antager dog ikke Associationslæren ganske i den Form,
hvori hine ældre Forskere fremstillede den; hans egne Undersøgelser
førte ham paa flere afgjørende Punkter ud over
den, som man navnlig ser af hans Noter til Faderens
Værk. Dette bringer da ikke faa Inkonsekventser ind i
hans Lære, idet han ikke er bleven sig bevidst, at de Berigtigelser,
han foretager ved den ældre Theori, ofte ligefrem
føre til fuldstændig at ophæve dennes Forudsætninger,
hvilke han dog mener at kunne fastholde. Saa klar og
aaben en Aand, som denne Tænker er, saa stor en Lyst
og Glæde, han finder i Debatten med andre Standpunkter,
og saa uforbeholdent han anerkjender det, naar han er
bleven sig Vanskeligheder bevidst, hans Theori ikke kan
løse, saa mærker han dog ikke, at han her har forladt sine
oprindelige Forudsætninger. Han omtaler i sin Autobiographi,
hvilken Kamp han havde at bestaa for at frigjøre
sig fra de ensidige Synspunkter, Faderens stærke Paavirkning
allerede tidlig havde bibragt ham; og vi ville finde
Sporene af denne Kamp i de antydede indre Modsigelser i
hans Lære. Associationspsykologien er det mægtige Vaaben,
han retter imod sine Modstandere, men den er tillige en
Vægt, der tynger paa ham selv og gjør ham ude af Stand
til at fuldende, hvad han stræber efter at naa. Ganske
% --- p. 23 ---
\setcounter{footnote}{0}%
\opage{23}i Almindelighed kan hin Modsigelse siges at fremkomme
ved, at Associationstheorien hviler paa en udvortes og
mekanisk Opfattelse af Sjælelivet, der vel er berettiget som
empirisk Grundlag, men som overser en Rest, der her er
Sagens egentlige Kjerne, nemlig det egentlige Selv, den
subjektive Energi, som maa være tilstede, om end i rudimentær
Form, i den dunkleste Fornemmelse, og som er
den Magt, der ene kan bringe Enhed og Sammenhæng i
Forestillingskredsen. Ideassociationen er en af de vigtigste
reale Faktorer i Sjælelivet, men kan selv kun opfattes ret
ved at ses som en elementær Form for Selvets Energi.
Naar man, som vi ville see Stuart Mill gjøre, indrømmer,
at denne Energi udgjør Aandens egentlige Væsen, og at
den ikke kan forklares ved Associationslæren, saa er det
inkonsekvent vedblivende at betragte dennes Principer som
eneste Grundlag for Psykologien.

Dersom Stuart Mills Forhold til den tidligere engelske
Philosophi havde været mere bekjendt, vilde man ikke saa
ofte have set ham betegnet som Positivist og som Discipel
af Auguste Comte. Han har selv udtrykkelig afvist denne
Betegnelse i sit Skrift: »Auguste Comte and Positivism.»
Ikke blot var hele hans Anskuelse grundlagt, før Comtes
Skrifter udkom, men han afviger ogsaa i meget væsentlige
Punkter fra den store franske Philosoph. Som dennes
Hovedfortjeneste fremhæver Mill hans Opfattelse af Menneskeaandens
Udvikling gjennem tre Stadier, det theologiske, det
metaphysiske og det positive, og den Maade, hvorpaa han
derved begrunder Videnskabernes Klassifikation. Men til
denne Anerkjendelse føier Mill nogle eiendommelige Indskrænkninger.
Fordi det indrømmes, at vi kun erkjende
Phænomenerne og deres Love, medens det Absolute ligger
ud over vor Rækkeevne, følger deraf ikke, at det theologiske
Stadium skulde være aldeles overvundet i den Forstand,
at der skulde komme en Tid, da Menneskeslægten ganske
holdt op med at spørge efter den oprindelige Aarsag, det
absolute Princip til Alt. Den »positive» Philosoph kan
herom danne sig den Mening, han anser for den rette; den
% --- p. 24 ---
\setcounter{footnote}{0}%
\opage{24}vil da blive Gjenstand for den positive Kritik; men det er
ikke nødvendigt, at alle Tænkere stemme overens i disse
Sager. »Den, som betragter alle Begivenheder som Led i
en fast Orden, idet enhver af dem er den uforanderlige
Følge af en foregaaende Betingelse eller af en Kombination
af foregaaende Betingelser, antager fuldstændig den positive
Maade at tænke paa, hvad enten han anerkjender eller
ikke anerkjender en universel Grund, hvis oprindelige Følge
hele Naturens System er, og hvad enten denne universelle
Grund opfattes som Aand eller ikke.» Fordi fremdeles det
positive Stadium skal afløse det metaphysiske, følger det
ikke, at det ikke længere skulde være nødvendigt at anstille
Undersøgelser og Analyser af de Forudsætninger og
Principer, hvorpaa vor Erkjendelse hviler. Men Betydningen
af saadanne Undersøgelser er Comte tilbøielig til at overse,
ligesom han slet ikke har givet Logiken og Psykologien
Plads i sit System (smlgn. ovenfor S. 9-10). Den almindelige
Videnskabslære skal løse to Opgaver: den skal
vise os de Veie, ad hvilke den menneskelige Tanke naaer
Erkjendelsen, og de Midler, ved hvilke Erkjendelsens Sandhed
kan godtgjøres. Men Comte har kun behandlet den
første Opgave, ikke den anden; han har ikke indset Verifikationens
Betydning, som den induktive Logik paaviser.
Dermed hænger det sammen, at Comte betragtede Hypotheserne
som Midler til at tilfredsstille vor intellektuelle
Trang til Klarhed og Sammenhæng mellem Phænomenerne,
i Stedet for som Veie til at finde Phænomenernes reale
Aarsager. Han erklærer endog, at »det mest philosophiske
Synspunkt tilsidst fører til at opfatte Studiet af Naturlovene
som bestemt til at fremstille den ydre Verden for os, med
Tilfredsstillelse af vor Forstands væsentlige Tilbøieligheder»,
til hvilke han regner Trangen til Orden og Harmoni, til
Skjønhed og Idealitet. Tyngdepunktet truer altsaa her med
at flytte sig fra det Objektive til det Subjektive, skjøndt
Positivismen efter Comtes Lære netop skulde udelukke alt
Subjektivt. Med Rette finder Mill heri Spiren til Comtes senere
mystiske Spekulationer, hvor den »subjektive Syn%
% --- p. 25 ---
\setcounter{footnote}{0}%
\opage{25}these» skulde træde i Stedet for den objektive. Mill er
her mere positivistisk end Comte, men han forlader Positivismen
ved at hævde Logiken som særegen Videnskab. ---
Ligesom Comte oversaa den logiske Analyses Betydning,
saaledes negtede han ogssa Psykologiens Realitet som % PRINTED AS IS (ogssa: blind read b1-48-11)
selvstændig Videnskab, fordi den hviler paa Selviagttagelse,
der efter hans Anskuelse kun fører til subjektive Illusioner.
I Stedet for Psykologien vil han sætte Phrenologien. Mill
hævder Psykologiens Selvstændighed som en Videnskab, der
har sine særegne Principer og Love, nemlig Ideassociationens
Love, der ere for Psykologien, hvad Tyngdeloven er for
Astronomien, og hvad Vævenes elementære Egenskaber ere
for Physiologien. --- Philosophien har saaledes for Mill en
dybere Grund og et videre Omfang end for Comte. Mill
selv har som Tænker en videre Horizont og en større Bevægelighed
end Comte; men han staaer paa den anden
Side tilbage for den franske Philosoph i streng Konsekvents
og systematisk Afsluttethed. ---

Stuart Mills Hovedværk er hans Logik. Han har her
med Mesterskab skildret den induktive Forskens Methode
i dens fineste Nuancer og er mere end nogen anden Logiker
kommen Ansigt til Ansigt med Erkjendelsens virkelige,
praktiske Funktioner. Det vil være vanskeligt at finde
noget Sidestykke til dette Værk, naar man ikke gaaer tilbage
til Aristoteles. Hvad denne gjorde for Syllogismen,
for den deduktive Logik, har Mill gjort for Induktionen,
for den experimentale Forskens Logik. Allerede Baco og
Descartes polemiserede mod Syllogismen og negtede dens
Betydning som Erkjendelsesmethode; men den nye Logik,
Baco lovede, kunde dengang endnu ikke fremstilles. Ligesom
Aristoteles først kunde opstille Theorien om Syllogismen,
efterat Reflexionen og Dialektiken var bleven øvet og udviklet
af Sokrates og Plato, saaledes kunde Induktionens
Theori først gjennemføres, efterat den experimentale Forsken
havde arbeidet længe og paa forskjellige Omraader.
Og da viste det sig ogsaa, at den formelle (syllogistiske)
Logik ikke var nok, men at der behøves en Lære om den
% --- p. 26 ---
\setcounter{footnote}{0}%
\opage{26}objektive Sandhed og Vished, en »philosophy of evidence»,
der ikke blot er analytisk og sikrer mod Selvmodsigelse,
men virkelig fører Erkjendelsen videre, ad Synthesens Vei
beriger med nyt Indhold.

Ved at slutte forstaaer man i Almindelighed at udlede
en Sætning af en anden. Men dette er ifølge Mill ikke
Slutningens oprindelige Form. Oprindelig slutter man
hverken fra det Almindelige til det Enkelte eller fra det
Enkelte til det Almindelige, men fra et Enkelt til et andet
Enkelt. Sandsningen viser os to Phænomener (a og b)
forbundne enten ved Samtidighed eller umiddelbar Succession,
og ifølge Associationens Lov fremkalder Opfattelsen af
det ene Phænomen (a) siden Forestillingen om det andet (b).
Dette er den mest elementære Slutning, som er rent instinktmæssig
og mekanisk, og som foretages selv af Dyr
og smaa Børn. Naar Barnet engang har brændt sig (b)
paa den varme Kakkelovn (a), vil Synet af Kakkelovnen
(a) siden fremkalde Forestillingen om Smerten ved at brænde
sig (b), og derfor »skyer brændt Barn Ilden.» --- Men naar
de Erfaringer, hvoraf Slutningen skal drages, ere mangfoldige
og adskilte i Tid, maa de fastholdes og sammenholdes
ved et Fællesnavn; kun derved bliver det muligt at
gjøre sig skarpt og bevidst Rede for den Operation, man
foretager. I Stedet for den elementære Slutning, der gik
fra det ene Enkelte til det andet Enkelte, træder altsaa nu
en mere indviklet Proces. Først opstilles ved Induktion
(i snevrere Betydning) en almindelig Sætning som Udtryk
for de givne Erfaringers Sum (ai + a2 -f- a3 = A), og af
denne almindelige Sætning (A) udlede vi saa det andet
Enkelte (b) ved syllogistisk Deduktion. Men denne indviklede
Proces er kun en anden Form for hin elementære
Slutning: dels er det Almene kun Summen af det Enkelte,
og dels have vi ligesaa megen Ret til at drage en Slutning
om en hel Klasse som om et enkelt Phænomen, forsaavidt
det besidder Klassens Attributer. Syllogismen, Slutningen
fra det Almene til det Enkelte, er kun en Prøve,
vi anstille, det Almene selv er kun Summen af den tid%
% --- p. 27 ---
\setcounter{footnote}{0}%
\opage{27}ligere Erfaring, et Slags Register, en kort Formel, som ikke
selv kan være oprindeligt Udgangspunkt.

Vi se saaledes Associationspsykologien ligge til Grund
for Mills logiske Theori. Men hvorledes komme vi da ud
over det Tilfældige i dette Grundlag? Hvorledes kan en
virkelig Indsigt og Erkjendelse udvikle sig af den blinde,
instinktmæssige Succession af Forestillinger? Stuart Mill
siger selv, at hvad Menneskene kunne begribe, og hvad de
ikke kunne begribe, for en stor Del er tilfældigt og afhænger
af deres Erfaring og Vaner. Ved særlig at kultivere
visse Ideassociationer kunne Menneskene gjøre sig aldeles
uskikkede til at begribe mange Ting, medens de paa
den anden Side kunne begribe Meget, som atter synes
Andre ubegribeligt. Hvorledes kunne vi da være sikre paa,
at vi pleie de \emph{rigtige} Ideassociationer? Det maa Logiken
nødvendigvis belære os om, da den ellers ikke vilde være,
hvad den efter Mills Definition skal være, »thi science of % PRINTED AS IS (thi science: O-B07: thi for the)
the conditions of correct thinking» Men det ses let, at
Mill kun ved en Inkonsekvents kan opstille et Begreb om
det Rigtige, hvorved dette bliver forskjelligt fra det faktisk
Antagne. De virkelig opstaaede Ideassociationer kunde
ikke have dannet sig anderledes, og lige saa lidt vil man
ifølge Mills Forudsætninger egentlig kunne forandre og
»kultivere» dem, thi hvorfra skal denne Dannelse udgaa,
naar der i Bevidstheden ikke er andre Faktorer end de
enkelte Forestillinger, og disse kun virke efter Associationens
uvilkaarlige Love?

Ifølge den almindelige Opfattelse gjælder det at fastholde
de Forbindelser mellem Phænomenerne, der have
Nødvendighedens Karakter, eller med andre Ord, at ordne
Phænomenerne ifølge Aarsagens Princip. Men Mill er ikke
tilfreds med dette Svar. Ved Nødvendighed, siger han,
tænker man sig i Almindelighed et mystisk Baand, der
forener Tingene, ved Aarsag Noget, der ikke blot \emph{stedse}
\emph{gaaer forud for} noget Andet, men som ligefrem \emph{frembringer}
noget Andet. Intet af disse Begreber kan
ifølge Mill den induktive Logik bruge. Den finder kun en
% --- p. 28 ---
\setcounter{footnote}{0}%
\opage{28}Række af Phænomener, af hvilke nogle bestandig indtræde
i Forening med hverandre, --- den eneste Sammenhæng,
den kjender. Om Nødvendighed kan her kun være Tale,
forsaavidt man derved forstaaer \emph{ubetinget} Forbindelse, idet
det ene Phænomen aldrig er tilstede, uden at det andet
ogsaa indtræder. Om Aarsag kan her kun være Tale,
forsaavidt man derved forstaaer Totalsummen af de positive
og negative Præcedentser, ved hvis Indtræden et vist
Phænomen uforanderlig ogsaa indtræder. Aarsagen er
derfor altid phænomenal, og Spørgsmaalet er, hvorledes vi
opnaa Erkjendelsen af de phænomenale Aarsager. --- Man
ser let, at Mill i sin Iver for at skyde den indre, i sig
grundede Nødvendighed tilside, ikke mærker, at han kun
henlægger Vanskelighederne et andet Sted, nemlig \emph{mellem}
Phænomenerne i Stedet for i deres Indre. En ubetinget
Forbindelse kan den empirisk-induktive Betragtning, kan i
hvert Tilfælde vor paa Ideassociationer beroende Erkjendelse
ikke naa. Naar Mill alligevel statuerer et saadant Ubetinget,
saa er det ikke blot en logisk Tilsnigelse; det er ogsaa
fordi der her, som det senere vil vise sig, ubevidst virker
et Princip, der fører ud over hans bevidste Forudsætninger.
---

Det er navnlig ved Angivelsen af de Veie, ad hvilke
de phænomenale Aarsager kunne erkjendes, at Mill har
ført Induktionens Logik videre. Baco var ikke naaet ud
over den tilfældige Induktion, som bestaaer i at samle og
opsummere de Tilfælde, der tale for, og dem, der tale imod
(de positive og de negative Instantser); Experimentets
Theori var det forbeholdt Mill at opstille. Og han har
gjort det med en saadan Klarhed og Fuldstændighed, at en
Naturforsker som Liebig har følt sig fremmet i sine Undersøgelser
ved Studiet af hans Logik.

Der er ifølge Mill fire Methoder for den empiriske
Forsken. De ere dels Opdagelsesmethoder, dels Bevismethoder.
Vi give her en ganske kort Karakteristik af
dem. 1) Overensstemmelsesmethoden: Naar to eller flere
Tilfælde af det Naturphænomen, der skal udforskes, kun
% --- p. 29 ---
\setcounter{footnote}{0}%
\opage{29}have én eneste Omstæadighed fælles, saa kan kun denne % PRINTED AS IS (Omstæadighed: O-B09)
Omstændighed være den søgte Aarsag (eller Virkning).
2) Uoverensstemmelsesmethoden: Naar et Tilfælde, i hvilket
Naturphænomenet indtræder, og et Tilfælde, i hvilket det
ikke indtræder, have alle Omstændigheder fælles paa én
nær, der kun forekommer i det første Tilfælde, saa er
denne den søgte Aarsag (eller Virkning). Denne Methode
fører os ved direkte Experiment med Vished til Aarsagen.
Hvor Experimentet ikke er muligt, maa vi lade os nøie
med Overensstemmelsesmethoden, der ikke fører os til
egentlige Kausallove, men foreløbig kun til phænomenale
Ensformigheder. 3) Restmethoden: Vi drage fra Phænomenet
den Del, der ved tidligere Induktioner er erkjendt som
Virkning af visse Betingelser; Resten af Phænomenet maa
da være Virkning af Resten af de forhaandenværende Betingelser.
4) De hverandre ledsagende Forandringers Methode:
Naar et Naturphænomen forandres ved en Forandring
i et andet Naturphænomen, maa det staa i Kausalsammenhæng
med dette. 

Det er ikke alle Opdagelser, vi umiddelbart kunne
gjøre ved Hjælp af disse Methoder. De forudsætte alle Experimenter
eller dog Iagttagelser. Men Naturphænomenerne
ere ofte for indviklede til, at man direkte kan skride til
at undersøge dem. De virkende Betingelser kunne være
saaledes blandede mellem hverandre i den virkelige Verden,
at en ved Hjælp af de experimentale Methoder fremadskridende
Forsken aldrig vilde kunne blive færdig med dem.
Her viser den deduktive Methode sig som et nødvendigt
Supplement til den induktive. Deduktionen forudsætter dog
stedse en forudgaaende Induktion. Deduktionens Opgave
er at finde Loven for en Virkning ud fra de forskjellige
Faktorer, hvis Produkt den er. Men saa maa enhver af
disse Faktorer først særlig være undersøgt ved Iagttagelse
eller Experiment. Af Lovene for de enkelte Faktorers
Virkning udledes saa Lovene for Virkningen af deres Kombination.
Denne Deduktion er en Syllogisme, der ogsaa
kan have mathematiske Grundsætninger til sine Præmisser,
% --- p. 30 ---
\setcounter{footnote}{0}%
\opage{30}saa at Virkningen af Kombinationen findes ved ligefrem
Beregning. Men Deduktionens Resultat maa, for at kunne
anerkjendes som videnskabelig Sandhed, verificeres ɔ: den
maa ved omhyggelig Sammenligning vise sig at stemme
med Fakta, som ere fundne ved direkte Iagttagelser. Saaledes
maatte Newtons Tyngdelære verificeres ved, at den ved
Deduktion førte til de keplerske Love.

At forklare et Naturphænomen eller en Sammenhæng
mellem Naturphænomener vil sige at føre det tilbage til
et eller andet bekjendt Phænomen. Men Forklaringen
sætter nu egentlig kun en ny Gaade i Stedet for den tidligere;
ud over Phænomenerne kunne vi aldrig komme.
Naturens Væsen vedbliver at være lige hemmelighedsfuldt,
hvor langt vi end rykke frem. Den høieste Opgave, Erkjendelsen
her kan stille sig, er, ifølge Mill, at naa de
færreste Antagelser, af hvilke Naturens Orden og Sammenhæng
kan forklares. Naturforklaringen finder i Besvarelsen
af dette Spørgsmaal sin bestemte Grændse. Paa den ene
Side kunne Naturens sidste Love ikke være mindre talrige
end de kvalitative Forskjelligheder mellem vore Fornemmelser
og aandelige Tilstande. Ere vi naaede saa vidt, at vi have
paavist en særegen Art af Aarsag for enhver af de uopløselige
Bevidsthedstilstande, saa kunne vi ikke komme
videre. Men foruden denne subjektive Grændse er der paa
den anden Side ogsaa en objektiv Grændse. Der er sidestillede
Aarsager i Naturen, som vi ikke kunne føre tilbage
til nogen fælles Lov. Vi kjende ingen Lov for Fordelingen
af de oprindelige naturlige Kræfter og Substantser
eller for det bestemte Forhold, der faktisk finder Sted
imellem dem. De elementære Kræfter staa i forskjellige
Forhold til hverandre, uden at vi formaa at opdage nogen
Regel deri. Vi kjende Lovene for Faktorernes Virken, men
ikke for deres oprindelige Forening. Saaledes er der i et Kaleidoskop
altid en regelmæssig Orden mellem Glasstykkerne,
og dog beroer Anordningen hver Gang kun paa en tilfældig
Rystelse.

Ifølge den almindelige Opfattelse er der i det Mindste
% --- p. 31 ---
\setcounter{footnote}{0}%
\opage{31}én Lov, som gjælder ud over  al Erfaring, nemlig den, at
enhver Ting har sin Aarsag, Kausalitetens Lov. Selve
Induktionen, i Betydning af en Slutning fra enkelte Tilfælde
til en almen Lov, hviler efter denne Opfattelse forsaavidt
ogsaa paa Kausalitetens Princip, som det, der giver
Tanken Ret til at slutte fra det Enkelte til det Almene,
kun kan være Forudsætningen om, at der overalt i Naturen
er en bestemt, lovmæssig Sammenhæng. Mill opfatter nu,
som vi have set, Induktionen væsentlig som en Slutning
fra et Enkelt til et andet Enkelt, der oprindelig foregaaer
som en aldeles mekanisk og instinktmæssig Ideassociation.
Erfaringen er altsaa ikke Andet end disse Ideassociationers
Sum med Udskillelse af de tilfældige og forbigaaende. Han
kan altsaa ikke anerkjende Kausalitetsloven som oprindelig
virkende Princip, men betragter den som en abstrakt
Generalisation af al den Erfaring, vi (hidtil) have gjort.
Et Absolut kunne vi ad denne Vei aldrig naa. Vi finde
Ensformigked og lovmæssig Sammenhæng i alle de enkelte % PRINTED AS IS (Ensformigked: O-B13)
Kredse af Phænomener, vi undersøge, og drage deraf den
Slutning, at der i \emph{hele} Naturen maa findes en saadan
Lovmæssighed. Egentlig have vi kun Ret til at sige, at
den findes i den os bekjendte Del af Naturen. »Jeg har den
Overbevisning,» siger Mill, »at Enhver, som er vant til Abstraktion
og Analyse, og som oprigtig bruger sine Evner, efter
tilstrækkelig at have øvet sin Indbildsningskraft ikke vil finde % PRINTED AS IS (Indbildsningskraft: unattested-word sweep, zoom z10)
nogen Vanskelighed ved at forestille sig, at der, f. Ex. i et
af de mange Firmamenter, i hvilke Astronomien nu inddeler
Universet, kan følge Phænomener efter hverandre paa
tilfældig Vis og uden nogen bestemt Lov; heller ikke findes
der Noget i vor Erfaring eller aandelige Organisation, som
kan afgive nogensomhelst Grund til at tro, at det ikke er
Tilfældet.» Det er en petitio principii, en Antagen af,
hvad man netop skulde bevise, naar man grunder Erfaringens
Mulighed paa Kausalitetsloven i Stedet for at grunde
Kausalitetsloven paa Erfaringens Virkelighed. For at gjøre
den Slutning: »Ild brænder,» behøver man ingen Forudsætning
om, at hele Naturen følger ensformige Love; det
% --- p. 32 ---
\setcounter{footnote}{0}%
\opage{32}er nok, naar der blot hersker Ensformighed i Ildens Virkninger
--- men denne Ensformighed forudsætte vi jo ikke;
den lære vi netop af Erfaringen. Og som det gaaer med
hver enkelt Kreds af Erfaringer, gaaer det med Erfaringen
som Helhed.

Den Induktion, hvorpaa Kausalitetsloven grunder sig,
er ikke experimental, men bestaaer blot i, hvad Baco kalder
inductio per enumerationen simplicem, en ligefrem Sum af % PRINTED AS IS (enumerationen: O-B14)
de enkelte Tilfælde, hvori Regelen har holdt Stik, uden
Deduktion og Verifikation. Hvis det var et specielt Faktum,
det dreiede sig om, vilde dette være en høist utilstrækkelig
Grund at bygge paa; men det Usikre i Bevismaaden forringes
ved, at vi her have at gjøre med den mest abstrakte
Generalisation, som omfatter vor Erfarings hele
Omraade. Hvad der gjælder om Kausalitetsloven, gjælder
ogsaa om de mathematiske Grundsætninger; ogsaa de ere
kun Induktioner fra Erfaringen. Men de angaa saa simple
og saa hyppig forekommende Forhold, at man ikke selv
mærker, at man foretager Induktionen. Alligevel vilde
Geometrien have været en physisk Videnskab, dersom dens
Gjenstand ikke lige saa godt havde kunnet betragtes i et
Billede eller ideel Konstruktion, som i selve Virkeligheden,
og dersom den ikke havde havt sin demonstrative Karakter,
en Karakter, som iøvrig ogsaa andre Grene af Naturvidenskaben
efterhaanden antage. Der er altsaa i Mathematiken
ingen høiere eller strengere Nødvendighed end andre Steder
indenfor vor Erfaring. A=C, dersom A=B og B=C; ---
to rette Linier, som engang have skaaret hinanden, ville
bestandig divergere; --- saadanne Sætninger ere induktive
Sandheder, hvis Gyldighed ligesom Kausalitetslovens kun
beroer paa, at de stedse ere befundne rigtige og ikke en
eneste Gang falske. En absolut Umulighed kjende vi ikke;
Grunden til, at vi ikke kunne antage, at to rette Linier,
som skjære hinanden, kunne konvergere, er derfor kun, at vi
aldrig have set det, Et fuldkomnere Bevis for almindelige
Sandheder have vi ikke.

Om disse Hovedpunkter, om Kausalitetsprincipets og
% --- p. 33 ---
\setcounter{footnote}{0}%
\opage{33}de mathematiske Axiomers aprioriske, ,eller aposterioriske
Gyldighed, har Mill ført en lang og interressant Strid, med % PRINTED AS IS (interressant: O-B17)
den Mand, hvis Værk over de induktive Videnskabers Historie
efter Mills egen Erklæring for en stor Del gjorde
det muligt for ham at fuldende sin Logik, William Whewell.
Ifølge Whewell, som i det Væsentlige hylder Kants Lære,
forudsætter enhver Erkjendelse ikke blot, at Phænomenerne
fremstille sig paa en vis Maade, men ogsaa at den Erkjendende
medbringer den rette Ide, det rette Synspunkt til
deres Opfattelse. Al Erfaring forudsætter derfor visse Ideer
eller Grundbegreber. Men, siger Mill, disse Ideer ere selv
kun Resultater, der erhverves ved Abstraktion af Erfaringen.
Naar man vil udtrykke en Iagttagelse i Ord, maa man
henføre den til tidligere vundne Begreber. Gives der nu
ikke et Begreb, hvorunder den kan henføres, saa maa man
danne et helt nyt Begreb som Udtryk for de nye Kjendsgjerninger.
Det var en Tilfældighed, at Kepler iforveien
havde Ideen om en Ellipse, da han studerede Marsbanen;
derfor kunde han beskrive og betegne den som en Ellipse;
ellers maatte han have dannet denne Ide ved Sammenligning
af Phænomenerne, for saa at anvende den til deres
Fortolkning.

De logiske Grundsætninger maa dele Skjæbne med de
mathematiske; de ere lige saa lidt som disse oprindelige
Principer, men de høieste Generalisationer af vor Erfaring.
Der kan altsaa ikke tilkomme dem en Nødvendighed, som
gaaer ud over, hvad Erfaringen lærer. En nødvendig Sætning
maatte være en saadan, som vi ikke kunne Andet
end tænke, hvis Modsætning vi ikke kunne fatte. Men
hvad vil det sige, at Noget er os ufatteligt? Det vil kun
sige, at der i forveien har dannet sig uadskillelige Ideassociationer
i vor Bevidsthed, der ikke tillade os at tænke det.
Enhver tilsyneladende universel og absolut Sætning grunder
sig paa en saadan fast Association. Kun ved at være opmærksom
paa disse Ideassociationers Betydning kan man
undgaa at hildes i subjektive Fordomme, som alle Tiders
Religion og Spekulation afgive Exempler paa. Derfor var
% --- p. 34 ---
\setcounter{footnote}{0}%
\opage{34}det en Velgjerning mod Menneskeheden, at James Mill i
sin »Analysis» gjennemførte den uimodstaaelige Ideassociation
som Princip for Forklaringen af alle sjælelige Phænomener.
Der gives Fakta, som i al vor Erfaring vise sig
forbundne, og som derfor ogsaa i vor Bevidsthed maa bevirke
en fast og ikke ved nogen Tvivl rokket Antagelse.
Men paa den anden Side indeholder Videnskabens Historie
Exempler nok paa, at hvad der til en Tid blev betragtet
som umuligt, senere har vist sig virkelig at være sandt.
Man antog f. Ex. i Middelalderen, at det var umuligt, at
der kunde gives Antipoder, og at Stjernerne for ikke at
falde ned maatte være befæstede til faste Sphærer. Naar
vi stadig se, at »ufattelige» Ting i Tidens Løb blive fattelige,
saa kan altsaa Ufatteligheden af det Modsatte af en
Sætning, ikke være et Kriterium for denne Sætnings
Sandhed.

Whewell bemærker hertil, at naar en Tid forkaster
som utænkeligt, hvad en anden Tid anser for sandt, eller
omvendt, da kommer det af, at Bevidstheden stadig udvikler
sig fra det Ufuldkomnere til det Fuldkomnere, en
Udvikling, under hvilken der stadig opstaa nye Ideer og
Synspunkter; men det rokker ikke den Sandhed, at det
sidste, mest abstrakte Kriterium, vi have, er Umuligheden
af det Modsatte. Ogsaa Lewes og Spencer optræde her mod
Mill, skjøndt de ellers i det Væsentlige ere enige med ham.
Spencer fremhæver navnlig (se ovenfor p. 14), at uden
et saadant Kriterium vilde vor Erkjendelse ikke kunne
foretage et eneste gyldigt Skridt. --- Mills sidste Ord er,
at selv om der kunde paavises visse Forhold, som i og
for sig ere utænkelige for os ifølge vor Natur, vilde dette
Utænkelige derfor godt kunne existere. At Noget er utænkeligt,
kan kun have rent subjektiv Betydning, med mindre
man, som Schelling og Hegel, vilde antage, at Tankens og
Tilværelsens Muligheder falde sammen, at der er en oprindelig
Overensstemmelse mellem Mikrokosmen og Makrokosmen;
men det er en Antagelse, som kun kunde begrundes
ved en overnaturlig Aabenbaring. Saalænge, siger
% --- p. 35 ---
\setcounter{footnote}{0}%
\opage{35}Mill tilsidst, endnu saa fremragende Tænkere som Spencer
og Lewes kunne betragte en spekulativ Nødvendighed, en
Nødvendighed for Tanken, som Bevis for Tilværelsen af en
real Nødvendighed i Universet, saa længe maa det siges, at
den gamle Metaphysik endnu hersker i Philosophien. Have
vi Myndighed til at sige, at der er Noget, Gud eller Naturen
ikke kan gjøre? Kjende vi da alle Naturens Kombinationer?
Der gives ingen Nødvendighed udenfor Naturens
egen Virken, og denne Virken lære vi kun at kjende gjennem
Erfaringen.

Vi kjende ikke det Værende i sig selv; derfor er al
Ontologi umulig. Den spekulative Philosophi søger forgjæves
at fastsætte Love for Væren i sig selv. Metaphysiken er
ikke philosophia prima men philosophia ultima, det sidste
Resultat af al vor Erfaring. Vore almindelige Ideer indeholde
Intet, som ikke er lagt ind i dem ved vore Erfaringer
eller vore Tankevaner. Der gives ingen abstrakte Ideer.
Ved sammensatte Ideer, som svare til existerende, individuelle
Gjenstande, ere vi i Stand til udelukkende at henvende
vor Opmærksomhed paa enkelte af Karaktermærkerne. Det
kommer derved til at se ud, som om vi kunde opfatte dem
adskilte fra Resten, især da vi ved Tegn og Ord kunne
fastholde det Lignende og oftere Indtrædende. Naar vi f.
Ex. se en hvid Hest, ere vi i Stand til at fæste vor Opmærksomhed
blot paa Karaktermærket »hvid», der saa gjenfindes
ved andre Gjenstande; og da vi have dannet et
særegent Ord for det, sige vi, at det er »den samme Egenskab»
i alle Tilfælde. Men det er dog stadig Forestillingerne
om konkrete Phænomener, som udgjøre Begrebernes Indhold,
og vi tænke ikke ved Hjælp af Begreber, men ved Hjælp
af konkrete Forestillinger og af de Navne, hvormed vi fastholde
disse Forestillinger. --- Mill nærmer sig altsaa her
den extreme Nominalisme, for hvilken Begreberne blive Et
med blotte Navne. ---

Efterat man har gjort sig bekjendt med Mills Logik,
vil man ikke uden Forundring kunne høre ham udtale, at
dette Værks Resultater slet ikke staa i nogen  nødvendig
% --- p. 36 ---
\setcounter{footnote}{0}%
\opage{36}Forbindelse med nogen af de to philosophiske Hovedretningers
eiendommelige Lærdomme. Thi ganske vist gjælder dette
for saadanne Partiers Vedkommende som det særdeles vigtige,
der fremstiller »den experimentale Forskens Methoder»;
men hele det Lys, hvori de logiske Processer stilles, hele
den Forklaring, der gives af deres Natur og Grund, bekræfter
dog i fuldt Maal, hvad vi fra Begyndelsen af fremhævede,
at det er en ganske bestemt, fra den tidligere empiriske
Skole arvet Lære, Associationspsykologien, der ligger
til Grund fra først til sidst. Associationspsykologien viser
igjen tilbage til Lockes Lære om Aanden som en ubeskreven
Tavle; kun naar Aanden opfattes som fuldkommen passiv,
kan dens Væsen udtømmende forklares ved Ideassociationens
Love. Andre Steder regner Mill sig ogsaa til den philosophiske
Skole, hvis Hovedrepræsentant Locke er, medens
den modstaaende repræsenteres af Kant. Men han negter,
at Hovedforskjellen mellem de to Skoler skal bestaa i, at
den ene betragter Aanden som rent passiv i Erkjendelsen,
medens den anden tillægger den en oprindelig Aktivitet.
Ogsaa den empiriske Skole anerkjender ifølge ham et oprindeligt
mentalt Element i vore Ideer og betragter ikke
Begreberne som blotte Kopier af Indtrykkene men som
Produkter af aandeligt Arbeide. Men i samme Afhandling,
hvori denne Udtalelse findes (Dissertations and discussions
III.), paaviser Mill det netop som Bains Fortjeneste at have
ført Associationspsykologien ud over den passive Determinisme
til Anerkjendelse af en oprindelig Aktivitet i Aanden.
Heri ligger den Indrømmelse, at den ældre empiriske Skole
\emph{ikke} lærte Aandens Aktivitet. Og det er et stort Spørgsmaal,
om ikke en af de væsentlige Eiendommeligheder ved
denne Skole dermed opgives. Det er i ethvert Tilfælde
umuligt baade at fastholde hin oprindelige Aktivitet i Selvets
Væsen og at betragte Ideassociationens Love som de eneste,
der gjælde for Aandens Liv. Den oprindelige Energi maa
jo dog præge ethvert Indtryk paa en eiendommelig Maade
og maa ligeledes være bestemmende for de Forbindelser, i
hvilke Forestillingerne træde med hverandre indbyrdes. Aner%
% --- p. 37 ---
\setcounter{footnote}{0}%
\opage{37}kjendelsen af en oprindelig Aktivitet fører os egentlig fra
Locke til Kant.

Stuart Mill sætter Hovedforskjellen mellem de to
philosophiske Skoler deri, at den ene (den aprioriske, den
intuitive eller den introspektive Skole) gaaer ud fra nogle
af de mest sammensatte Aandsphænomener som oprindelige
og mener i dem at have Grundforudsætninger om Tingenes
Natur og Væsen, medens den anden (den experimentale eller
den aposterioriske) søger at reducere de tilsyneladende oprindelige
Principer ved fortsat Analyse. Det er klart, at
der herved ikke er angivet nogen principiel Forskjel; thi
begge Skoler ville dog vist være enige om, at man skal
føre Analysen saa vidt, som man i det Hele taget kan føre
den. De Analyser, vi finde hos Spinoza og Leibnitz, hos
Kant og Hegel, kunne ogsaa i Klarhed og Konsekvents
maale sig med dem, vi træffe hos den anden Skoles Tænkere.
Og Mill indrømmer selv, at Analysen kun \emph{kan} føres
til et vist Punkt. Vi standse tilsidst ved Bevidsthedstilstande,
der ikke kunne føres tilbage til simplere og mere
elementære. Navnlig ligger Fornemmelsen af Lighed og
Ulighed til Grund for alle andre Fornemmelser. Indviklede
Ligheds- og Ulighedsforhold kunne vi føre tilbage til simple
og usammensatte; men selve Evnen til at opfatte Ligheds- og
Ulighedsforhold maa bestandig forudsættes. Ogsaa her
støde vi paa et Element, der viser ud over det Empiriske,
et Bidrag til Erkjendelsen fra Aandens Side. Opfattelse af
Lighed beroer paa en Sammenlignen, en aandelig Virksomhed.
Vi vises  saaledes tilbage til Aandens Aktivitet.
Kant saa denne bestaaende i en Sammenfatten af det for
Fornemmelsen givne Mangfoldige til Anskuelsens og Tankens
Enhed. Og den sammenlignende Virksomhed viser i
samme Retning; Sammenligning er Betingelsen for Sammenfattelsen.
Al Ideassociation forudsætter, som William
Hamilton har vist\footnote[1]{Dette anerkjender Mill. Se hans Note til James Mills „Analysis“ I p. 112 f.}, denne Sammenlignen; thi for at en
% --- p. 38 ---
\setcounter{footnote}{0}%
\opage{38}Forestilling skal kunne fremkalde en anden, med hvem den
pleier at følge, maa den først gjenkjendes, ses som en, vi
før have havt; og Gjenkjendelsen sker ved Sammenlignen.
Deraf følger igjen, at selve Ideassociationen kun bliver en
elementær Form for Selvets Aktivitet. Men saa er det
Princip, Mill har lagt til Grund for sin logiske Theori, kun
et afledet Princip, og den sidste Forklaring er ikke given;
det vilde kun være Tilfældet, naar de logiske Love førtes
tilbage til hin oprindelige Aktivitet, i hvilken de maa have
deres Spire. Det er Mill selv, der ikke har gjennemført
Analysen. Og naar Sagen forholder sig saaledes, ere vi
ikke blot komne ud over Empirismen, men ogsaa over
Skepticismen. Vor Erkjendelse ses da at hvile paa Forudsætninger,
der ere givne i vor aandelige Organisation, i
vor Bevidstheds Væsen, og som det er os umuligt at ophæve.
Selve det at tænke eller forestille sig Noget medfører
stedse, at Tankens og Forestillingens Indhold ordnes
og forbindes paa en vis Maade; heraf den Forudsætning,
at der overalt maa være lovmæssig Sammenhæng, en Forudsætning,
som kun kan opgives, naar selve Tænkningen
opgives. Den Nødvendighed, vi erkjende, kan indtræde som
Led i en større Sammenhæng, som vi ikke kunne overse;
men derved ophæves den ikke.\footnote[1]{Naar Stuart Mill mener, at man godt kunde tænke sig en Verden, hvor to og to ikke var fire men fem, da har man dertil rigtig svaret, at det i saa Fald ikke vilde være 2 + 2, der var = 5, men 2 + 2 + 1, idet 1 skyldes den Omstændighed, hvorved hin Verdens Indretning er forskjellig fra vor. Se Mansel: The philosophy of the conditioned p. 131.} Denne Forudsætning om
en lovbestemt Sammenhæng findes derfor hos Alle og i
alle Tider, om den end først ved den nyere Naturvidenskab
har faaet sin fulde Bekræftelse og sikre Udførelse. Ogsaa
i de mythologiske Forestillinger, selv i den vildeste Overtro
ligger Overbevisningen om Orden og Sammenhæng til
Grund.

Ved at gaa tilbage til Forudsætningerne have vi saa%
% --- p. 39 ---
\setcounter{footnote}{0}%
\opage{39}ledes fundet Grundfeilen hos Mill. Paa denne Maade faaer
Kritiken større Betydning, end naar man, som det for det
Meste er Tilfældet, argumenterer fra de skeptiske Konsekventser,
som kunne drages af Mills Lære, Konsekventser,
som Mills smidige Dialektik desuden vilde kunne gjøre
uskadelige. Ved at rette Kritiken mod Forudsætningerne
vinde vi tillige den Fordel, at det bliver os muligt at hævde
den store Betydning af Mills Logik (smlgn. ovenfor p. 25 og
28). Den gjør ikke Fyldest som udtømmende Fremstilling af
Erkjendelsesprocessen; men naar det Ensidige og Relative
i Synspunktet er indset, maa man beundre den Energi og
den Skarpsindighed, hvormed Opgaverne behandles. En
Fremstilling, der ængstelig vilde stræbe efter Korrekthed
og Fuldstændighed med Hensyn til de første Principer,
vilde maaske neppe kunne trænge saaledes ind i den
konkrete Erkjendelses Væsen og Former, som Mill fra sit
Standpunkt har formaaet det. ---

Hvad der bringer Mill til at forlade Associationspsykologiens
oprindelige Forudsætninger, er, at han har faaet
Øie paa Selvhedens, Subjektivitetens Begreb. Idet dette
Begreb trænger sig frem, bringer det Uro og Modsigelse
ind i de empiriske Forudsætninger. I sine særlig psykologiske
Undersøgelser kommer Mill dette Begreb endnu
nærmere end i sin Logik. Mod disse vende vi os nu, idet
vi først fremdrage hans Opfattelse af Psykologien i Almindelighed.\footnote[1]{I Logiken omhandles Psykologiens Methode; de andre Undersøgelser findes især i „Examination of Sir William Hamiltons philosophy“ og i Noterne til James Mills „Analysis“.}


Mill hævder, som allerede omtalt, Psykologiens Selvstændighed
som Videnskab. Den er ikke blot en Del af
Physiologien. At enhver aandelig Tilstand har en vis Nervetilstand
til sin umiddelbare Betingelse og nærmeste Aarsag,
kan ikke siges at være godtgjort i samme strenge Betydning,
hvori det er godtgjort for Sandsefornemmelsernes
% --- p. 40 ---
\setcounter{footnote}{0}%
\opage{40}Vedkommende. I ethvert Tilfælde kjende vi ikke disse
Nervetilstandes Væsen, ei heller deres Forhold til de aandelige
Tilstande; vi ere derfor indskrænkede til at studere disse
sidste direkte. Psykologien har til Gjenstand de oprindelige
eller afledede Love, ifølge hvilke en aandelig Tilstand
følger efter en anden. Ganske vist kan her i det Enkelte
ikke være Tale om exakt Indsigt; vi kunne finde Grundlovene
og Grundbetingelserne, men disse modificeres i de
enkelte Tilfælde ved de individuelle Forhold. Mill sammenligner
den psykologiske Viden med Læren om Ebbe og Flod;
ogsaa her staa Hovedaarsagerne (Solens og Maanens Tiltrækning)
fast og kunne forklares og forudsiges; men der
er en Mængde sekundære Faktorer (Vindens Retning, lokale
Omstændigheder, Beskaffenheden af Oceanets Bund etc.),
der hver for sig afhænge af særegne Love, og som øve
betydelige Modifikationer. --- Vi vide allerede, at Ideassociationens
Love ifølge Mill ere de psykologiske Grundlove.
For at finde disse, maa man ligesom i Naturvidenskaben
anvende en dobbelt Methode: dels maa Kjendsgjerningerne
føres tilbage til Love, disse igjen til mere almindelige
Love, og saa fremdeles; dels maa de sammensatte
Phænomener opløses i enkelte og elementære. Den første
Methode har sin Analogi i Newtons Opdagelse af Tyngdeloven
som Forklaring baade af Himmellegemernes Bevægelse og
af Faldbevægelsen paa Jorden; den anden har sit Forbillede
i den chemiske Analyse, ved hvilken Stofferne opløses i
deres Bestanddele. Men naar man saaledes har ført et
sjæleligt Phænomen tilbage til simplere og mere elementære
Phænomener, ved hvis Kombination og Association det er
opstaaet, da er det ikke, som Hartley og James Mill antog,
nok at studere disse elementære Phænomener; man kan
ikke af dem deducere det høiere, mere sammensatte Phænomens
Natur, lige saa lidt som man kan deducere Vandets
Egenskaber af Iltens og Brintens. --- Den egentlige Psykologi
er en Erfaringsvidenskab. Men af de almindelige
Love, den fører til, kan man ad Deduktionens Vei finde,
hvorledes den menneskelige Karakter vil forme sig under
% --- p. 41 ---
\setcounter{footnote}{0}%
\opage{41}visse bestemte physiske og moralske Betingelser. Psykologien
gaaer da over i Ethologien, Læren om den menneskelige
Karakter. Der kan naturligvis ikke være Tale om at forudsige
bestemte Fakta; men af de rigtig erkjendte psykologiske
Love kan man udlede, hvilke Tilbøieligheder og
Bestræbelser der ville fremtræde under visse Forhold. For
at faa Gyldighed maa denne Deduktions Resultater verificeres,
eftervises som stemmende med den virkelige Iagttagelse.
--- Til Psykologien og Ethologien støtter sig endelig Sociologien,
Læren om det menneskelige Samfund. Ogsaa
denne er en deduktiv Videnskab. Efter en Undersøgelse
af enhver især af de Faktorer, der gribe ind i Samfundslivet
og dér øve deres Indflydelse gjennem de almindelige
psykologiske og ethologiske Love, skal alle disse Faktorers
forenede Virken bestemmes. Kombinationen kan her ofte
blive saa indviklet, at Deduktionen ikke kan gjennemføres,
ligesom Pkysiken ikke kan løse Problemet om den gjensidige % PRINTED AS IS (Pkysiken: O-B22)
Indvirkning af tre mod hverandre graviterende Legemer.
Sociologien maa derfor undertiden være omvendt
deduktiv ɔ: i Stedet for at begynde med at deducere fra
de psykologiske og ethologiske Love, maa den begynde med
at udlede en empirisk Lov af historiske og sociale Fakta,
for derpaa at finde en dybere Begrundelse af denne empiriske
Lov i den menneskelige Natur. Ideen om Menneskehedens
Fremskridt og Udvikling er et Exempel paa en saadan
empirisk Lov, der for at blive en rationel Lov, en
Kausallov, maa ses i sit Udspring af selve den menneskelige
Natur.

Man kunde savne Ethiken i denne Oversigt over
Aandsvidenskaberne, og man har undertiden troet, at Mill
ikke anerkjendte anden Ethik end Ethologien, Læren om
Karakterernes Udvikling. Men dette er en Misforstaaelse.
Ifølge Mill er Ethiken væsentlig en Kunst, en Del af den
Kunst, der svarer til Videnskaben om den menneskelige
Natur og det menneskelige Samfund. Det Eiendommelige
for Kunsten i Modsætning til Videnskaben er, at den taler
i Form af Regler og Forskrifter, ikke i Form af Udsagn
% --- p. 42 ---
\setcounter{footnote}{0}%
\opage{42}om Fakta. Kunsten sætter sig et Øiemed som ønskeligt,
definerer det og overgiver det til Videnskaben, for at denne
kan undersøge dets Aarsager og Betingelser. Læren om
Formaalene og de almindelige Grundsætninger, som kunne
udledes af dem, udgjør Teleologien eller hvad de tydske
Philosopher kalde den praktiske Fornufts Grundsætninger,
en Livets Kunst, som bestaaer af tre Dele: Ethik, Politik
og Æsthetik, --- Læren om det Rette, om det Hensigtsmæssige
og om det Ædle og Skjønne. Ved Hjælp af denne
Kunst vindes en Rettesnor for Bedømmelsen af alle Øiemed
og Gjenstande for vore Drifter.

Den Forskjel, Mill saaledes sætter mellem Ethiken og
Ethologien, forudsætter nødvendigvis en Faktor i Aanden,
hvorved denne kan hæve sig op over den psykologiske
Mekanisme, Associationspsykologien lærer. For at forstaa
Overgangen fra Psykologi til Ethik maa vi betragte Selvets
og Villiens Begreber.

Det er en Forudsætning, Mill har tilfælles med Positivismen,
at vi kun erkjende Phænomener; vi kjende Intet
i og for sig, Aanden i sig selv lige saa lidt som Materien.
Den skotske Skole gik ud fra den almindelige Menneskeforstand
og dens instinktmæssige Antagelse af en ydre
Gjenstand som vore Fornemmelsers Aarsag. Men den »almindelige
Menneskeforstand» er noget altfor Kombineret,
til at Associationspsykologien kan erkjende den for noget
Oprindeligt, der ikke kan underkastes Analyse. I Antagelsen
af ydre, af os uafhængige Gjenstande ligger der egentlig
ikke Andet end Antagelsen af, at der foruden vore virkelige,
nærværende Sandsefornemmelser gives en real Mulighed
til mange andre Fornemmelser, som kunne indtræde
under visse Betingelser. Tingene ere ikke Andet end stadige
Muligheder af Sandseindtryk. For os staa de som
Mere end Muligheder, fordi vi uvilkaarlig anvende Substants-
og Aarsagsbegrebet, der kun har Gyldighed indenfor
Phænomenernes Række, paa selve denne Rækkes Mulighed
og Grund. Og Forestillingen om denne reale Substants
vinder i Styrke, naar vi mærke, at ei blot vi selv, men
% --- p. 43 ---
\setcounter{footnote}{0}%
\opage{43}ogsaa Andre gaa ud fra Forudsætningen om dens Existents.
--- Stuart Mill slutter sig altsaa (ligesom James Mill) til
Berkeleys Idealisme og betragter det som en Illusion, naar
man mener, at en materiel Substants er nødvendig til at
bære Phænomenerne, ligesom Inderne mente, at Jorden
trængte til et Støttepunkt og derfor hvilede paa en Elephant.
Som han følger Berkeley med Hensyn til Materiens Begreb,
saaledes synes han at følge Hume med Hensyn til Aandens
Begreb. Aanden er kun Rækken af vore virkelige og
mulige Fornemmelser. Den adskiller sig fra den Række,
som udgjør Materien, derved, at den ogsaa indbefatter de
Forestillinger, som ikke ere stadfæstede ved Erfaring, og
alle vore Tanker, Følelser og Beslutninger, men navnlig
ved, at denne Række er indskrænket til os selv, hver for
sig, og ikke som hin fælles for alle Mennesker under samme
Forhold.

Denne Fornyelse af Berkeleys og Humes Phænomenisme
er en nødvendig Konsekvents af Associationspsykologien.
Ligesom denne Lære i forrige Aarhundrede fremkaldte
Reids Appel til den sunde Menneskeforstand og den
almindelige Moral, saaledes har den ogsaa nu vakt Angsten
for den absolute Skepticisme, og Mill har anset det for
nødvendigt at godtgjøre, at Beviserne for Guds Tilværelse
og Sjælens Udødelighed ere «lige saa gode» ifølge hans
Lære som under de sædvanlige Forudsætninger. For os
har det størst Interesse at se, hvorledes Mill ogsaa her aldeles
bestemt gaaer ud over sine egne Forudsætninger og
godtgjør deres Utiltrækkelighed. Strax efter (i Skriftet om % PRINTED AS IS (Utiltrækkelighed: O-B24)
Hamilton) at have givet den angivne Forklaring af Jegets
Begreb erklærer han nemlig, at denne Forklaring ikke
slaaer til, at den navnlig ikke omfatter Phænomener som
Erindring og Forventning. «Naar vi altsaa», siger han,
»tale om Aanden som en Række af Fornemmelser, maa
vi føie til, en Fornemmelsesrække, der er sig selv bevidst
som forbigangen og tilkommende. Og vi staa da ved det
Dilemma enten at antage, at Jeget eller Aanden selv er
noget fra de virkelige og mulige Fornemmelsers Række
% --- p. 44 ---
\setcounter{footnote}{0}%
\opage{44}Forskjelligt, eller ogsaa at antage det Paradox, at en Række
af Fornemmelser kan være sig selv bevidst som en Række.»
Jeget og Erindringen, siger han i sine Noter til James
Mills «Analysis» (hvor man i det Hele vil finde det klareste
og skarpeste Udtryk for hans psykologiske Anskuelse), ere
kun to Sider af samme Kjendsgjerning: uden Erindring
intet Jeg, og omvendt: uden Jeg ingen Erindring. Med
Rette kalder Mill det et Paradox, at en «Række» skulde
kunne være sig selv bevidst; men naar dette er saa, hvad
bliver der da af den Definition af Aanden som Mill og
Associationspsykologien antage, at den er en Række af
Fornemmelser? Er det saa ikke nødvendigt at gribe det
første Alternativ og antage, at Aanden eller Jeget ikke
falder absolut sammen med Fornemmelsernes Række? Dette
synes Mill ogsaa at gjøre, naar han erklærer, at Baandet,
som sammenholder Rækken er lige saa realt som selve de
enkelte Fornemmelser, der danne Rækken, og ikke er et
blot Tankeprodukt, og at vi her have et oprindeligt Element,
der ikke har Noget tilfælles med nogensomhelst anden Ting
og derfor maa betegnes med sit eget Navn. Naar han
ogsaa erklærer, at vi her staa ved en uløselig Gaade, da
maa man dog bemærke, at han selv paa en kunstig Maade
for en stor Del har frembragt dette Gaadefulde ved at slaa
en Definition af Aandens Væsen fast, om hvilken han strax
efter maa indrømme, at den ikke rammer den egentlige
Kjerne. --- Det er ikke let at forstaa, hvorledes Mill har
kunnet blive beskyldt for Materialisme; thi det er klart, at
forsaavidt han lader sig fortrænge fra sit empirisk-skeptiske
Standpunkt, maa han netop gaa i idealistisk Retning. Udgaaende
fra den empiriske Psykologi viser han os Nødvendigheden
af en Psykologi, der, som det siden Kant har
været den tydske Philosophis Opgave, udvikler Sjælelivets
centrale Begreber, navnlig Selvbevidsthedens Begreb. Og
en indre Forbindelse mellem de to Undersøgelsesrækker vil
blive mulig, dersom det kan lykkes at vise, at hele det
System af mekaniske Processer, som Associationspsykologien
% --- p. 45 ---
\setcounter{footnote}{0}%
\opage{45}fremstiller, er tjenende Grundlag og Middel for Selvbevidsthedens
Udvikling. ---

I nær Forbindelse med Selvets Begreb staaer Villiens.
Med Hensyn til Opfattelsen af dette Begreb har Mill en
Støtte i sin logiske Theori, som han anvender her paa en
interessant Maade. Han negter, som vi have set, Nødvendigheden
og sætter i dens Sted fast og ensformig Følge.
Det forekommer ham, at man i Almindelighed ved Nødvendighed
forstaaer en mystisk Magt, som udenfra overvælder
den individuelle Villie. Opfattes Nødvendigheden
paa denne Maade, bliver Nødvendighedslæren (Determinismen)
til Fatalisme, som hos Muhamedanerne, ifølge hvis Tro det
Forudbestemte vil ske, hvad man saa gjør for at hindre
det, og som hos Owen, ifølge hvis Lære Menneskets Karakter
dannes for ham og ikke ved ham. Mill vil derfor
helst afskaffe Udtrykket Nødvendighed. Enhver Handling
forudsætter som sin Betingelse tidligere Bevidsthedstilstande
(Tanker, Følelser og Drifter), uden hvilke den ikke vilde
være bleven til. Virkning følger her efter Aarsag lige saa
regelmæssig som paa det physiske Omraade; men om indre,
mystisk Nødvendighed er der her lige saa lidt Tale som
hist. Vi kunne selv udvikle vor Villie ved at anvende
Ideassociationens Love til at skaffe de ædlere Motiver voxende
Magt. Oprindelig bestemmes Handlingen ved Lyst og
Ulyst; men ved Ideassociationens Magt bliver efterhaanden
Midlet Gjenstand for vort Ønske uden Hensyn til det oprindelige
Formaal: Handlingen villes da i sig selv, ganske
afset fra, at vi finde Behag i, hvad den forskaffer os. Forskjellen
mellem en slet og en god Mand bestaaer ikke i,
at den sidste handler imod sine stærkeste Tilbøieligheder,
men deri, at hans Stræben efter at handle ret og hans
Modbydelighed for det Urette ere stærke nok til at holde
enhver anden Drift Stangen og bringe den til Taushed.
Da nu en saadan aandelig Tilstand er mulig, ere Menneskene
skikkede til Selvopdragelse og moralsk Ledelse. Selvopdragelsen
er Enhvers moralske Pligt. Et nødvendigt
Middel til den er, at Tanken danner sig en klar Maalestok
% --- p. 46 ---
\setcounter{footnote}{0}%
\opage{46}for Ret og Uret, for at moralsk Drift og Afsky kunne virke
paa rette Maade. --- Mills Lære viser tydelig nok ogsaa
paa dette Punkt ud over hans oprindelige Forudsætninger.
Thi hvorledes kan Begrebet Selvopdragelse faa Realitet, naar
Selvet ikke opfattes som en oprindelig Energi, der formaaer
at gjennemtrænge og udvikle de givne Betingelser (Indtryk
og Forestillinger) og lede Ideassociationen i en vis bestemt
Retning? ---

Spørgsmaalet om Villiens Frihed var et af de
Punkter, i hvilke Mill adskilte sig fra den skarpsindige
skotske Philosoph, mod hvem han har rettet et af sine
vigtigste Skrifter, William Hamilton, i hvis Lære den frie
Villie i Betydning af absolut Valgfrihed spiller en Hovedrolle,
idet den, og den af den betingede moralske Ansvarlighed,
efter hans Overbevisning var det eneste Grundlag,
hvorpaa Troen paa Guds Tilværelse kunde bygges. Men
der er ogsaa andre Punkter af Mills Polemik, som vi her
maa fremdrage til Karakteristik af hans Standpunkt. Hamilton
hævder energisk Erkjendelsens Relativitet; det Absolute
og det Uendelige kunne vi ikke tænke, fordi vor
Tanke ikke kan opfatte det, som er afsluttet i sig selv, og
som ingen Grændser har. Skjøndt Mill naturligvis maa være
enig med Hamilton i Hævdelsen af Erkjendelsens Relativitet,
saa polemiserer han dog mod denne Motivering.
Det er, mener han, rene Abstraktioner, Hamilton kæmper
imod. \emph{Det} Absolute og \emph{det} Uendelige er i sig meningsløst.
Hvis det Absolute virkelig existerede, maatte det jo
besidde alle mulige Prædikater i den fuldeste Udstrækning:
det maatte være absolut godt og absolut ondt, absolut vist
og absolut dumt, absolut mægtigt og absolut svagt, kort
sagt, det vilde være selvmodsigende. Naar Hegel alligevel
har opstillet et saadant Begreb, saa har han dermed kun
indlagt sig den Fortjeneste at have opløst den spekulative
Metaphysik ved en Række reductiones in absurdissimum.
--- Det er klart, at Mill aldeles ikke har forstaaet Hegels
Tanke eller i det Hele havt Øie for det Problem, der ligger
i Begrebet om det Absolute. Hans Styrke ligger i det
% --- p. 47 ---
\setcounter{footnote}{0}%
\opage{47}Hele i Behandlingen af enkelte, konkrete Spørgsmaal; han
samler ikke de enkelte Philosophemer til en Totalanskuelse.\footnote[1]{„Det er det fulde, det hele Syn, han mangler; hans Aands store Dygtighed udfolder sig ved Gjenstande, der efter hans Behandling af dem vedblive at være lige saa isolerede, som de vare i Forveien; de bringes ikke til Enhed og Harmoni med hverandre; Philosophens Haand behøves endnu stedse med Hensyn til dem.“ Quarterly Review. July 1872. p. 81.}
Han kan uden at føle noget intellektuelt Savn lade Erkjendelsen
ende i en tilfældig Flerhed, og for saa vidt han
opkaster Spørgsmaalet om Gudsideens Realitet, staaer dette
Spørgsmaal ikke i Forbindelse med hans egentlige Philosopheren.
Den intellektuelle Drift, der ikke hviler, før den
kan faa de spredte Traade forbundne, eller dog i det
Mindste se dem konvergere, denne Drift er ham fremmed.
Hans empiristiske og prosaiske Tankegang, der, som man
har bebreidet ham,\footnote[2]{Ibid. p. 117.} ikke føler »Livets Mysterium», anser
det for muligt, om han end ikke selv synes at følge denne
Udvei, at slaa sig til Ro i en Deisme, hvor det Overnaturlige
henlægges til »Verdens Begyndelse», og hvor Guddommen
ligesom hos Epikur lever i ligegyldig Fjernhed fra
Verden. Derimod har det været Philosophiens Hovedproblem
lige fra dens første Dage at søge et ubetinget Princip,
som sidste Forklaring og Forsoning af Tilværelsens Modsætninger.


Mill mener, at man i Stedet for \emph{det} Absolute, \emph{det}
Uendelige, skal sætte noget Absolut, \emph{noget} Uendeligt (et
absolut og uendeligt Væsen), det Konkrete i Stedet for det
Abstrakte, og han anser det ikke for umuligt at idealisere
de høieste Egenskaber, vi kjende, og derved naa en relativ
Erkjendelse af Guds Væsen. I det Mindste hævder han
denne Mulighed overfor Hamilton for at afskjære de theologiske
Konsekventser, denne drog af Umuligheden af at
erkjende det Ubetingede. Man mærker, at Mill her er
kommen ind paa et Omraade, hvor han ikke er hjemme.
% --- p. 48 ---
\setcounter{footnote}{0}%
\opage{48}Fra Hamiltons Standpunkt har man træffende svaret ham,\footnote[1]{\emph{Henry Mansel}: The philosophy of the conditioned. p. 115.}
at »noget Absolut», »et absolut Væsen» er et selvmodsigende
Begreb for os, da ethvert »Noget», ethvert
»Væsen», vi kjende, er relativt og endeligt, saa at det tiltilføiede Absolut »absolut» fører os ud over vor Erkjendelses % PRINTED AS IS (tiltilføiede: til- repeated at the line turn; blind read b1-31-07)
Grændser.

Det er efter Mills egen Erklæring i hans Autobiographi
ikke saa meget theoretiske som praktiske Motiver, der
førte ham i Kamp mod Hamilton, i hvem han saa den
betydeligste Repræsentant i England for den hans egen
modsatte philosophiske Retning. »Forskjellen mellem de
to philosophiske Skoler, den, som støtter sig til Intuition,
og den, som støtter sig til Erfaring og Association, er
ikke blot den abstrakte Spekulations Sag; den er fuld af
praktiske Konsekventser og ligger ved Roden af alle de
største Modsætninger i praktiske Anskuelser i en Fremskridtets
Tidsalder. Den praktiske Reformator maa bestandig fordre
Forandringer i Ting, der støttes af mægtige og vidt udbredte
Følelser, eller betvivle den tilsyneladende Nødvenvendighed % PRINTED AS IS (Nødvenvendighed: O-C01: the syllable doubled at the line turn)
og Urokkelighed af bestaaende Kjendsgjerninger;
og det er ofte en uafviselig Del af hans Indlæg at godtgjøre,
hvilken Oprindelse disse mægtige Følelser have, og
hvorledes disse Kjendsgjerninger kom til at synes nødvendige
og urokkelige. Derfor er der et naturligt Fjendskab
mellem ham og en Philosophi, der tilsidesætter Forklaringen
af Følelser og moralske Kjendsgjerninger ved Omstændigheder
og Ideassociationer og foretrækker at behandle dem
som oprindelige Elementer i den menneskelige Natur, en
Philosophi, som har Tilbøielighed til at hævde yndede Lærdomme
som intuitive Sandheder og betragte Intuitionen
som Naturens og Guds Stemme, der taler med en høiere
Autoritet end vor Fornufts.» Man tager neppe feil, naar
man i denne Udtalelse finder det dybeste Motiv til Stuart
Mills hele Standpunkt. Men hvad selve Sagen angaaer,
% --- p. 49 ---
\setcounter{footnote}{0}%
\opage{49}maa det bemærkes, at Mill her karakteriserer den »intuitive»
Skole, som om den skotske Skole var en værdig
videnskabelig Repræsentant for den, og at man blot behøver
at nævne Navne som Spinoza, Kant og Fichte for strax
at godtgjøre det Uberettigede i Mills Antagelse, at kun den
empiriske Skole skulde repræsentere Aandsfrihedens og
Fremskridtets Philosophi.

Men Mill bekæmper i Hamilton ikke blot den »intuitive»
Philosophi, ogsaa Theologien, Supranaturalismen.
Hamiltons Polemik mod det Ubetingedes Philosophi (Schelling
og Hegel) havde nemlig til Baggrund en moralsk-religiøs
Tro, der ikke skulde kunne finde noget Udtryk eller
nogen Begrundelse gjennem Erkjendelsen, fordi dens Gjenstand
er det Ufattelige, det over alt Menneskeligt absolut
Ophøiede. Imod denne Lære udtaler Mill sig saaledes:
»Dersom man i Stedet for det glade Budskab, at der
existerer et Væsen, i hvem alle de Fortrin, den høieste
menneskelige Tanke kan begribe, existere i en for os ubegribelig
Grad, forkynder mig, at Verden styres af et Væsen,
hvis Attributer ere uendelige, uden at vi kunne forstaa,
hvad de ere, lige saa lidt som hvilke Principerne for hans
Styrelse ere, undtagen at »»den høieste menneskelige Moralitet,
som vi ere i Stand til at tænke os»», ikke godkjender
dem, --- overbevis mig derom, og jeg vil bære min
Skjæbne saa godt jeg kan. Men naar man fortæller mig,
at jeg skal tro dette og til samme Tid kalde dette Væsen
ved de Navne, som udtrykke og bekræfte den høieste
menneskelige Moralitet, saa siger jeg med rene Ord, at jeg
ikke vil. Hvilken Magt et saadant Væsen end kan have
over mig, én Ting er der, som han ikke skal kunne: han
skal ikke tvinge mig til at dyrke ham. Jeg vil ikke kalde
noget Væsen godt, som ikke er, hvad jeg mener, naar jeg
anvender dette Prædikat paa mine Medskabninger; og dersom
saadant et Væsen kan dømme mig til Helvede, fordi jeg
ikke vilde kalde ham saaledes, saa vil jeg gaa til Helvede.»\footnote[1]{Examination of Sir William Hamiltons philosophy p. 103.}

% --- p. 50 ---
\setcounter{footnote}{0}%
\opage{50}Det er først gjennem Læsningen af Autobiographien,
at man faaer en fuldkommen Forklaring af den Energi og
den Bitterhed, hvormed Mill udtaler sig. Man ser af den,
at han af sin Fader tidlig førtes til den Overbevisning, at
Troen paa et Forsyn ikke blot var theoretisk uholdbar,
men ogsaa moralsk forkastelig, fordi den nødvendigvis tilslører
alt det Onde og al den Fordærvelse, som findes i
Verden, ja endog fører til Tilbedelse af et Væsen, som med
Vidende og Villie har skabt en saadan Verden! --- Den
eneste religiøse Tro, der er mulig under saadanne Forudsætninger,
maa være en praktisk Tro; den maa have sin
Gjenstand forud for sig og ikke bagved sig, i Muligheden,
ikke i Virkeligheden. Stuart Mill synes, skjøndt han intetsteds
tydelig udtaler sig i sit eget Navn om disse Sager,
nærmest at slutte sig til dem, der, »skjøndt de ikke kunne
tro, at Verden kan have en Skaber og Styrer, der besidder
absolut Magt og fuldkommen Godhed, dog have det, der
udgjør alle mulige Religioners egentlige Værd, et idealt
Begreb om et fuldkomment Væsen,  mod hvilket de bestandig
fæste deres Blik som deres Samvittigheds Leder;
og dette Ideal af det Gode staaer i Almindelighed Fuldkommenheden
nærmere end deres objektive Guddom, som
anse sig forpligtede til at finde absolut Godhed hos Skaberen
af en Verden saa opfyldt af Lidelse og saa vansiret af
Uretfærdighed som vor.» --- Det er et Spørgsmaal, om Mill
slipper fra Vanskelighederne ved at stille en subjektiv
Gudsidé i Stedet for en objektiv, og om i det Hele en
ideal praktisk Stræben, saaledes som han forlanger den, kan
være mulig, hvis den Virkelighed, fra hvilken enhver Stræben
jo dog maa gaa ud, og i hvilken enhver Mulighed maa
have sit Tilknytningspunkt, er saa ideforladt, at den manichæiske
Dualisme, efter hvad James Mill lærte, skulde være
den eneste moralsk forsvarlige Anskuelse. Men i ethvert
Tilfælde viser Mill her hen til et Ideal, som ikke kan forklares
ved Ideassociationens Virken men peger tilbage til
dybere Kilder, Kilder, som Mill i saa mange Antydninger
har aabnet for os, men kun for strax at tilstoppe dem igjen.
% --- p. 51 ---
\setcounter{footnote}{0}%
\opage{51}Interessant er det at se, hvorledes han i sin praktiske
Opfattelse af Guddommen mødes med I. G. Fichte, hvis
Gudsidé var den moralske Verdensorden, der realiseres ved
den frie Handlen. Fichte er maaske den af de tydske Philosopher,
Mill mest sympathiserer med; men Fichte er konsekventere
end han, idet han har set, at den praktiske
Idealitet nødvendigvis forudsætter den theoretiske og ikke
kan naaes ad rent empirisk Vei. ---

Som Mill i det Hele i sin philosophiske Anskuelse
bygger paa den tidligere engelske Forskning, medens han
dog ofte fjerner sig mere fra den, end han selv mærker,
saaledes gaaer det ogsaa, hvad Ethiken angaaer. Mills
Ethik er Utilitarianisme, »Moral grundet paa Nytte- eller
Lykkeprincipet.» Dette Princip kommer i den engelske
Philosophi allerede frem hos Baco, hvem Macaulay som bekjendt
har fremstillet som det Nyttiges Prophet. Som den
egentlige Grundlægger af Utilitarianismen synes\footnote[1]{Kfr. \emph{Hallam}: Introduction to the litterature of Europe. (7th ed.). IV. p. 160.}% / PRINTED AS IS (litterature: O-C04)
% CHECK p. 51: star lost in the layer: paired with the mark by order
 Richard
Cumberland (i det 17de Aarhundrede) at maatte regnes;
ifølge ham kunne alle moralske Love reduceres til den ene:
Stræben efter det for alle fornuftige Væsener fælles Gode,
hvorved vort eget Bedste fremmes som en Del af det Hele,
saa at den Enkeltes Vei til sit Vel er Alles Vei til det
fælles Vel. Dog er hos Cumberland det individuelle Hensyn
aldeles underordnet Hensynet til det Almene, til Samfundet.
Da Lockes Philosophi var trængt igjennem, og al Erkjendelse
opfattedes som beroende paa Sandsning, blev ogsaa
Ethiken grundlagt paa Lyksalighedsdriften, paa den egne
Interesse. Whewell (History of Moral philosophy in England)
bemærker, at denne Lære stemmede med den nationale
Tænkemaade, og Tocqueville viser, at den er den naturlige
for et demokratisk Handelsfolk som det amerikanske.
Hertil kommer, at den meget vel kan bringes i Forbindelse
med theologiske Synspunkter. Saaledes hos William Paley,
% --- p. 52 ---
\setcounter{footnote}{0}%
\opage{52}hvis ethiske Værk (udk. 1785) lige ind i dette Aarhundrede
lagdes til Grund ved Universitetet i Cambridge. Paley,
der i sin Ungdom havde vundet Prisen for en Afhandling,
i hvilken han stillede den epikuræiske Moral over den stoiske,
lærte, at man skal gjøre det Nyttige, fordi det er Guds
Villie, og fordi det er nødvendigt til Salighed. Den individuelle
Interesse støttede han ved det Haab og den Frygt,
Religionen fremkalder. Men Nyttetheorien blev verdsliggjort
og derved tillige først ret gjennemført af den berømte Retslærde
Jeremy Bentham († 1832). Benthams Betydning
ligger vistnok nærmest paa den praktiske Reforms Omraade,
og han har særlig Fortjeneste ved sin Polemik mod de
mange Subtiliteter og det traditionelle Formelvæsen, der
gjorde den engelske Ret til et fuldstændigt Chaos, hvor de
skjændigste Misbrug kunde trives. Men selve hans Reformtendents
havde sin Grund og sin Støtte i hans energiske
og methodiske Tænkning, der af Nytteprincipet søgte at
udlede Normer til Bedømmelse af de bestaaende Love. Det
er fra Bentham, at Mills ethiske og sociale Anskuelser have
deres Udspring; han gik i sin Ungdom i den Grad op i
Benthamismen, at han omfattede den med en Sekterers
Iver. Vel har han senere berigtiget flere Punkter i denne
Lære, men han er ikke naaet til nogen Klarhed over
Forskjellen mellem sin og Benthams Opfattelse.

Der ligger i Utilitarianismen Muligheden af et dobbelt
Udgangs- og Tyngdepunkt: man kan lægge Vægten paa
den egoistiske Interesse og vise, hvorledes denne for sin
egen Skyld maa træde i et harmonisk Forhold til Andres
Interesser; og man kan gaa ud fra Sympathien og den
Enkeltes Forhold til det større Hele og vise, hvorledes den
egoistiske Interesse kun er et underordnet Element indenfor
dette. Som allerede antydet gaaer Cumberland den
sidste af disse Veie, Locke og Paley den første. Forholdet
mellem Mill og Bentham beroer ligeledes paa, at de hver
gaa i sin Retning, opfatte det samme Princip fra to forskjellige
Sider. Ifølge Bentham følger ethvert Menneske
i enhver af sine Handlinger sin egen Lyst og Interesse.
% --- p. 53 ---
\setcounter{footnote}{0}%
\opage{53}Selvopofrelse bestaaer kun i at opgive en ringere Lyst for
en større. Men den Enes Lykke kan ikke sondres fra den
Andens; Enhvers velforstaaede Interesse maa derfor lede
ham til at arbeide for Andres Lykke. Herved fremkommer,
hvad Bentham kalder det Ethiske. »Drøm ikke om,» siger
han (Works I. p. 27), »at Folk ville bevæge deres lille
Finger for at tjene dig, naar deres Fordel derved ikke staaer
klart for dem. Men de ville ønske at tjene dig, naar de
indse, at de ved at gjøre det kunne tjene sig selv. I disse
gjensidige Tjenester bestaaer Dyden.»\footnote[1]{Man kommer her til at tænke paa den Opgave, Carlyle forelagde Utilitarianerne: „Given en Verden fuld af Slyngler, at vise, hvorledes Dyden resulterer af deres forenede Bestræbelser.“} Bentham antager
vel Sympathi for en oprindelig Faktor i den menneskelige
Natur; men han anser den ikke for stærk nok til at bære
det Ethiske. Stuart Mill adskiller sig netop fra Bentham
ved den Vægt, han lægger paa de sympathetiske Følelser
som Vehikel for det Ethiskes Udvikling. Hos Bentham, der
væsentlig holdt sig til den egoistiske Interesse, blev Dyden
Et med Klogskab og ydre Retskaffenhed. Han fremhæver
heller ikke ønsket om Selvopdragelse, om ethisk Uddannelse
af den egne Personlighed, men betragter Mennesket kun
som Forretningsvæsen, forstaaer kun »the business part of
human affairs.» Hos Mill er den egoistiske Interesse derimod
ikke blot underordnet Forholdet til Andre man ogsaa den
ideale Trang til Selvudvikling. Allerede i Slutningen af
Logiken udtaler han sig paa en Maade, der synes at vise
udover »Nyttetheorien»: »Jeg indrømmer ganske, at Uddannelsen
af en ideal Adel i Villie og Handling skal være et
Formaal, for hvilket den særlige Stræben efter egen og
Andres Lykke (for saa vidt den ikke er indbefattet i denne
Ide) i ethvert Kollisionstilfælde skal træde tilbage. Men
jeg troer, at netop det Spørgsmaal: hvori bestaaer denne
Ophøiethed i Karakteren? --- maa afgjøres ved at gaa tilbage
til Lykken som Maalestok. For den Enkelte skal
Karakteren selv være det høieste Maal, fordi en saadan ideal
% --- p. 54 ---
\setcounter{footnote}{0}%
\opage{54}Ophøiethed eller en Tilnærmelse til den mere end alt Andet
vil bidrage til at gjøre det menneskelige Liv lykkeligt, saa
vel i den forholdsvis beskedne Betydning af Fornøielse
og Frihed for Smerte, som ogsaa i den høiere Betydning,
at Livet ikke bliver, hvad det nu næsten almindelig er,
barnagtigt og betydningsløst, men bliver et saadant, som
menneskelige Væsener med høiere udviklede Evner kunne
ønske at føre.»

Det er er urigtig Forestilling, man gjør sig om Nyttemoralen, % PRINTED AS IS (er er: O-C07: doubled)
naar man mener, at den nødvendigvis maa være
en lav og endelig Betragtningsmaade. Bentham og Mill
ere Enthusiaster, lidenskabelig begeistrede for Menneskehedens
Fremskridt. De anse en Tilstand for mulig, hvor
almindelig Lykke hersker, og de betegne herved den fuldstændigste
Modsætning til Schopenhauers Pessimisme. Men
deres Optimisme er ikke, som Optimismen saa ofte er, fad
og overfladisk; den har et dybt Kjendskab og en inderlig
Tillid til den menneskelige Natur til sin Forudsætning; ---
den er heller ikke konservativ, mener ikke, som den vulgære
Optimisme, at Alt er godt, som det er, men bestaaer
væsentlig i Troen paa, at Alt kan blive godt, hvor langt
det end nu er fra at være det, eftersom de nitten Tyvendedele
af Menneskeslægten endog i de Dele af den nu existerende
Verden, der stikke mindst dybt i Barbari, leve i
Ulykke. Den nuværende elendige Opdragelse og elendige
Samfundsorden ere de eneste virkelige Hindringer for et
lykkeligt Liv for hele Menneskeheden, ɔ: »en Existents,
som bestaaer af faa og forbigaaende Lidelser, mange og
forskjelligartede Glæder med en bestemt Overvægt af
de aktive Glæder over de passive, og som til Grundlag
for det Hele har dette: ikke at vente Mere af Livet end
Livet kan byde.» De høiere intellektuelle og moralske Nydelser
ville mere og mere blive tilgængelige for Alle. Og
i ethvert Tilfælde ligger der i selve Kampen mod Ulykken
og i de Interesser og Følelser, denne Kamp fremkalder, en
ædel Tilfredsstillelse, som staaer høit over en hvilkensomhelst
egoistisk Nydelse.

% --- p. 55 ---
\setcounter{footnote}{0}%
\opage{55}Det Eiendommelige for Nyttetheorien, saaledes som
Mill opfatter den, er det, at den forudsætter det Ethiske i
Stedet for at begrunde og udvikle det. Ifølge Mill er
Stridspunktet ikke, om man skal handle ret, men hvad der
er det Rette; og det, mener han, kan kun Nyttetheorien
lære. Men det egentlig Ethiske er dog ikke et »Hvad,»
men netop et »At;» selve Forpligtelsesforholdet er det
Punkt, hvor alle enkelte Opgaver, alle Dyder og Pligter
maa have deres Tilknytningspunkt. Kun ved sin Forbindelse
med dette centrale Punkt bliver Noget godt i ethisk
Forstand. Nyttetheorien derimod holder sig til det Peripheriske,
til Handlingernes Indhold og Virkning i Stedet
for til deres indre, personlige Karakter. Mill siger endog,
at deri stemmer han ganske med Bentham, at Bevæggrunden
ikke har Noget at gjøre med en Handlings moralske
Værd, skjøndt den har saare Meget at sige med Hensyn til
den Handlendes Værd. Men kan man virkelig saaledes
betragte Bedømmelsen af Handlingerne og af de handlende
Personer som to forskjellige Ting? Handlingen kan være
nyttig, uden at Personen derfor er sædelig. En god og
ædel Karakter er ikke Noget, der kan betragtes som Tillæg
til det Ethiske i Handlingen; tvertimod, uden Forstaaelse
af Karakteren og de personlige Egenskaber kan der slet
ingen Dom fældes om Handlingens Moralitet; --- derfor er
moralsk Bedømmelse saa overordentlig vanskelig. Paa den
anden Side vise Nytten og Lykken, idealt opfattede, tilbage
til den personlige Karakterudvikling; ja, der gives,
som vi ogsaa have hørt Mill selv udtale, men som han i
sine egne Begrebsbestemmelser for det Meste glemmer, ingen
sand Nytte eller Lykke udenfor denne. Dersom Mill ikke
i Grunden opfattede Sagen paa denne Maade, --- og altsaa
her, ligesom i Logiken og Psykologien, gik ud over sine
egne Forudsætninger, --- vilde der reise sig en farlig Fjende
for ham i Schopenhauers og Hartmanns Pessimisme, som er
den sande Konsekvents af Sondringen mellem Lykke og
Opgave, mellem »Hvad» og »At.» Men baade Pessimismen
og Utilitarianismen adskille, hvad der ikke skal adskilles.
% --- p. 56 ---
\setcounter{footnote}{0}%
\opage{56}»Individets og Slægtens Berettigelse beroer kun paa deres
eiendommelige Formaal ag Opgave, gjennem hvis Løsning % PRINTED AS IS (ag: O-C08: for og)
deres eget Væsen udvikles, og kun i selve denne Udvikling
kan Lykken og Tilfredsstillelsen søges\dots{} Det Eudæmonologiske
og det Ethiske staa i en indre Sammenhæng, og
Pessimismens Betydning ligger i indirekte at godtgjøre Nødvendigheden
af denne Enhed.»\footnote[1]{„Philosophien i Tydskland efter Hegel.“ p. 285-286.}

De forskjellige Synspunkter, der saaledes gjøre sig
gjældende hos Mill, maa naturligvis medføre Uklarhed og
Vaklen i Begrebsbestemmelserne. Han udtaler saaledes, at
»Sjælen ikke befinder sig i en sund Tilstand, ikke i den
Tilstand, som stemmer med det Nyttige, naar den ikke
elsker Dyden for dens egen Skyld,» --- men tilføier i samme
Aandedræt, at Handlinger og Planer kun ere dydige, fordi
de fremme et andet Formaal end Dyd. Paa det Spørgsmaal:
er Dyden Middel eller Formaal? har Mill intet afgjørende
Svar. Snart hedder det, at den bør attraaes for
sin egen Skyld; snart, at Nyttemoralisterne »stille Dyden
i Spidsen for de Ting, der ere gode som Middel til det
sidste Formaal;» snart, at den er eller kan blive en af
»Lykkens Bestanddele,» --- »en Del af Maalet,» ja endogsaa »en vigtigere Del af det end nogen af de Gjenstande,
hvortil det er Middel»!\footnote[2]{\emph{Stuart Mill}: Moral grundet paa Nytte- eller Lykkeprincipet. Oversat og indledet af G. Brandes. p. 46-48.}

Af Autobiographien ser man, at disse dialektiske
Vanskeligheder i Mills eget Liv førte til en Krise, som
han tilskriver stor Betydning for sin aandelige Udvikling.
Som ivrig, efter sin egen Indrømmelse endog fanatisk
Benthamianer sværmede han i sin Ungdom for Udbredelsen
af Nytte og Lykke i Menneskeheden. Men --- kom han
til at spørge sig selv --- vilde du virkelig føle dig lykkelig,
dersom alle dine Idealer virkeliggjordes? --- og han maatte
besvare dette Spørgsmaal med Nei. Lang Tids Tungsind
% --- p. 57 ---
\setcounter{footnote}{0}%
\opage{57}og Fortvivlelse fulgte paa dette Nei. Omsider fik han nyt
Mod og ny Kraft ved den Overbevisning, at Lykken, skjøndt
den er Livets Maal, dog kun kan naaes ved ikke at gjøre
den til direkte Maal; kun de ere lykkelige, der gjøre noget
Andet end Lykke til deres ideale Formaal; Lykken finde
de da underveis. »Saasnart du spørger dig selv, om du
er lykkelig, vil du ophøre at være det.»\footnote[1]{Autobiography p. 142.}

Tydeligere kan det ikke godtgjøres, at Lykkeprincipet
viser ud over sig selv, kun er et relativt Princip i ethisk
Henseende, og at det altsaa ikke kan lægges til Grund
for Ethiken. Mill har forvexlet det psykologiske Udgangspunkt,
det opdragende Middel med det ethiske Formaal.
Men hans Styrke ligger nu paa den anden Side ogsaa netop
i at vise, hvorledes dette opdragende Middel virker.
Følelsen af Lyst og Ulyst ere de oprindelige Motiver for
al menneskelig Virken og Handlen, og Mill viser nu, ved
en træffende Anvendelse af Associationspsykologien (en Anvendelse,
hvori han har Hartley og James Mill, ja, egentlig
allerede Platon og Aristoteles, til Forgængere), hvorledes
Mennesket ved selve disse Motivers Virken Skridt for Skridt
føres til at stille sig høiere og ædlere Formaal, til at forlade
den egoistiske Begrændsning og optage de universelle,
humane Synspunkter i sine Motiver. Samfundets Interesse
fortrænger lidt efter lidt de individuelle Interesser, og der
vil efterhaanden endog kunne danne sig en saa stærk Følelse
for det Almene, at man vil kunne nære Frygt for, at den
skal hæmme den individuelle Udvikling. Ideassociationen
mellem Pligtfølelsen og Nytten, støttet af Sympathiens
naturlige Kraft, vil kunne frembringe en psykologisk
Magt og en social Virkekraft, som overgaaer endog den
største Indflydelse, der er bleven udøvet af en positiv
Religion.

Mill begynder i sin Ethik som Epikuræer og ender
som Stoiker. En Mulighed til at forene de stridende Syns%
% --- p. 58 ---
\setcounter{footnote}{0}%
\opage{58}punkter ligger i den Idé, der  paa saa mange Steder dukker
frem hos ham, men som han ikke lægger til Grund, Udviklingens
Idé. Kun ved den kan hans Philosophi befries
fra sine indre Modsigelser. I Grunden er det jo saaledes
ikke Lykken, der er Mills ethiske Princip, men Lykken
udviklet til sit Ideal, hvor den er Et med den ethiske Karakter.
Det er selve denne Udvikling, som er den ethiske
Opgave, og hvis Nødvendighed begrunder Pligtforholdet.
Det ethiske Motiv er Stræben efter stedse høiere Udvikling,
en Udvikling, hvorunder de individuelle og de kulturhistoriske
Opgaver staa i inderlig Forbindelse og Vexelvirkning
med hverandre.\footnote[1]{En Antydning i denne Retning, i Retning af at lade Nyttens Ide gaa op i Udviklingens Ide, findes hos Mill i hans Skrift „On liberty“ p. 24: „Jeg betragter Nytte som sidste Maalestok for alle ethiske Spørgsmaal; men det maa være Nytte i \emph{mest omfattende Betydning, grundet paa Menneskets blivende Interesser som fremadskridende Væsen}.“}

Utilitarianismen er Empirismen paa det ethiske
Omraade. I Erkjendelseslæren fandt Empirismen sit Udtryk
i Opfattelsen af Menneskeaanden som oprindelig passiv;
med Anerkjendelsen af den oprindelige Aktivitet i Individet
opløstes derfor Empirismen og Associationspsykologien. I
Ethiken finder den et tilsvarende Udtryk, idet Attraaen,
»den passive Modtagelighedstilstand,» som Mill bestemmer
den, tages til Udgangspunkt, medens dog Villien, »det aktive
Phænomen,» betragtes som »et Skud fra Attraaens
Stamme.» Men hvorledes kan der udvikle sig en Aktivitet
af en passiv Tilstand, og maa hin Anerkjendelse af den
oprindelige Aktivitet ikke ogsaa faa Konsekventser med
Hensyn til Grundlaget for det Ethiske, saa at Stoikerne,
ogsaa hvad det oprindelige Udgangspunkt angaaer, faa Ret
overfor Epikuræerne? ---

Det er tidligere blevet antydet, at Mills hele theoretiske
Standpunkt har sin Grund i hans praktiske Begeistring,
hans Iver for Fremskridt og Udvikling. Denne Iver
% --- p. 59 ---
\setcounter{footnote}{0}%
\opage{59}viser sig ogsaa tydelig i hans Ethik; det er maaske hans
Utaalmodighed efter af komme til de praktiske Opgaver, % PRINTED AS IS (af komme: O-C11: for at komme)
der gjør, at vi i hans ethiske Begrebsbestemmelser ofte
have ondt ved at gjenkjende den store Logiker. Men
navnlig træder naturligvis den praktiske Begeistring frem,
naar vi komme til selve Samfundsvidenskaben. Hos Faa
eller Ingen i vort Aarhundrede har der rørt sig en stærkere
Interesse for Humanitet og Fremskridt, forbunden med en
saa overlegen Intelligents og saa forskjelligartede Kundskaber.
Mill er ikke blot Philosoph, men har ogsaa et
Navn som Nationaløkonom. Hans ethiske Idealer svæve
ikke i Luften, men begrundes ved ædruelig Undersøgelse
af de virkelige Forhold. Den skarpeste, prosaiske Forstandighed
har hos ham ført til Konsekventser, der synes den
tilvante Tænkemaade overspændte og phantastiske. Uden
her at kunne gaa nærmere ind paa de praktiske Problemer,
med hvilke Mill har beskjæftiget sig, maa vi dog i Korthed
vise Sammenhængen mellem hans Behandling af dem og
hans almindelige philosophiske Standpunkt.

Den sociale Videnskab er ifølge Mills Logik de samme
Methoder underlagt som Naturvidenskaben: Induktionens
Resultater danne Grundlaget for en Deduktion, hvis Resultater
igjen bekræftes ved direkte Erfaring. Mill gaaer derfor
ikke ud fra dogmatiske Forudsætninger, men søger at
finde bestemte Udgangspunkter i Erfaringen. Han udtaler
sig her endog ofte aldeles skeptisk. Er der f. Ex. Tale om
Kvindens Stilling i Samfundet, saa gaaer han ikke ud fra
visse Forestillinger om Kvindens Natur, men erklærer
endog, at vi aldeles ikke kjende denne. Kvinderne have
hidtil levet under visse Indskrænkninger, ja endog under
ligefrem Tvang; hvad man kalder deres Natur, er i høieste
Grad et Kunstprodukt, Resultatet af et forceret Tryk i en
Retning og en unaturlig Ophidselse i en anden Retning.
Visse af deres Evner ere stedse blevne udviklede ved Forholdenes
Magt, medens andre ere henvisnede og tilsyneladende
uddøde af Mangel paa Næring og Pleie. Saalænge
Ethologien endnu er i sin Barndom, saalænge man endnu
% --- p. 60 ---
\setcounter{footnote}{0}%
\opage{60}har saa lidet Kjendskab til Lovene for Karakterens Dannelse
og for de Maader, hvorpaa de ydre Betingelser her gribe
ind, saa længe i det Mindste kan man endnu ikke føre
noget positivt Bevis for, at de intellektuelle og moralske
Forskjelligheder mellem Mænd og Kvinder ere grundede i
Naturen. Kun det, der ikke kan forklares ved Opdragelse
og ydre Omstændigheder, har man Ret til at kalde naturligt;
hvor indgribende disse Omstændigheder kunne virke,
véd man endnu slet ikke\footnote[1]{Kvindernes Underkuelse. Oversat af G. Brandes. p. 37-40.} Paa den anden Side bestaaer
hele den historiske Udvikling i en fremskridende Ophævelse
af den Stærkeres Ret, som i saa lange Tider var eneste
og ubetinget Regel i alle Livsforhold. Ved denne Frigjørelse
for ydre Tvang, som er den moderne Verdens største
Eiendommelighed, bliver det muligt for de Enkelte at udvikle
deres Individualitet uden at standses af kunstige
Skranker. Kvindens Udelukkelse fra de Rettigheder og de
Veie, der staa aabne for Mændene, er kun en Levning
af hin brutale Magt, som nu bekæmpes paa de fleste Omraader.


Med Ophævelsen af den Stærkeres Ret vil efterhaanden
enhver ydre Autoritet falde, enhver Magt, der vil stille sig
som Formynder overfor Individerne i Stedet for at lade
disse kæmpe Livets Kamp med egne Hænder og paa eget
Ansvar. Ifølge den gamle Anskuelse troede man, at Individet
vilde fare vild, naar det blev overladt til sig selv.
Men den moderne Anskuelse lærer, at de Anliggender, i
hvilke Individet alene er interesseret, ikke ville gaa godt,
uden naar man overlader dem til hans egen Ledelse, og
at Autoritetens Mellemkomst kun gjør Skade, naar det ikke
netop gjælder om at beskytte Andres Rettigheder. I Afhandlingen
»om Friheden» (et af Mills klareste og smukkeste
Skrifter, om hvilket han ogsaa selv mente, at det vilde
holde sig længst af alle hans Værker) udvikles denne
Tanke videre. En sund intellektuel og moralsk Udvikling
er kun mulig under Forudsætning af Frihed; naar Sam%
% --- p. 61 ---
\setcounter{footnote}{0}%
\opage{61}fundet, det være nu demokratisk eller aristokratisk i sin
Forfatning, vil udstrække sin Indflydelse baade over det
aandelige og det legemlige Liv, naar det søger at lede
Individerne ind paa visse Livsretninger, indpode dem visse
bestemte Meninger, saa hæmmes Karakterudviklingen, og
høiere Selvstændighed og Eiendommelighed bliver umulig.
Derfor fordrer Mill absolut Frihed for Tanke og Ord og for
Anlægget af Livsplanen i det Hele. Kun med Hensyn til
Handlinger, der kunne indvirke paa andre Menneskers
Velfærd, kunne Skranker for den individuelle Frihed retfærdiggjøres.
Følger man ikke denne Regel, vil Resultatet
blive en gjennemgaaende Middelmaadighed; Mændene ville
blive Dverge, og man vil snart erfare, at med smaa Mænd
kan intet Stort udføres. Faren er navnlig stor i demokratiske
Stater, efterdi den offentlige Mening, der i dem
bliver den herskende Magt, griber langt dybere og umærkeligere
ind i det aandelige Liv end ydre Lovbud og
Skranker, der sættes ved Magt. Mills Bestræbelser gaa
derfor ud paa at finde politiske Midler til at hævde
Minoriteternes Ret. Men det kraftigste Middel søger han
dog i en Reform af den offentlige Mening, en absolut,
og ikke blot ydre og formel Anerkjendelse af Personlighedens
Ret.

Man maa ikke heraf slutte, at Mill betragter Statens
Rolle som udspillet, saa at Alt nu skulde kunne overlades
til den individuelle Stræben og Konkurrence. Han bestrider
saadanne Anskuelser som Buckles, der udlede alt Ondt og
Forkert i Historien af Regeringernes Indflydelse. Han
hævder i sin »politiske Økonomi,» at der er en stor Mængde
Opgaver, som angaa Menneskehedens og de fremtidige Generationers
Interesser eller ogsaa de Medlemmer af Samfundet,
som behøve ydre Hjælp, --- Opgaver, som ikke kunne
løses uden af Statsmagten. Staten skal ikke blot sørge
for Retssikkerhed og øve Kontrol og Voldgift i de Enkeltes
indbyrdes Anliggender; den skal ogsaa gribe ind, hvor der
enten kan finde Forsømmelse Sted fra Individernes Side,
saaledes i Opdragelsen og Undervisningen (kun at den ikke
% --- p. 62 ---
\setcounter{footnote}{0}%
\opage{62}gjør denne til et Monopol), eller hvor de Enkeltes Kræfter
ikke slaa til (Fattigvæsen, Kolonisation). Der kan efter
hans Opfattelse her ikke fastsættes bestemte Grændser;
man maa holde sig til den almindelige Regel, at Statsmagtens
Indgriben kun er tilstedelig, hvor Nytten deraf
er aabenbar. Og tillige bør Statsmagten altid virke paa
en saadan Maade, at de individuelle Kræfter ikke trænges
tilbage, men tvertimod udvikles og opdrages gjennem dens
Forholdsregler, saa at Regeringsunderstøttelsen bliver »et
Lærekursus for Folket i den Kunst at løse store Opgaver
ved individuel Virksomhed og frivillig Samvirken.»

Opdragen og uddannet i Benthams og James Mills
Skole var Stuart Mill fra sin første Fremtræden en begeistret
Demokrat. Hans politiske Tro undergik senere
en dobbelt Modifikation. Dels kom han snart til at tænke
paa, at den demokratiske Udvikling vilde kunne medføre
Fare for den individuelle Udvikling, at Minoriteterne let
vilde blive undertrykkede; denne Frygt udtaler han allerede
i sin Artikel i Edingburgh Review om Tocquevilles Bog % PRINTED AS IS (Edingburgh: O-C15)
»La démocratie en Amérique,» og han kom senere, som
allerede vist, tilbage til den i Afhandlingen om Friheden.
Men dels saa han ogsaa, at det sociale Spørgsmaal
var vigtigere end det politiske. Tidligere havde han
med den ældre politiske Skole betragtet Eiendomsretten
som noget Fast og Absolut; men efterhaanden førtes han
mere og mere til at se, at den, som enhver Ret, er betinget
af Samfundsforholdene. Eiendomsretten kunde heller
ei ifølge hans hele philosophiske Anskuelse være et Dogme,
hvorfra man nødvendigvis skulde gaa ud. Den er en Institution,
som er opstaaet i Medfør af Tingenes naturlige
Gang, og som har sine store Fortrin, men som dog i og
for sig ikke kan gjøre Fordring paa absolut Gyldighed.
Derfor dømmer Mill gunstigere om de kommunistiske og
socialistiske Systemer (navnlig Fouriers), end Nationaløkonomerne
pleie at gjøre, ligesom han finder det naturligt,
at kommunistiske Ideer have udbredt sig en Tid som
vor, »hvor en almindelig Prøvelse af alle første Principer
% --- p. 63 ---
\setcounter{footnote}{0}%
\opage{63}erkjendes for uundgaaelig, og hvor for første Gang i Historien
de mest lidende Klasser i Samfundet have en
Stemme med ved Undersøgelsen.» De fleste Vanskeligheder,
Kritiken kan paavise ved de socialistiske Systemer ramme
efter hans Mening maaske i endnu høiere Grad de nuværende
Forhold. Naar man f. Ex. siger, at under en socialistisk
Forfatning vil Enhver skyde saa meget Arbeide som
muligt over paa Andre, fordi han ei er særlig interesseret
deri, saa svarer Mill, at det Samme maa gjælde om ni
Tiendedele af det nuværende Arbeide: hvor ringe en Del
af, hvad der produceres i England, produceres af Mænd,
som selv høste Udbyttet deraf? En Fabrikarbeider har
mindre personlig Interesse ved sit Arbeide end et Medlem
af en socialistisk Association. Desuden hævder Mill, at
Menneskeheden er istand til en langt høiere Grad af Almenaand
end vor Tidsalder har vænnet sig til at anse for
mulig. Og i ethvert Tilfælde, mener han, kunde man
vanskelig faa en Fordeling, som var uligere og uretfærdigere
end den nuværende. --- Afgjørelsen af, hvilket af de to
Systemer, Kommunismens eller Privateiendommens, Fortrinet
maa gives, maa bero paa, hvilket af dem der bedst kan
forenes med menneskelig Frihed og Udvikling. Mill finder
ingen Nødvendighed af at forlade Privateiendommens System,
naar det kun befries fra alle de Misligheder, som
ikke uundgaaelig ere forbundne med det. Naar Lovene kun
gjorde lige saa Meget for at mildne de sociale Uligheder
og begunstige en større Fordeling af Formuen, som de nu
gjøre for at forhøie hine Uligheder og koncentrere Besiddelsen,
saa vilde det vise sig, at Privateiendommens Princip ikke
staaer i nogen nødvendig Sammenhæng med de legemlige
og aandelige Lidelser, som Socialisterne anse for uadskillelige
derfra. Ligeledes er han uenig med Socialisterne i deres
Polemik mod den frie Konkurrence, idet han gjør opmærksom
paa, at man kun har Valget mellem Konkurrence
og Monopol, og at »Beskyttelse mod Konkurrence vil sige
det Samme som Beskyttelse af Ladhed og Dovenskab, som
% --- p. 64 ---
\setcounter{footnote}{0}%
\opage{64}en Fritagelse for Nødvendigheden af at være lige saa intelligent
og virksom som andre Mennesker.»

Mill regner sig i Autobiographien til Socialismens
Tilhængere. Dog er han, som man ser, uenig med Socialisterne
i væsentlige Punkter. Det Socialistiske i hans Anskuelse
beroer paa, at han forlanger Statsmagtens Indgriben
for at løse det sociale Spørgsmaal. Arbeidernes Vilkaar
kunne ikke forandres uden en gjennemgribende Ændring i
deres hele Tænkemaade, og denne Ændring forudsætter igjen
en Forbedring af deres ydre Vilkaar, som kan give dem
Lyst til og Smag for ædlere Nydelser og høiere Kultur. Kun
derved vil den Sløvhed og Letsindighed kunne udryddes,
som nu gjør det umuligt for Proletariatet at hæve sig, idet
det ved sin stærke Formering bestandig holder sig selv
nede. De Forholdsregler, der skulle tages, maa altsaa virke
dels til Udvikling af Intelligents, dels til Ophævelse af
Fattigdommen. Det Første opnaaes ved en national Børneopdragelse
og ved Folkeundervisning, og som Midler til det
Andet foreslaaer Mill store Kolonianlæg og Udstykning af
Jordegodset. Naar en Generation af Arbeidere herved blev
stillet under nye og bedre Vilkaar, vilde der være
tilveiebragt gunstige Betingelser for Udviklingen af en ny
Generation med høiere Formaal. og større Energi. --- Vi
kunne ikke her indlade os paa nogen Bedømmelse af disse
Forslag, men indskrænke os til at fremhæve det Mærkelige
i, at den samme Mand, der stærkere end nogen Anden har
talt den frie Individualitets Sag, saa indtrængende fordrer
Statsmagtens Indgriben paa det sociale Omraade. Der
kunde synes at være en Modsigelse i at appellere til disse
modsatte Principer. Men man ser let, at Mill statuerer
en indre Forbindelse mellem dem. Samfundsmagtens høieste
Formaal er at udvikle den individuelle Selvstændighed og
Personlighed, at sikre Individets Virken paa de Punkter,
hvor Alt kan overlades til dets frie Spil, og at nære og
pleie den individuelle Frihed der, hvor den endnu ikke har
kunnet udfolde sig. Det er naturligvis langt fra, at Mill
har overvundet alle Vanskeligheder. Men han har den
% --- p. 65 ---
\setcounter{footnote}{0}%
\opage{65}store Fortjeneste at have vist, hvor urigtig man opfatter
Frihedens Væsen, naar man mener, at al Indgriben af
Statsmagten paa de sociale Opgavers Omraade skulde være
forbi, og at Alt nu kunde overlades til den »naturlige»
Gang. Der er endnu en lang Vei at gaa, førend den i
Sandhed naturlige Udvikling kan indtræde. Staten, ved
hvis Hjælp Mennesket er fort ud af den dyriske Naturtilstand,
kan ikke slippe ham, førend den har ført ham til
den humane Naturtilstand. Mill har endelig ogsaa den
Fortjeneste at have paavist de ideelle og ethiske Elementers
Betydning for de sociale Spørgsmaal og derved ført ud over
den snevre Synskreds, som i Almindelighed gjør sig gjældende
ved deres Behandling.

Som Resultat af den sociale Udvikling tænker Mill
sig en Tilstand, hvor saavel Kapitalen som Befolkningen
er omtrent stationær. Det at erhverve kan ikke være den
sidste Opgave; Menneskeheden maa dog engang komme saa
vidt, at den sætter Formaalet over Midlet. Man arbeider
da ikke længer paa at forøge Kapitalen; alle Opfindelser
gaa derimod ud paa at forkorte Arbeidstiden for at vinde
Tid til friere og ædlere Sysler; det intellektuelle, æsthetiske
og moralske Liv spiller Hovedrollen og lægger Beslag paa
den største Del af Energien. Befolkningen er stationær,
fordi man indser, at en for stærk Forøgelse af den vilde
hindre den høiere og ædlere Udvikling. Selv om der var
Næring og Klæder nok til Flere endnu, vilde der dog være
Overbefolkning, naar Menneskene boede saa tæt, at det blev
umuligt at nyde Ensomheden, og naar den frie Natur ganske
var tagen i den menneskelige Nyttes Tjeneste. »En Verden,
fra hvilken Ensomheden var banlyst, vilde være et fattigt
Ideal. Ensomheden er en væsentlig Betingelse for Dybde i
Tænkning og i Karakter; og Ensomhed i Forbindelse med
Naturskjønhed er Vuggen for Tanker og Indskydelser, der
ikke blot komme det enkelte Menneske til Gode, men som
heller ikke Samfundet godt kan undvære.» --- Selv om
det Billede, Mill her udkaster, skulde være en Utopi, har
det sin Betydning og Interesse, ikke blot som anskueligt
% --- p. 66 ---
\setcounter{footnote}{0}%
\opage{66}Udtryk for en human Bestræbelses Maal, men ogsaa som
Vidnesbyrd om det ædle, friske og menneskekjærlige Sind
hos den store Tænker, vi --- trods al den Kritik, vi have
forsøgt at anlægge paa hans Lære --- ikke kunne forlade
uden Beundring.

II. \emph{Alexander Bain} (f. 1818) er født i Skotland
og kan siges at have fortsat den skotske Skoles Bestræbelser
paa den empiriske Psykologis Omraade. Hans Lære kan
kaldes »den skotske Psykologi, saaledes som den vilde se
ud i det 19de Aarhundrede. »\footnote[1]{\emph{Ribot}: La psychologie Anglaise contemporaine. Paris 1870. p. 222.} Men i Stedet for de ældre
skotske Philosophers Ængstelighed, deres Appel til »den
sunde Menneskeforstands» Udsagn og deres Tilbøielighed til
at standse ved visse saakaldte Bevidsthedskjendsgjerninger,
som de ansaa for fundamentale, træder hos deres moderne
Efterfølger en sikker og omfattende Anvendelse af den induktive
Methode, et omhyggeligt Studium af Physiologi og
Anatomi som Grundlag for Psykologien, og en skarpere
Kritik og Analyse. Dog sporer man endnu den ældre
Skoles Eftervirkning i den empirisk-deskriptive Methode, i
en Mangel paa principiel Gjennemførelse af de ved Induktion
og Analyse vundne Begreber og paa synthetisk Udvikling
og Fremstilling, hvorved der kunde vindes en Totalanskuelse
af Aandslivet. Det er dette, som og at han i sit hele
Standpunkt staaer Mill saa nær, der berettiger til at henføre
ham blandt den rene Empirismes Repræsentanter. I
England selv regnes han til Mills Skole. Hans Hovedværker
ere »The senses and the intellect» og »The emotions and
the will,» i hvilke han, følgende Botanikens og Zoologiens
beskrivende Methode, giver en Aandens Naturhistorie med
en Mangfoldighed af Enkeltundersøgelser og skarpsindige
Bemærkninger, som gjør disse Bøger høist lærerige og interessante,
skjøndt Oversigten ofte hindres netop af Stoffets
% --- p. 67 ---
\setcounter{footnote}{0}%
\opage{67}Mangfoldighed. I et mindre Skrift »On the study of
character» har han givet et Udkast til en Ethologi efter
Mills Principer.

Vi have allerede havt Leilighed til at fremhæve, at
Bain har korrigeret den sensualistiske Opfattelse af Aanden
som fuldstændig passiv. Det er et afgjørende Bevis for,
at den engelske Philosophi i vor Tid ikke er en blot Fornyelse
af det forrige Aarhundredes Sensualisme (maaske
hellere Sensationisme). Det Eiendommelige for denne var,
at den reducerede Tænkning til Sandsning og opfattede
Sandsningen som en rent passiv Proces, en blot Liden og
Modtagen. I Modsætning hertil opfattede Sensualismens
Modstandere gjerne Tænkningen som noget rent Aktivt og
sondrede skarpt mellem den og den passive Sandsning.
Denne Sondring hviler egentlig paa samme Forudsætning
som Sensualismen. Der maa være et indre Forhold mellem
Sandsning og Tænkning. Naar man erkjender, at vi vel
ikke have medfødte Begreber, men at Grundformerne i vor
aandelige Organisation maa være oprindelige, at vi, som
Leibnitz udtrykte det, ere os selv medfødte, saa maa man
ogsaa erkjende, at den aandelige Energi, der ytrer sig i
Erkjendelsens høiere Former, maa være tilstede i de første
aandelige Livsytringer, i de elementære Sandsefornemmelser,
og at den under Udviklingens Gang fremgaaer af disse og
bliver sig selv bevidst. Dette lærer nu netop Bain: »Sandsning
er allerede en Funktion af Forstanden; den er den
rent ydre Aabenbarelse af de ideale Processer» (Senses and
Intellect p. 339). Og disse ideale Processer have deres
Spire i selve Sjælelivets Begyndelsespunkt.

Selv i de laveste sjælelige Livsytringer findes der,
siger Bain, Elementer af Aand eller Tanke (intelligence).
Den primitive Kjendsgjerning, som udtrykker Aandens Væsen,
er Evnen til at skjelne, til at opfatte Forskjelle (discrimination).
Dette er en oprindelig aandelig Virksomhed, som
ikke nærmere kan forklares, men hvorpaa al Bevidsthed hviler.
Den har til sin Forudsætning Evnen til at fastholde et
Indtryk, selv naar dets Kilde ikke længer er tilstede,  saa
% --- p. 68 ---
\setcounter{footnote}{0}%
\opage{68}at det kan blive sammenlignet med et andet Indtryk. Ved
den maa derfor den psykologiske Analyse standse. Allerede Thomas
Hobbes har indskærpet dette Grundfaktum, idet
han lærte, at uden Forandring og Bevægelse var hverken
Fornemmelse eller Bevidsthed mulig.\footnote[1]{Per sensionem vulgo intelligimus aliqvam de rebus objectis per phantasmata judicationem; phantasmata scilicet comparando et distingvendo. . . . Adeo sentire semper idem et non sentire, ad idem recidunt. \emph{Hobbes}: De corpore IV, 25, 5.} En Virken uden
Forandring frembringer intet Indtryk paa os. Vi mærke
ikke Jordens Bevægelse i Verdensrummet, men derimod
Skibets Bevægelse, hvorpaa vi seile, og endnu mere Vognens
Bevægelse, hvorpaa vi kjøre. Sandsen for Forskjel er saaledes
det oprindelige Aandselement. Al Bevidsthed beroer
egentlig paa Modsætning og Overraskelse og er des mere
intensiv, jo større Modsætningen og som Følge deraf Overraskelsen
er. Efter et dybt Mørke fornemmes Lyset stærkere
end ellers.\footnote[2]{I denne Kjendsgjerning tage Webers og Fechners interessante Undersøgelser om Fornemmelsens (Bevidsthedens) Grader deres Udgangspunkt. Smlgn. „Philosophien i Tydskland“ p. 253---255.} Og det gaaer med Erkjendelsen som med
Bevidstheden i det Hele; jo flere Forskjelle, der opfattes,
des klarere og bestemtere er Erkjendelsen. Ingen Karakteristik
af en Gjenstand er mulig uden gjennem Bestemmelse
af dens Forskjelligheder fra andre Gjenstande. Selv i den
mest elementære Sandsefornemmelse er der en saadan Opfattelse
af det Forskjellige, og for saa vidt en Dom, skjøndt
man ved Dømmen i Almindelighed tænker paa de høiere
udviklede Former af Skjelneevnen. Bain stemmer altsaa
med William Hamilton og Skotterne\footnote[3]{Ogsaa Lotze lærer, at Sandseopfattelsen er en Dom. Se anf. Skr. p. 295.} i den Lære, at
der i enhver Bevidsthedsakt indeholdes en Dom. Sandsning
forudsætter altsaa Forstand (intellect). Den gamle
Sætning: Nihil est in intellectu, qvod non erat in sensu
--- er aldeles urigtig. Der er, som snart skal vises, en
% --- p. 69 ---
\setcounter{footnote}{0}%
\opage{69}Mængde Begreber. som ikke have deres Oprindelse fra
Sandsningen.

Den første Fornemmelse er ikke en Sandsning af
noget Ydre, saaledes som den tidligere sensualistiske Skole
antog. »Bevægelse gaaer forud for Sandsning og er fra
Begyndelsen af uafhængig af enhver Impuls udenfra. Virksomhed
hører inderligere og mere uadskillelig til vor Natur,
end nogen af vore Sandsefornemmelser gjør, og er faktisk
medvirkende i enhver af vore Sandsninger, idet den giver
dem Karakteren af sammensatte Phænomener, medens den
selv er en simpel og elementær Egenskab.» Men hvorfra
kommer da den oprindelige Bevægelse? En Række af Fakta
viser, at Nervecentrene kunne udsende Impulser til Musklerne,
blot ifølge deres indre Spænding, som gjør en Udladning
nødvendig. Naar Organismen har faaet tilbørlig Næring,
maa Nervesystemet have en Afløber for sin forøgede Energi.
Vi bringes herved til at se den aandelige Kraft i dens Forbindelse
med de andre kosmiske Kræfter. Gjennem Nervekraften
staaer den i Forbindelse med de chemiske Kræfter,
med Magnetisme, Elektricitet, Varme og mekanisk Kraft,
som alle vise tilbage til samme Oprindelse og alle kunne
gaa over i hverandre.

Det første sjælelige Phænomen er den spontane Bevægelse,
som vi allerede træffe hos Fosteret i Moders Liv,
men som ogsaa ytrer sig hos de Nyfødte og hos ældre
Individer i den Livlighed, som fremkaldes ved god og
rigelig Næring, og i den spontane Opvaagnen af Søvnen.\footnote[1]{Bain støtter sig her til den berømte Physiolog Joh. Müller.}
Denne spontane Bevægelse blive vi os bevidst ved Muskelfornemmelsen,
som altsaa er den første Fornemmelse, en
Fornemmelse, der adskiller sig fra den egentlige Sandsning
ved, at den ikke har sin Kilde i en ydre, men i en indre
Impuls. Vi kunne i Muskelfornemmelsen, som i enhver
Fornemmelse, adskille to Elementer, et Følelses- og et
Erkjendelseselement. Disse to Elementer staa i et omvendt
% --- p. 70 ---
\setcounter{footnote}{0}%
\opage{70}Forhold til hinanden. Jo mere jeg føler den Lyst eller
Smerte, som Muskelbevægelsen fremkalder, des mindre er
jeg i Stand til den Skjelnen og Sondren, hvorpaa Erkjendelsen
beroer.

Den oprindelige Bevægelse er uden Maal og Med, er
en uvilkaarlig Udledning af Nervekraften. Men Bevægelsens
Retning bestemmes efterhaanden ved de Indtryk den fremkalder;
den ledes ved dem i bestemte Baner. Man kan
iagttage dette paa nyfødte Lam. De prøve uvilkaarlig paa
at bevæge Lemmerne, men denne Bevægelse sker snart til
Siden, snart fremad; kun efterhaanden lære de, at Bevægelsen
fremad er den naturligste. Naar Hyrden fører dem
hen til Moderfaaret og lader dem berøre det, gjør Varmeindtrykket,
de derved modtage, at de søge at bevare eller
igjen at indtage denne Stilling. Naar deres Mund bringes
hen til Moderens Yver, og de uvilkaarlige Sugebevægelser
ere begyndte, vil den Nydelse, de føle herved, anspore til
Fortsættelse eller Gjentagelse. Snart kjende de Moderen i
Frastand og kunne foretage de Bevægelser, som udfordres
for at komme til hende. Den oprindelige, spontane Bevægelse
er altsaa Forudsætningen for, at de egentlige Sandsefornemmelser
kunne lære Dyret Noget.

Fra Muskelfornemmelserne stamme Begreber som Modstand,
Inerti, Vægt og Naturens mekaniske Egenskaber.
De maales efter den Grad, i hvilken vi maa anstrænge os.
Forskjellem mellem 4 og 6 Punds Vægt opfatte vi ved % PRINTED AS IS (Forskjellem: O-D05)
Følelsen af forøget Anstrængelse, som vi maa anvende, naar
vi i Stedet for 4 faa 6 Pund at bære paa. Begreberne om
Varighed og Udstrækning have deres Oprindelse fra Muskelvirksomheden
i Forbindelse med Følesandsen.  Kun ved
Hjælp af denne Sands vinde vi bestemte Punkter til at
maale Bevægelsen med, baade hvad Rummet og Tiden angaaer.
Kombinationen af Muskelfornemmelserne med Syn,
Hørelse og de andre Sandser giver os de mere sammensatte
Sandsninger. Ved Hjælp af disse udvikler Erkjendelsen sig,
men dens Spire er fra først af tilstede.

Allerede ved det første Skridt ud over, den første
% --- p. 71 ---
\setcounter{footnote}{0}%
\opage{71}nærmere Bestemmelse af den oprindelige spontane Bevægelse,
aabenbarer det aandelige Livs Grundlove sig, Ideassociationens
Love, i Kraft af hvilke sjælelige Phænomener indgaa en
saa nøie Forbindelse, at det ene fremkalder det andet, ---
Love, som gjælde for hele Aandslivets Omraade fra de laveste
Trin til de høieste, og som, efter Stuart Mills Udtryk, ere
for den aandelige Verden, hvad Tyngdeloven er for den
ydre Verden. Disse Love kunne føres tilbage til to: Kontiguitetens
(Berøringens) Lov og Similaritetens (Lighedens)
Lov. Sjælelige Phænomener, som indtræffe samtidig med
eller umiddelbart efter hinanden, stræbe at fremkalde hinanden,
naar blot et af dem senere indtræffer. Og ligeledes
stræbe sjælelige Phænomener, der ligne hinanden, indbyrdes
at fremkalde hinanden. Ældre Psykologer, navnlig James
Mill, forsøgte at forklare alle sjælelige Phænomener ved
Kontiguitetens Lov. Men denne Lov forudsætter Lighedens
Lov, saaledes som William Hamilton har paavist. Thi for
at et Phænomen skal kunne fremkalde Forestillingen om
et andet Phænomen, med hvilket det før er indtruffet
sammen, maa hint første selv gjenkjendes, erkjendes som
det samme, der før indtraf sammen med det andet. Lighedens
Lov er altsaa en nødvendig Betingelse for, at Berøringens
Lov skal kunne virke. Men hvor Ligheden er
fuldstændig, hvor det kan siges, at det er det samme Phænomen,
der kommer igjen, der behøver man ikke særlig at
fremhæve den. Omvendt danner Berøringen den stiltiende
Forudsætning for, at Lighedens Lov kan virke; thi for at
Ligheden mellem to Phænomener skal kunne opfattes, maa
de først bringes sammen for Bevidstheden. At der virkelig
er to forskjellige Love, skjøndt de saaledes kunne synes at
gaa over i hinanden, ses bl. A. af, at afgjørende aandelige
Forskjelligheder bero paa, hvilken af de to Love der er
den fremherskende. Association ved Berøring i Rum eller
Tid (Kontiguitet) spiller især en Rolle ved Opsamlen af
Kundskaber og ved Udførelse af praktiske Forretninger.
Vi erindre da Tingene i den samme Orden, i hvilken de
pleie at fremstille sig; det er Vane og Routine, der herske.
% --- p. 72 ---
\setcounter{footnote}{0}%
\opage{72}Association ved Lighed er derimod fremherskende hos Opdageren,
Tænkeren og Digteren, og betinger i det Hele
Produktiviteten og Originaliteten, medens Berøringens Lov
især betinger Receptiviteten.

Vi betragte først Berøringens eller Kontiguitetens
Lov. At Bevidsthedstilstande, som engang ere forekomne
sammen, senere fremkalde hinanden, kan ikke videre forklares.
Det er et sjæleligt Grundfaktum, som igjen viser
tilbage til en Grundeiendommelighed ved Nervevirksomheden.
Ligesom Forestillingerne saaledes maa ogsaa de tilsvarende
Nervetilstande forbindes. Forestillingen om en Sandsning
er kun en svagere Sandsning, der beroer paa samme Nervestrømning
som den oprindelige Sandsning, kun at ogsaa
denne Strømning er svagere end ved den oprindelige
Sandsning. Gjentagelsen danner ligesom Kanaler for disse
Strømninger, saaledes at de lettere kunne bevæge sig i
visse Retninger end i andre. En levende Forestilling om
en vis Bevægelse kan uvilkaarlig fremkalde denne Bevægelse;
Forestillingen om en Afgrund kan vække Svimmelhed.
Ja, Tanken selv er egentlig kun en tilbagekaldt, en standset
Bevægelse. Hjernen er derfor ikke et Magasin, hvor Tanker
og Forestillinger skulde ligge ophobede. Der kan heller
ikke gives medfødte Ideer; alle Ideer bero paa Funktion
og Bevægelse og ere ikke til uden dem, altsaa heller ikke
før dem (smlgn. Schelling; se ovenfor p. 12 f.).

Association ved Berøring kan enten finde Sted mellem
Bevidsthedstilstande af samme Art (f. Ex. en Lyd og en
anden Lyd, en Bevægelse og en anden Bevægelse) eller af
forskjellig Art (f. Ex. Farve med Modstand, Bevægelse med
Afstand o. s. v.). Af første Art er navnlig Associationen
mellem forskjellige enkelte Bevægelser af Vigtighed. Næsten
enhver betydningsfuld Livsytring beroer hos Dyr og Mennesker
paa meget sammensatte Bevægelser, der først efterhaanden
komme i Stand. Et af de vigtigste Exemplar paa Association
mellem Bevidsthedstilstande af forskjellig Aart er den
Proces, ved hvilken Forestillingen om en ydre Verden opstaaer.
I Almindelighed betragter man denne Forestilling
% --- p. 73 ---
\setcounter{footnote}{0}%
\opage{73}som naturlig og instinktmæssig og aner ikke, at den har
en meget indviklet og sammensat Oprindelse. Mange forskjellige
Bevægelser og Fornemmelser maa ved hyppig Gjentagelse
indgaa en fast Forbindelse, for at hin Forestilling kan
opstaa. Det er Syns- og Følefornemmelserne i Forening
med Muskelfornemmelserne, som frembringe Forestillingerne
om Soliditet og Udstrækning, der ere de væsentlige
Elementer i Forestillingen om ydre materielle Ting. Navnlig
lægger Bain her Vægt paa Fornemmelsen af vore egne
Bevægelser. Naar vi bevæge vore Lemmer og som Følge
deraf modtage visse Sandsefornemmelser, udlede vi disse fra
ydre Ting. Dersom vor Sandsning var fuldkommen passiv,
vilde vor Opfattelse af Universet være som en Drøm, hvor
det ene Billede afløser det andet. Vor Sandsning er nu
derimod, som i det Foregaaende paavist, snarere aktiv end
passiv. Og vi tro paa en ydre Virkelighed, fordi vi vide,
at naar vi foretage en vis Bevægelse, ville vi derefter faa
en vis Sandsefornemmelse. »Vor ydre Verden er Totalsummen
af alle Leiligheder til at udvikle Energi.» Antagelsen
af en ydre Verden er altsaa en paa specielle Erfaringer
grundet Generalisation eller Abstraktion, ved hvilken
vi sammenfatte det Forbigangne og forudsige det Tilkommende;
men ved en Illusion sætter vor Bevidsthed saa
denne Abstraktion som en oprindelig og uafhængig Kilde
og Aarsag til hine enkelte Erfaringer, omtrent paa samme
Maade som Platon statuerer Ideerne om det Gode og det
Skjønne som selvstændige Realiteter udenfor de enkelte
gode og skjønne Ting, af hvilke de ere abstraherede. Men
en saadan absolut selvstændig Yderverden have vi ingen
Grund til at antage. »Verden kan kun erkjendes i sit
Forhold til Aanden. Erkjendelsen er en aandelig Tilstand;
Begrebet om en materiel Ting er en aandelig Ting. Vi
ere ude af Stand til at undersøge Tilværelsen af en uafhængig
materiel Verden; et Forsøg derpaa vilde være selvmodsigende.
Vi kunne kun tale om den Verden, som
fremstiller sig for vor Aand. Ved en Illusion, som Sproget
fører med sig, tro vi os i Stand til at betragte en Verden,
% --- p. 74 ---
\setcounter{footnote}{0}%
\opage{74}der ikke indtræder i vor egen aandelige Tilværelse. Men
selve denne Betragtning er jo en aandelig Virksomhed.»

Der er altsaa hos Bain den samme Tilnærmelse til
Berkeley, som hos James og Stuart Mill. Dog kan man
hos ham lige saa lidt som hos disse drage den Konsekvents,
at Alt maa opløse sig i Skin og Bedrag; thi han fastholder
en gjennemgaaende Dobbelthed i Verdensbetragtningen,
beroende paa Modsætningen mellem Subjekt og
Objekt, Forestilling og Virkelighed. Vor Erkjendelses
Grundvirksomhed er jo ifølge Bain Skjelnen og Sammenlignen.
Den forudsætter derfor bestandig Flerhed og Forandring.
Vi erkjende kun Forhold, Forskjel og Lighed.
En Ting kan kun opfattes, naar der er en Baggrund, en
Folie, i Forhold til hvilken den kan bedømmes og karakteriseres.
Med Opdagelsen af nye Folier, af nye Modsætninger
og Forskjelle rykker Erkjendelsen fremad. I sidste
Instants kunne alle Modsætninger reduceres til Modsætningen
mellem Bevidsthed og Existents (Aktualitet), Subjekt og
Objekt. En Ophævelse af denne Modsætning er umulig; et
absolut Subjekt og et absolut Objekt ere lige utænkelige.
Berkeleys Feil bestod i at lade det Aktuelle gaa op i det
Ideelle, det Objektive i det Subjektive. Denne Modsætning
maa man fastholde, selv om man forkaster Antagelsen af
en fra Aanden absolut adskilt Verden. Aktualitet og
Idealitet ere i deres Væsen beslægtede; Forestilling og virkelig
Bevægelse bero paa samme Nervefunktion. Der kan
ogsaa paavises en hel Række af Overgange fra den objektive
Pol til den subjektive; men selve Modsætningen mellem de
to Poler kan ikke ophæves. »Udstrækningen er det mest
udprægede objektive Faktum, Lyst og Smerte de mest udprægede
subjektive Tilstande. Mellem Bevidstheden om
Udstrækningen og Bevidstheden om Lyst er der den største
Forskjel, som den menneskelige Erfaring kan statuere i
hele den existerende Verden. Det er det extreme Objekt
og det extreme Subjekt; og den gjennemførte Analyse viser
os, at det extreme Objekt beroer paa Fornemmelsen af
anvendt Muskelenergi.»

% --- p. 75 ---
\setcounter{footnote}{0}%
\opage{75}Ideassociation paa Grund af Lighed spiller især en
Rolle i den høiere intellektuelle Virken, hvorved der bringes
Orden og Sammenhæng i det ved Kontiguitetens Virken
ophobede Stof, og hvorved dette finder sin Forklaring og
opfattes paa nye Maader. Den identificerende (Lighed og
Enhed søgende) Kraft virker i al Tænkning, men kan være
tilstede i meget forskjellige Grader. Dens Styrke viser sig
ved Størrelsen og Mængden af de Forskjelligheder, gjennem
hvis Dække Overensstemmelsen maa søges. Paa den beroer
Klassifikation, Abstraktion og Definition, hvorved man identificerer
flere Gjenstande ved at fremhæve et fælles Træk
eller en fælles Gruppe af Egenskaber som særlig Gjenstand
for Betragtning. Jo mere Lighedens Lov virker, des fuldkomnere
bliver Abstraktionen og Klassifikationen. Hvilket
Fremskridt er der f. Ex. ikke mellem Aristoteles's Inddeling
af Dyrene i Vanddyr, Landdyr og Luftdyr og Cuviers
System, som er grundet paa Dyrenes indre Bygning og efter
denne deler dem i Hvirveldyr, Leddyr, Bløddyr og Zoophyter.
I Stedet for de ydre og tilfældige Ligheder træde de indre,
væsentlige Overensstemmelser. Ved Klassifikation, Abstraktion
og Definition danne vi almindelige Begreber. Almindelige
Domme danne vi ad Induktionens Vei. Vi vende
her ikke, som ved Abstraktionen, vort Blik mod en enkelt
Egenskab eller en enkelt Gruppe af Egenskaber, men vi
fremhæve et vist Sammentræf, en vis Forbindelse af to
eller flere Egenskaber og Fakta. Naar vi udsige to
Kjendsgjerningers bestandige Koexistents, udtale vi en
Naturlov.

Kontiguitetens og Similaritetens Love ere Ideassociationens
Grundlove. Men der gives ikke blot enkelte, ogsaa
sammensatte Ideassociationer. Naar jeg f. Ex. betragter
et Bjerg, kan dette føre min Tanke ind paa mange forskjellige
Veie. Ifølge Kontiguitetens Lov kan jeg ledes til
at tænke paa Nabobjergene, paa Dale og Landsbyer i Nærheden,
paa de Mineralier, hvoraf Bjerget bestaaer, paa historiske
Minder, som knytte sig til det. Ifølge Similaritetens
Lov kan jeg komme til at tænke paa lignende Bjerge,
% --- p. 76 ---
\setcounter{footnote}{0}%
\opage{76}jeg har set eller læst om, paa Malerier og Digterværker
o. s. v. Hvilken af alle disse Retninger Tanken vælger,
beroer paa, hvilken af dem der styrkes ved Forbindelse og
Sammensætning med andre Forestillinger. Naar jeg f. Ex.
er Historiker af Profession, virke de historiske Minder
stærkest; er jeg derimod Geolog, vil Bjergets physiske Beskaffenhed
fængsle mig mest o. s. v. Associationer, som
hver for sig ere for svage, kunne faa Magten, naar de virke
sammen med flere andre. Oftest spiller her Følelsen og
Villien en vigtig Rolle, enten hindrende eller fremmende.
Naar de sammensatte Associationer ere af en rent tilfældig
Karakter, tjene de netop til at lede Tanken bort fra det
Væsentlige. Heraf er al Selvmodsigelse og Inkonsekvents
at forklare. Kun naar Lighedens Kraft er tilstrækkelig udviklet,
vil den befri Aanden fra Modsigelser ved at stille
det Modstridende nær sammen og gjøre opmærksom paa
Uoverensstemmelserne. Naar En f. Ex. har to Meninger,
der staa i Modsigelse med hinanden, og mellem hvilke der
altsaa maa vælges, er det Lighedens Lov, der begrunder
Valgets Nødvendighed. Men denne Virken hindres af Tilbøieligheders
og Stemningers Indflydelse, som have knyttet
sig til visse Forestillinger. Heraf kommer det, at absolut
Konsekvents saa sjeldent findes i de menneskelige Anskuelser
og Karakterer.

De Love, som hidtil ere paaviste for Ideassociationerne,
angaa tidligere Forestillingers Gjenopvækkelse. Men
Aanden er ikke bunden til det tidligere Modtagne; den
kan danne nye Forbindelser. Da have vi, hvad Bain kalder
konstruktiv Association. Phantasiens Skaben beroër, som
allerede Hobbes har antydet, herpaa; vi have Forestillingen
om Guld og Forestillingen om et Bjerg, og Eventyrdigteren
danner nu deraf Forestillingen om et Guldbjerg. Ikke blot
Kunst og Poesi bero paa konstruktiv Association, ogsaa
al Naturopfattelse forudsætter den. Thi den virkelige Natur
forstaa vi kun, naar vi forene i en Totalsammenhæng, hvad
Analysen og Abstraktionen have adskilt. Det Samme
gjælder, naar vi ville sætte os ind i Andres Stemninger og
% --- p. 77 ---
\setcounter{footnote}{0}%
\opage{77}Lidenskaber; vi forstaa dem kun, naar vi konstruere deres
aandelige Tilstand ved Sammenfatten af de enkelte Elementer
(Frygt, Vrede, Attraa o. s. v.), I alle saadanne
Konstruktioner, selv de mest originale og skaberiske, er
det intet helt nyt Princip, der virker; der behøves kun en
kraftig Virken og Samvirken af Kontiguitetens og Similaritetens
Love. ---

Bain følger den gamle Tredeling af de sjælelige Phænomener
i Erkjendelse, Følelse og Villie, ikke i Betydning
af særlige Evner eller Kræfter, men som Grupper af sammenhørende
Phænomener, som naturlige Ordener, der hver for
sig igjen indbefatte flere Underafdelinger. Erkjendelseslivet
er forholdsvis lettest at forstaa; Analysen kan snarere
gjennemføres paa den klare Bevidstheds Omraade, hvor de
enkelte Elementer i Regelen ere bekjendte. Vanskeligere
bliver Opfattelsen af Følelses- og Villieslivet, fordi den
dunkle Naturgrund og det ubevidste Liv her spille saa stor
Rolle. Her slaaer Bains beskrivende og analyserende
Methode heller ikke til; den fremdrager kun de Phænomener,
der bevæge sig paa Overfladen, men paaviser ikke
den dybere Sammenhæng. Spencer har, i sin Kritik af
Bains Værk »Emotions and Will,» fremhævet, at ikke blot maa
Analysen gjennemføres til mere primitive Former end Bain
har gjort det, men der maa til Grund for hele Studiet af
Følelseslivet lægges en komparativ Undersøgelse af dets
Fremtræden hos de forskjellige Dyrearter og Menneskeracer
og hos det samme Individ i forskjellige Aldre, for at man
kan lære enhver Følelses Betingelser og Udviklingsgang at
kjende. Dog har Spencer selv kun aldeles i Forbigaaende
behandlet dette vanskelige psykologiske Problem.

Til Grund for sin Inddeling af Følelserne lægger Bain
tre Kjendemærker, nemlig foruden selve Bevidstshedstilstandenes % PRINTED AS IS (Bevidstshedstilstandenes: O-D13)
Karakter deres physiologiske Virkninger og deres æsthetiske
eller moralske Ytringsformer. Ad denne Vei kommer han
til at opstille følgende Klasser. 1) Naar Indtrykkene
stemme overens, harmonere, saaledes at Følelsen i det Hele
kan have sit frie Løb, opstaaer der Lyst; i modsat Til%
% --- p. 78 ---
\setcounter{footnote}{0}%
\opage{78}fælde en Følelse af Tvang og Smerte. 2) Maaske den
mest elementære Følelse er Forundringen, Følelsen af Overraskelse;
den er, som vi allerede have hørt, knyttet til selve
Bevidsthedens Grundlov, ifølge hvilken Relation, Modsætning
og Forandring ere nødvendige for al Opfattelse. 3) Frygt
og alt hvad dertil hører, Overtro, Ængstelse o. s. v., samt
dens Modsætning, Modet i dets forskjellige Arter. 4) Følelsen
af Ømhed og Velvillie, som er af en overveiende
passiv Karakter, og snarere sætter Kirtlerne end Musklerne
i Bevægelse. Ømhedsfølelsen (the tender affection) maa
ikke forvexles med Sympathien, der er en meget sammensat
Følelse, som kan udvikle sig af Ømhedsfølelsen, men
ikke nødvendig forudsætter den. Sympathien forudsætter,
at man ganske har sat sig ind i en Andens Titstand og % PRINTED AS IS (Titstand: O-D14)
betragter den som sin egen; men Ømhedsfølelsen kan godt
indskrænke sig til Ens egen Person, og vi faa da 5) Selvkjærligheden
og Selvtilfredsheden, der have deres Grund i
den passive Nydelse og Glæde ved det egne Jeg. 6) Følelsen
af vor Magt og af den overlegne Virken, vi udfolde,
paa hvad Omraade det saa er (Politik, Rigdom, kunstnerisk
og videnskabelig Dygtighed o. s. v.). Denne Følelse kan
smelte sammen med Selvkjærligheden, men de ere dog i
deres Væsen forskjellige. Der findes Mænd, især haarde og
strenge Naturer, hos hvem Magtbevidstkeden spiller en stor % PRINTED AS IS (Magtbevidstkeden: O-D15)
Rolle, uden at de dog kunne siges at være Egoister eller
Selvbeundrere. 7) Følelsen af Vrede, hvor det ikke er en
Lyst, der fremkalder en anden Lyst, men en Smerte, der
bevirker en anden Smerte. 8) Den Følelse af Lyst, som
bevirkes ved ivrig at forfølge et Formaal (i Alvor eller i
Leg), ved Anstrengelser, hvor man ganske er optaget af én
Gjenstand (emotion of pursuit; plot-interest). 9) Den intellektuelle
Nydelse, den Følelse af Lyst, som Tankens
Virksomhed fremkalder, navnlig ved Opdagelsen af Lighed,
Overensstemmelse og Konsekvents; herpaa beroer Glæden
og Interessen ved videnskabelige Opdagelser, digteriske
Billeder og industriel Kløgt. Paa den anden Side fremkalde
Inkonsekvents og Disharmoni en pinlig Fornemmelse især
% --- p. 79 ---
\setcounter{footnote}{0}%
\opage{79}hos tænkende Naturer.\footnote[1]{„Der er en Art Smerte, som er karakteristisk for intellektuelle Naturer, og som saare lidet føles af Andre; det er Smerten ved inkonsekvente Forestillinger, Meninger og Principer. Stærk Følelse forenet med mindre udviklet Forstand fører let ud over det Anstød, Modsigelsen vækker; men paa et modsat begavet Sind virker Inkonsekvents som en skurrende Lyd og bliver en Spore til intellektuel Anstrengelse, der forhøier den naturlige Tilbøielighed til at lade Aandens herskende Evne virke. Denne Smerte maa vi tilskrive en vigtig Funktion, nemlig Ansporingen til at stræbe efter Sandhed og ikke blot efter mangfoldige intellektuelle Billeder og Forestillinger. . . . Heri ligger et stort Motiv til at gjennemføre den videnskabelige Undersøgelse og udrydde populære Vildfarelser.“ \emph{Bain}: On the study of character. p. 272 f.} Disse Følelser blive af Vigtighed
ogsaa i moralsk Henseende ved at lede til Sandhed og Retfærdighed.
--- Til disse elementære Følelser slutte sig endelig
de mest sammensatte af alle, de æsthetiske og moralske
Følelser.

Allerede ældre engelske Forskere, Alison og efter ham
James Mill, have anvendt Ideassociationens Love til at forklare
Følelsen for det Skjønne og det Ophøiede. Coleridge
gjorde, efter Kant, opmærksom paa Vanskeligheden ved at
skjelne mellem det Skjønne og det Behagelige, naar man
holdt sig til de blotte Ideassociationer. Den engelske
Æsthetiker John Ruskin lader Ideassociationerne tjene som
Midler til at fremkalde høiere Ideer; skjønt er efter ham,
hvad der minder om guddommelige Ideer og Attributer
(Uendelighed, Enhed, Fred, Harmoni, Renhed o. s. v.), der
finde deres Udtryk i Verden.\footnote[2]{\emph{Milsand}: L’esthétique Anglaise. Etude sur Mr. John Ruskin. Paris 1864. p. 149. Kfr. \emph{Stuart Mills} Note 48 til James Mills „Analysis“ II.} Endnu mere end Ruskin
viser Bain her tilbage til Kant. For at et Indtryk i det
Hele skal være af æsthetisk Beskaffenhed, siger han, maa
det ikke være rent sandseligt, da det saa faaer en interesseret
og egoistisk Karakter, som gjør det rent individuelt
og umuliggjør en ren og fælles Nydelse. Derfor ere Syn
og Hørelse de mest æsthetiske Sandser; ved disse er det
lettest muligt at se bort fra den egoistiske Drift. Det
% --- p. 80 ---
\setcounter{footnote}{0}%
\opage{80}egentlig Skjønne beroer paa Harmoni. Det Ophøiede beroer
paa Sympathi. Vor Tanke gaaer uvilkaarlig i samme Retning
som Indtrykkene: det Store i Naturen vækker Tanken
om det Store i os, især i den menneskelige Aands høieste
Former. Det sande Ophøiede er den menneskelige Magt:
det Ophøiede, vi finde i Naturen, opfatte vi kun ved at
iføre denne menneskelige Former og fortolke den efter
menneskelige Idealer.

De moralske Følelser ere ifølge Bain endnu mere
sammensatte end de æsthetiske. Han polemiserer mod de
ensidige Forklaringsforsøg, som udlede alle Handlinger af
en enkelt Drift eller Side i Menneskets Væsen. Han
polemiserer tillige mod Antagelsen af en universel moralsk
Lov, som skulde kunne tjene til Regel og Maalestok i alle
enkelte Tilfælde. Hvor finde vi en saadan Lov? Vi kunne
stille vore Uhre efter Observatoriet i Greenwich; men hvor
faa vi et fælles Signal til at regulere vore Handlinger? Der
gives lige saa lidt nogen almindelig moralsk Lov, som der
gives en almindelig Fornuft; hver Enkelt har sin Fornuft
og fælder sine Domme efter sin individuelle Maalestok
Der gives lige saa lidt noget i sig Godt, som der gives
noget i sig Skjønt eller Sandt; disse Ideer existere ikke
udenfor de individuelle Bevidstheder: selv om Individer
stemme overens i deres Domme, kan denne Overensstemmelse
dog lige saa lidt udledes af en i sig existerende
Sandhed, som man kan tillægge et Sangerchor en almindelig
Stemme, forskjellig fra de enkelte Syngendes Stemmer. ---
Derimod behøve vi blot at kaste Blikket omkring os og
læse Historien for at finde de Magter, der virkelig raade i
den sædelige Verden. Det er ikke et usynligt moralsk
Ideal, men reale Autoriteter, der virke. Vi finde ingen oprindelig
Følelse af Pligt og Retfærdighed; de oprindelige
Drivfjedre ere Frygt for Straf, Respekt for Autoriteten,
Ømhedsfølelse og Sympathi, Egenkjærlighed o. s. v. Af
disse vil Egoismen stedse vinde Overhaand, dersom ikke
selve den egoistiske Drift bringes i Pligtens Tjeneste,
hvilket sker ved Forholdet til en straffende Myndighed. For%
% --- p. 81 ---
\setcounter{footnote}{0}%
\opage{81}pligtelse er det moralske Grundphænomen. Men Forpligtelse
er der kun Tale om, hvor Overtrædelse af en Lov
medfører Straf. Ansvarlighed er altsaa Strafbarhed. Og
Straffen forudsætter igjen en ydre Magt, en virkelig Autoritet,
fra hvem Buddet udgaaer, og som kan revse dets
Overtrædelse. Men ved Ideassociationens Magt tabes det
oprindelige Motiv efterhaanden af Sigte; de Dyder, der
først kun øvedes af Tvang, blive efterhaanden Formaal i
sig selv, ligesom Penge blive et Formaal i sig selv, medens
de oprindelig kun attraaedes som Middel til visse Fornødenheders
Tilfredsstillelse. Ideassociationerne have en Tendents
til at frembringe nye Tyngdepunkter, hvor ofte mægtigere
Følelser kunne koncentreres, end de oprindelige Motiver
formaaede at fremkalde. Saaledes emanciperer Samvittigheden
sig efterhaanden fra Frygten for Straf, især ved
Hjælp af Erkjendelsens og Sympathiens Virken. Den Underkastelse
og Følelse af Forpligtelse, som Individet har lært
i Forholdet til den ydre Magt, overføres nu paa Forholdet
til den indre, ideelle Magt. Derved bliver Mennesket aandelig
frit, hæver sig over den slaviske Tænkemaade.

Som Bidrag til Menneskets ethiske Udviklingshistorie
er denne Skildring fortræffelig. Langt fra, som Bain mener,
at modbevise den ethiske Lovs Selvgyldighed, bekræfter den
netop denne. Thi den godtgjør, at der gives en Udviklingslov,
i Kraft af hvilken Mennesket hæver sig fra lavere
til høiere og ædlere Aandsformer. Denne Lov er en Naturlov,
men en Naturlov, som den selvbevidste Aand kan gjøre
til sin Livsregel, hvorefter den kan danne sit Ideal; her
er da Overgangen fra Naturlov (i snevrere Betydning) til
Morallov. Det er i Kraft af denne Lov, at Mennesket kan
vende sig med Kritik og Protest mod Autoriteter, overfor
hvilke han maaske udadtil maa bukke under; det er i Kraft af
denne Lov, at Autoriteterne kæmpe indbyrdes, og den ene
fortrænger den anden. Under saadanne Konflikter viser
det sig klart, at den egentlige Magt er den indre, den
ideelle Magt, der bruger hine ydre Kræfter som Midler for
sig egen Udfoldelse. Ifølge Bains Lære maatte man slutte,
% --- p. 82 ---
\setcounter{footnote}{0}%
\opage{82}at Moraliteten ophører, naar Autoriteterne tabe deres Magt.
Denne Konsekvents er ogsaa dragen af Kirchmann, en
skarpsindig tydsk Forfatter,\footnote[1]{I. H. v. \emph{Kirchmann}: Die Grundbegriffe des Rechs und der Moral. Berlin 1869.} som gaaer ud fra lignende % PRINTED AS IS (Rechs: O-D19)
Forudsætninger som Bain. Det Sædelige er da selv kun
et Middel, som falder bort, naar Maalet er naaet. Men
dersom det Sædelige har sin Grund i Menneskets Væsen
som anlagt til Udvikling, saa vil der aldrig komme den
Tid, hvor det skulde kunne afløses af andre Ideer, --- selv
ikke i »den stationære Tilstand.»

Ifølge Bain maa der bestemt skjelnes mellem Pligt og
Fortjeneste. Hvad man er forpligtet til, er kun hvad man
vilde blive dadlet eller straffet for at undlade; hvad man
gjør ud over dette, er Fortjeneste. Dette er ikke blot
konsekvent, men tillige i høi Grad karakteristisk for den
empiriske Moral, der begrunder det Sædelige paa faktiske
Tilstande (Autoritet og Drift) og derfor ikke kan omfatte
den ideale Stræben, som gaaer ud over disse, --- og denne
Stræben er dog det Sædeliges høieste Form. Paa dette
Punkt er Empirismen enig med Katholicismen, der har en
Regel for menige Christne og en anden for Helgene. Mill
beundrer ogsaa i denne Henseende »de kloge og klart seende
Mænd, der skabte den katholske Kirke.» Den protestantiske
Ethik har nedlagt en energisk Indsigelse mod denne
Modsætning mellem Pligt og Fortjeneste. Vi ere forpligtede
til at indlægge os al den Fortjeneste, vi kunne; --- selv
naar vi gjøre det, have vi dog kun gjort vor Pligt. Der
er i disse Begreber en Dialektik, som godtgjør den empiriske
Morals Utilstrækkelighed. Dette er et af de Punkter,
hvor den franske Positivisme afviger fra den engelske Empirisme.\footnote[2]{Kfr. \emph{Littrés} Polemik mod Utilitarianismen: La science au point de vue philosophique. 2éd. p. 340 f.}
Den energiske og af Enthusiasme glødende
Auguste Comte, der ellers beundrer den middelalderlige
Katholicisme, stemmer i dette Punkt med Protestantismen
% --- p. 83 ---
\setcounter{footnote}{0}%
\opage{83}og fordrer af Enhver en absolut Hengivelse i Menneskehedens
Tjeneste, en fuldstændig Opgivelse af Egoismen. Den empiriske
Skole, der med Bentham og Mill bygger paa »den
velforstaaede Interesse,» kan ikke anerkjende en saadan Forpligtelse,
ja betragter vel endog Selvopofrelsen og den uinteresserede
Handlen med mistænksomt Blik.\footnote[2]{„Uinteresserethed er, strengt taget, en Art Ufornuft eller Afsindighed med Hensyn til Individets eget Selv, en Modsigelse mod det følende Væsens inderste Natur, som bevæger sig hen efter Lysten og bort fra Smerten  Naar vi opgive dette fundamentale Princip for vor Konstitution, staa vi stedse paa Grændsen af Urimelighed og Lovløshed. Thi naar en Mand har ophørt at sørge for sit eget Vel eller at handle i Retning af sin egen Lyst, hvorfor skulde han da sørge for nogensomhelst Andens Velfærd? Han har fornegtet Lyksalighed som Formaal og maa, hvis han er konsekvent, udstrække samme Maalestok til Alle omkring ham.“ \emph{Bain}: On the study of character. p. 292.} % PRINTED AS IS (the mark in the text is **), the note is labelled *); zoom z3)

Bains ethiske Principer godtgjøre det Mangelfulde i
hans psykologiske Undersøgelser. Han har ikke gjennemført
Analysen af Følelserne saaledes, at de førte til at opdage
Grunddriften i den menneskelige Natur, men bliver
staaende ved en Flerhed af primitive Følelser En videnskabelig
Ethik maa have sit reale Grundlag i hin Grunddrift.
Det Vaklende i Utilitarianismens Begreber har netop
sin Grund i en Uklarhed med Hensyn til det oprindelige
psykologiske Udgangspunkt. Dernæst har Bain heller ikke
fremhævet den centrale Betydning, Erkjendelsen har for
den ethiske Bevidstheds Lutring; det er ved Tankens
Arbeide, at der opgaaer for Bevidstheden en Forestilling
om et Ideal, som viser ud over det Endelige, Faktiske og
Egoistiske. ---

Den oprindelige Aktivitet, som ifølge Bain er Sjælelivets
Udgangspunkt, traadte i Baggrunden ved Undersøgelsen
af Følelseslivet, men den maa ifølge Sagens Natur
atter blive af Betydning ved Undersøgelsen af Villiens
Væsen. Den tidligere Associationspsykologi maatte egentlig
ganske negte Villiens Existents. Ifølge Bain har Villiens
Udvikling to Forudsætninger: den spontane Bevægelse, som
% --- p. 84 ---
\setcounter{footnote}{0}%
\opage{84}gaaer forud for al Fornemmelse, og det Baand, som lader
en vis Bevægelse følge efter en vis Fornemmelse. Villiens
Fødselsstund er det Øieblik, hvor en Fornemmelse første
Gang bestemmer en Handling. Den uvilkaarlige Bevægelse
bliver ved Association med Fornemmelsen af Lyst eller
Ulyst en vilkaarlig Handling. Ad denne Vei faaer Individet
Evne til mere sammensatte Bevægelser og opdrages til
Herredømme over sine Lemmer. Efterhaanden sker Overgangen
fra Fornemmelse til Handling ikke øieblikkelig.
Ind imellem begge skyder sig Forestillingen først om Handlingens
direkte, dernæst ogsaa om dens indirekte Følger.
Der opstaaer en Kamp mellem de forskjellige Forestillinger
og de Impulser, de give; overfor den nærværende eller nære
Lyst og Smerte stiller sig Forestillingen om mulig indtrædende
stærkere Fornemmelser. Denne Debat, under hvilken
Handlingen suspenderes, er Overveielsen. Den er oprindelig
fremkaldt ved den Erfaring, at hurtig Handlen førte Smerte
med sig. Ved den udvides Villiens Omraade, og en høiere
Udvikling, en Forædlen af Motiverne bliver mulig. De
frivillige Handlinger adskille sig fra Reflexvirkningerne ved,
at Bevidstheden træder ind imellem Indtrykket og Handlingen.
Herved bliver et Herredømme muligt ei blot over
vore ydre Handlinger, men ogsaa over vore Tanker og
Følelser, der jo selv egentlig ere Bevægelser. Men dette
Herredømme kan kun udøves indirekte, ad en Omvei; vi
hemme Følelsen ved at holde dens ydre Tegn tilbage eller
ved at vende vor Opmærksomhed mod visse Forestillinger,
der kunne fremkalde andre Følelser.

Den ældre skotske Skole (Reid og hans Disciple, ja
endnu W. Hamilton) antog Villien i Betydning af det frie
Herredømme over Lemmerne for en oprindelig Evne.
Bain erklærer dette for en Illusion. Villien har sin Spire
i den oprindelige Aktivitet, men den maa gjennemgaa en
lang Proces for at blive til, hvad man i Almindelighed
kalder fri Villie. Om Frihed kan der her ikke være Tale
i den Forstand, at en fra de enkelte sjælelige Phænomener
forskjellig Magt, et indre Jeg skulde gribe ind; Associa%
% --- p. 85 ---
\setcounter{footnote}{0}%
\opage{85}tionernes og Motivernes Række kan ikke afbrydes. Bain
behandler dette Spørgsmaal lige som Stuart Mill og polemiserer
lige som denne saa vel mod Frihedslæren som mod
Nødvendighedslæren. Navnlig afviser han paa en træffende
Maade Paaberaabelsen af Bevidsthedens Udsagn, af Følelsen
af Frihed. Den umiddelbare Bevidsthed, siger han, kan
ikke lære os Mere om de psykologiske Love, end den
umiddelbare Betragtning af Himmelen kan lære os om de
astronomiske Love. Bevidstheden lærer os de enkelte
sjælelige Phænomener at kjende, men lærer os ikke deres
Natur og Sammenhæng, som kun kan findes ved Undersøgelse
og Analyse. I ethvert Tilfælde have vi ingen ufeilbar
Autoritet i den. Ufeilbarhed og Erkjendelse udelukke
hinanden. Villien er, som det har vist sig, et meget
sammensat Phænomen; ved Kombinationen af de mange
Enkeltheder, der høre med til at opfatte den, kan der derfor
begaaes altfor mange Feil, til at vi skulde kunne bygge
paa en umiddelbar Forvisning. ---

Til Bains Theori om Villien slutter sig nøie hans
Opfattelse af Begrebet Tro eller Antagelse (kfr. allerede
ovenfor p. 73). Hvad hører der til, for at vi skulle tro
Noget? James Mill, til hvem her Herbert Spencer slutter
sig, forklarede Troen som en fast og uopløselig Association
mellem to Forestillinger. Men ifølge Bain, med hvem Stuart
Mill har erklæret sig enig, kan Troen ikke forklares ved at
betragtes som et rent intellektuelt Phænomen. Det er
Forbindelsen med et praktisk Formaal, med en Handlen,
der adskiller Troen fra en blot Phantasi. Den mest primitive
Form for Tro er Forventningen om, at der efter en
vis Handling vil følge en vis Virkning. Den oprindelige
Aktivitet har en uimodstaaelig Trang til at udlade sig, og
Udladningen sker naturligvis i den Retning, i hvilken den
tidligere har vist sig at medføre Tilfredsstillelse. Det sultne
Rovdyr holder sig til den Vei, paa hvilken det før har
truffet sit Bytte. Synet af Bækken lokker den Tørstige,
fordi han der venter Tilfredsstillelse af sin Trang. Ogsaa
de høiere Former af Tro have deres Kjendemærke og Maale%
% --- p. 86 ---
\setcounter{footnote}{0}%
\opage{86}stok i den Handlen, de føre til; de, som ikke haabe, frygte
og handle i Overensstemmelse med deres religiøse Tro,
indbilde sig kun at tro, men tro i Virkeligheden ikke.

Af den oprindelige Aktivitet udspringer der derfor en
stærk Tendents til at tro. Her ligger Spiren til alle Generalisationer,
alle almindelige Sætninger: de første Erfaringer
ophøies af den primitive Drift til almindelige Love. Saaledes
opstaa de intuitive eller aprioriske Antagelser; og saaledes
er navnlig Forestillingen om det Uendelige og det
Ideale opstaaet. Men den fremskridende Erfaring berigtiger
Skridt for Skridt de oprindelige Intuitioner. Begyndelsen
er den spontane Antagelse, den intuitive Anskuelse; men
snart reiser Tvivlen sig ved de stridende Tilfælde, Erfaringen
paaviser. Erfaringen er derfor den eneste sikre Erkjendelseskilde,
den virkelige Autoritet, efter hvilken vi maa maale
alle vore Antagelser. Vore Anskuelser voxe altsaa ikke en
rolig og jevn Væxt, men danne sig under en stadig Kamp
mellem den spontane, generaliserende og idealiserende Intuition
(Bacos antecipatio naturæ) og den overveiende og
prøvende Erfaring. Vi have en stærk Tendents til først at
tro og saa siden efter bevise, og det vil sandsynligvis vare
længe, inden de Skuffelser, som denne Fremgangsmaade
medfører, ville faa en Ende.\footnote[1]{„Mennesket har ganske vist en ubegrændset Higen, og den metaphysiske Lære om det Uendelige svarer paa en Maade til den. Men i det virkelige Liv finde vi saare faa af vore Ønsker fuldt tilfredsstillede, end ikke dem, der ere mest ærefulde for den menneskelige Aand, som Erkjendelsesdriften, Attraa efter Selvudvikling, Driften til at gjøre vel. Hvorledes kunne vi da forvisse os om, hvilken Higen der fører sin egen Tilfredsstillelse med sig? Især da de intuitive Bestræbelser gaa i meget forskjellig Retning hos Menneskene, og alle ikke kunne være lige sande“. \emph{Bain}: The emotions and the will. London 1859. p. 587.}

Denne sidste Udvikling karakteriserer Bains hele Standpunkt.
Den oprindelige Aktivitet, som er hans psykologiske
Opdagelse, og ved hvis Anerkjendelse han afgjort forlader
den ældre sensualistiske Skole, bruger han kun som et
Begyndelsespunkt, fra hvilket han igjen stræber at befri
% --- p. 87 ---
\setcounter{footnote}{0}%
\opage{87}sig, idet han lægger Hovedvægten paa Erfaringerne og Ideassociationerne.
Foruden hin oprindelige Aktivitet, som
endnu tilhører det forbevidste Liv, antager han endnu et
oprindeligt aandeligt Element, der ligger til Grund for alle
Bevidsthedsytringer, nemlig den oprindelige Sands for Forskjel;
uden den kan der heller ikke være Tale om Erfaring
eller Ideassociation. Men vi standse saaledes ved en Modsætning
mellem de aktive og de passive Aandsphænomener,
som dog i enhver virkelig aandelig Funktion skulle gaa
sammen, uden at der gjøres noget Forsøg paa at paavise,
hvorledes Forholdet imellem dem er. Sandsningen kan ei
foregaa rent passivisk, naar den har den spontane Aktivitet
og den elementære Forstandsfunktion til sin Forudsætning.
Og lige saa lidt kan Ideassociationen være en rent passivisk
og mekanisk Akt. Navnlig viser Lighedens Lov afgjort
tilbage til den oprindelige Energi, med hvilken Aanden
stræber efter Enhed i sin Anskuelse og sin Tænkning. Det
viser sig paa alle Punkter, at Modtageligheden forudsætter
Virksomheden, Hengivelsen Friheden, ligesom paa den
anden Side det Modtagne forhøier Virksomheden, Hengivelsen
lutrer og forædler Friheden. Vi komme da ogsaa
her til det Resultat, at Empirien fører ud over sig selv,
baade ved de oprindelige ideelle Forudsætninger, den hviler
paa, og ved det ideelle Formaal, i hvis Tjeneste den arbeider.
Ideerne om det Uendelige og Absolute have, netop
ifølge Bains Psykologi, deres Kilde i Menneskets inderste
Væsen. Saa kan det da kun være tilsyneladende, at den
\emph{menneskelige} Erfaring modbeviser dem. Erfaringen
kan gjendrive de betingede og historiske Former, hvorunder
Ideerne optræde; men denne Gjendrivelse kan kun blive
Indledning til Arbeidet paa et nyt og fuldkomnere Udtryk,
svarende til de nye Forudsætninger. Bains træffende Ytring,
at Ufeilbarhed og Erkjendelse udelukke hinanden, gjælder
intet Sted mere end ved Erkjendelsen af det Ideale og
Uendelige; --- og dog synes han selv at gaa ud fra, at vi
ikke kunne naa en saadan Erkjendelse, fordi vi intet Ufeilbart
kunne naa!

% --- p. 88 ---
\setcounter{footnote}{0}%
\chapterhead{\opage{88}Tredie Kapitel.}{Den kritiske Skole i England.}


Empirismen fører, som det har vist sig, baade paa
negativ og positiv Maade ud over sig selv. Negativt ---
thi ved at stille Erfaring og Ideassociation i Modsætning
til den indre Nødvendighed afskjærer den sig selv Muligheden
af enhver Begrundelse, af at komme ud over det
Tilfældige. Positivt --- thi dens egne Undersøgelser føre til
Anerkjendelse af det menneskelige Væsens oprindelige Aktivitet
og af den menneskelige Tanke som primitivt og irreduktibelt
Element i al Erkjendelse. Der er en nøie Forbindelse
mellem begge Resultater, det positive og det negative.
Thi den sidste Begrundelse, som den menneskelige
Erkjendelse formaaer at vinde, kan kun opnaaes ved at gaa
til al Erkjendelses sidste Forudsætning, der ikke kan søges
andet Steds end i selve det erkjendende Væsens egen Natur
og oprindelige Organisation. Og her staa vi da netop ved
Kants Grundtanke. At vi ikke, som Mill antog, kunne
tænke os en absolut Uorden, en ubetinget Tilfældighed, men
at selv Phantasien vil negte sin Tjeneste hertil, kommer af,
at hint oprindelige Tankeelement bringer en Sammenhæng
tilstede i al Sandsning og Erkjendelse. Og paa den anden
Side bliver det nu klart, hvorledes vi kunne fælde Domme,
der have Gyldighed for al mulig Erfaring og i alle Tilfælde
(som de logiske og mathematiske Sætninger); det er i
Kraft af det samme, til Grund for al Erfaring liggende og
derfor af Erfaringen relativt uafhængige Element. Vi bestemmes
kun ved Erfaringen, idet vi selv bestemme Erfaringen.


Af Kants kritiske Philosophi, hvis store Betydning
beroer paa at have fremdraget denne Betragtningsmaade af
Erkjendelsen, ere de Mænd stærkt paavirkede, der i Eng%
% --- p. 89 ---
\setcounter{footnote}{0}%
\opage{89}land staa som Empirismens Modstandere. Det er paa den
anden Side interessant at se, at den engelske Philosophis
almindelige Retning mod Induktionen og det Reale ogsaa
tydelig raader hos dem. Ogsaa de anerkjende, at Erfaring
og Induktion ere vore væsentlige Erkjendelsesveie; de vende
sig med skarp Polemik mod de tydske spekulative Systemer
og hævde det Relative og Empiriske som vor Erkjendelses
egentlige Omraade. Men de hævde, at Erfaringen ikke kan
være sin egen Grundvold, men forudsætter visse Principer,
paa hvilke dens egen Mulighed og Gyldighed beroer.\footnote[1]{Dog have ogsaa den tydske Spekulations Ideer vundet Indgeng i England. Ikke blot gives der Hegelianere i Oxford, men allerede \emph{Coleridge} udviklede selvstændig flere Schellingske Ideer, navnlig paa det æsthetiske Omraade (Phantasiens Begreb). Og særlig maa her fremhæves en nyere Forfatter, den bekjendte Historiker \emph{Thomas Carlyle} (f. 1795), der i sin „Sartor resartus“ (forfattet 1831, udkommen nogle Aar senere, først i Nordamerika, hvor Carlyle er meget læst) i humoristisk Form har fremstillet en human Livsanskuelse, der er udsprungen af alvorlig og dybsindig Tænkning. Hans Verdensanskuelse er en Idealisme, for hvem „alle Gjenstande kun ere Vinduer, gjennem hvilke det philosophiske Øie skuer ind i selve Uendeligheden.“ „Hele det ydre Univers og Alt, hvad det indbefatter, er kun en Klædning; og al Videnskabs Væsen ligger i Klædernes Philosophi.“ Alt er Symboler; men det høieste Symbol er i de heroiske Personligheder, Menneskehedens Forbilleder. --- Stuart Mill udtaler i sin Autobiographi, hvor meget han skylder Carlyle, som ofte ad sin intuitive Vei udviklede Ideer, Mill først lang Tid efter naaede til gjennem Induktionen. Saaledes vil man bl. A. i det anførte Skrift af Carlyle (Book III. chap. IV.) finde Mills sociologiske Anskuelse in nuce. Carlyle opfatter heller ikke Philosophien blot som Spekulation; foruden den ovenfor anførte Definition giver han ogsaa en anden, der viser Theoriens praktiske Betydning: „Hvad er Philosophi Andet end en stadig Kamp mod Vanen, end en stedse fornyet Anstrengelse for at gaa ud over (to transcend) den blinde Vanes Sphære og derved blive transcendental?“ Denne Definition er træffende overfor Empirismen, ifølge hvilken al Erkjendelse beroer paa Vane, saa at Forsøget paa at komme ud over Vanen er selvmodsigende, uagtet Mill har troet i Empirismen at have det eneste brugelige Middel til at bekæmpe overleverede Vaner og Fordomme. --- Dog er for Carlyle ikke Tanken, men Phantasien det Høieste; „Forstanden er Vinduet, men Phantasien er Øiet, der ser igjennem det.“ De philosophiske Ideer ere for ham tjenende Grundlag for Mystik og Moral.} % PRINTED AS IS: „Indgeng“
En afgjørende Modsætning er der mellem de Forskere,
% --- p. 90 ---
\setcounter{footnote}{0}%
\opage{90}vi skulle beskjæftige os med i dette Kapitel, og deres
Modstandere. Den empiriske og realistiske Skoles Repræsentanter
have deres Livsanskuelse i deres Philosophi. For
Whewell og Hamilton er Philosophien derimod kun et Udenværk;
den egentlige Helligdom, selve Livsanskuelsen hviler
hos begge disse Mænd paa den positive Religion. Dette
maa naturligvis ogsaa drage Følger med sig i philosophisk
Henseende, især med Hensyn til Konsekvents og Gjennemførelse
af Principerne. De philosophiske Problemer fastholdes
og debatteres ikke i deres Renhed, naar en anden
Interesse spiller ind med baade hos Angribere og Forsvarere.
Skjøndt Kampen kun undtagelsesvis har strakt sig ind paa
det religiøse Omraade, maatte den Følelse dog ofte opstaa,
at det var dette, den egentlig gjaldt. Paa den anden
Side har denne Omstændighed givet Kampen en forhøiet
Interesse, idet Principspørgsmaalene, som Englænderne ellers
helst skjule under Enkelthedernes Masse, derved ere traadte
klarere frem. ---

I. \emph{William Whewell} (f. 1795, d. 1866). ---
At ogsaa andre Forskere end de, der tilhøre den empiriske
Skole, kunne tilkjende Erfaringen principiel Betydning og
indlægge sig Fortjeneste af den, derpaa have vi et Exempel
i Whewell, der ikke blot selv begyndte som naturvidenskabelig
Fagmand (Mineralog), men hvis »History of the
inductive sciences,» ifølge Stuart Mills egen Erklæring var
ham et uundværligt Grundlag ved Udarbeidelsen af hans
Logik, som uden dette Værk ikke havde kunnet gjennemføres.
Studiet af Naturvidenskaberne førte Whewell til at
undersøge deres Methode og Principerne og Betingelserne
for deres Udvikling. Disse Undersøgelser førte ham mere
og mere til en Opfattelse af Erkjendelsens og Erfaringens
Væsen, som i det Hele stemmer med den kantiske, skjøndt
han ogsaa, ikke altid til Fordel for Bestemtheden og Konsekventsen,
er paavirket af den skotske Skole. Vi finde
ogsaa hos ham denne Skoles Tilbøielighed til at appellere
% --- p. 91 ---
\setcounter{footnote}{0}%
\opage{91}til den almindelige Bevidsthed og lade sig nøie med de
overleverede Forestillinger i Stedet for at gaa ind paa
strengere philosophiske Analyser.

Efterat have fuldendt »de induktive Videnskabers Historie»
(1837) udgav Whewell sit andet Hovedværk »Philosophy
of the inductive sciences, founded upon their history»
(1840, omarbeidet 1847), i hvilket han ved Hjælp af Videnskabernes
historiske Udvikling søger et efterspore Forholdet % PRINTED AS IS (et efterspore: O-E04: for at)
mellem Erkjendelsens Elementer. Til dette Værk ville vi
i det Følgende især holde os. ---

Ikke blot Naturvidenskaberne, ogsaa Aandsvidenskaberne
(Psykologi og Ethik) ere induktive Videnskaber. Men
der er et Element i Induktionen, man ofte har overset.
De store Opdagelser og Fremskridt i Videnskabens Historie
ske ikke blot gjennem Paavisningen af nye Fakta; en
væsentlig Betingelse er de Ideer og Former, ved hvis Hjælp
disse Fakta opfattes og bestemmes; ja, selve Paavisningen
af Fakta bliver ofte kun mulig ved, at Forskeren ledes af
den rigtige Betragtningsmaade. De enkelte Fakta skulle
ikke blot bringes sammen, men der føies ved selve den
aandelige Aktivitet, hvorved de kombineres, et nyt Element
til; der dannes en almindelig Sætning, som ikke existerede
i noget af de iagttagne Fakta i og for sig. Der maa
findes en Mellembestemmelse, i hvilken de bekjendte og de
ubekjendte Fakta kunne mødes. Saaledes opfattede Kepler
Planeternes Baner som Ellipser, og herved blev det ham
muligt at opstille sine Love.\footnote[1]{„Kun de,“ siger Whewell (Gesch. der induct. Wissensch. übers. von Littrow I. p. 426) i Anledning af Kepler, „opnaa Besiddelsen af nye Sandheder, der kunne foretage et livligt Spring fra et Punkt indenfor deres Erkjendelse til andre, som ofte ere langt fjernede derfra, og bringe en Forbindelse i Stand mellem dem, --- ikke de, der forsigtig staa stille paa ethvert Sted og vente, indtil de udenfra drives til at gaa videre.“ Smlgn. \emph{Reuschle}: Kepler und die Astronomie. Frankfurt 1871. p. 141: „Paa det Spørgsmaal, hvorledes Kepler har fundet disse overordentlige Sætninger (de tre berømte Love), er vel det korte Svar, at han har regnet dem ud af Tycho Brahes Iagttagelser, idet disse vare saa gode, at de kun fyldestgjordes ved den \opage{92}elliptiske Theori, og han havde mathematisk Kraft og Udholdenhed nok til ikke at hvile, førend han i alle Henseender havde gjort Fyldest for hine Iagttagelser. Men vi maa nu stige dybere ned i hans Genis Værksted og forfølge Gangen i hans Forsken ud fra hans Grundtanke, for at se, hvorledes hans objektive, physiske Værks Historie paa det Inderligste hænger sammen med hans subjektive, metaphysiske Systems Historie ...... Til Grund for alle Keplers mathematiske Undersøgelser og metaphysiske Forskninger ligger den Ide, hvis Opfattelse alene vilde være nok til at grunde hans Storhed, at Solsystemet, hvis raa ydre Omrids Kopernikus havde givet, kun da er et virkeligt System, naar der hersker Lov og Orden deri, saa at de Love, der beherske Bevægelserne, have deres Aarsag i bestemte indre Kræfter, som ere væsentlige for Systemets og dets Legemers Natur.“}% the note runs over from p. 91 to p. 92 (joined by draft.py)
 Paa samme Maade op%
% --- p. 92 ---
\setcounter{footnote}{0}%
\opage{92}fattede Newton Tyngden som Tiltrækning og denne som
universel Verdenskraft. Kun under Forudsætning af en
saadan Ides Indførelse bliver en Opdagelse mulig. Men i
Almindelighed opfatter man Induktionen mere som Bevismiddel
end som Opdagelsesmethode.

Begreb og Kjendsgjerning staa i et Vexelforhold til
hinanden. Begrebet kan aldrig være uafhængigt af Kjendsgjerningen,
men paa den anden Side maa Kjendsgjerningen
bestandig drages hen imod Begrebet. Erkjendelsen forudsætter
en Kombination af Tanke og Gjenstand. Uden
Tanken vilde der ingen Forbindelse være mellem Gjenstandene,
uden Gjenstand vilde der ingen Realitet være. Enhver
enkelt Erkjendelsesakt indeslutter begge Elementer. Forskjellen
mellem dem kan kun være relativ. Hvad der fra
ét Synspunkt er Theori, er fra et andet Synspunkt Faktum.
Hvad der i Forhold til en forudgaaende Undersøgelse er
Theori, er i Forhold til en følgende Kjendsgjerning. At
Planeterne bevæge sig i Ellipser, var hos Kepler en Theori,
men bliver hos Newton et Faktum. En sand Theori er
et Faktum; et Faktum er en Theori, med hvilken man er
fortrolig. Begge ere kun Slutninger, vi drage fra Phænomenerne;
hvad vi kalde Theori, er en bevidst Slutning;
hvad vi kalde Faktum er en ubevidst Slutning.

Lad os i denne Henseende betragte selve Sandsningen. Et sammenhængende Billede af en sandselig 
Gjen%
% --- p. 93 ---
\setcounter{footnote}{0}%
\opage{93}stand danner sig først, naar de enkelte Fornemmelser
kombineres og sammenfattes til Enhed. Der ligger en
ubevidst Slutning til Grund for alle vore Sandsninger.
Begreber som Rum, Afstand, Soliditet o. s. v. udvikle sig
først ved Fornemmelsen af vore Muskelbevægelser, og vi
overføre dem saa uvilkaarlig paa Gjenstandene. Vi se en
Bevægelse foregaa og slutte os til Kraften som dens Grund.
Vi se forskjellige Planter eller Dyr, men kun fordi vi anvende
vore Forestillinger om Lighed og Ulighed paa dem. Ligheden
eller Uligheden se vi lige saa lidt, som vi se Kraften
eller Rummet; eller rettere, vi se dem kun, for saa vidt
vi lægge dem ind i Phænomenerne. Heraf ses, hvor meningsløst
det er at betegne Ideerne eller Begreberne som
blotte omdannede Sandsefornemmelser. Selv den mest elementære
Fornemmelse kan ikke komme i Stand uden et
Bidrag fra det fornemmende Subjekt. Vor Naturopfattelse
er fra først til sidst en Fortolkning af Naturen, en Oversættelse
af dens Indhold paa vort eget Sprog.

Subjektets Bidrag til Phænomenernes Erkjendelse eller
Fortolkning bestaaer i en Række af Ideer eller Grundbegreber,
som bringe Sammenhæng og Enhed i Erfaringsindholdet.
Ikke som om disse Ideer skulde ligge fuldfærdige
i Aanden fra først af. De maa ikke sondres fra selve den
opfattende Virksomhed, men ere netop denne Virksomheds
Former og Love; kun ved Reflexion over hin Virksomhed
lære vi disse Former og Love at kjende, og kun ved Erfaringen sættes hin Virksomhed i Bevægelse. Bevidstheden
vilde altsa ingen Ideer og Grundbegreber have, hvis den % PRINTED AS IS: „altsa“
ikke stod i Vexelvirkning med Rummets, Tidens og Materiens Verden; men vi kjende jo heller slet ingen Bevidsthed,
som ikke staaer i et saadant Vexelforhold. Vi kunne
ikke tænke os noget Trin af menneskelig Existents, hvor
hin Virksomhed og dens Grundlove ikke allerede træde i
Kraft og udvikles ved Erfaringen. Formerne og Ideerne
virke længe før vi blive os selv bevidste, og udvikles og
modificeres paa mange Maader ved Begivenhedernes og Phænomenernes Indflydelse. Hvor vidt denne Modifikation strækker
% --- p. 94 ---
\setcounter{footnote}{0}%
\opage{94}sig, kunne vi ikke bedømme; men den er sikkert af væsentlig
Betydning for den Proces, der gaaer for sig i vor
Bevidsthed.

Grundbegreberne erhverves i Løbet af Erfaringen, men
de læres ikke af Erfaringen i samme Betydning som specielle
Kjendsgjerninger læres. De enkelte Kjendsgjerninger ere
for os mere eller mindre tilfældige; en Tidlang ere de ubekjendte,
til en anden Tid bekjendte; Nogle kjende dem,
Andre ikke; de kjendes med større eller mindre Vished.
Men Grundbegreberne ere de Love, efter hvilke al Opfattelse,
al Erfaring foregaaer, Det er denne deres almindelige
og nødvendige Gyldighed, hvorpaa deres Forskjel
fra Erfaringsbegreberne beroer. Vi kunne lære almindelige
Sandheder, f. Ex. Tyngdeloven, af Erfaringen, Men hvis
disse Sandheder virkelig have universel Gyldighed, saa
skyldes denne de Ideer, som Forstanden anvender paa det
Givne.

De to Elementer, som gjennemtrænge hinanden i vor
Erkjendelse, betinge saaledes en Grundmodsætning, der
aabenbarer sig under en Række af Former: Ide og Faktum,
Begreb og Gjenstand, Theori og Kjendsgjerning, Slutning
og Iagttagelse, men som har sit skarpeste Udtryk i Modsætningen
mellem nødvendig (apriorisk) og experimental
Sandhed. Hin udspringer af Tanken, denne af Iagttagelsen.
Vi kunne ikke opfatte Grundbegreberne i sig selv; de finde
deres Udtryk i visse Grundsætninger og Begrebsbestemmelser,
som Tanken kommer til, naar den reflekterer over sin egen
Virken. Hver enkelt Videnskab har sine oprindelige Grundsætninger
eller Axiomer, som udspringe af det Grundbegreb,
den Ide, hvorpaa den hviler. Saaledes udspringe Mathematikens
Axiomer af Rummets og Tidens Ideer, som ligge
til Grund for al Opfattelse af det Givne, Mekanikens
Axiomer udspringe af Aarsagens Ide, der ligeledes er en
oprindelig Lov for Tankens Virken. Af det almindelige
Axiom, at enhver Ting har sin Aarsag, følge de mekaniske
Axiomer, at Aarsager maales ved deres Virkninger, og at
Tilbagevirkningen er lige stor med Indvirkningen. Bevæ%
% --- p. 95 ---
\setcounter{footnote}{0}%
\opage{95}gelsens Love hente deres Form fra Aarsagens Ide, deres
Stof fra Erfaringen. Axiomerne ere naturligvis abstrakte
og maa fortolkes ved Studium af den virkelige Natur; men
de udgjøre den faste Kjerne, om hvilken Erkjendelsen af
denne kan udvikle sig. At de anførte Ideer ere nødvendige,
vil sige, at det er os umuligt at tænke det Modsatte af
dem, og at enhver Indvending mod deres Gyldighed netop
forudsætter denne. Vi kunne tage feil ved Forklaringen i
det Enkelte; men at der maa være en Forklaring, en
Aarsag, det er en Forudsætning, vi ikke kunne Andet end
gaa ud fra. Vi kunne f. Ex. ikke tænke os en Sten svævende
i Luften i Stedet for at falde til Jorden, uden at
vi maa forudsætte en Aarsag til denne Afvigelse fra den
almindelige Lov.

Ligesom Rummets og Tidens Ideer ligge til Grund
for Mathematiken og Aarsagens Ide for Mekaniken, saaledes
er Hensigtsmæssighedens eller Formaalets Ide nødvendig
for at opfatte de organiske Phænomener. Heller ikke denne
Ide kan udledes af Erfaringen; den udvikler sig nødvendig
som den eneste Form, under hvilken vi kunne opfatte
Livets Phænomener. Harveys og Cuviers Opdagelser bleve
kun mulige ved den. Vel har man ment, at efterhaanden
som den videnskabelige Erkjendelse rykker fremad, indskrænkes
Hensigtsaarsagens Omraade og de mekanisk virkende
Aarsager træde i Stedet. Men dette er en Illusion.
Formaalets Ide maa nu opfattes paa en anden Maade end
før, men kan ikke opgives. Vi se ikke Hensigtsmæssigheden
i enkelte særlige Phænomener, men i den hele
Sammenhæng. Vi antage ikke, at der sendes en Sky, for
at det kan regne paa vor Mark, men vi finde Formaalets
Ide udtrykt i Sammenhængen mellem Varmens, Luftens og
Vædskernes Love og det vegetabile Livs Udvikling. Formaalets
Ide faaer derved netop en mere storartet Karakter:
i Stedet for nogle faa Traade, som forbinde særskilte
Gjenstande, faa vi en total, hele Virkeligheden omfattende
Sammenhæng.

% --- p. 96 ---
\setcounter{footnote}{0}%
\opage{96}Ogsaa paa Ethikens Omraade\footnote[1]{\emph{Whewell}: History of Moral philosophy in England. 1852.} have vi en oprindelig
virkende Ide, der lægger sig for Dagen i det Axiom, at vi
skulle gjøre, hvad der er ret. Selv om man er uenig om,
\emph{hvad} der er ret, hersker der dog ingen Strid om, \emph{at} det
Rette skal gjøres. Denne Spire til det moralske Liv udvikles
nu videre ved Erfaringen paa lignende Maade, som
paa det theoretiske Omraade Grundbegrebet udvikles
gjennem Iagttagelserne. Menneskenes Drifter føre dem
sammen og gjøre faste Love for deres Samliv nødvendige;
herved opdrages den moralske Bevidsthed til efterhaanden
ogsaa at anvende Lovens Begreb paa sit eget Liv. Den
høieste Lov for et Menneskes Handlinger maa være en
saadan, der omfatter hele hans Væsen, Tanke og Drift
saavel som Handling. Fra den ydre, legale Pligtopfyldelse
føres Mennesket derfor til den indre, ethiske, og vinder
efterhaanden en høiere Maalestok, efter hvilken han kan
bedømme Samfundets Love. --- I Modsætning til denne
ethiske Opfattelse, der gaaer ud fra et indre Princip som
Maalestok og Forudsætning for Bedømmelsen af alt Ydre,
den uafhængige Ethik, Principernes Ethik (the morality of
principles), staaer den ethiske Lære, som gaaer ud fra
Handlingens Følger (the morality of consequences), fra
Nytten (Utilitarianismen), hvis vigtigste Repræsentanter ere
Epikur i Oldtiden, Hobbes, Paley og Bentham i nyere Tid.
Denne ethiske Empirisme overser ligesom den theoretiske
Empirisme de oprindelige Forudsætninger. ---

Man vil i den her fremstillede Anskuelse let gjenfinde
Kants Grundtanker, men man vil gjenfinde dem i en noget
vag og ubestemt Form. Whewells Fortjeneste bestaaer i,
at han har gjennemført Videnskabernes Historie med særligt
Hensyn til Methoden og Erkjendelsesproblemet, og at han
har hævdet Tankens Oprindelighed og Aktivitet i Erkendelsesprocessen
overfor den empiriske Skole. Men han
giver ikke det ved den historiske og kritiske Prøvelse
% --- p. 97 ---
\setcounter{footnote}{0}%
\opage{97}Fundne en dyberegaaende philosophisk Behandling. Han
lader sig for det Meste nøie med for hvert enkelt videnskabeligt
Omraades Vedkommende at slutte tilbage til en
apriorisk Idé, men undersøger ikke, hvorvidt disse Ideer
ikke skulde kunne reduceres til én eller dog nogle faa, ei
heller, hvorledes de forholde sig til Bevidsthedens Væsen,
eller hvorledes de have udviklet sig. Denne Undersøgelse
er saa meget des mere nødvendig, som ikke blot Ideerne
eller Grundbegreberne træde i en Modsætning til det Empiriske,
som ikke kan deduceres af dem, men ogsaa i en
indbyrdes Modsætning til hverandre, Aarsagens (ɔ: den mekaniske
Aarsags) Begreb f. Ex. til Formaalets. Flere af
»Ideerne» optager han endog saa godt som i ganske populær
Form, f. Ex. Formaalets Idé. Han citerer her Cuviers Ord:
»Zoologien har et for den eiendommeligt Princip, som den
med stor Fordel anvender ved flere Leiligheder, og det er
Existentsbetingelsernes Princip, som man sædvanligvis ogsaa
kalder Hensigtsaarsagernes Princip.» Men ser man nærmere
til, føre disse Ord netop ud over den populære Forestilling
om et Formaal. Thi hvad Cuvier mener, er egentlig ikke
Andet, end at et levende Væsens Bygning og Organer maa
være saaledes beskafne, at dets Existents er mulig; derfor
kan Analysen af Existentsbetingelserne (conditions de l'existence)
ofte føre til Love, som stemme med Erfaringen.
Men om en forudgaaende Afpassen af Individet efter Betingelserne
behøver der derfor ikke at være Tale; Overensstemmelsen
kan ganske simpelt have sin Grund i de
levende Væseners Vexelvirkning med de ydre Betingelser.

Whewells Mening er, som det ses af Fremstillingen,
ikke, at hine Grundbegreber fra først af ligge færdige i
Bevidstheden. Han antyder, som ligeledes anført, et enkelt
Sted, at de Former, under hvilke vi opfatte Tingene, allerede
danne sig, før vi blive os selv bevidste, saa at vi
nu kun kjende dem i den Skikkelse, hvortil det ubevidste
Livs Begivenheder har bragt dem. Denne Antydning viser
i en rigtig Retning, ligesom Whewell jo i det Hele hævder,
at Ideerne udspringe af Aandens Aktivitet, ikke have nogen
% --- p. 98 ---
\setcounter{footnote}{0}%
\opage{98}selvstændig Væren. Men des mere paafaldende er det, at
han kan lade sig nøie med en blot og bar Opregnen af
dem uden nærmere Kritik og Analyse. I denne Henseende
har hans Modstander Stuart Mill et stort Fortrin
for ham.

Medens det ikke kan negtes, at Whewell med Hensyn
til den almindelige Opfattelse af Erkjendelsens Væsen har
set rigtigere end Mill, kan han dog ikke maale sig med
denne i den skarpere Angivelse af den induktive Forskens
Gang. Mill fastholder det exakte Synspunkt, fordrer det
strenge logiske Bevis. Derved ledes han til at formulere
de specielle induktive Methoder. Whewell fremhæver den
aandelige Produktivitet, som lægger sig for Dagen i Opdagelserne.
Rent logisk set har Mill Ret, naar han siger,
at vi ved Whewells Bestemmelse af Induktionen ikke komme
ud over Hypothesen, fordi den exakte Verifikation mangler.
Metaphysisk og historisk har Whewell, logisk har Mill
Retten paa sin Side. ---

Whewell iagttager ikke de Grændser, som Kant afstak
for Grundbegrebernes (Kategoriernes) Anvendelse. De Ideer,
vi nødvendigvis maa anvende ved Betragtningen af Naturen
og Menneskelivet, ere ifølge hans Lære de eneste naturlige
Midler, ved hvilke vi kunne danne os et Begreb om
Verdens Skaber og Styrer. Lige saa vel som vi slutte, at
enhver enkelt Ting maa have sin Aarsag, lige saa vel maa
vi slutte, at hele Verden maa have en Aarsag. At der
maa være en første Aarsag, ansér han endog for et Axiom.
Vi kunne ikke standse ved en uendelig Række af endelige
Aarsager; vi maa statuere en Begyndelse; for at Phænomenerne
kunne blive forstaaelige for Tanken maa denne
hæve sig over den endeløse Succession. Og de moralske
Ideer føre os til at tænke hin første Aarsag som aandelig
Forsynsmagt. --- Whewell har forfattet en af de saakaldte
Bridgewater-Afhandlinger,\footnote[1]{Saa kaldte efter en 1829 afdød Greve af Bridgewater, der udsatte en Sum til saadanne Skrifters Udgivelse.} som gik ud paa »ad natur%
% --- p. 99 ---
\setcounter{footnote}{0}%
\opage{99}videnskabelig Vei» at bevise Skaberens Visdom og Godhed
som aabenbaret i Verden. Hans ukritiske Sammenblanding
af Tænkning og Erfaring og deraf følgende Uklarhed med
Hensyn til Erkjendelsens rationelle Karakter lægger sig
navnlig for Dagen i hans theologiske Forestillinger og deres
Forbindelse med de philosophiske Begreber. Lige saa let
som det falder ham at gaa til Forestillingen om en første
Aarsag, lige saa let falder det ham at forbinde Bibelens
og Videnskabens Lære. De moralske Ideer skulle ikke blot
føre os til Troen paa et Forsyn; de skulle ogsaa begrunde
den Konsekvents, at naar det sædelige Ideal ei er fyldestgjort,
maa der være truffet overnaturlige Foranstaltninger
for at virkeliggjøre det. Til disse Foranstaltninger høre
navnlig Bibelen selv, hvis Fortælling om Verdens Skabelse
Videnskaben, der ikke naaer op til Tidens Begyndelse,
hverken kan bekræfte eller forkaste. --- Ogsaa her finde vi
en Ubestemthed i Begreberne, en Mangel paa dybere gaaende
Kritik og paa Bevidsthed om de afgjørende Forskjelle, som forringer Whewells Fortjenester af Erkjendelsesproblemet.


II. \emph{William Hamilton} (f. 1788, d. 1856). ---
Dette Navn fører os til Skotland, hvorfra den første Reaktion
i Brittanien selv mod Empirismen udgik i Slutningen
af forrige Aarhundrede. Thomas Reids Lære har siden
den Tid været herskende ved de skotske Universiteter.
Den synes at stemme med Folkekarakteren. Til sin Baggrund
har den en fast og alvorlig moralsk og religiøs
Overbevisning, saaledes som man kan vente det i John
Knox's Fædreland; den holder sig besindig og rolig indenfor
den sædvanlige Bevidstheds Grændser og søger blot at
fremstille denne Bevidstheds »Kjendsgjerninger;» men den
lægger dog helt igjennem Vægten paa Erfaringen og vil
anvende Bacos og Newtons Principer, kun at den ikke med
Locke og Condillac indskrænker Erfaringen til det Sandselige.
Den søger at finde Bevidsthedens primitive og nød%
% --- p. 100 ---
\setcounter{footnote}{0}%
\opage{100}vendige Udsagn, der ligge til Grund for al anden Erfaring.
Den skotske Philosophi er væsentlig empirisk Psykologi,
kun at den søger at paavise Bevidsthedskjendsgjerninger,
som den sædvanlige Erfaring efter dens Mening overser.
Den staaer paa Empirismens Grund, skjøndt den vil bestride
dens Konsekventser. Den mener umiddelbart at
kunne finde primitive Fakta og standser ved en Række af
Principer, som den anser for de sidste, der kan naaes.
Derved bliver den en Slags »systematisk Lettroenhed» og
synder mod den sande empiriske Methode, der først anerkjender
et Faktum, naar videnskabelig Undersøgelse godtgjør
det, lige saa vel som mod den philosophiske Analyses
Princip, ifølge hvilket et Faktum ikke ligefrem kan forklare
et andet, men begge maa forklares af en høiere Lov.

Den skotske Skole viser saaledes i to Retninger.
Hovedvægten kan lægges paa den empiriske Tendents, og
vi komme da ved Skærpelse af den videnskabelige Methode
til den moderne engelske Psykologis Standpunkt. Paa
denne Vei slog allerede Thomas Brown (d. 1820) ind, lige
som Alexander Bain herigjennem hænger sammen med sine
skotske Forgængere (se ovenfor p. 66). Men man kunde
ogsaa fremdrage det kritisk-spekulative Element, og derved
nærme sig til Kant. Ad denne Vei har Hamilton
bragt den skotske Skole ud over dens vage og uvidenskabelige
Holdning, Han medbragte til denne Opgave et dybt
og alsidigt Kjenskab til den græske og den tydske Philosophi,
overhovedet en forbausende og i England sjælden
Lærdom i Philosophiens Historie, hvilket ikke lidet bidrager
til den store Interesse og Belæring, hans Skrifter yde. Men
han var desuden en skarp Logiker og Dialektiker, en Forsker,
der fandt den philosophiske Tænknings Maal og Udbytte i
selve Tankens Aktivitet og i den Bevægelse, de høieste
aandelige Kræfter derved sættes i. Han har ikke skrevet
Meget, men virkede i mange Aar som philosophisk Professor
i Edingburg. Selv udgav han kun nogle Afhandlinger, som % PRINTED AS IS (Edingburg: O-E14)
han føiede til sin Udgave af Reid, og en Samling af sine
Artikler i Edingburgh Review under Titel af »Discussions % PRINTED AS IS (Edingburgh: O-E15)
% --- p. 101 ---
\setcounter{footnote}{0}%
\opage{101}on philosophy» (1852); efter hans Død ere hans Forelæsninger
over Metaphysik og Logik udgivne af Mansel og
Veitch i fire Bind.

Det første Arbeide, som gjorde Hamilton berømt, var
en Afhandling i »Edingburgh Review» 1829, rettet imod
Schellings og Cousins Lære om det Absolute. Det er et
Forsøg paa at vise, at der ikke kan gives nogen Philosophi
om det Ubetingede. Vi kunne kun erkjende det Betingede
og det Begrændsede. Det Uendelige ɔ: det som ubetinget
er uden Grændse, og det Absolute ɔ: det som er ubetinget
afsluttet i sig selv, kunne vi ikke fatte, ikke konstruere i
vor Tanke. Vor Tanke er fra først til sidst betinget, ja,
det at tænke er at gjøre betinget (to think is to condition).
Tanken kan ei overskride Bevidstheden, lige saa lidt som
Ørnen kan flyve ud over Atmosphærens Grændser; og Bevidstheden
har sine bestemte Betingelser, som ethvert Bevidsthedsindhold
er underlagt. Vi opfatte Existentsen i Tid
og Rum, ikke i og for sig. Men Tid og Rum forudsætte
Forhold mellem forskjellige existerende Elementer; det i Tid
og Rum Existerende er derfor relativt, og Tid og Rum
kunne ikke anvendes paa det Ubetingede. Bevidstheden
forudsætter fremdeles en Modsætning mellem et Erkjendende
og et Erkjendt, mellem Subjekt og Objekt; men denne Modsætning
kan ikke anvendes paa det Ubetingede, som derved
vilde blive relativt. Den eneste Udvei vilde her være med
Schelling at antage, at det Ubetingede er hævet over Modsætningen
mellem Subjekt og Objekt, er begges Indifferents,
og at vi for at erkjende denne Indifferents maa hæve os
over Bevidstheden og Reflexionen, der stedse forudsætte
Flerhed og Relativitet. Men hvorledes kunne vi hæve os
over vor egen Bevidsthed? Og med hvad Ret kunne vi
statuere en Evne, der strider mod selve vor Bevidstheds
Love?

Man ser, hvorledes Kants Lære optages af Hamilton
og spiller en Hovedrolle i hans Anskuelse. Dog skiller han
sig fra Kant i et afgjørende Punkt. Kant havde vel indskrænket
Anvendeligheden af vore Anskuelses- og Tanke%
% --- p. 102 ---
\setcounter{footnote}{0}%
\opage{102}former til det Phænomenale og Relative, men han erkjendte
dog det Ubetingede som et Fornuftprincip, der vel ikke
kunde lægges til Grund som konstruktivt Princip, men som
har regulativ Betydning, tjener til at føre ud over det
relativt Givne og stille ideale Opgaver. Heri finder Hamilton
Spiren til den tydske Spekulations paa det Ubetingedes
Idé byggede Systemer. Grundfeilen er efter hans
Mening, at man antager det Ubetingede for et positivt
Begreb. Kun det Betingede kunde vi positivt erkjende og
begribe, omfatte og fremstille med vor Tanke; det Ubetingede,
Negationen af det Betingede, er uerkjendeligt og
ubegribeligt, er er rent negativt Begreb, ligesom det Ufattelige. % PRINTED AS IS (er er: O-E16: doubled)
Vilde man nu med Kant sige, at vi alligevel besidde
det Ubetingede som positivt, regulativt Princip, saa vilde
der være indrømmet en Selvmodsigelse i vore høieste aandelige
Evner, som nødvendigvis vilde føre os ud i Skepticisme.
Løsningen er derimod ganske simpel, naar man
fastholder, at vi kun kunne opfatte det Betingede, og at
det Ubetingede for os er en ren Negation Det Betingede
ligger stedse mellem to Extremer, som have Karakteren af
det Ubetingede; vi kunne ikke fatte noget af dem, og dog
maa et af dem være det Sande. Verden maa f. Ex. engang
have en Begyndelse eller være evig; men vi kunne
lige saa lidt fatte en absolut Begyndelse som en uendelig
Evighed. Vi lære heraf, at vi ikke kunne gjøre Tanken
til Maalestok for Existentsen. Vi kjende kun betinget
Existents, ei Existentsen i og for sig. Hvad der ligger ud
over det Betingede, er Gjenstand for Tro. Der maa være
en Existents i og for sig; ellers kunde der ingen betinget
Existents være. »Og ved en vidunderlig Aabenbaring indgydes
der os saaledes ved selve Bevidstheden om vor
Mangel paa Evne til at fatte Noget, der ligger ud over det
Relative og Endelige, en Tro paa noget Ubetinget hinsides
al den Virkelighed, vi kunne fatte.»

Dersom Philosophien skulde have det Ubetingede til
Gjenstand, vilde den være umulig. En absolut Sandhed
er os utilgængelig. Men paa en saadan beroer Philosophiens
% --- p. 103 ---
\setcounter{footnote}{0}%
\opage{103}Værd og Betydning heller ikke. Det er ikke færdige Resultater,
den gaaer ud paa. En fuldstændig Løsning at
Problemerne, som standsede al videre Discussion, vilde være
en Lammelse af den philosophiske Stræben og føre os tilbage
til et intellektuelt Barbari. Vort Væsen og vor Lykke
bestaaer i Virksomhed, og Erkjendelsen er kun af Værdi,
for saa vidt den er Virken og Stræben, ikke død Besiddelse.
Spekulationens Betydning og Nødvendighed maa hovedsagelig
søges i den Øvelse af vore høieste aandelige Kræfter, den
bevirker. Lessing har Ret: for os er den evige Stræben
efter Sandhed bedre end Besiddelsen af Sandheden. En
vaagen Vildfarelse er bedre end en sovende Sandhed. Det
Høieste, Philosophien kommer til, er Anerkjendelsen af vor
Uvidenhed med Hensyn til Existentsen i og for sig selv.
Men der er to Slags Uvidenhed. Vi philosophere for at
undgaa den blinde, tankeløse Uvidenhed, der ikke kjender
Bevidsthedens og Sandhedens Betingelser, og naar vi have
erkjendt disse Betingelser, ende vi i den høiere, den bevidste
Uvidenhed (docta ignorantia).

Hamilton statuerer altsaa en Modsætning mellem Viden
og Tro. Kun denne sidste naaer det Ubetingede. Men
hvilke Motiver har den da til at søge ud over det Endelige
og hævde det Ubetingedes positive Realitet i Stedet for at
lade det staa hen som en ren Mulighed, som en »Ding
an sich?» Disse Motiver ligge i Menneskets Væsen som
moralsk Personlighed. Mennesket er sig en ubetinget
Pligtlov bevidst, og den moralske Ansvarlighed forudsætter
igjen Villiens Frihed. Menneskets Væsen har sit eiendommelige
Udtryk i Kampen mellem den moralske Frihed og
den physiske Nødvendighed. Og kun som frit moralsk
Væsen kan han tro paa en vis og hellig Verdensstyrer.
Vistnok er en ubetinget Frihed ubegribelig, men vi vide jo,
at Tankens Grændser ikke ere Mulighedens og Virkelighedens
Grændser. Ved Anerkjendelsen heraf baner Philosophien
Vei for Anerkjendelsen af Ansvarlighed og Frihed,
som ogsaa den christelige Aabenbaring forudsætter. Philosophiens
Ende er Theologiens Begyndelse.

% --- p. 104 ---
\setcounter{footnote}{0}%
\opage{104}Vi have her naaet Kjernepunktet i Hamiltons Livsanskuelse.
Erkjendelsens Relativitet giver her Plads for
den moralske og religiøse Tro paa det Ubetingede. Menneskets
aandelige Væsen, saaledes som det godtgjøres ved
den moralske Ansvarlighed og Frihed, var for ham det
Høieste, lige som for Kant. Hans Valgsprog var: »Paa
Jorden er der intet Stort uden Mennesket, og i Mennesket
er der intet Stort uden Aanden.» Stuart Mill opponerer
mod dette Valgsprog, som har en afgjort spiritualistisk
Tendents. I dette Kjernepunkt fremtræder Modsætningen
mellem de to Philosopher, Repræsentanterne for de to stridende
Skoler, tydelig. Mill fastholder lige som Hamilton,
at vi kun erkjende det Relative og Betingede, men mener
med Rette, at vi saa ogsaa maa opfatte Menneskets Væsen
som betinget, at den Methode, der fører os til Indsigt i
Naturens Love, ogsaa maa lære os Menneskelivets Love
at kjende, eftersom Menneskelivet selv er et Led i Naturlivet.
Mills Styrke ligger netop paa dette Punkt; hans
Opfattelse af Aarsagsbegrebet hjælper ham ud over Friheds- og Nødvendighedslærens Abstraktioner. Han har tillige
klart set, at Hamiltons Philosophi bestemmes i sin hele
Eiendommelighed ved Ideen om den ubetingede Frihed.
For Hamilton maatte Mills Lære staa som Fatalisme, og
Fatalismen førte efter hans Overbevisning nødvendigvis til
Materialisme og Atheisme. De Videnskaber, som særlig
hvile paa Nødvendighedens Begreb, Mathematik og Physik,
var han derfor ikke fri for at betragte med mistænkelige
Øine. Denne hans Fjernhed fra Naturvidenskaberne stiller
ham i Modsætning til hele den nyeste Philosophis Stræben,
som netop gaaer i Retning af en nøiere Tilslutning til
denne.

Mill var, som vi have set, heldigere med de Projektiler,
han kastede ind i Fæstningen selv, end med sine Angreb
paa Udenværkerne. Hamiltons Fortjeneste beroer paa den
Klarhed og Originalitet, hvormed han har opfattet Kants
Lære om Erkjendelsens Relativitet, og Mansel har paavist,
at Mill, hvad Hovedpunkterne i hans Angreb paa denne
% --- p. 105 ---
\setcounter{footnote}{0}%
\opage{105}Lære angaaer, ikke har forstaaet Hamilton. Desuden kunde
Mill ifølge sin empiristiske Retning ikke angribe Hamilton
paa det egentlig afgjørende Punkt. Dette Punkt finde vi i
det Forhold mellem det Ubetingede og det Betingede,
hvorfra Hamilton gaaer ud. Det Ubetingede bestemmer
han som Negationen af det Betingede. Hvis man gaaer
ud fra denne Bestemmelse, behøver man ikke først at paavise
Umuligheden af at opfatte det Ubetingede; thi saa er
dette Begreb fra første Færd selvmodsigende. Dersom
nemlig det Ubetingede har det Betingede udenfor sig som
sin Modsætning, saa bliver det derved selv betinget og relativt,
ophæves altsaa i samme Øieblik det sættes. Men
hvad Ret har man til at gaa ud fra denne Forudsætning
om Modsætningen mellem det Ubetingede og det Betingede?
Det er en Forudsætning, som ganske vist ligger i den almindelige
Bevidsthed og indeholdes i de overleverede religiøse
Forestillinger. Men den philosophiske Tankes første
Skridt maa være at hæve sig over denne Dualisme.\footnote[1]{Den optræder i Tænkningens Historie i en dobbelt Form. For Grækerne var det Endelige og Begrændsede det Fuldkomne, det Uendelige det Lavere og i sidste Instants det Onde. Endnu Origenes opfattede derfor Gud som begrændset. I Middelalderen (kfr. Augustinus og Anselmus) vendtes Forholdet om; Gud blev den Uendelige, evig hævet over Endelighedens Verden. Strax ved den nyere Tids Begyndelse søgte derimod dybsindige Tænkere (Cusa, Bruno, Spinoza) den indre Enhed af det Uendelige og det Endelige. Kfr. \emph{Weisse}: Das philosophische Problem der Gegenwart p. 195 f. R. \emph{Nielsen}: Mathematik og Dialektik. p. 117 f. --- Hamilton har med sit Kjendskab til Philosophiens Historie klart set dette; han finder netop allerede hos Cusa den nyere spekulative Philosophis Ide om det Absolute, som han polemiserer imod.}

Det sande Ubetingede maa være det, der faaer sin
nærmere Bestemmelse fra det Betingede, --- det, i hvis
Væsen alle Betingelser have deres sidste Grund; det sande
Absolute maa staa i et indre Forhold til det Relative. Vel er
derfor en umiddelbar og absolut Erkjendelse af det Ubetingede
umulig; for saa vidt har Hamilton Ret. Men vi kunne
gjennem Erkjendelsen af det Betingede vinde afgjørende
% --- p. 106 ---
\setcounter{footnote}{0}%
\opage{106}Væsensbestemmelser af det Ubetingede, i Kraft af det indre
Forhold imellem dem. Her ligger en Række af Problemer,
som den engelske Philosophi kun lidt beskjæftiget sig med,
men som vi i de store Systemer paa Fastlandet, navnlig
hos Spinoza og Hegel, finde epokegjørende Bearbeidelser af,
Problemer, der uundgaaelig trænge sig frem igjen, hvorofte
de end skydes til Side

Hamilton kan selv ikke gjennemføre Modsætningen
mellem Existentsen i sig selv og Phænomenerne, det Ubetingede
og det Betingede. Han udtaler (Lectures I p. 125),
at vi fra Phænomenernes eiendommelige Beskaffenhed kunne
drage mange gyldige Slutninger angaaende det ubekjendte
Princip, hvis Aabenbaringer de ere. Skjøndt derfor Existents
kun er aabenbaret for os i Phænomenerne, og vi kun
kunne have en relativ Erkjendelse af Aand og Materie, saa
kunne vi dog ved Slutning og Analogi hæve os over de
blotte Tilsyneladelser, som Erfaring og Iagttagelse vise
os. Men hvorledes kunde dette være muligt, naar der
ikke var et indre Forhold mellem Existentsen i sig selv og
Phænomenerne?

Ifølge Hamilton er Philosophiens Ende Theologiens
Begyndelse. Men Spørgsmaalet bliver nu, med hvilke
Midler Theologien skal bygge sin Lære. Se vi tilbage til
Theologiens Historie, finde vi, at den stedse har benyttet
philosophiske Begreber eller i det Mindste Begrebsformer.
Men naar nu disse Begreber, naar al vor Tanke og Erkjendelse
kun angaaer det Betingede, hvorledes bliver saa
en theologisk Lære om det Ubetingede mulig? Hamilton
siger selv, at ingen Vanskelighed opstaaer i Theologien, som
ikke allerede viser sig i Philosophien. Men saa maa man
ogsaa omvendt sige, at de Vanskeligheder, der reise sig i
Philosophien, ogsaa maa gjælde for Theologien. Det er da
inkonsekvent, naar Hamilton et Sted udleder den engelske
Theologis Ufuldkommenhed og Ubetydelighed af dens uphilosophiske
Karakter. Konsekvent er det derimod, naar han
indskærper, at Bibelens Udsagn skulle gjælde uden Omtydning,
at Theologens Hovedhverv er Bibelfortolkning, og
% --- p. 107 ---
\setcounter{footnote}{0}%
\opage{107}philologisk Lærdom derfor den eneste nødvendige Betingelse
for Theologien. Konsekvent er det ogsaa, naar hans
begavede Discipel \emph{Henry Mansel} af hans Principer har
udledet Umuligheden af en videnskabelig Theologi.

For det saakaldte Bampton-Legat holdes der i Oxford
aarlig Forelæsninger til Forsvar for de christelige Lærdomme.
Saadanne »Bampton-lectures» har ogsaa Mansel holdt, og
han har her »anvendt sin korte Optræden paa et theologisk
Katheder til at bevise, at der ingen Theologi gives.»\footnote[1]{\emph{Rémusat}: La théologie naturelle en France et en Angleterre. p. 75.}
Han erklærer sig først saa vel imod Dogmatismen som imod
Rationalismen. Dogmatismen vil deducere videnskabelig
udformede Dogmer af den religiøse Tro; Rationalismen vil
reducere de religiøse Forestillinger til symbolske Udtryk for
abstrakte Sandheder. Begge gaa ud paa at naa en høiere
Erkjendelse af de guddommelige Ting. Man mener, at naar
man ophæver det Billedlige og Anthropomorphistiske, vil
man naa en sandere Erkjendelse af Gud. Men de abstrakte
Former, der sættes i Stedet for de saakaldte billedlige, ere
betingede af vor Bevidstheds og vor Erkjendelses Væsen,
som er endeligt. Vi formaa ikke at fatte det Ubetingede
som personligt Væsen, hvad vi dog af moralske og religiøse
Grunde maa gjøre. Vor Personlighed er endelig og betinget
og kan derfor ei være Form for Opfattelsen af det
Ubetingede. Ethvert Forsøg paa at udvikle et rationelt
Gudsbegreb, fører derfor, som Philosophiens Historie viser,
til Pantheisme, der er uforenelig med den moralske og
religiøse Tro. Derfor kan der ingen videnskabelig Theologi
gives. Men da Tanken ikke udgjør hele Menneskets Væsen,
da han som moralsk og religiøs Personlighed finder sig i
et Afhængigheds- og Forpligtelsesforhold til et høieste
personligt Væsen, Gjenstanden for hans Dyrkelse og Kilden
til Loven for hans Handlen,\footnote[2]{„At anvende Guds Navn i abstrakt Spekulation blot som en vedtagen Benævnelse for det høieste Led i Tankens Kjede, og \opage{108}at tro paa ham af praktiske Grunde i Ærefrygt og Lydighed, er to meget forskjellige Ting; og kun til det Sidste, ikke til det Første er Begrebet om Gud som Person nødvendigt. Dersom Mennesket var et rent Forstandsvæsen, vilde Problemet om det Ubetingede miste sin Hovedvanskelighed; men han er ogsaa et Væsen med religiøse og moralske Evner, og ogsaa disse fordre at tilfredsstilles ved en gyldig Løsning af Problemet.“ \emph{Mansel}: The philos. of the cond. p. 13.} saa maa de religiøse 
% the note runs over from p. 107 to p. 108 (joined by draft.py)
Lær%
% --- p. 108 ---
\setcounter{footnote}{0}%
\opage{108}domme vende sig mod alle Sider af Mennesket, appellere
til hans Følelse og Villie, vel anvende Analogier og Sandsynlighedsslutninger
fra det Endelige, men ikke gruude deres % PRINTED AS IS: „gruude“
Gyldighed udelukkende paa Tankens Anerkjendelse. Vi
maa altsaa skarpt adskille Viden og Tro. Vi kunne ikke
begribe, hvorledes det Ubetingede og det Personlige ere
forenede i det guddommelige Væsen; men vi maa tro, at
de paa en eller anden Maade ere forenede deri. Vor Tanke
kan derfor ikke tages til Maalestok for den religiøse Tro;
vi kunne ikke danne os noget videnskabeligt Gudsbegreb,
efter hvilket vi skulde kunne bedømme de religiøse Lærdomme.
--- Dette er Mansels Lære om »den religiøse
Tankes Grændser,» som har vakt saa megen Opsigt i England
i de senere Aar.

Det er et Tidernes Tegn, at der fra saa mange forskjellige
Sider i vore Dage rettes Angreb paa Theologien.
Det nu skildrede, som udgik fra selve den engelske Orthodoxis
gamle Bulværk, Oxforduniversitetet, har det tilfælles
med R. Nielsens Polemik hos os, at det ikke udgaaer fra
et fjendtligt Forhold til den positive Religion, men paa
engang vil hævde den religiøse Tros Eiendommelighed og
den videnskabelige Erkjendelses Love. Dog er der én
væsentlig Forskjel mellem Mansels og Nielsens Standpunkt,
og i et Punkt, hvor man maa indrømme,, at den engelske
Philosoph er mere konsekvent end den danske. Nielsen
hævder, hvad Mansel negter, en Philosophi om det Ubetingede.
Udgaaet fra Hegels absolute Idealisme fastholder
han ogsaa i sin senere Lære Muligheden af en philosophisk
Bestemmelse af den absolute Aands Væsen, en Mulighed,
% --- p. 109 ---
\setcounter{footnote}{0}%
\opage{109}som for ham bliver Et med Erkjendelsens egen Mulighed
(Grundideernes Logik I p. 53). Her reiser sig da den
Vanskelighed at bestemme Forholdet mellem denne philosophiske
Gudside og den Opfattelse af Guddommen, som det
religiøse Livs Krav skal føre til; der kan dog ikke være to
absolute Ideer, og Striden mellem Philosophi og Religion
begynder da paa Ny ved Spørgsmaalet om, hvilken af dem
der skal afgjøre, hvad det \emph{sande} Ubetingede er. Dette
har Mansel, og allerede Hamilton, klart set. De negte, at
Philosophien kan erkjende det Absolute og Ubetingede, og
derved staaer Pladsen da aaben for de religiøse Forestillinger.
Mansel bemærker, at de græske Philosopher kunde have en
philosophisk Ide om det Absolute, om det til Grund for
Alt liggende Princip, uden at behøve at gjøre det til Et
med Folkereligionens Guddom; de kunde betragte Folkereligionens
Guder som afledede personlige Former for hint
Absolute, ligesom i Theogonien de høiere Guder have udviklet
sig af Chaos og Natten. Men indenfor Christendommen maa det absolute Princip, man antager, nødvendigvis
gjøres til Et med den Gud, man troer paa.

Saa meget man end maa give Mansel Ret heri, saa
er det dog et Spørgsmaal, om Modsætningen mellem Tro
og Viden kan fastholdes, selv naar man opfatter den saaledes,
som han gjør. Enhver positiv Religion vil jo dog
netop lære, hvad Gud og det Guddommelige er, og navnlig
den christne Kirke opstiller en Række af bestemte Lærdomme,
den forlanger anerkjendte; ved disse Lærdommes
Udvikling er det, hvad Dogmehistorien ogsaa viser, umuligt
at holde Tanken og Erkjendelsen ude. Troen kan altsaa
ikke være uden Theologi, og det vil derfor vise sig at være
en Illusion at mene, at man kan omstyrte al Theologi og
dog hævde den positive Tro.

Det er kun, naar man sætter en dualistisk Sondring
mellem det Ubetingede og det Betingede, at man kommer
til en saadan Modsætning mellem Viden og Tro. Vil man
hævde det Overnaturlige, da er det ganske vist konsekvent
at anerkjende Umuligheden af en Begrundelse ad Tænk%
% --- p. 110 ---
\setcounter{footnote}{0}%
\opage{110}ningens Vei. Men det overnaturlige Princip, som i sig
indeslutter en Fornegtelse af Tankens Gyldighed, er som
saadant fremmed for Philosophien. Denne kan kun opfatte
det Absolute som iboende, i dets indre Enhed med det
Relative, og den ser heller ikke i Menneskets moralske og
religiøse Væsen nogen Nødvendighed af at forlade dette
rationelle, humane Standpunkt, eftersom vi ingen andre
Midler have til at \emph{erkjende} den sande Moralitet og den
sande Religiøsitet end netop Tanken, Erkjendelsen selv.
Dersom der, som Mansel erklærer, skulde gives »retmæssige
Kilder til Overbevisning, forskjellige fra dem, der konstituere
den særlig saakaldte Erkjendelse,» saa maatte denne
Overbevisning udvikle sig efter en Logik, der er forskjellig
fra den, vi kjende og anvende, saa at vi fik to Slags
Logik --- og den gamle Strid begyndte da igjen forfra. ---

Vi vende nu tilbage til Hamilton. Stuart Mills Kritik
af Hamiltons Philosophi gik dels ud paa at bekæmpe Forsøget
paa at hævde en af Viden uafhængig Tro; --- og i
dette Punkt i det Mindste var Mill ikke just heldig, skjøndt
han kæmpede for Philosophiens Sag; det dreiede sig her
om Problemer, som han og hans Skole altid have skudt
til Side, og det var derfor et fremmed Omraade, han her
vovede sig ind paa. Dels var den rettet mod Hamiltons
logiske og psykologiske Lære, som vi nu skulle karakterisere.


Logiken er ifølge Hamilton en rent formel og apriorisk
Videnskab. Den skal slet intet Hensyn tage til den virlige % PRINTED AS IS (virlige: O-E29: vir-|lige)
Existents og dens Forhold, men beskjæftiger sig kun
med Tankens egne Love og Betingelser. Om Sætningernes
Sandhed eller Falskhed i sig selv afgjør Logiken Intet;
for den kan Alt være sandt, hvad der ikke indeholder en
Selvmodsigelse. Logiken har altsaa kun at gjøre med den
Nødvendighed, som beroer paa Begrebernes indbyrdes Forhold,
men kan Intet afgjøre om den reale Nødvendighed,
som udspringer af Tingenes virkelige Natur. Hamilton
bebreider endog Aristoteles at have krænket Logikens for%
% --- p. 111 ---
\setcounter{footnote}{0}%
\opage{111}melle og abstrakte Karakter ved at indblande Hensyn til
den reale Sammenhæng.

Indenfor den formelle Logiks Omraade har Hamilton
lagt megen Iver og Dygtighed for Dagen. Han har forsøgt
at opstille en ny Theori om Slutningerne, som vi dog
ikke her kunne gaa nærmere ind paa.\footnote[1]{Hos \emph{Trendelenburg}: Logische Untersuchungen. 2te Ausg. II p. 304---307 findes en kort Karakteristik og træffende Kritik af denne „New analytic of logical forms.“ Navnlig vises, hvorledes Hamiltons egen Theori fører ud over den formelle Logik.} Men med Hensyn
til den formelle Logik i det Hele maa vi bemærke, at
Mill har Ret, naar han hævder, at den vel kan tjene til
at gjøre Prøve paa, om vore Slutninger ei ere selvmodsigende,
men ikke til at vise, hvorledes vi vinde nye Erkjendelser.
Den Udvidelse, Mill har givet Logiken, er
derfor et virkeligt videnskabeligt Fremskridt. Den formelle
Logik prøver Tankens Overensstemmelse med sig selv, men
løser ikke den vigtigste Opgave, Tankens Overensstemmelse
med Virkeligbeden, det egentlige Erkjendelsesproblem. % PRINTED AS IS (Virkeligbeden: O-E30)

Sit almindelige philosophiske Udgangspunkt har Hamilton
fælles med den tidligere skotske Philosophi, idet
han betragter Philosophien som den videnskabelige Fremstilling
og Udvikling af Bevidsthedens Kjendsgjerninger.\footnote[2]{Philosophien i egentlig Forstand faaer altsaa kun med det Aandelige, slet ikke med det Materielle at gjøre; kfr Hamiltons Valgsprog. Philosophien bliver kun „philosophy of mind;“ „natural philosophy“ er ikke, hvad vi forstaa ved Philosophi.}
Bevidsthedens oprindelige Udsagn, vore instinktmæssige Antagelser
maa enhver Philosophi gaa ud fra, der ikke vil
gjøre sig selv til Intet. Hvad Bibelen er for Theologen,
det er Bevidstheden for Philosophen. At hine oprindelige
Udsagn ere sande, maa man antage, »hvis ikke selve Roden
til vor Natur er en Løgn.» Og det vigtigste af disse Udsagn og oprindelige Fakta er Modsætningen mellem Bevidstheden
og dens Gjenstand, mellem Jeg og Ikke-Jeg.
Idealismen, som ophæver Ikke-Jeget og gjør det til et
% --- p. 112 ---
\setcounter{footnote}{0}%
\opage{112}blot subjektivt Phænemen, opløser Alt til Drøm og Skygger. % PRINTED AS IS (Phænemen: O-E31)


Dette er den samme Argumentation, som Reid anvendte
overfor Hume. Men i to Punkter afviger Hamilton
fra Reid. Han betragter ikke som denne Bevidstheden
som en særegen Evne, forskjellig fra Iagttagelsen og Opmærksomheden
, men som Et med selve den umiddelbare
Erkjendelse. Vi erkjende en Gjenstand umiddelbart, ɔ: vi
ere os dens Tilværelse bevidste. Vi erkjende os selv
umiddelbart, ɔ: vi ere os vor egen Tilværelse bevidste. Hamilton
sætter dernæst som en nødvendig Forudsætning, at
Bevidsthedens primitive Udvagn ikke ere selvmodsigende. % PRINTED AS IS (Udvagn: O-E32)
Kun under denne Forudsætning have disse primitive Udsagn,
den almindelige Menneskeforstands Principer, Gyldighed
fremfor alle afledede Bevidsthedsphænomener, hvis Gyldighed
beroer paa hines.

Bevidstheden selv kan ikke defineres nærmere, da den
ligger ved Roden af al Erkjendelse. Alt hvad vi kunne
forestille os og begribe, forestille vi ved Bevidsthedens
Hjælp. Denne forudsættes bestandig. Hvis den nu selv
skulde kunne belyses og begribes, saa maatte vi jo have
en anden Bevidsthed, der kunde opfatte hin første, og
saaledes i det Uendelige. Men skjøndt Bevidstheden ei
kan defineres, kan den dog analyseres. Vi kunne fastholde
og sammenligne Bevidsthedsphænomenerne og derved
slutte os til Bevidsthedsakternes almindelige Betingelser.

Philosophien (som philosophy of mind) falder i tre
Hoveddele. I den første, Phænomenologien eller den empiriske
Psykologi, fremstilles Bevidsthedsphænomenerne saaledes,
som de faktisk forefindes. De ordne sig i tre Grupper:
Erkjendelse, Følelse og Villie. Den anden Del, Nomologien
eller den rationelle Psykologi, fremstiller de Love, som
gjælde for disse Phænomener, og deler sig derfor i Logik,
Æsthetik og Ethik. Den tredie Del, Ontologien eller »inferential
psychology,» uddrager saa endelig heraf Resultater
med Hensyn til Aandens Væsen i sig selv; navnlig behandles
her Problemerne om Guds Tilværelse og Sjælens
% --- p. 113 ---
\setcounter{footnote}{0}%
\opage{113}Udødelighed. --- Hamilton har kun naaet at behandle Phænomenologien
og Logiken. ---

Egentlig maa denne Anskuelse strande paa de samme
Skjær som Empirismen. Den gaaer ud fra den empiriske
Psykologi og kan kun ved en Inkonsekvents naa ud over
dette empiriske Grundlag og opstille Love for det Sande,
Skjønne og Gode. I ethvert Tilfælde vilde den hele Anskuelse
være rent subjektiv, eftersom den kun bygges paa
de subjektive Bevidsthedsphænomener, --- dersom der ikke
begyndtes med den postulerede Forudsætning, at disse
Phænomener have objektiv Gyldighed, da »Roden til vor
Natur ikke kan være en Løgn.»

Dog viser det sig ved nærmere Betragtning, at Hamilton
egentlig har et andet Udgangspunkt end hint empirisk-psykologiske,
at han ligesom Kant gaaer ud fra Bevidsthedens
oprindelige Aktivitet. Herpaa beroer det
kritisk-spekulative Element hos ham, hvorved han afgiver
interessante Korrektiver til den empiriske Skoles Lærdomme.

Al Bevidsthed forudsætter, siger han, en Dobbelthed,
en Modsætning. Derfor er al vor Erkjendelse betinget.
Men Modsætning og Dobbelthed kan ikke opfattes uden
ved Sammenligning. Derfor ligger der i enhver Erkjendelses 
akt en Skjelnen og Dømmen. Den sammenlignende Virksomhed
spiller en Rolle i den simpleste Iagttagelse. Den
første Bevidsthedsakt er Bekræftelsen af, at Noget existerer;
men det forudsætter en Adskillelse mellem dets Existents og
dets Ikke-Existents; altsaa en Sammenligning af disse to
Tilstande. Dernæst skjelnes der mellem Selvet og den ydre
Gjenstand, mellem de forskjellige indre og ydre Phænomener
o. s. v. Fra de simpleste til de høieste og mest sammensatte
Bevidsthedstilstande spore vi den sammenlignende
Virksomhed som gjennemgaaende Betingelse. Under denne
Virken følger Bevidstheden visse Love, som angive de Betingelser,
uden hvilke den ikke kan fungere. Disse Love
ere derfor vor Erkjendelses Grundlove (de logiske og mathematiske
Grundsætninger); de ere aprioriske, i deres Princip
uafhængige af Erfaringen, hvilket følger af deres Nødven%
% --- p. 114 ---
\setcounter{footnote}{0}%
\opage{114}dighed; en virkelig Nødvendighed kan Erfaringen ikke lære
os Først ved hin Bevidsthedsaktivitet og de Love, den
følger, bliver Erkjendelsen til. Har Erkjendelsen derfor end
sin Begyndelse i Erfaringen, saa har den dog sin egentlige
Oprindelse i Aanden selv.

Naar Hamilton har en saadan Opfattelse af Erkjendelsens
Væsen, kan han ikke med de engelske Psykologer
betragte Ideassociationens Love som de eneste, der gjælde
paa Aandslivets Omraade. Han kan heller ikke opfatte dem
som Udtryk for rent passive Tildragelser indenfor Bevidstheden,
men maa søge at føre dem tilbage til Bevidsthedsenergien som deres sidste Grund. Som der ytrer sig en
saadan Energi allerede i Opfattelsen og Iagttagelsen af
Gjenstandene, saaledes skyldes det ogsaa Selvets Energi, at
Forestillingen om et vist Indtryk fastholdes og bevares
saaledes, at den senere han vækkes igjen. Og denne samme
Energi virker ogsaa i Ideassociationen, hvis almindelige Lov
kan betegnes som Fuldstændiggjørelsens Lov (the law of
redintegration): Forestillinger, som høre sammen, ɔ: udgjøre
et Hele for Tanken, enten fordi de tidligere have berørt
hinanden som successive eller samtidige, eller fordi de
ligne hinanden, have en Tendents til at fremkalde
hinanden. Denne Lov, der i sig indeslutter de speciellere
Love for Ideassociationen, synes at have sin Grund i det
aktive Selvs Enhed, der ogsaa stræber at bevirke Enhed og
Sammenhæng mellem Forestillingerne indbyrdes. --- Paa
denne Maade vikler, som allerede ovenfor (p. 37 f.) vist,
Associationspsykologien, og med den Empirismen, sig op og
gaaer ind under et høiere Princip.

Associationspsykologien troede særlig at kunne forklare
Kausalitetsprincipet ud af Erfaringen ved at betragte det
som Exempel paa en Ideassociation, der paa Grund af den
fuldstændige Bekræftelse, den har vundet i Erfaringen, er
fast og uopløselig. Men Ideassociationen kan, indvender
Hamilton, kun forklare, hvorledes vi komme til at henføre
et vist Phænomen til en vis Aarsag, derimod ikke, hvorledes
vi i det Hele komme til at spørge om en Aarsag.
% --- p. 115 ---
\setcounter{footnote}{0}%
\opage{115}De fleste Philosopher, der ikke hylde den rene Empirisme,
pleie at opstille Kausaliteten som et særligt Princip, en
særegen Lov for vor Erkjendelse. Ifølge Hamilton er dette
ikke nødvendigt. Kausalitetens Lov er ifølge hans Opfattelse
kun en særegen Form for det Betingedes Lov, ɔ:
den Lov, ifølge hvilken vor Erkjendelse kun omfatter det
Betingede og Relative. Da vi ikke kunne begribe et Absolut
eller konstruere det i vor Tanke, saa kunne vi heller
ikke tænke en absolut Begyndelse eller en absolut Ende.
Naar en Ting opstaaer, er dette for os kun tilsyneladende;
vi forudsætte uvilkaarlig, at den har været til iforveien
under en anden Form, ɔ: vi udlede den af en vis Aarsag,
eller af Aarsager, thi der er altid flere. Vi opfatte ikke
Existentsen absolut eller i sig selv, men Existentsen under
Tidens Form. Fra Existentsen i det Hele kan der ikke
drages et Atom, lige saa lidt, som der kan føies et Atom
til den; vi kunne hverken begribe, hvorledes Summen af
det Tilværende skulde kunne forøges eller formindskes. Ex nihilo
nihil. Omnia mutantur, nihil interit. Naar Noget begynder
eller ender, kan det kun være som Phænomen. Selv naar
vi forsøge at tænke os Verdens Skabelse, tænke vi stedse
Verden som indeholdt efter Muligheden i Guds Væsen, saa
at Skabelsen kun er en Overgang fra Mulighed til Virkelighed.
--- Det er altsaa ifølge vor Aands Begrændsning og
Endelighed, ifølge vor Erkjendelses Relativitet, at vi ere
nødte til at tænke os Aarsager til alle Ting.

Denne Opfattelse af Kausalitetsprincipet har den Fortjeneste,
at den stræber at føre dette Princip tilbage til
selve vor Bevidstheds Væsen i Stedet for at sætte det som
et eget isoleret Princip. Den betoner særlig den Side ved
Aarsagsforholdet, at dette angiver en Sammenhæng mellem,
en Enhed af Phænomenerne. Nu betragter Hamilton det
vistnok som en Følge af vor Begrændsning, at vi ikke
kunne Andet end henføre enhver Begivenhed til en tidligere,
hvoraf den har udviklet sig; men naar denne Henføren
skal betyde, at vi opfatte Phænomenerne som Et i deres
Grund og kun forskjellige i deres Fremtræden,  da ligger
% --- p. 116 ---
\setcounter{footnote}{0}%
\opage{116}det lige saa nær at betragte det som en Fuldkommenhed,
at vor Bevidsthed søger tilbage til Enheden og ikke bliver
staaende ved den phænomenale og tilfældige Flerhed. I sin
Konsekvents fører, som man rigtig har bemærket, denne
Lære til Ideen om en evig Substants, hvis Aabenbaringer i
Tiden de phænomenale Aarsager og Virkninger ere, ---
altsaa til den Pantheisme, som Hamilton frygtede saa meget.
Hamiltons Lære danner her en fuldstændig Modsætning til
Empirismen, der i Aarsagsforholdet blot ser et Successionsforhold,
en Forbindelse mellem forskjellige Phænomener
uden indre, nødvendig Sammenhæng. Allerede Kant har
vist, at vore Domme om Tingene kun have objektiv Gyldighed,
kun angaa reale Forhold og ikke Drømmebilleder,
naar Aarsagsforholdet er Mere end tilfældig Tidsfølge.
Kausalitetsloven er Udtryk for en af de Former, hvorunder
vor Tanke sammenfatter det Givne til Enhed. En saadan
Enhed maa vi ifølge vor Organisation forudsætte i Tilværelsen;
vi maa i sidste Instants antage en evig Substants
som Grund til al nødvendig Sammenhæng mellem Phænomenerne.
 Det er da klart, at Hamilton her, lige som Kant,
selv viser ud over det Betingedes Philosophi. Den Enhed
og Sammenhæng, vi finde i Tilværelsen, beroer paa, at det
Ubetingede er den indre Grund til det Betingede, er det
alt Betinget Betingende. --- Og her have vi da en naturlig
Overgang til den Tænker, som netop paa den nu antydede
Maade har udviklet Hamiltons Lære saa vel som Empirismen
videre, som paa original Vis søger ud over begges
Ufuldkommenheder, og som derfor kan siges at betegne
Høidepunktet af de philosophiske Bestræbelser i England i
vor Tid.

\chapterhead{Fjerde Kapitel.}{Udviklingsphilosophien.}


Det følger af den engelske Philosophis empiriske
Karakter, at den sjelden har stræbt efter at samle sine
% --- p. 117 ---
\setcounter{footnote}{0}%
\opage{117}Lærdomme til et afsluttet System. Den har altid havt
skarpest Blik for det Endelige og Betingede i den menneskelige
Viden, og ethvert System synes jo dog at forudsætte
et ubetinget Princip, der kan bringe Afslutning og Enhed
i Verdensopfattelsen. De Retninger vi i det Foregaaende
have beskjæftiget os med, gjøre ingen Undtagelse herfra.
Stuart Mill har sin Styrke i Behandlingen af de enkelte,
isolerede Problemer; disse staa hos ham Side om Side og
bringes ikke i indre Forbindelse, skjøndt det naturligvis er
den samme Grundanskuelse, vi møde ved hans Behandling
af dem hver især. Mill har intet indre Princip; hans Princip
er jo, at der intet Princip er, men at alle ydre og indre
Phænomener gruppere sig efter Love, der ikke besidde nogen
i sig gyldig Nødvendighed. William Hamilton kunde heller
ikke stræbe efter et System, eftersom hans Hovedlærdom jo
var, at det Ubetingede er os aldeles ufatteligt; og skjøndt
han betragtede Nødvendigheden som Mere end blot ydre
Forbindelse i Rum og Tid, var det dog efter hans Opfattelse
kun en Følge af vor Begrændsning, at vi anvende Aarsagsbegrebet
paa Tingene; for det Ubetingedes Vedkommende
havde dette Begreb ingen Gyldighed.

Hvad der giver \emph{Herbert Spencer} enestaaende Betydning
i den engelske Philosophi, er derfor ikke mindst
hans systematiske Stræben, hans Stræben efter overalt at
samle de spredte enkelte Træk til en Helhedsanskuelse.
Han er ikke tilfreds med at behandle de enkelte Problemer.
Skjøndt han besidder en forbausende Viden paa saa godt
som alle særlige Videnskabers Omraade, betragter han den
dog blot som det Materiale, hvoraf en ny, Alt omfattende
Bygning skal opføres. Han er et sjeldent Exempel paa en
Forening af analytisk og synthetisk Begavelse. Hos de andre
engelske Philosopher er Analysen det Fremherskende; det
Givne og Foreliggende prøves og opløses i sine Elementer,
føres tilbage til sine simpleste Former. Derimod stræbe de
ikke efter en saadan Rekonstruktion af det Givne, at dettes
Udvikling kan komme til at aabenbare dets Natur i dens
hele Sammenhæng. Af de simpleste Former, de usynlige
% --- p. 118 ---
\setcounter{footnote}{0}%
\opage{118}Spirer søger den synthetiske Tænker Skridt for Skridt
at udlede de høieste og betydningsfuldeste Funktioner, idet
han griber den for hele Udviklingen gjældende Lov.

Det er dette paa engang omfattende og dybtgaaende
Blik, der gjør, at Spencer mere end de andre fra den ældre
engelske Skole udgaaede Tænkere har Øie for det Vendepunkt,
hvortil Empirismen i vore Dage er naaet. Han
lader sig ikke nøie med at komme tilbage til de Forudsætninger
og Principer, ved hvilke Empirismen ophæver
sig selv, men drager med Bevidsthed deres Konsekventser,
ser dem i deres Betydning som Grundlag for hele Verdensanskuelsen.
Det er Paavirkningen af Hamilton og den
kritiske Skole, der lader ham se skarpere paa det metaphysiske
Omraade end Empirikerne formaa. Den dybere
gaaende philosophiske Begrundelse, som kun kan vindes ved
at føre Tankens og Virkelighedens Love tilbage til et
Princip, der ligger til Grund for begge, denne Begrundelse,
som Empirismen hverken kunde søge eller finde, stræber
Spencer at give.

Men trods denne metaphysiske Tendents vil Spencer
ikke forlade Erfaringens Omraade. Han regner sig endnu
til Erfaringsphilosophien, skjøndt han polemiserer mod den
rene Empirisme, der tilsidst maa ende i det Principløse.
Han adskiller sig ikke blot fra Empirikerne ved den dybere
Begrundelse, han søger, men ogsaa ved den nye Belysning,
han paa det saaledes vundne Grundlag giver de faktiske
Bestemmelser, hvorved Empirismen slaaer sig til Taals. Vi
have tidligere set, at denne Standsen ved det Positive og
Faktiske netop er det, der bringer Modsigelse ind i den
empiriske Anskuelse, idet de enkelte Bestemmelser staa
faste og ubevægelige overfor hverandre, --- at der vel paa
nogle Steder gjøres Tilløb til at bringe Bevægelse og Udvikling
ind, men at Empirismens atomistiske Karakter
hindrer denne Stræben, og at Bevægelse og Udvikling i
hvert Tilfælde ikke lægges til Grund som herskende Synspunkt.
Anderledes hos Spencer. De faste Former smelte,
Liv og Bevægelse opstaa i de døde Masser. Udviklingen
% --- p. 119 ---
\setcounter{footnote}{0}%
\opage{119}bliver ikke et isoleret Faktum, man paa et eller andet Sted
støder paa, men Udtrykket for en Verdenslov, som raader
i alle Phænomener. Og herved banes der Vei for en Art
Positivitetens Dialektik, der overvinder og forsoner Modsætningerne
ved at indordne dem som Led i den uendelige
Bevægelse.

Spencers Bestræbelser have i de senere Aar vundet
fortjent Anerkjendelse. Det ser man ikke blot af hans
Skrifters store Udbredelse i England og Nordamerika og af
Oversættelserne paa fremmede Sprog (Fransk og Russisk),
men ogsaa af Udtalelser af fremragende Videnskabsmænd.
George Lewes opkaster i sin Philosophiens Historie det
Spørgsmaal, om der nogensinde er opstaaet en mere fremragende
Tænker i England. Og skjøndt Stuart Mill paa
Grund af Spencers Kritik af Empirismen altid omtaler ham
med et vist Forbehold, erklærer han ham dog for Auguste
Comtes Jevnbyrdige, en ikke ringe Roes i hans Mund. Den
berømte Naturforsker Huxley siger om Spencers Hovedværk:
»Den eneste fuldstændige og methodiske Fremstilling
af Udviklingens Theori, jeg kjender, findes i Herbert
Spencers »philosophiske System,» et Værk, der med Omhu
bør studeres af alle dem, som ønske at gjøre sig bekjendte
med den videnskabelige Bevægelses Tendentser.»\footnote[1]{Disse Udtalelser findes anførte hos \emph{Ribot}: „La psychologie Anglaise contemporaine,“ som ogsaa selv giver en træffende og sympathetisk Karakteristik af Spencer (p. 145---152).} --- Men
det er ogsaa først nu, da hans store Værk nærmer sig sin
Fuldendelse, at Spencer er bleven anerkjendt. Hans Stræben
var fremmed for den engelske Aandsretning; han er aaben
og resolut i sin Tankegang og sine Udtalelser og viser
ikke den empiristiske Besindighed, der lader selv Mænd
som Stuart Mill og Darwin lempe sig efter de herskende
Forestillinger; --- derfor kunde han ikke gjøre Regning paa
Sympathi, men maatte føle sig tilovers i det engelske
Samfund. »Midt under den almindelige Ligegyldighed er
Spencer forbleven Philosophien urokkelig tro. Han har
% --- p. 120 ---
\setcounter{footnote}{0}%
\opage{120}sikkert maattet besidde et heroisk Mod og en sjelden
Uafhængighedsfølelse for at kunne hellige sig til strenge
Studier, der kun kunde skaffe ham enkeltstaaende og ubemærkede
Bifaldsytringer. Med det Talent, den aandelige
Frugtbarhed og den næsten Alt omfattende Mangfoldighed
af Kundskaber, hvorom hans Skrifter vidne, vilde han, saafremt
han var gaaet ind paa at følge de slagne Veie, med
Lethed have opnaaet alle de Gaver, som det engelske
Samfund i rigt Maal skjænker dem, der tjene det saaledes,
som det selv vil tjenes. Han har foretrukket at indvie
sig til Fattigdom og til --- hvad der er endnu vanskeligere
at bære --- Uberømthed. Men Spencer fortjener Mere end
tomme Tilkjendegivelser af Sympathi» etc.\footnote[1]{Aug. \emph{Laugel}. Revue de deux mondes. 15. févr. 1864.} Baade i
Karakter og Verdensanskuelse synes Spencer at have meget
tilfælles med Spinoza. ---

For strax at give en Forestilling om Spencers Lære
ville vi fremdrage den Hovedtanke, som behersker den fra
først til sidst, Udviklingens Ide. Denne Ide er for Spencer
Resultatet af en omhyggelig Analyse af Erfaringens Kjendsgjerninger;
den angiver en Lov, som raader overalt i
Naturens og Aandens Verden. Og paa den anden Side
kan man kun ved at lægge denne Ide til Grund synthetisk
fremstille vor Erkjendelses hele Indhold. Spencer bliver
ikke træt af at fremstille sin Ide under begge disse Synspunkter
og vise dens Sandhed paa alle Omraader. Sit
Udgangspunkt havde han oprindelig paa det ethiske og
sociale Omraade; her saa han først Loven om den af
naturlige Aarsager fremkaldte Udvikling gjælde. Men en
fast Form, en større Sammenhæng kom der først i hans
Anskuelser, da han lærte den Sandhed at kjende, »som
først dunkelt antydedes ved Harveys embryologiske Undersøgelser,
som siden efter opfattedes klarere af Wolff, og
som blev bragt i en bestemt Form af von Baer --- den
Sandhed, at al organisk Udvikling er en Overgang fra en
% --- p. 121 ---
\setcounter{footnote}{0}%
\opage{121}ensformig til en forskjelligartet Tilstand.» I et af sine
første Skrifter, »Social Statics» (1850), fremstillede han ikke
blot Menneskets og Samfundets Udvikling, men udtalte
ogsaa, at det i Principet var den samme Lov, der gjaldt
for den organiske og den sociale Udvikling. Siden førtes
han saa til at underkaste denne Lov en nøiere Undersøgelse,
som baade førte til en rationel Begrundelse af den og til
at opstille flere Former for Udviklingens almindelige Lov,
end han fra først af havde gjort.\footnote[1]{\emph{Spencer}: The classification of the sciences. 3d. ed. London 1873. p. 46.} Denne Undersøgelse
indeholdes i hans »First Principles» (1862). Men førend
vi gaa over til den, ville vi gjøre nogle almindelige Bemærkninger
om den Udstrækning og den Belysning, hvori
Spencer opfatter Udviklingens Lov.

Astronomien viser os efter Laplaces Hypothese, hvorledes
den oprindelige Taagemasse efterhaanden artikuleredes
og omdannedes, saa at det nuværende Solsystem fremgik
deraf. Geologien viser, hvorledes Jorden fra sin oprindelige
glødende Tilstand, der var en Følge af den ved Koncentration
af Taagemassen udviklede Varme, har opnaaet sin
nuværende Tilstand gjennem en Række af Forandringer:
Jordskorpens Afkjøling og Størkning, Fordelingen af Land
af Vand, Leiring af de forskjellige Lag o. s. v. Det organiske % PRINTED AS IS (af Vand: O-F07: for og Vand)
Liv viser os en lignende Proces baade indenfor Individets
og indenfor Arternes Historie. Det er stedse de
høiere, rigere og mangfoldig udviklede Former, der opstaa
senest. Hvad Mennesket angaaer, kunne vi paavise en hel
Række af Trin lige fra den halv dyriske Vilde til den
dannede Europæer. Ogsaa med Hensyn til enhver enkelt
menneskelig Bestræbelse gjælder det Samme, hvad enten vi
forfølge Sprogets, Kunstens, Videnskabens, Handelens eller
Industriens Historie. Alle Phænomener vise hen til samme
Lov. Men deraf følger paa den anden Side ogsaa, at denne
Lov kun er en Hypothese, hvortil vi med Nødvendighed
komme; vi kunne ikke tilskrive den absolut Gyldighed, da
% --- p. 122 ---
\setcounter{footnote}{0}%
\opage{122}vi ikke kjende det Absolute og altsaa ei kunne vide, hvorvidt
denne Lov har Anvendelse paa det. De store Problemer
angaaende Tingenes Væsen og Grund løses ikke
herved; det sidste Mysterium rykkes kun et Skridt tilbage,
men forbliver i sig selv lige saa dunkelt, som det altid
har været.

I Almindelighed opfatter man Udviklingens og Fremskridtets
Ide paa ensidig teleologisk Maade, som om alle
Trin i Udviklingen kun skulde have til Formaal at fremme
den menneskelige Lyksalighed. Men dette er et endeligt
og rent subjektivt Synspunkt. Den philosophiske Forsken
spørger derimod om, hvilke Love der gjælde for Phænomenerne,
for de bestandige Forandringer og Fremskridt,
aldeles afset fra vore Interesser, --- spørger om, hvad Udviklingen
er i sig selv, og hvilke dens Aarsager ere. Og
som Udviklingen selv kun er Udtryk for en phænomenal
Lov, saaledes kunne ogsaa dens Aarsager kun være phænomenale
og relative; den absolute Aarsag er os utilgængelig.

Som Udviklingens Lov kun angaaer det Relative og
Phænomenale, saaledes er den i sig selv et Udtryk for
Relativitetens Princip. Naar Alt er i Udvikling, saa er
Alt ogsaa kun relativt. Derfor betoner Spencer ogsaa ofte,
hvor umuligt det er at give fast afpælede Definitioner. De
absolute Grændsers Tid er forbi. Man kan ikke paavise
noget Punkt, hvor det aandelige Liv begynder; der finder
en successiv Overgang Sted fra de rent legemlige Funktioner
til de sjælelige. Lige saa lidt kan man fastslaa nogen
Grændse mellem de organiske og de uorganiske Phænomener,
mellem det Levende og det Livløse. Den nyere Physik
viser os, hvilken indre Sammenhæng der er mellem de
forskjellige Kræfter, Elektricitet, Magnetisme, mekanisk
Kraft, Varme o. s. v.; den ene af disse kan omsættes i
den anden. Videnskaben selv staaer ikke i nogen absolut
Modsætning til den populære Bevidsthed; den raader ikke
over Midler og Kræfter, som ikke allerede ere tilstede i
denne Bevidsthed, men betegner et høiere Udviklingstrin
af disse samme Kræfter, en Udvidelse af Synskredsen, en
% --- p. 123 ---
\setcounter{footnote}{0}%
\opage{123}konsekventere og nøiagtigere Gjennemførelse af Synspunkterne.


Udviklingens Lov er Udtrykket for, at Relativitetens
Verdensanskuelse er bleven sig selv bevidst. Den falder
ikke længere ud i det Principløse og Usammenhængende,
men samler sig i et Princip. Ligesom Kritiken af Materiens
Begreb i den nyeste Naturphilosophi har ført til at sætte
Materialiteten i Stedet for Materien, saaledes kan man sige,
at Spencer har sat Relativiteten i Stedet for det Relative.
Men hvor der er et Princip, der er der forsaavidt ogsaa et
Absolut, et Ubetinget. Idet Spencer opstiller et Princip og
paa Basis af det udvikler et System, gaaer han derfor ud
over Empirismen og han er sig dette bevidst. Ved konsekvent
Gjennemførelse og Gjennemtænkning af den empiriske
realistiske Skoles Anskuelser er han naaet ud over den
og har grebet Tanker, der ligge til Grund for den spekulative
(idealistiske) Retning. Han sætter det ogsaa udtrykkelig
som sin Opgave at forsone Locke og Kant.

Men vi kunne allerede her angive, hvor det vanskeligste
Punkt i Spencers Philosophi ligger. Thi det Princip,
hvortil det Relative føres tilbage, og hvori Loven for Phænomenerne
har sin Grund, kan ikke søges andet Steds end
i det Absolute og Ubetingede; --- ellers vilde det jo ikke
være det sande Princip, men der maatte søges et andet.
Det Relative kan kun samles under et principielt Synspunkt,
naar man forudsætter, at det Relative og Phænomenale
staaer i indre Forhold til det Absolute og Ubetingede;
kun fra dette kan al Enhed og Sammenhæng udledes.
Enhver virkelig Totalanskuelse maa betegne et Fremskridt
i Erkjendelsen af Virkeligheden selv; den nominalistiske
Tanke, at kun det isolerede Faktum har Realitet, opløser
sig selv i den Erkjendelse, at netop Enheden og Helheden
er den sande Realitet. Og henimod en saadan Erkjendelse
er det ogsaa, at Spencers egen philosophiske Anskuelse
viser.

Relativitetens Princip ophæver, som vi have set, de
faste Grændser. Men naar alle Grændser falde, hvorledes
% --- p. 124 ---
\setcounter{footnote}{0}%
\opage{124}skulde da Grændsen mellem det Relative og det Absolute
kunne blive staaende? Maa ikke denne Grændse, som for
Spencer er Et med Grændsen mellem det Erkjendelige og
det Uerkjendelige, selv blive relativ og flydende, ɔ: underkastes
Udviklingens Lov? --- Der synes ingen anden Udvei
at være, hvis Spencer ikke vil falde tilbage til den »rene
Empirisme,» han bekæmper. Dualismen mellem det Absolute
og det Relative vilde gjøre Spencers egne Bestræbelser efter
en »synthetisk Philosophi» frugtesløse.

Spencer hævder, som vi ville faa at se, udtrykkelig
Nødvendigheden af det Absolutes Ide som sidste Princip,
sidste Forudsætning for al Tænkning og Erfaring. Ja, han
antyder endog, at de væsentlige Former og Love for Tilværelsen,
vi kunne erkjende, ogsaa maa være Udtryk for
det Ubetingedes Væsen; han viser, hvorledes Forestillingen
om dette i Tidernes Løb, under Religionens og Videnskabens
Udvikling, har antaget stedse høiere og fuldkomnere Former;
--- men heri ligger da, at det ikke kan være absolut uerkjendeligt.
Hvis man vel antog et absolut Princip, men
holdt det Relative udenfor al Berøring med det, saa vilde
dette Princip ikke blot være et uerkjendeligt, men et
meningsløst Begreb.

Den Usikkerhed, vi finde hos Spencer paa dette afgjørende
Punkt, er naturligvis medført ved de dobbeltsidige
Indvirkninger, han har modtaget, --- fra den empiriské og
fra den kritiske Skole, og den kunde neppe undgaaes hos
en Tænker, der har sat sig et saa storartet Maal som det
at forene de to modsatte Strømninger i den nyere Philosophi.
Det Spørgsmaal kunde nu opkastes, hvilken af de
to Strømninger der faaer Overhaanden i Spencers Anskuelse.
Dersom jeg opfatter denne ret, synes Resultatet i sidste
Instants at være, at Spencer i endnu høiere Grad, end han
selv er sig det bevidst, har overskredet Empirismens Grændser
og nærmet sig de ideale Opfattelser, som have arbeidet sig
frem gjennem Fastlandets philosophiske Systemer. Begrundelsen
heraf skal den følgende Fremstilling give. ---

Spencer har udkastet og tildels allerede udført en
% --- p. 125 ---
\setcounter{footnote}{0}%
\opage{125}storartet Plan. I Aaret 1860 udstedte han en Subskriptionsindbydelse
paa et philosophisk System, en »synthetisk
Philosophi,» der var anslaaet til ti Bind. Af disse foreligge
de fem første. Det første Bind indeholder under Navn af
»First principles» Grundlaget for hele Anskuelsen, idet det
undersøger Forholdet mellem Religion og Videnskab, opstiller
de Forudsætninger, paa hvilke al Tænkning og Erfaring
hviler, og de almindelige Love, der fremgaa som
sidste Resultater af de særlige Videnskabers Undersøgelser,
samt paaviser, hvorledes disse Love kunne udledes af hine
Forudsætninger. Det er altsaa en Metaphysik. De to
næste Bind give under Navn af »Biologiens Principer» en
Fremstilling af det organiske Livs Betingelser og Udvikling og
viser, hvorledes Livets Former og Funktioner danne sig
under Vexelvirkningeu med Omverdenen. »Psykologiens % PRINTED AS IS (Vexelvirkningeu: O-F11)
Principer,» ligeledes i to Bind, fremstiller paa tilsvarende
Maade Sjælelivets Udvikling under Kampen med Omgivelserne.
Af de fem Bind, som staa tilbage, skulle de tre
(»Principles of sociology») indeholde Samfundsvidenskaben,
de to sidste (»Principles of morality») Ethiken. --- Foruden
de her nævnte Skrifter har Spencer udgivet flere mindre
Arbeider. Den allerede omtalte »Social Statics» udvikler
nogle af hans ethiske Ideer; »The classification of sciences»
giver foruden det i Titelen angivne Emne en Bestemmelse
af hans Forhold til Positivismen; desuden har han samlet
en Del spredte, meget interessante Afhandlinger i to Bind
»Essays, scientific, politic and speculative.» --- I vor Fremstilling
følge vi hans egen systematiske Orden. Kun derved
kunne vi faa en Forestilling om den Enhed, der gaaer
igjennem hans Anskuelse, og om den Energi, med hvilken
han søger at faa Erkjendelsens forskjellige Strømme til at
forene sig til én fremadilende Flod. ---

Den skarpeste Modsætning, hvori forskjellige Anskuelser
kunne staa til hinanden, er den, vi finde mellem Religion
og Videnskab. Den skriver sig fra den allerførste Iagttagelse
og Opdagelse af en Lov, en regelmæssig Sammenhæng i
Naturen, og den har gjentaget sig paa ethvert senere Trin
% --- p. 126 ---
\setcounter{footnote}{0}%
\opage{126}af Religionens og Videnskabens Udvikling. Dog er det
ikke nogen uforsonlig Strid. Her som overalt maa der være
Noget, som begge de stridende Parter ere eller kunne blive
enige om. Selv i det Falske er der »en Sjæl af Sandhed;»
selv den vildeste Overtro maa indslutte et berettiget Element.
Ellers vilde man ikke kunne forstaa, hvorledes
Religionen har kunnet hævde sig trods Videnskabens atter
og atter gjentagne Angreb paa dens Dogmer. Selv om man
vil føre Dogmerne tilbage til den religiøse Følelse som deres
Udspring, saa staaer dog denne Følelse selv tilbage som
saa væsentligt et Faktum, at enhver philosophisk Theori, der
ikke forklarer det, er ufuldstændig. Hvis man antager, at
denne Følelse umiddelbart er skabt, saa maa man ogsaa
antage, at den svarer til sit Ophav; hvis man antager at den
lige som alle sjælelige Evner har udviklet sig under Individernes
Vexelvirkning med deres Omgivelser, saa maa man
ogsaa antage, at der er bestemte reale Betingelser, hvortil
den svarer. I begge Tilfælde føres vi altsaa til at betragte
den religiøse Følelse med Respekt. --- Hvad alle Religioner
have tilfælles, kan kun være den Sandhed, der ligger til
Grund for dem alle; og ligeledes kan det, Religion og Videnskab
have tilfælles, kun være en abstrakt Sandhed, i
hvilken de begge mødes. Lykkes det at eftervise denne
Sandhed, hvori den menneskelige Bevidstheds to Poler
mødes, saa vil det vise sig, at Striden mellem Religion og
Videnskab har været af samme Beskaffenhed som Kampen
mellem de to Riddersmænd om Farven af et Skjold, af
hvilket de hver kun havde set sin Side.

De høieste religiøse Ideer angaa Verdens Væsen og
Oprindelse. Dette Spørgsmaal er blevet besvaret paa tre
forskjellige Maader. Ifølge Atheismen bestaaer Verden ved
sig selv og er ikke frembragt; ifølge Pantheismen har den
frembragt sig selv; ifølge Theismen er den frembragt af en
ydre Aarsag. Den atheistiske Opfattelse er uholdbar, da vi
ikke kunne danne os en Forestilling om Noget, der ingen
Oprindelse har; Existents i og ved sig selv er for vor
Tanke noget Ufatteligt. I ethvert Tilfælde giver denne Op%
% --- p. 127 ---
\setcounter{footnote}{0}%
\opage{127}fattelse os intet Svar paa Spørgsmaalet, ingen Forklaring
af Verdens Tilværelse. Ogsaa Pantheismen er uholdbar:
hvad der frembringer sig selv, maa i Forveien have været
til under en anden Form, --- og Spørgsmaalet vender da
tilbage i det Uendelige. Paa samme Maade gaaer det
Theismen: den fører fra Verden tilbage til Skaberen og
antager Skaberen at bestaa ved sig selv; --- men saa ende
vi atter her ved det for os utilgængelige Begreb om et i
sig Bestaaende.

Den kritiske Betragtning godtgjør saaledes Uholdbarheden
af alle de Hovedformer, den religiøse Ide har antaget.
Vi kunne, som Hamilton og efter ham Mansel have
vist, ikke danne os noget Begreb om det Absolute og
Ubetingede, derfor er ethvert Gudsbegreb fuldt af Modsigelser.


Heraf lære vi, at Sandheden i Religionerne ikke kan
bestaa i den Maade, hvorpaa de løse Problemet om Verdens
Væsen og Oprindelse. Vi maa derimod søge deres Sandhed
i selve Anerkjendelsen af, at her er et Problem, at der gives
et Mysterium, der danner Tilværelsens Baggrund og ligger
ud over vor Erkjendelses Rækkeevne. Alle Religioner anerkjende
et Mysterium, men søge dog, hver paa sin Maade,
at give en Løsning af Gaaden, en Aabenbaring, hvorved
Mysteriet da atter faktisk ophæves. Dog viser det sig, at
jo høiere en Religion staaer, des høiere hæves det Ubetingede
over Erkjendelsens Omraade. Lavest staaer Fetischdyrkeren,
der forestiller sig sin Gud under de mest dagligdags
og endelige Former; Polytheismen viser os allerede
Guddommen under mere ideale Skikkelser; Monotheismen
reiser Altre for den ubekjendte og uerkjendelige Gud, og erkjender
tilsidst, at en Gud, vi kunne begribe, ikke vilde
være nogen Gud, og at det vilde være bespotteligt at tænke,
at Gud er saaledes, som vi kunne tænke ham. Mysteriet
er ikke, som Religionerne mene, et relativt, men et absolut
Mysterium. Naar den religiøse Bevidsthed har hævet
sig til dette Punkt, vil den være forsonet med Videnskaben.

Det gaaer nemlig med de høieste videnskabelige Ideer,
% --- p. 128 ---
\setcounter{footnote}{0}%
\opage{128}som med de høieste religiøse Ideer; de ere selvmodsigende,
naar de betragtes som absolute. Hvad Kraft eg Materie, % PRINTED AS IS (eg: O-F16: for og)
Rum og Tid ere i sig selv, kunne vi ikke vide Noget om.
Paa intet Punkt naa vi en Afslutning. Vi kunne finde
almindelige Love, men disse vise stedse tilbage til endnu
mere almindelige Love, og saa fremdeles. Vor Tænkning
er, som Hamilton har vist, betinget ved Modsætning og
Forskjel; hvad der ingen Modsætning har og ikke kan betegnes
ved sin Forskjel fra Andet, kunne vi ikke tænke.
Paa den anden Side forudsætter al Erkjendelse ogsaa et
Lighedsforhold; vi erkjende en Ting ved at henføre den til
en anden lignende; men det Absolute, det i sig selv Bestaaende
kunne vi ikke sammenligne med noget Andet end
det selv. Endelig er al vor Erkjendelse subjektiv, bestaaer
i vore Bevidsthedstilstandes Overensstemmelse med de ydre
Phænomener; derfor er den nødvendigvis ogsaa relativ, naaer
ikke Tingenes Væsen i sig selv.

Det Spørgsmaal fremstiller sig da: er vor Bevidsthed
indskrænket til de blotte Phænomener, til det Endelige og
Relative? Rent logisk set synes Sagen at forholde sig saaledes,
at det Absolute, som Hamilton siger, er et rent negativt
Begreb. Men for en nærmere Betragtning stiller
Sagen sig anderledes. Der maa dog være Andet i Bevidstheden
end selve Bevidsthedens Grændser og Betingelser; der
maa være Noget, som begrændses og betinges. Enhver
bestemt udformet Bevidsthedstilstand omfatter kun et Relativt;
men den forudsætter et ubestemt, uformet Indhold,
et Ubetinget og Uendeligt, som er det nødvendige Materiale
for Tankens Arbeide. Selve Begrebet om det Relative forudsætter
som nødvendigt Korrelat det Absolutes Realitet;
ellers vilde det Relative jo være absolut. Enhver virkelig
Tanke er endelig, men viser ud over sig selv; ethvert
Existerende, vi tænke, viser ud over sig selv. Enhver ny
Form, Tanken vil give det Absolute, falder for Kritiken,
men det Absolute skaber sig saa atter andre Former, og
saaledes i det Uendelige. Dersom det Absolute ikke var
i Bevidstheden, vilde Tanken mangle Indhold, ligesom den
% --- p. 129 ---
\setcounter{footnote}{0}%
\opage{129}vilde mangle Form, naar kun det Absolute var. Der er
Noget, som er betinget, Noget, som staaer i de relative
Forhold, og som vi kun kunne udtrykke ved successivt at
ophæve alle Betingelser, alle Grændser. Vi komme da til
den positive skjøndt ubestemte Bevidsthed om det, som ligger
ud over al bestemt Bevidsthed, om Existents i og for sig,
om en Substants bag alle Phænomener, et Værende bag alle
Tilsyneladelser. Denne Bevidsthed er det nødvendige Modstykke
til vor Selvbevidsthed.

»Foruden den \emph{bestemte} Bevidsthed, hvis Love Logiken
formulerer, er der ogsaa en \emph{ubestemt} Bevidsthed,
der ikke kan formuleres. Foruden fuldstændige Tanker, og
foruden de ufuldstændige Tanker, der kunne fuldstændiggjøres,
er der Tanker, som det er umuligt at fuldstændiggjøre
(gjennemføre), og som dog ogsaa ere reale i den Forstand,
at de ere normale Fornufttilstande» (normal affections
of the intellect). --- »Da vor Bevidsthed om det
Ubetingede bogstavelig er det Samme som den ubetingede
Bevidsthed eller det raa Tankemateriale, som vi ved Tænkningen
give bestemte Former, saa følger, at en stedse
nærværende Fornemmelse (sense) af real Existents er selve
vor Aands Grundlag.»

Religion og Videnskab høre derfor nødvendig sammen,
som to hinanden supplerende Bevidsthedsformer. Vor Bevidsthed
om Naturen er fra en Side set Religion, fra en
anden Side set Videnskab. Videnskaben angriber ethvert
enkelt overnaturligt Faktum, som Religionen statuerer, men
samtidig indrømmer den, at dens egne Resultater kun ere
relative, ikke naa det absolute Mysterium.

Men ere vi da ikke nødte til at fremstille det Absolute
for os under en eller anden Form? --- Auguste Comte
antog, at det theologiske Begreb om et Forsyn og det
metaphysiske Begreb om Naturen havde udspilt deres Rolle,
og at det Ideal, der nu stod tilbage, var den positive Videnskabs:
at føre alle enkelte Fakta tilbage til et eneste
universelt Faktum. Herimod indvender Spencer, at den
høieste religiøse, metaphysiske og videnskabelige Ide i sig
% --- p. 130 ---
\setcounter{footnote}{0}%
\opage{130}selv er én og den samme. Naar Religionen har frigjort
sig for al Anthropomorphisme, bliver den Troen paa et absolut
Væsen, der iboer alle Phænomener; og Videnskaben
fører os tilsidst til Anerkjendelsen af en universel Kraft, én
Existents, der som universel og absolut er os ubegribelig.
Og naar Comte mente, at den religiøse Følelse herefter vil
have Menneskeheden til sin Gjenstand, saa strider dette
mod Loven for Religionens Udvikling, der viser en stadig
Fremgang fra det Endelige henimod det Uendelige; Bevidstheden,
der har hævet sig til Ideen om det Ubetingede, kan
ikke atter finde sin høieste Gjenstand i et endeligt Væsen.\footnote[1]{The classification of the sciences. p. 35-36. 41.}
--- Ifølge Hamilton og Mansel er det en Pligt, som med
Nødvendighed udspringer af vor Natur, at tænke sig Gud
som personligt Væsen. Men Pligt kan det hverken være
at bekræfte eller at benegte Guds Personlighed. Vor Pligt
er kun at erkjende vor Begrændsning og ikke gaa ud over
vor Tankes Skranker. Der kan ikke være nogen Strid
mellem vore intellektuelle Evner og vore moralske Forpligtelser.
Naar vi vægre os ved at tro paa Guds Personlighed,
da er det ikke fordi det er for høi en Bestemmelse,
men netop fordi det er for lav, endelig og menneskelig en
Bestemmelse af det Ubetingede. Skulde der ikke kunne
gives en Tilværelsesform, der stod lige saa høit over Bevidsthed
og Villie som disse over den rent mekaniske Bevægelse?
Vi kunne vistnok ikke danne os noget Begreb om
en saadan Tilværelsesform; men det er ikke en Grund til
at betvivle dens Existents; --- vi have jo set, at det Betingede
ikke kan fatte det Ubetingede. Og dog vil Tanken
maaske stedse vedblive med sine Forsøg paa at finde Udtryk
og bestemte Former for det absolute Væsen, der er dens
egen Baggrund og Basis. Der vil heller ikke være noget
Urigtigt deri, saa længe saadanne bestemte Former udtrykkelig
anerkjendes som blotte Symboler uden indre Overensstemmelse
med det, de betegne. Under den bestandige
% --- p. 131 ---
\setcounter{footnote}{0}%
\opage{131}Uddannelse af saadanne Symboler og derpaa følgende Forkastelse
af dem som utilstrækkelige holdes Bevidstheden
levende om Modsætningen mellem det Ubetingede og det
Betingede. De positive Religioners Tilhængere klamre sig
fast til de endelige Former for det Ubetingede, som om de
kunde være absolute, og ere derfor mindre fromme end
Tænkeren, som véd, at alle slige Former kun ere utilstrækkelige
Symboler. Og selv om han staaer ene med
denne Overbevisning, véd han dog, at den udvikler sig med
Nødvendighed og vil have Fremtiden for sig. »Han maa
erindre at lige som han (med alle sine Evner, Bestræbelser
og Overbevisninger) er et Barn af Fortiden, saa er han
Fader til Fremtiden, og at hans Tanker ere som Børn, der
ere fødte ham, og som han ikke maa lade dø af Mangel
paa Pleie. Han kan, som ethvert andet Menneske, med
Rette betragte sig som et af de utallige Organer, gjennem
hvilke den ubekjendte Aarsag virker, og naar den ubekjendte
Aarsag frembringer en vis Overbevisning i ham, saa bemyndiges
han derved til at udtrykke den i Tale og Handling.
Thi --- for at anvende Digterens Ord i deres høieste
Betydning --- »»Naturen kan ikke forbedres ved noget
Middel, uden at jo Naturen selv har frembragt dette Middel;
over den Kunst, som du siger føier Noget til Naturen, staaer
der en Kunst, som Naturen selv frembringer.»» Derfor vil
den Vise ikke betragte den Overbevisning, der rører sig i
ham, som tilfældig. Han vil uden Frygt udtale den høieste
Sandhed, han kan se, fordi han véd, at hvad saa Følgen
bliver, spiller han derved sin rette Rolle i Verden; --- fordi
han véd, at \emph{kan} han bevirke den Forandring, han stræber
efter, godt! --- kan han \emph{ikke,} ogsaa godt, skjøndt ikke
\emph{saa} godt!»\footnote[1]{\emph{Spencer}: First Principles. 3d ed. p. 123.}

Man ser, at Spencer ikke gjør sig Illusioner med
Hensyn til den Forsoning mellem Religion og Videnskab,
han mener at have fundet. Det er ogsaa klart, at de po%
% --- p. 132 ---
\setcounter{footnote}{0}%
\opage{132}sitive Religioner ikke kunne gaa ind paa denne Forsoning
uden netop at opgive det, der gjør dem positive, nemlig
Troen paa en overnaturlig, under visse bestemte Former
aabenbaret Magt. Ifølge Spencer er det Absolute den Forudsætning,
til hvilken al Undersøgelse af Bevidsthedens og
Tilværelsens Væsen med Nødvendighed fører tilbage. Det
staaer ikke i dualistisk (overnaturlig) Sondring fra Phænomenerne,
men i immanent Forhold til dem, er »den ubekjendte
Aarsag,» »den universelle Kraft,» som ligger til
Grund for dem. Derfor er det ikke klart, hvorledes Spencer
kan vedblive at betegne det Absolute som det Uerkjendelige;
det vilde det kun være, dersom Phænomenerne ikke stod i
indre Forhold til det, saa at de ikke kunde afgive i det
Mindste indirekte Bidrag til Opfattelsen af det. Konsekvent
maatte Spencer opfatte Phænomenerne som relative Former
og Udtryk for det Absolute, som Indholdsbestemmelser, i
hvilke dette udtrykker sit Væsen, saaledes at vi gjennem
dem kunne vinde en real, skjøndt betinget og relativ Erkjendelse
af det Absolute. Dette viser sig klart, naar vi
fremdrage Spencers Bestemmelse af Philosophiens Opgave
og første Forudsætninger.

Ved Philosophi forstaaer han den til Enhed og Helhed
udviklede Erkjendelse (completely-unified knowledge). Den
søger de almindelige Love, som indbefatte de særlige Love,
der ere de enkelte Videnskabers Gjenstande. Idet Philo
sophien overalt søger Enhed og Overensstemmelse, forudsætter
den Gyldigheden af den Bevidsthedsfunktion, hvorved
vi opfatte Tingene som hverandre lige eller ulige. Endog
den mest skeptiske Tanke maa anerkjende denne Gyldighed;
thi Tvivlen selv kan kun hævdes ved Hjælp af hin Bevidsthedsfunktion,
som ligger til Grund for enhver Slutning,
baade Benegtelse og Bekræftelse. Men foruden denne
Funktion selv forudsættes ogsaa et Resultat, hvortil den
har ført. Vore Forestillinger sondre sig i to store Grupper,
som ikke kunne gjøres til én: en Række af livlige Forestillinger,
som vi sammenfatte under Navn af Objekt og
henføre til ydre Gjenstande, og en Række svagere Fore%
% --- p. 133 ---
\setcounter{footnote}{0}%
\opage{133}stillinger, som vi sammenfatte under Navn af Subjekt eller
Jeg. Denne Sondring er Resultatet af en Række Bevidsthedsakter,
der statuere Ligheder og Uligheder mellem Forestillingerne.
Det Subjektive og det Objektive, Jeg og Ikke-Jeg
staa saaledes som Grundmodsætninger indenfor det
Relative, som de to Hovedformer for det Uerkjendeliges
Aabenbaring. Den relative Realitet, vor Erkjendelse omfatter,
maa vi nødvendigvis tænke i indre Forbindelse med
den absolute, for os utilgængelige Realitet De Grundbetingelser,
under hvilke Realiteten fremtræder for os, ---
Kraft og Materie, Tid og Rum, Bevægelse --- Former, som
ligge til Grund for al Opfattelse og Erfaring, maa vi derfor
tilskrive real Gyldighed. Vel have vi ingen Ret til at antage,
at de have absolut Realitet, ere Former for det Absolute
selv; men da disse Former for os have en ubetinget
og uforanderlig Nødvendighed, saa maa vi ogsaa antage,
at den ubekjendte, absolute Kraft virker uforanderlig paa
os: uden en saadan Uforanderlighed vilde hine Former være
utænkelige. I sidste Instants kunne alle hine Former føres
tilbage til Kraften: Materie og Bevægelse ere kun forskjellig
betingede Ytringer af Kraft, og Rum og Tid ere de Betingelser,
hvorunder disse Ytringer fremtræde for os. Kraft
er ogsaa den egentlige og oprindelige Gjenstand for al Bevidsthed.
Bevidsthedens Grundkarakter bestaaer jo i Opfattelse
af Lighed og Ulighed, altsaa i en Sammenligning,
men denne forudsætter Forskjel mellem Indtryk; vi opfatte
altsaa det, der bevirker Forandring, --- men vi kunne kun
opfatte det ved Hjælp af den Fornemmelse, vi have af vor
egen Kraft til at bevirke Forandringer. Denne Bevidsthed
om Kraft bliver altsaa det sidste Punkt, hvortil Analysen
fører os. Det er naturligvis ikke den absolute Kraft, vi
her opfatte, men kun en betinget Virkning af den ubetingede
Aarsag.

Kraftens Bestaaen er nu den nødvendige Forudsætning
for al Videnskab og Philosophi Saa længe der tænkes,
saa længe Bevidstheden fungerer, saa længe maa der ogsaa
finde Overgange og Forandringer Sted: men derved forud%
% --- p. 134 ---
\setcounter{footnote}{0}%
\opage{134}sættes da, at Kraften bestaaer, at Materien ikke tilintetgjøres,
at Bevægelsen stadig fortsættes. Holdt Kraften op
at virke, blev Materien til Intet, standsede Bevægelsen absolut,
saa vilde vi faa en Bevidsthedsakt, som bestod i et
Forhold mellem Noget og Intet, ɔ: et Forhold, hvor det ene
Led manglede, ɔ: et Forhold som ikke er noget Forhold,
en Selvmodsigelse. En absolut Ende kunne vi lige saa lidt
tænke som en absolut Begyndelse, ja, vi kunne ikke tænke
os nogen Bevidsthed, for hvem dette kunde forholde sig
anderledes. Enhver enkelt Form af Kraft kunne vi tænke
os omsat i en anden, men selve Kraften kunne vi ikke ophæve.
Det er ikke nogen enkelt, ingen relativ Kraft, der
bestaaer, men Kraften i det Hele. Vi vinde da her et
nyt Synspunkt for den Sandhed, som ovenfor hævdedes, at
vi nødvendigvis maa forudsætte et Ubetinget som Baggrund
og Bærer for alt Relativt og Phænomenalt. Den Kraft,
hvis Bestaaen vi lære, er den absolute Kraft, om hvilken
vi have en ubestemt Bevidsthed som det nødvendige Korrelat
til den Kraft, vi kjende. At hævde Kraftens Bestaaen
er altsaa kun en anden Form for Hævdelsen af en ubetinget
Realitet uden Begyndelse eller Ende.

Hvad der saaledes følger af vor Bevidstheds Væsen,
a priori, bekræftes a posteriori ved den moderne Physiks og
Chemis Resultater, der fastsætte Kraftens Vedvaren og
Materiens Indestruktililitet som Erfaringskjendsgjerninger. % PRINTED AS IS (Indestruktililitet: O-F23)
Ja, ethvert enkelt Experiment forudsætter egentlig denne
Kjendsgjerning; thi saa snart der skal veies eller maales,
maa der lægges en Krafténhed til Grund, som maa antages
konstant; men dersom Kraften ikke bestod, vilde man ikke
have nogen Sikkerhed for denne Enheds Konstants, og
al Maaling, al exakt Experimenteren, vilde altsaa være
umulig.

Den eneste Sandhed, som gaaer  ud over Erfaringen,
idet den ligger til Grund for Erfaringen, er Kraftens Bestaaen.
Som Erfaringens Grundlag maa den ogsaa være
Grundlag for enhver videnskabelig Organisation af Erfaringen.
% --- p. 135 ---
\setcounter{footnote}{0}%
\opage{135}Til den fører en gjennemført Analyse os tilbage, og paa den
maa en rationel Synthese bygges.

Af denne Grundforudsætning om Kraftens Bestaaen er
det da igjen en nødvendig Følge, at der maa være en fast
Lovmæssighed i Naturen. Naar der i to givne Tilfælde
finder fuldstændig Lighed Sted mellem de forudgaaende
Betingelser, saa maa Følgerne ogsaa være de samme; ---
ellers vilde vi komme til den Konsekvents, at der enten
maatte være opstaaet en hel ny Kraft, eller at en Kraft
havde ophørt at existere; men i begge Tilfælde vilde vi
komme i Strid med Sætningen om Kraftens Bestaaen, som
afskjærer absolut Begyndelse og absolut Ende.\footnote[1]{Man ser let, at denne Forudsætning ligger til Grund forMills „fire Methoder for den experimentale Forsken“ (ovenfor p. 28 f.).} % PRINTED AS IS: „forMills“

Ligeledes følger det af Kraftens Bestaaen, at de forskjellige
enkelte Kræfter kunne gaa over i hverandre indbyrdes,
saaledes at en vis Grad af den ene Kraft er ækvivalent
med en vis Grad af den anden Kraft. Den enkelte
Krafts tilsyneladende Ophør vil kun sige, at den omsættes
i en anden Form. Saaledes kunne mekanisk Bevægelse,
Varme, Lys, chemisk Tiltrækning, Elektricitet etc. omsættes
indbyrdes, ja de kunne omsættes til Fornemmelse, Følelse
og Tanke, ligesom disse kunne omsættes til hine. \emph{Hvorledes}
denne Omsætning sker, vide vi ikke. Hvorledes
Svingninger i Luften kunne fremkalde Lydfornemmelser,
chemiske Forandringer i Hjernen aandelige Tilstande, det
er Mysterier, vi ei kunne fatte, men ikke dybere Mysterier
end Omsætningen mellem de physiske Kræfter indbyrdes.
Disse Omsætninger ere ikke blot empiriske Fakta; det følger
af selve Bevidsthedens Grundlov, at enhver Kraftytring maa
være en Følge af en foregaaende Kraftytring, og at den ved
sit Ophør maa gaa over i en anden Art af Kraftytring;
hvis ikke, ophæves Sætningen om Kraftens Bestaaen og
med den selve Bevidstheden. --- Paa lignende Maade kan
det bevises, at Bevægelsen stedse maa gaa i den Retning,
hvor der ydes mindst Modstand, og at Bevægelsen paa
% --- p. 136 ---
\setcounter{footnote}{0}%
\opage{136}Grund af, at forskjellige Kræfter samtidig virke, maa blive
rhythmisk ---

Hvad vi hidtil have fundet, er kun de abstrakte Love,
som gjælde for de enkelte Faktorer. Men den Enhedserkjendelse,
som Philosophien søger, fuldendes først, naar vi vende os
til de konkrete Phænomener, der svare til og kunne udledes
af hine abstrakte Love. De konkrete Phænomener i sig
selv har Philosophien ikke at gjøre med; de undersøges i
de specielle Videnskaber. Men naar Geologen har forklaret
os Jordens Historie, naar Chemikeren har lært os Stoffernes
Kombinationer og Forandringer at kjende, naar Physiologen
har fundet Loven for Organismens Væxt, saa opstaaer det
Spørgsmaal, om der ikke er Noget, som er fælles for alle
disse forskjellige Processer, om der ikke kan paavises Elementer,
som komme igjen i alle disse konkrete Phænomeners
Historie.

De konkrete Phænomener have Materie og Bevægelse
til deres Elementer. Loven for deres Historie maa derfor
være en saadan, der nærmere angiver Formen for Materiens
og Bevægelsens bestandig fornyede Fordeling Nu undergaaer
enhver Existents, vi kjende, to forskjellige Processer.
Den bliver til for os derved, at Elementer, som tidligere
befandt sig i spredt og forvirret Tilstand, samle sig til et
Hele, idet Materien koncentreres og de enkelte Dele tabe
deres Bevægelse. Dette er, hvad vi kalde Udvikling, som
altsaa bestaaer i Opsamlen af Materie og Adspredelse af
Bevægelse. Den anden Proces, Phænomenet undergaaer, er
en Opløsning, som omvendt bestaaer i Absorberen af Bevægelse
og Adspredelse af Materien. Udvikling og Opløsning
betegne de to Hovedperioder i ethvert Phænomens
Historie.

Udviklingen, hvis almindelige Karakter er angivet som
Opsamling af Materie og Adspredelse af Bevægelse, maa
nu bestemmes nærmere ad Induktionens og Erfaringens Vei.
Vi maa i Astronomien og Geologien, i Biologien, Psykologien
og Sociologien søge de speciellere Eiendommeligheder
ved Udviklingen, saaledes som den hersker gjennem
% --- p. 137 ---
\setcounter{footnote}{0}%
\opage{137}hele Universet, gjennem alle Aabenbarelser af den uerkjendelige
Grundkraft.

Den første af disse nærmere Bestemmelser bestaaer i,
at Udviklingen gaaer fra en mindre sammenhængende til en
mere sammenhængende Tilstand. Ifølge Laplaces Hypothese
har Solsystemet fra en uhyre Taagemasse fortættet sig til
de enkelte Kloder, hvoraf det nu bestaaer; og enhver af
disse Kloder er selv undergaaet en Fortætningsproces.
Enhver Plante voxer ved at koncentrere Elementer i sig, der
før existerede i adspredt Tilstand; ethvert Dyr voxer
ved atter at koncentrere Elementer i sig, der før existerede
spredt i Planter eller andre Dyr. Den menneskelige
Historie viser os en stedse nøiere Sammenslutning mellem
Slægter og Stammer, indtil der nu synes at danne sig et
civiliseret Samfund, som omfatter hele Jorden; og sammenligne
vi en civiliseret Stat med en Flok Vilde, ville vi strax
se, hvilken Forskjel der er i Koncentration.

Men Overgangen fra adspredt til koncentreret Tilstand
er tillige en Overgang fra en ensartet til en uensartet Tilstand.
Den oprindelige Taagemasse blev under sin Koncentration
efterhaanden artikuleret: der dannede sig Ringe,
som afløstes fra Hovedmassen, og som igjen deltes i Kloder,
indtil vort nuværende Solsystem med alle dets indre Forskjelligheder
uddannede sig. Jordkloden var først en glødende
Masse; saa størknede Skorpen, og der blev en Modsætning
mellem den og den indre flydende og glødende Kjerne; paa
Jordens Overflade udviklede sig Forskjelle mellem Land og
Hav, Bjerge og Sletter, klimatiske Forskjelligheder etc. Fra
sin første ensartede Spire udvikler Organismen sig til en
rig Mangfoldighed af Dele og Organer. Det menneskelige
Samfund viser os det samme Phænomen i, hvad man
kalder Arbeidets Deling; oprindelig maatte Enhver selv
frembringe Alt, hvad han behøvede; efterhaanden ere de
forskjellige Arbeider blevne fordelte, saa Enhver faaer sin
Bestilling.

Overgangen fra det Ensartede til det Uensartede medfører
en Overgang fra det Ubestemte til det Bestemte. Over%
% --- p. 138 ---
\setcounter{footnote}{0}%
\opage{138}alt træder Orden i Forvirringens Sted. Tydeligst ses dette
i Organismens Udvikling. Dens Bygning og dens Organer
blive bestandig mere udformede og artikulerede. Fra en
geléagtig Masse udvikler den sig til et bestemt begrændset
Hele med skarpt udprægede Dele.

Føie vi nu hertil, at det ikke blot er Materien, men
ogsaa Bevægelsen, der undergaaer disse Forandringer, ---
parallelt med Organismens Bygning udvikler sig saaledes
dens Funktioner, --- saa faa vi følgende fuldstændige Definition:
Udvikling er en Opsamlen af Materie og samtidig
Adspredelse af Bevægelse, i Løbet af hvilken Materien gaaer
over fra en ubestemt og usammenhængende Ensartethed
til en bestemt og sammenhængende Uensartethed, og i
Løbet af hvilken Bevægelsen undergaaer en parallel Omdannelse.


Erfaringen godtgjør, at der finder en saadan Udvikling
Sted overalt. Himmellegemer, Organismer og menneskelige
Samfund have i det Væsentlige samme Historie.
Men hvad Erfaringen saaledes lærer os, maa vi give en
dybere Begrundelse; vi maa paavise Udviklingens Aarsager
(naturligvis relative og phænomenale Aarsager; den absolute,
noumenale Aarsag er os utilgængelig), og disse Aarsager
maa vi udlede fra den høieste Sandhed, der ligger til Grund
for al Erfaring og derfor gaaer ud over al Erfaring, Sætningen
om Kraftens Bestaaen. --- Det er klart, at i den
ensartede Masse, som stedse danner Udgangspunktet for
Udviklingen, maa de enkelte Dele modtage forskjellige Indvirkninger
fra Omgivelserne. Der maa først danne sig en
Modsætning mellem det Indre og det Ydre, og dernæst
mellem det Ydres og det Indres forskjellige Dele. Det
Ensartede er altid ustadigt og gaaer derfor med Nødvendighed
over til Forskjellighed. Derfor maa f. Ex. en Plante- eller
Dyreart variere paa Grund af de forskjellige Forhold,
hvorunder de forskjellige Individer leve. I modsat Fald
vilde forskjellige Kræfter have samme Virkninger, hvilket
strider mod Sætningen om Kraftens Bestaaen. --- Men naar
nu saaledes forskjellige Dele udformes, saa maa ogsaa en
% --- p. 139 ---
\setcounter{footnote}{0}%
\opage{139}og samme Kraft have forskjellige Virkninger paa de forskjellige
Dele. Kraften bliver uensartet, ligesom det Aggregat
den virker paa. Og det vil igjen sige, at en og samme
Aarsag kan have mangfoldige Virkninger; Virkningen er
altid mere sammensat end Aarsagen. Lyset modificeres ved
de Virkninger, det fremkalder ved at falde paa forskjellige
Gjenstande, og saaledes opstaa de forskjellige Farver.
Ogsaa denne Lov følger af Kraftens Bestaaen: forskjellige
Ting vil jo ikke sige Andet end Ting, som reagere paa forskjellig
Maade; hvis altsaa en og samme Kraft ikke modificeres
paa forskjellig Maade ved de forskjellige Ting, vilde
en (reagerende) Kraft være tilintetgjort, hvilket er umuligt.
--- En tredie Aarsag til Udviklingen maa søges i Udsondringen
af det Ligeartede. Ligesom Vinden samler de visne
Blade, saaledes virker i det Hele taget en og samme Kraft
ensartet paa det, som i alle Henseender er indbyrdes ligt.
Ting og Personer, som ere samme Betingelser underkastede,
ville ogsaa komme til at ligne hverandre. Herved bevirkes,
hvad Darwin kalder den naturlige Udvælgelse, som omfatter
de Væsener, der ere i Stand til at kæmpe Livets Kamp
uden at bukke under. Beslægtede Væsener bringes derved
sammen, og større Bestemthed i Form og Udtryk blive
mulige. ---

Efterat nu Udviklingens Væsen og Aarsager ere paaviste,
reiser det Spørgsmaal sig: vedbliver Udviklingen
uafbrudt i det Uendelige, eller finder den sin Afslutning i
en bestemt vedvarende Tilstand? Enhver Bevægelse vedvarer,
indtil der er opnaaet Ligevægt, og jo større Modstanden er
udenfra, des hurtigere opnaaes Ligevægten. Naar den ydre
Modstand ikke er betydelig, kan den indre Bevægelse
mellem Delene indbyrdes vedvare, saa at det Hele befinder
sig i en bevægelig Ligevægt (eqvilibrium mobile). Men
enhver saadan bevægelig Ligevægt har dog en Tendents til
at gaa til absolut Ligevægt. De ydre Kræfter høre aldrig
ganske op at virke, selv om deres Virken er forholdsvis
umærkelig. --- Solsystemet kan siges at være i bevægelig
Ligevægt, idet enhver Bevægelse eller Perturbation opveies
% --- p. 140 ---
\setcounter{footnote}{0}%
\opage{140}ved en Bevægelse i modsat Retning. Planetbevægelsen kan
betragtes som en Svingning mellem Aphelion og Perihelion.
Men af Formindskelsen i den Enckeske Komets Omløbstid
har man som bekjendt sluttet, at Himmelrummet er opfyldt
af en Æther, der hemmer Klodernes Bevægelser og
tilsidst maa kunne standse dem; da ville alle de mindre
Kloder igjen samles til en uhyre Masse. Endnu vigtigere
er det dog, at Solen engang vil have udstraalet al sin Varme
i Verdensrummet; og da Solvarmen er Kilden til al Bevægelse
og alt Liv, vil dette medføre en fuldstændig Hvile,
en absolut Ligevægt. --- Hvad Menneskelivet angaaer, da
vil der engang komme det Tidspunkt, da den menneskelige
Aand har arbeidet sig frem til fulstændig Harmoni med Naturen. % PRINTED AS IS (fulstændig: O-F29)
Er dette Maal end uendelig langt borte, saa følger det
dog af Bevægelsens og Livets Natur, at det er det, henimod
hvilket Bevægelsen gaaer; med andre Ord, Udviklingen
maa ende med Tilveiebringelsen af den høieste Fuldkommenhed
og Lyksalighed. ---

Ogsaa dette følger af Grundloven om Kraftens Bestaaen.
Da vi ikke kunne tænke os Kraften i dens absolute
Enhed, men kun under Form af Vexelvirkning mellem forskjellige
Kræfter, saa foregaaer al Bevægelse under Modstand.
Det Bevægede afgiver derfor bestandig en Del af
sin Bevægelse til det, der gjør Modstand, det Medium,
gjennem hvilket det bevæger sig; tilsidst maa det da have
mistet al sin Bevægelse.

Naar Ligevægten er opnaaet, gaaer Udviklingen over
i Opløsningen. Medens Udviklingen bestaaer i Opsamlen
af Materie og Adspreden af Bevægelse, bestaaer Opløsningen
omvendt i Opsamlen af Bevægelse og Adspreden af Materie.
At Opløsningen indtræder, følger af de ydre Kræfters fortsatte
Virken. Naar Ligevægt er opnaaet, er der jo ikke
Kraft tilovers til at møde disse ydre Kræfter med. Derfor
se vi overalt, baade i de menneskelige Samfunds og i den
enkelte Organismes Historie, baade i Chemi, Geologi og
Astronomi, at en Udvikling paa sit høieste Punkt gaaer
over i Opløsning.

% --- p. 141 ---
\setcounter{footnote}{0}%
\opage{141}Have vi da her naaet den absolute Ende, og vil Universet
tilsidst stivne i en evig Død, blive en uhyre Dynge
af udslukte Solsystemer? Det følger af sig selv, at vi ikke
kunne finde noget absolut Svar paa et saadant Spørgsmaal.
Vi kjende kun Udviklingen indenfor vor Del af Universet
og vide ikke, om der gives Stjernesystemer overalt
i det uendelige Rum, saa at nye Bevægelser stedse kunne
opstaa ved Indvirkning fra dem. Det Uendelige kunne vi
ikke konstruere. Holde vi os derimod til Grundloven om
Kraftens Bestaaen, da er det en nødvendig Antagelse, at
Bevægelsen aldrig kan ophøre. Ligevægten og Opløsningen
kunne kun danne Indledningen til en ny Udviklingsperode. % PRINTED AS IS (Udviklingsperode: unattested-word sweep, zoom z10)
At Bevægelsen er rhythmisk paa Grund af Kræfternes Flerhed,
er en Lov, der ogsaa gjælder for Universet i dets
Helhed. I én Periode har den tiltrækkende Kraft Overhaand,
Masserne koncentrere sig, og Bevægelsen spredes; i
en anden Periode have de frastødende Kræfter Overhaand,
Bevægelsen koncentreres, og Masserne spredes. Det vil
altsaa sige, at Udvikling og Opløsning vexle i det Uendelige
Udviklingens Resultat er de Betingelser, der fremkalde Opløsningen,
og Opløsningens Resultat er de Betingelser, der
fremkalde Udviklingen. Forud for vor Verden ligger der
altsaa en uendelig Række af Udviklings- og Opløsningsperioder,
og efter vor Verdens Ophør vil der følge en ny
Uendelighed af saadanne skiftende Perioder. Principet i
disse er det samme, skjøndt de specielle Resultater aldrig
ville være de samme.

Dette er den Slutning, hvortil vor begrændsede Erkjendelse
fører os. Lige saa lidt som vi kunne tænke os
nogen Grændse for Tid og Rum, lige saa lidt kunne vi
tænke os nogen Grændse for den uerkjendelige Kraft, som
aabenbarer sig i Universet. Selv uendelig maa den virke i
det Uendelige. --- Da vi ikke kjende det Absolute, kjende
vi heller ingen absolut Nødvendighed, men kun relativ Nødvendighed,
Nødvendighed for vor Tanke. Lige saa lidt
kjende vi det indre Forhold mellem Phænomenerne og den
absolute Kraft, som ligger til Grund for dem (the connection
% --- p. 142 ---
\setcounter{footnote}{0}%
\opage{142}between the phenomenal order and the ontological order).
Alle vore Grundbegreber (Materie, Bevægelse, Kraft) ere kun
Symboler for en ubekjendt Realitet. Naar der derfor er
blevet reist den Anklage, at Udviklingsphilosophien er materialistisk,
fordi den fører Alt, selv Livets, Aandens og det
menneskelige Samfunds Phænomener, tilbage til Materiens,
Bevægelsens og Kraftens almindelige Love, saa oversér man
netop, at selve disse Grundbegreber ere Symboler for det Absolute,
der aabenbarer sig i alle Phænomener, i de saakaldte
materielle saa vel som i de aandelige, og ikke kan finde sit
fuldstændige Udtryk i nogen enkelt Klasse af Phænomener.
Holder man dette fast, da taber Striden mellem Spiritualismen
og Materialismen sin Betydning, og man vil indsé,
at Udviklingsphilosophien ikke drager det Høiere ned ved
sine Forklaringer, men netop omvendt hæver det saakaldte
Lavere. ---

Spencers »First Principles» er et af de mest epokegjørende
Værker i vort Aarhundrede. Intetsteds finder man
som her alle Ideer, alle videnskabelige Resultater sammenarbeidede
til en fast, organisk Enhed. Rigdommen i det
induktive Grundlag, Skarpheden og Sikkerheden i den deduktive
Begrundelse og i den psykologiske Analyse, som danner
Forudsætningen saa vel for Induktionen som for Deduktionen,
ere lige beundringsværdige. Spencer staaer for saa vidt paa
Empirismens Standpunkt, som han fastholder, at det høieste
Punkt, vi kunne naa, er den mest almindelige phænomenale
Lov, Erfaringen lærer os kjende. Men han standser ikke,
som Mill og Comte, ved empiriske Almensætninger; han
fører de empiriske Love tilbage til en eneste Lov, Loven
om Kraftens Bestaaen, og søger saa at udlede denne Lov
af selve Tankens og Bevidsthedens Væsen. Derved kommer
der en indre Sammenhæng mellem de forskjellige Led i
hans Anskuelse. Philosophien som Enhedsopfattelse af Tilværelsen
maa gaa ud fra en oprindelig Forudsætning,  der
ikke kan søges andet Steds end i selve Tankens, Erkjendelsens
Væsen. Denne Forudsætning og Nødvendigheden af at føre
al anden Erkjendelse tilbage til den har ingen Tænker i
% --- p. 143 ---
\setcounter{footnote}{0}%
\opage{143}vor Tid fastholdt og gjennemført med en saadan Energi
som Spencer.

I det Enkelte forelaa der mange Forarbeider. I philosophisk
Henseende er Spencer særlig paavirket af Hamilton.
Han optager dennes Lære om Erkjendelsens Relativitet, men
paaviser, at Relativiteten med Nødvendighed forudsætter et
absolut Princip, hvis Aabenbarelse og Symbol den er. Ligeledes
har han optaget og gjennemført Hamiltons Lære om
Kausalitetsprincipet som grundet i selve vor Bevidstheds
Væsen. Den Sætning, at Alt har sin Aarsag, var ifølge
Empirismen hentet fra Erfaringen; men Spencer viser, at
den kun er en anden Form for Loven om Kraftens Bestaaen
og lige som denne har sin Grund i vor Bevidstheds inderste
Væsen. Han skiller sig fra Hamilton ved den nøie Forbindelse,
han saaledes fra først af sætter mellem Physikens
og Psykologiens Resultater. Hamilton var streng Spiritualist
og havde ingen Sands for Naturvidenskaben; for
ham vilde en saadan Synthese, hvorved der opstilles en
fælles Grundlov for Aandens og Naturens Verden, staa som
Materialisme og Atheisme, og Mansel har da ogsaa anklaget
Spencer for at have presset Hamiltons Lære i denne Retning
ved kun at antage den negative (kritiske) Del af hans
Philosophi og udelade den positive (moralsk-religiøse, eller
rettere: theologiske). Ifølge Mansel er Spencers Lære en
Kombination af Pantheisme og Positivisme (Philos. of the
cond. p. 39---40). Men Sandheden er, at Spencer ved
konsekvent Gjennemførelse af de Hamiltonske Principer er
kommen til sin eiendommelige Lære, der trods den Vaklen,
vi allerede have paavist i den hvad et afgjørende Punkt
angaaer, dog ikke er en eklektisk Kombination, men en
helstøbt Udvikling af en dybt begrundet Hovedtanke. Man
kunde maaske sammenligne hans Forhold til Hamilton med
Spinozas til Cartisius. % PRINTED AS IS (Cartisius: O-G03)

Foruden philosophiske Forarbeider danne ogsaa Naturvidenskabens
Resultater et nødvendigt Grundlag for hans
Anskuelser. Som Momenter af særlig Vigtighed fremhæve
vi: Laplaces Theori om Solsystemets Udvikling; Læren om
% --- p. 144 ---
\setcounter{footnote}{0}%
\opage{144}Naturkræfternes Overgang til hverandre indbyrdes, saaledes
som den klarest er fremstillet af Joule og Helmholtz; v. Baers
Definition af Livet som en Udvikling fra det Ensartede til
det Uensartede; Darwins Hypothese om Arternes Oprindelse.
Dog havde Spencer tildels selv foregrebet flere af disse
Resultater, og i ethvert Tilfælde bestaaer hans Originalitet
i Bearbeidelsen og Sammenføiningen af alle disse Enkeltheder
til et i sig afsluttet Hele.

I visse Henseender vilde man kunne sammenligne
Spencers Udviklingsphilosophi med Hegels System, hvis
Grundtanke er den dialektiske Bevægelse fra det Umiddelbare
gjennen Modsætningen (Reflexionen) til Ideen. Som % PRINTED AS IS (gjennen: O-G06)
hos Spencer Udviklingen begynder med en ensartet Tilstand,
der paa Grund af sin Ensartethed er i en ustadig Ligevægt,
hvis Ophævelse fører til Artikulation og Individualisation,
saaledes ligger Umiddelbarheden ifølge Hegel aaben for
Modsigelsen, som med Nødvendighed opløser den i Reflexion,
indtil Spliden og Sondringen hæves i Ideen (Spencers »bevægelige
Ligevægt»). Men hvad der hos Hegel er en abstrakt
logisk Proces, hvor Virkelighedens Indhold er presset
ind i Tankebevægelsens Former, det er hos Spencer en real
Naturproces, som omfatter det virkelige Univers i alle dets
enkelte Dele. Vel fører ogsaa Spencer Udviklingens Ide
tilbage til Bevidsthedens Væsen, men denne metaphysiske
Begrundelse danner kun Forudsætningen, afgiver ikke en
tvingende Form for det reale Indhold. Medens Hegels Methode
er absolut deduktiv, støtte hos Spencer Induktion og
Deduktion hinanden, idet Induktionens Resultater deduceres
ud af de sidste Forudsætninger for al Erkjendelse og Bevidsthed.
--- Hegel havde fundet Motivet til sin Lære i en
Analyse af Tankebevægelsen, som viste tre forskjellige Stadier,
af hvilke det ene med logisk Nødvendighed førte til det
andet. Hos Spencer er det derimod, som vi have set, en
ethisk-historisk Ide, der har afgivet Motivet; denne Ide
bragtes saa i Forbindelse med Livets Ide i Almindelighed,
og derved banedes Veien til videre Opfattelse og dybere
% --- p. 145 ---
\setcounter{footnote}{0}%
\opage{145}Begrundelse, hvorved det ethisk-historiske Synspunkt kun
blev et Led indenfor en større Sammenhæng.

Hovedmodsætningen mellem Spencer og Hegel beroer
dog paa, at denne mente at give en Theori, en Begrebsudvikling
af det Absolutes Væsen, medens Spencer kun
giver sin Anskuelse som en Hypothese, der vel efter hans
Overbevisning er den eneste holdbare, fordi den er den
eneste, der paa engang stemmer med Bevidsthedens Grundlov
og med Erfaringens Resultater, --- men som dog ikke
kan lære os Noget om det Absolute i sig selv, da al vor
Erkjendelse kun angaaer det Relative og Phænomenale.
Ifølge Hegel udvikler det Absolute sig i Alt, kommer til
sig selv gjennem Verdensudviklingen; ifølge Spencer tør vi
ikke anvende Udviklingens Ide paa det Absolute, da enhver
Begrebsbestemmelse af hvad der ligger ud over Phænomenerne,
indvikler os i Selvmodsigelse.

Vi staa her atter ved Hovedvanskeligheden hos Spencer.
Naar det Relative og Phænomenale erkjendes som Former
for det Ubetingede,\footnote[1]{First Principles. p. 278: „The sphere of knowledge is co-extensive with the phenomenal --- co-extensive with all modes of the Unknowable that can affect consciousness.“} saa maa Erkjendelsen af disse Former
ogsaa føre til Væsensbestemmelser for det Ubetingede selv.
Det vilde netop være selvmodsigende, dersom de Bestemmelser,
der gjælde for Formerne, ikke havde i det Mindse en relativ % PRINTED AS IS (Mindse: O-G08)
Gyldighed for Væsenet selv, der finder sit Udtryk i disse
Former. »Man stiller det Uendelige overfor, over, hinsides,
udenfor det Endelige, ja anbringer endog en uhyre Skillevæg,
for at de ikke skulle komme hinanden for nær. Men
de ligge slet ikke udenfor hinanden. Det Endelige er det
Uendeliges Indhold. Derfor er det Uendelige heller ikke
ufatteligt; det han fattes i sit Indhold; kun omfattes kan
det ikke.» (Fechner). Udviklingen kan derfor ikke være en
nødvendig Grundidé indenfor det Endelige uden ogsaa at
have Gyldighed for det Uendelige. Som vi allerede have
sagt, det som for os er den høieste Nødvendighed, er maaske
% --- p. 146 ---
\setcounter{footnote}{0}%
\opage{146}kun en relativ Nødvendighed, men dog en virkelig Nødvendighed,
som vi maa lade staa, og som kun Phantasien
kan sætte sig ud over. Det Ubetingede vilde ikke fortjene
sit Navn, dersom dets Væsen ikke var Liv og Udvikling;
det er en Sandhed, som de metaphysiske Undersøgelser i
de idealistiske Systemer have ført til. Og den Udvikling,
der maa tilskrives det Absolute, kan ikke være en abstrakt,
ideel og tidløs, som hos Hegel, den Tænker, som med størst
Energi har hævdet Anvendelsen af Udviklingens Begreb paa
det Absolute; anerkjender man, med den efterhegelske Philosophi,
at Tid, Rum og Materie ikke kunne sættes udenfor
det Absolute, men at »Naturbegrebet knytter Verdensbevidstheden
til Gudsbevidstheden» (Weisse), saa maa dets
Udvikling ses som real, dets Fuldkommenhed ikke som Afsluttethed
men som ubegrændset Fremadskriden.\footnote[1]{Philosophien i Tydskland. p. 216 f. 262 f.} Ganske
vist mangle vi Midler til en Erkjendelse af det Absolutes
Udviklingsproces, og Hegel kunde kun ved en spekulativ
Illusion mene at give en saadan; men det kan ikke fritage
os for at gaa den Vei, den logiske Konsekvents fordrer,
især da Konsekventsen her er af saa stor Betydning for
hele Verdensanskuelsen.

Vi have allerede tidligere set, at Spencer ikke selv
fastholder denne Distinktion, naar han tillægger det Absolute
Bestemmelserne Kraft og Aarsag, som dog kun kunne hentes
fra det Relative. Men han gaaer endnu videre. Af de tre
religiøse Hovedanskuelser, hvis Uholdbarhed han mener at
have godtgjort, kan han egentlig kun siges at staa i afgjort
Modsætning til Atheismen og den vulgære Theisme:
han hævder i Mødsætning til den første, at den religiøse % PRINTED AS IS (Mødsætning: O-G10)
Bevidsthed er i sin Ret, hvad det Centrale angaaer, og
overfor den anden, at dette Centrale ikke kan finde sit
Udtryk i Personlighedens Form, som for ham er en endelig
Form.\footnote[2]{Vi ville senere komme til at omtale, at Spencer bestemt sætter Udvikling i Modsætning til Skabelse, saa at hin gjør denne utænkelig.} Forsaavidt kan han kaldes. Pantheist, skjøndt
% --- p. 147 ---
\setcounter{footnote}{0}%
\opage{147}han ikke ynder saadanne dogmatiske Udtryk. Og det maa
tillige fastholdes, at det Absolute for ham ikke er et Upersonligt,
men et Overpersonligt; han har set, at Debatten
mellem Pantheisme og Theisme har godtgjort det Urigtige
i Modsætningen mellem det Personlige og det Upersonlige,
naar Talen er om det Absolute.\footnote[1]{Anf. Skr. p. 161---165; 213---214; 279; 313; 324---325.}

I sit Ungdomsskrift opfattede Spencer Udviklingen
teleologisk og saa Harmoni og Lykke som Indholdet af
den guddommelige Ide, der skulde virkeliggjøres gjennem
den. Ogsaa senere fastholder han de velgjørende Resultater,
hvortil Udviklingen maa føre; men da hans Synspunkt nu
er objektivt og universelt, saa at den menneskelige Udvikling
kun bliver en Episode indenfor det store Verdensdrama (eller
Dramaer), saa kan der ikke længere være Tale om en saadan
Teleologi. Derfor har man klaget over det Trøstesløse i
denne Anskuelse, over at Spencer --- efterat have vist os,
hvorledes enhver Udviklingsproces ender med Ligevægt og
Opløsning, hvorefter nye Processer følge i det Uendelige, ---
»forlader os ængstelige, bøiede over Randen af den evige
Afgrund, idet vi med vor Tanke udspørge det uendelige
Mørke.»\footnote[2]{\emph{Caro} i Revue de deux mondes. 15 octbr. 1873 p. 774.} Dersom nogen Philosophi, nogen Videnskab
kan svare paa de Spørgsmaal, som her reise sig, kan Svaret
kun gaa ud fra den Erkjendelse, hvorpaa vi overfor
Spencer have lagt saa megen Vægt, at det Relative staaer
i et indre Forhold til det Absolute; den Værdi, Udviklingen,
dens Kampe og dens Resultater have fra det relative,
menneskelige Synspunkt, kan ikke gaa tabt i det
Absolute. ---

Fra Spencers »almindelige Philosophi» gaa vi nu til
hans »specielle Philosophi,» af hvilke kun Biologien og
Psykologien forelige fuldt udarbeidede. Dog ville vi i hans % PRINTED AS IS (forelige: O-G15)
»Essays,» i »Social Statics» og i den under dette Skrifts
Trykning udkomne Bog »The study of sociology» finde
Bidrag til hans ethiske og sociologiske Anskuelse. ---


% --- p. 148 ---
\setcounter{footnote}{0}%
\opage{148}Spencers hele System har paa en Maade sin Spire i
den nyere Physiologis Lære om Livet som en Udvikling fra
Ensartethed til Forskjellighed, en Artikulation og Individualisation.
Denne Lære anvendte han først paa det
menneskelige Samfundsliv og dernæst paa hele Tilværelsen,
og vandt derved en dybere og alsidigere Bestemmelse af
Udviklingens almindelige Love. Disse Love maa nu atter
finde deres klareste Udtryk i de levende Væseners Udvikling.
Livet bestaaer i en bestemt formet Kombination af uensartede
Forandringer, baade samtidige og successive, hvorved
de indre Forhold bringes i Overensstemmelse med de ydre.
Livet er i sin Almindelighed en Tillempelse af det Indre
efter det Ydre. Paa de livløse Gjenstande gjøre Omgivelserne
deres Indtryk uden at disse Indtryk have nogen
videre Følge; i ethvert levende Legeme foregaaer derimod
en sekundær Forandring som Følge af det modtagne Indtryk,
en Forandring, hvorved det levende Væsen bliver i
Stand til at møde en følgende Forandring i Omverdenen.
Herved sker Tillempelsen efter de ydre Forhold; Organismen
foretager indre Forandringer for at kunne møde nye Forhold.
Livet forudsætter altsaa et Vexelforhold mellem det
Indre og det Ydre. For at det levende Væsen skal kunne
bestaa under de ydre Forhold i deres Mangfoldighed og
Forskjellighed, maa der udvikles et System af tilsvarende
indre Forskjelligheder. Funktion, Virksomhed er saaledes
Hovedmomentet i Livets Ide. Organisk Funktion gaaer
forud for organisk Struktur. Funktionen udvikler Strukturen;
Strukturen giver paa den anden Side Funktionen
dens bestemtere udprægede Form. Et Exempel paa Liv
uden egentlig Organisation have vi i de saakaldte Rhizopoder,
som ikke vise nogen Adskillelse af Dele eller Organer, men
alligevel æde, sove og bevæge sig. De bestaa af en mikroskopisk
Slimmasse som idelig skifter Form, idet den
skyder sig ud i fine, uregelmæssige Fryndser, som snart
trække sig tilbage igjen; disse Bevægelser tjene baade til
Dyrets Ernæring og til at bringe det frem i Vandet. Her
% --- p. 149 ---
\setcounter{footnote}{0}%
\opage{149}have vi i den mest elementære Form den Afpassen af indre
Forhold efter ydre, hvori Livet bestaaer.

Hvad det levende Væsen optager fra Omverdenen,
omdanner det til en for det eiendommelig Form. Naar
Organismen er bleven beskadiget, har den til en vis Grud % PRINTED AS IS (Grud: for Grad; blind read b1-50-10)
Evne til at gjøre sig hel igjen. Hos lavere Dyr kan der
strax udvoxe et nyt Lem af samme Form som det tabte.
Forplantningen og Nedarvelsen af bestemte Egenskaber vise
os en lignende Evne til at forme Delene efter et vist
Mønster. Allerede den uorganiske Materie har en saadan
Evne. De materielle Atomer have deres bestemte Betingelser,
under hvilke de komme i Ligevægt. Jo mere sammensat
Molekulen er, des ustadigere er dens Ligevægt.
Og da nu de organiske Molekuler ere de mest sammensatte
af alle, saa er Ligevægten her den mest ustadige. En
ringe Indvirkning er nok til at bringe Forstyrrelse. Men
ligesom vi allerede hos Krystallet finde en Evne til at
forme de Atomer, der optages eller bringes ud af Ligevægt,
paa en med det hele krystalliserende Aggregats Struktur
stemmende Maade, saaledes maa vi antage en lignende
Evne hos Organismen. Naar Benet rives af et Firben,
fremvoxer strax Spiren til et nyt, som efterat have gjennemgaaet
samme Udviklingsproces som det tidligere, antager
en lignende Form. Her virker Organismens Struktur bestemmende
paa Udviklingen paa det enkelte Punkt. Men
dette kan kun ske, naar Formen for enhver Art af Organisme
bestemmes ved en Eiendommelighed i dets Elementer
(de physiologiske Enere), idet disse, ligesom de uorganiske
Atomer, have en særegen Struktur, efter hvilken de stræbe
at ordne sig. Den Kraft, ifølge hvilken uorganiske Atomer
ophobes i en for dem eiendommelig Krystallisationsform,
kalde vi deres Polaritet; i Analogi hermed kunde man kalde
den tilsvarende Evne hos de levende Legemers Elementer
deres »organiske Polaritet.»

Disse Elementer, de physiologiske Enere, ere paa den
ene Side forskjellige fra de chemiske Enere, der sammensætte
de organiske Legemer; thi disse ere omtrent de samme
% --- p. 150 ---
\setcounter{footnote}{0}%
\opage{150}hos alle Arter og forklare altsaa ikke Forskjellene mellem
disse; --- paa den anden Side fra de morphologiske Enere
eller Cellerne; thi dels er Cellen ikke nogen universel Form
(Rhizopoderne ere saaledes ikke cellulære), dels forudsætter
selve Cellens Tilblivelse netop hin organiske Polaritet. De
physiologiske Enere ere meget mere sammensatte end de
chemiske, saa at en ringe Forandring i deres Sammensætning
vil medføre en Forandring i den Form, som hele Aggregatet
antager. Herved bliver det muligt, at de ydre
Paavirkninger kunne tjene til det organiske Livs Udvikling.
De physiologiske Eneres store Plasticitet giver Indtrykkene
deres Magt, medens deres Polaritet bestemmer den Maade,
hvorpaa disse Indtryk bearbeides og formes, den Retning,
i hvilken de komme til at virke. De laveste organiske
Phænomener ere endnu aldeles uden Organisation; man kan
hos dem hverken tale om Form eller Type. Deres Skikkelse
forandrer sig fra det ene Øieblik til det andet, alt efter
som de ydre Vilkaar føre det med sig. En absolut Begyndelse
af det organiske Liv kunne vi ikke paavise. Chemiens
Resultater paa den ene Side og Biologiens paa den anden
Side synes at berettige os til den Antagelse, at de organiske
Sammensætninger efterhaanden have udviklet sig,
og at de efterhaanden have vundet større og større
Fasthed.

Dersom den Erstatning, det organiske Legeme formaaer
at skaffe sig for de Tab, dets Funktioner volde, kun gik
lige op mod det Tabte, vilde der ingen Udvikling eller
Lempelse efter Omverdenen ske. Men naar Funktionen og
det Tab, den volder, ikke overskrider visse Grændser, mere
end erstattes det Tabte ved hvad der optages. Bevægelse
og Virksomhed er baade for hele Organismen og for dens
enkelte Dele gunstigere for Ophoben af Cellevæv, end
Hvile er. Det er en bekjendt Sag, at Musklerne voxe ved
Øvelse. Og da nu de forskjellige Funktioner staa i Forbindelse
med hverandre, saa medfører en Stigen af den ene
ogsaa en Stigen af den anden. Muskelbevægelsen forudsætter
saaledes Blodomløbet, og begge maa derfor udvikles sammen.
% --- p. 151 ---
\setcounter{footnote}{0}%
\opage{151}Men heraf følger tillige, at Lempelsen efter Omverdenen
har sine Grændser, navnlig hos det udvoxede Individ, og
at der hører Meget til, at den skal forplante sig til Efterkommerne.
Dette sker kun, naar der er dannet en ny
Tilstand af Ligevægt paa Grund af Forandringerne. Da
kan Summen af alle tidligere Paavirkninger og Erfaringer
faa Overmagten over de enkelte, enestaaende Paavirkninger,
og den Modifikation, som først kun var en temporær og
individuel Lempelse, bliver faststaaende og arvelig.

Ad denne Vei maa Oprindelsen af nye Arter forklares.
Vel kan det ikke direkte bevises, at Arterne ere opstaaede
ved fortsatte Lempelser og disse Lempelsers Virkninger paa
Strukturen (den organiske Polaritet); men vi vide dog saa
Meget, at der paa denne Maade sker indgribende Modifikationer,
dannes Varieteter af samme Art, der kunne afvige
meget fra hverandre; ja, man kan ofte ikke afgjøre, om en
vis Forskjellighed konstituerer en Artsforskjel eller kun betegner
en Varietet. Dette Svævende i Artens Begreb viser
os, at det kun er et relativt Begreb.

Det følger derhos af den foregaaende Udvikling, at
Arterne maa antages opstaaede under de levende Væseners
Vexelvirkning med Omverdenen. Den Antagelse, at en Art
ikke skulde forandre sig ved i Millioner af Aar at være
udsat for forandrede Paavirkninger, vilde være en Negtelse
af Sætningen om Kraftens Bestaaen. Alle Organismer ere
bestandig udsatte for astronomiske, geologiske, meteorologiske
og organiske Paavirkninger. Enhver ny Færdighed, ethvert
Fremskridt i Organisation, som erhverves, medfører Muligheden
af nye Paavirkninger, saa at Forandringerne maa
stige i geometrisk Progression. Men Forudsætningen er
bestandig den, at Forandringerne ikke ere for store eller
for pludselige. Kun saadanne Forandringer udvikle, der
ikke ophæve den almindelige Harmoni mellem de paavirkende
Kræfter og Organismens Funktioner.

Spencer er saaledes ad sin egen Vei kommen til det
Resultat, Darwin hævdede i sit berømte Værk om »Arternes
Oprindelse.» Vel er Spencers Biologi senere end dette Værk,
% --- p. 152 ---
\setcounter{footnote}{0}%
\opage{152}men han havde allerede tidligere, i forskjellige Afhandlinger
(se hans »Essays») og i første Udgave af sin Psykologi,
hævdet Udviklingshypothesen, hvorfor Darwin ogsaa i Fortalen
til sit Værk nævner ham blandt sine Forgængere. Dog er
der stadig en Forskjellighed i Synspunktet mellem de to
Forskere. Darwin lægger som bekjendt Hovedvægten paa
det saakaldte Kvalitetsvalg, ɔ: paa den Kjendsgjerning, at
de organiske Væsener, som ikke kunne tage Kampen op
med Omverdenen, maa gaa under, og kun de Væsener bevares
og forplante sig, som ere stærke nok, men navnlig
de, som paa en eller anden Maade have erhvervet sig et
Fortrin fremfor de andre. Dette kalder Darwin ogsaa den
naturlige Udvælgelse: den Stærkeste lever, de Svage gaa
under. Derved opretholdes og udvikles Harmonien mellem
Livet og dets Betingelser. Men ifølge Spencers Anskuelse
har Darwin lagt for overveiende en Vægt paa denne indirekte
Udvikling. Selv tilskriver han den direkte Udviklingsmaade den største Betydning: Funktionernes Virken uddanner
Strukturen, og denne fortsatte Lempelse gaaer efterhaanden
over i det levende Væsens Natur, idet der danner sig en
ny Art af Ligevægt (organisk Polaritet) mellem de physiologiske
Elementer. Saalænge Organismen endnu er ufuldkommen
og passiv, spiller vistnok Kvalitetsvalget, den indirekte
Frembringelse af Ligevægt (indirect eqvilibration),
den største Rolle. Men efterhaanden som Aktiviteten udvikler
sig, gaaer den indirekte Proces over i den direkte.
Kvalitetsvalget vedbliver naturligvis at spille en stor Rolle,
men bliver mere og mere en underordnet Faktor. Naar
der har udviklet sig et Nerve- og Muskelsystem, og dermed
en Evne til at variere Handlingerne efter Omstændighederne,
saa bliver den direkte Frembringelse af Ligevægt (direct
eqvilibration) den vigtigste. I de civiliserede Menneskeracer
endelig rammer den indirekte Proces kun dem, Samfundsordningen
\emph{endnu} ikke formaaer at beskytte.

Forholdet mellem den direkte og den indirekte Proces
kan i Korthed angives saaledes: En ydre Paavirkning vil
dræbe Organismen, hvis denne er for svag, --- og i saa Fald
% --- p. 153 ---
\setcounter{footnote}{0}%
\opage{153}virker den indirekte Udviklingsmaade; --- hvis Organismen
derimod er stærk nok, vil den lempe sig efter Indtrykket,
og da virker den direkte Udviklingsmaade, idet Lempelsen
fører til en ny Ligevægtstilstand indenfor Organismen.“

Denne Theori, Evolutionstheorien, stiller Spencer i
skarp Modsætning til den sædvanlige Anskuelse, Kreationstheorien,
ifølge hvilken hver Art er skabt for sig efter en
fast Type. Spencer underkaster denne sidste Anskuelse en
knusende Kritik. For det Første indeholder Forestillingen
om en Skabelse en i sig selv umulig Tanke. Den statuerer
en absolut Begyndelse, sætter et Forhold mellem Intet og
Noget, et Forhold, som ikke er et Forhold. Af selve vor
Tankes Væsen og Grundlov følger saaledes Skabelseshypothesens
Ugyldighed. De, der tænke sig hver Art skabt for
sig, gjøre det uden virkelig at tænke, uden at oversætte
Ordene i Tanker; de tro det i Virkeligheden ikke, men tro,
at de tro det; thi Tro forudsætter, at man ideelt konstruerer
Sagen for sig, men det er her umuligt. Dernæst vilde det
følge af denne Hypothese, at Skaberen havde fastholdt den
anlagte Type, selv hvor denne var uhensigtsmæssig, og hvor
andre Former vilde have passet bedre. En og samme
Dyreklasses Medlemmer leve jo under mange forskjellige
Forhold. Saaledes er Sælhunden et Pattedyr, men lever i
Vandet. Antager man nu, at den er skabt til sit nuværende
Opholdssted, saa er det meningsløst, at dens Lemmer have
samme Grundform som Landpattedyrenes, da de dog skulle
være Svømmeredskaber, hvorimod det bliver let forklarligt,
naar den nedstammer fra et Pattedyr, der lidt efter lidt
blev Vanddyr i sin Levevis. Endelig vilde det følge af
Skabelseshypothesen, at visse Arter skabtes til at skade andre,
og at Smerte og Nød var en Del af den guddommelige
Plan. Ifølge Evolutionshypothesen frembringes der derimod
under Udviklingens Gang langsomt men sikkert en høiere
Grad af Lykke, idet alle Onder kun ere forbigaaende. Der
maa stedse tilveiebringes en større Overensstemmelse mellem
de levende Væsener og den ydre Verden. Ganske vist kan
man spørge, hvorfor alle Onder, al Nød og Smerte ikke
% --- p. 154 ---
\setcounter{footnote}{0}%
\opage{154}kunde være undgaaede; men der kan ikke længer være Tale
om, at de tilføies med Overlæg.

Evolutionshypothesen lærer os lige saa lidt som Kreationshypothesen
den egentlige Aarsag at kjende; Striden
mellem dem angaaer kun Maaden, hvorpaa den ubekjendte
Aarsag virker. Vi finde nu paa ethvert Omraade, at den
anthropomorphistiske Opfattelse af den ubekjendte Aarsag
forsvinder; man holder efterhaanden op med at udlede Phænomenerne
af en guddommelig Villies umiddelbare Indgriben,
skjøndt Mange, som paa andre Omraader have opgivet Troen
paa en saadan Indgriben, endnu fastholde den ved Spørgsmaalet
om Arternes Oprindelse; --- og dog har den her ikke
bedre Hjemmel end andre Steder. --- Udviklingshypothesen
er paa engang fremtraadt paa mange forskjellige Omraader:
indenfor Astronomien, Geologien, Physiologien, Politiken og
Philologien. Ligesom ifølge Laplace Solsystemet gjennem
en af naturlige Aarsager betinget Udviklingsproces har naaet
sin nuværende Skikkelse, saaledes har ifølge Lyell de endnu
daglig virkende Aarsager gjennem en Række umærkelige
Forandringer frembragt Jordoverfladens nuværende Udseende.
Physiologien ser ikke mere i Spiren en færdig Model af den
modne Organisme; man sander i Politiken Coleridges Ord,
at Forfatninger ikke laves, men voxe frem; og man betragter
ikke længer Sprogets Oprindelse som en overnaturlig
Akt. ---

Nu gives der vistnok et Standpunkt midt imellem
begge Hypotheser; det er det, hvorpaa \emph{Auguste Comte}\footnote[1]{Cours de philosophie positive. 2e éd. III p. 390 f.}
og med ham mange naturvidenskabelige Forskere have stillet
sig, som holde sig til det Faktiske, at en Overgang fra en
Art til en anden endnu ikke er paavist, og som derfor
slet ikke ville opstille nogen Theori om Arternes Oprindelse.
Nylig har \emph{Charles Robin} med stærkt Eftertryk
gjentaget Comtes Indsigelse mod saadanne Hypotheser.\footnote[2]{Anatomie et physiologie cellulaires. Paris 1873. p. XXIV: „Naar man ser bort fra en antividenskabelig Indblanding af theolo%
\opage{155}giske Forudsætninger ved de biologiske Problemer, saa er det ingenlunde Mangel paa philosophisk Dannelse hos vor Tids Forskere, som i Frankrig formindsker Antallet af dem, der slutte sig til Lamarcks og hans Efterfølgeres Lærdomme; ... Grunden er tvertimod den udviklede philosophiske Erkjendelse, at Gyldigheden af en spekulativ Anskuelse maa bedømmes efter Muligheden af en positiv Kontrol, og Manglen af et virkeligt Bevis paa dette Punkt, idet indtil nu et saadant Bevis ikke er givet for en eneste Klasse af Væsener, saa lidt som for en eneste Art af Celler ... Endnu er der til Bevis kun anført Muligheder, om hvilke faa Naturforskere ere enige, og ikke reelle Fakta.“ --- En anden af Comtes berømteste Disciple, \emph{Littré}, stiller sig mere imødekommende, men dog forsigtig afventende overfor Udviklingshypothesen (La science au point de vue philos. p. 539-559).} ---
% --- p. 155 ---
\setcounter{footnote}{0}%
\opage{155}Spencer betragter denne skeptiske og afventende Holdning
som en Tilstand af »ustadig Ligevægt.»  Uvilkaarlig maa
Tanken gaa i den ene eller i den anden Retning: Arternes
Oprindelse maa jo dog \emph{enten} være naturlig \emph{eller} overnaturlig,
et Tredie gives ikke. Og naar man i det Hele
hylder en rationel Verdensanskuelse, følger det af sig selv,
hvilket Led af dette Dilemma man maa gribe. Naar Comte
erklærer Problemet om Arternes Oprindelse for Frugten af
unyttig Spekulation, da udtaler Spencer i Modsætning hertil,
at »den Del af Biologien, som beskjæftiger sig med
Arternes Oprindelse er den vigtigste, i Forhold til hvilken
alle andre ere underordnede. Thi paa Biologiens Afgjørelse
af denne Sag beroer vor Opfattelse af den menneskelige
Natur i Fortid, Nutid og Fremtid, vore Theorier om den
menneskelige Aand og det menneskelige Samfund.»\footnote[1]{The classification of the sciences. p. 40.} ---

Alle morphologiske og physiologiske Phænomener maa
nu kunne finde deres Forklaring ifølge Udviklingshypothesen.
De organiske Former kunne ikke længere betragtes som
Typer, der have foresvævet Skaberen, eller som Midler, han
har villet anvende til Individets Vel; de maa betragtes som
nødvendige Resultater af Vexelvirkningen mellem Individerne
og Omverdenen. Den mekaniske Opfattelse træder altsaa i
Stedet for den typiske og den teleologiske. Spencer opstiller
for at godtgjøre denne Opfattelses Mulighed en Række de%
% --- p. 156 ---
\setcounter{footnote}{0}%
\opage{156}taillerede, tildels overordentlig skarpsindige Hypotheser, som
vi her ikke kunne dvæle ved. Hans Hovedtanke er, paa
Basis af den sammenlignende Anatomi og Physiologi, at
paavise en trinvis Udvikling parallelt med Betingelsernes
Udvikling. Selv de høieste og mest sammensatte Strukturer
og Funktioner maa kunne forklares ad denne Vei. Allerede
det strukturløse Protoplasma besidder f. Ex., i lavere Grad,
den samme Egenskab, som Nerven hos de høiere Dyr besidder.
En Rhizopodes geléagtige Masse og et ubefrugtet
Æg vise Bevægelser, som forudsætte, at et Indtryk er
blevet forplantet fra den ene Del af Massen til den anden;
og gjennem en Polyps nerveløse Krop se vi en ved Berøring
af en Føletraad fremkaldt Sammentrækning langsomt forplante
og sprede sig. Ved fortsat Funktion maa da det
særlige Organ for denne Funktion kunne udvikles.

Livets Grad beroer paa den Grad, i hvilken Overensstemmelsen
mellem Organismen og Omverdenen er omfattende
og nøiagtig. Idet der udvikles et System af indre Forskjelle,
svarende til de ydre, bliver Organismen stedse mere
sammensat; men ogsaa Forbindelsen mellem dens Dele blive
efterhaanden inderligere. I en saa udviklet Organisme som
den menneskelige kan et ringe Indtryk frembringe en
Mængde forskjellige Virkninger, kan fremkalde Muskelbevægelser
og Sammentrækninger og vække Ideer, Følelser og
Beslutninger. ---

Som der ikke kan paavises noget bestemt Punkt, hvor
det Uorganiske gaaer over til det Organiske, saaledes kunne
vi heller ikke statuere nogen fast Grændse mellem de physiologiske
og de psykologiske Phænomener. Begge Klasser
høre ind under Biologien og følge Livets almindelige Love.
Ligesom Geologien er en speciel Del af Astronomien, og
Biologien af Geologien, saaledes er Psykologien igjen en
Underafdeling af den almindelige Biologi. Den adskiller sig
fra de Videnskaber, hvorpaa den hviler, derved at den ikke,
som Physiken, betragter de ydre Phænomener i og for sig,
eller, som Physiologien, de indre i og for sig, men undersøger
Forbindelsen og Forholdet mellem det Indre og det
% --- p. 157 ---
\setcounter{footnote}{0}%
\opage{157}Ydre. Den objektive Psykologi studerer de Funktioner, ved
hvilke Organismen er i Stand til at lempe sig efter de ydre
Begivenheder og Forhold, og den Maade, paa hvilken disse
Funktioner fremtræde i Individernes Handlinger. Den
subjektive Psykologi staaer i Modsætning til alle andre Videnskaber.
Den undersøger selve Sandsningerne, Forestillingerne,
Følelserne og Villiesytringerne, som ledsage den
synlige Lempelse af de indre Forhold efter de ydre, og betragter
disse Bevidsthedsphænomener i deres Oprindelse og
indbyrdes Forbindelser. Dens Gjenstand, Bevidstheden selv,
og dens Methode, Analyse af selve Jeget, hvorved man ene
kan finde Lovene for Forandringerne i Bevidstheden, stiller
den i samme Forhold til alle andre Videnskaber, som det
Subjektive staaer i til det Objektive. Psykologien (baade
som objektiv og som subjektiv --- thi denne er en Forudsætning
for hin) har derfor endnu større Krav paa at være
en selvstændig Videnskab end nogensomhelst anden.

Aand og Bevidsthed er Noget, der slet ikke har Lighed
med andre Ting, er noget i sit Slags Enestaaende. Aandens
Substants er det, som bestaaer under alle Forandringer i
Bevidstheden, og hvis Modifikationer alle disse Forandringer
ere. Derfor har Hume Uret, naar han lærer, at der ikke
existerer Andet end Indtryk og Forestillinger. Enhver enkelt
Fornemmelse eller Forestilling kunde være fraværende,
men det, som forbinder Fornemmelserne og Forestillingerne
kan ikke tænkes borte; det er det, som opretholder Forestillingskredsens
Enhed trods alle Forsøg paa at dele den,
--- Aandens Substants i Modsætning til alle de skiftende
Former, den kan antage. Men vi kunne ikke erkjende
denne Substants. Enhver Bevidsthedstilstand er en Modifikation
af den, men i ingen Bevidsthedstilstand er den helt
tilstede. Desuden beroer al Erkjendelse paa Modsætningen
mellem Subjekt og Objekt; men her skulde det Samme
være Erkjendelsens Subjekt og Objekt, hvilket er umuligt.
Al Erkjendelse beroer fremdeles paa en Sammenligning,
men da Aanden er sui generis, have vi Intet at sammenligne
den med.

% --- p. 158 ---
\setcounter{footnote}{0}%
\opage{158}Psykologien kan  altsaa lige saa lidt gribe Aandens
Væsen i sig selv, som Chemien kan gribe Materiens Væsen
i sig selv. Dog kan den give en Forklaring af de aandelige
Phænomener, lige som Chemien af de legemlige. Som
Chemien søger de Enere, de uopløselige Atomer, hvoraf alle
Stoffer sammensættes, saaledes søger Psykologien de Bevidsthedsenere,
de sidste Elementer, paa hvis Kombination
Bevidsthedslivet beroer. Man kan danne sig en Forestilling
om en saadan psykologisk Ener (unit of consciousness) ved
at tænke sig den i Analogi med et stærkt og pludseligt
Indtryk. Ved et saadant (Slag, Lysvirkning eller Lyd) kan
man ikke i første Øieblik afgjøre, af hvad Art det er;
Fornemmelsens Beskaffenhed erkjendes kun, naar Indtrykket
har en længere Varighed. Af saadanne Stød eller Slag er
enhver Fornemmelse sammensat, ligesom en Tone er sammensat
af Luftsvingninger, som kun blive mærkelige hver
for sig, naar den Hurtighed, hvormed de følge efter hverandre,
ikke overgaaer 16 i Sekundet. Alle Forskjelle indenfor
Bevidsthedslivet beroer paa de forskjellige Maader og
Grader af disse psykologiske Eneres Kombination.

Undersøgelsen af den materielle Verden fører os tilsidst
til en objektiv Ener, et materielt Atom; Undersøgelsen
af den aandelige Verden fører os tilsidst til en subjektiv
Ener, et aandeligt Element. Spørgsmaalet opstaaer nu, hvorledes
de forholde sig til hinanden. Man vil forgjæves prøve
paa at reducere den ene til den anden. Vi kunne kun
tænke det Materielle ved Hjælp af det Aandelige. Den
mest elementære materielle Egenskab er Modstand, men den
forstaa vi kun ved at lægge vor egen Følelse og Kraft
over i de ydre Ting. Kun i vor egen Bevægelse, vor egen
Energi have vi en Maalestok for de ydre Phænomener; vi
kjende kun Tingene gjennem deres Reaktion paa os, men
denne Reaktion opfatte vi kun ved Hjælp af vor egen
Aktion. Uden Muskelspænding vilde vi ikke kunne opfatte
Materien. (Kfr. ovenfor p. 133). Dersom vi derfor havde
Valget mellem at reducere de aandelige Elementer til de
legemlige eller de legemlige til de aandelige, maatte vi vælge
% --- p. 159 ---
\setcounter{footnote}{0}%
\opage{159}det Sidste. Det er absurd at føre det Bekjendte tilbage
til det Ubekjendte; en forstaaelig Hypothese faa vi kun,
naar vi føre det Ubekjendte tilbage til det Bekjendte og
altsaa betragte de objektive, materielle Elementer som i
deres Væsen af samme Natur som vor Bevidstheds subjektive
Elementer.\footnote[1]{Denne Vei gaaer \emph{Lotze}. Se „Philosophien i Tydskland“ p. 300-302.} Men dog kan heller ei en saadan Hypothese
gjennemføres. Ligesom vi kun kunne opfatte Materien ved
Hjælp af aandelige Bestemmelser, saaledes kunne vi kun
opfatte Aanden ved Hjælp af materielle Bestemmelser. Vi
udtrykke de aandelige Phænomener ved Betegnelser, der ere
hentede fra de materielle Phænomener. Vi komme altsaa
aldrig ud over det Symbolske. Den betingede Form, under
hvilken Væsenet aabenbarer sig i Subjektet, kan lige saa
lidt som den betingede Form, under hvilken det aabenbarer
sig i Objektet, være selve det ubetingede Væsen, der ligger
til Grund for begge. Kun saa Meget se vi, at den samme
Udviklingens Lov, der gjælder indenfor den ene Række af
det Ubetingedes Aabenbaringsformer, ogsaa gjælder indenfor
den anden.

Enhver Fornemmelse maa, som vi have set, betragtes
som sammensat af elementære Indtryk. Men for at en
Fornemmelse kan danne sig, maa ikke blot de elementære
Indtryk, de psykologiske Enere, samles til en Enhed; denne
Enhed bliver kun Fornemmelse ved at være Led i en
Kjede af Bevidsthedstilstande; kun ved sit Forhold til
andre Fornemmelser, successive, samtidige eller begge
Dele, kan den enkelte Fornemmelse være til. Her
komme vi til en af de afgjørende Forskjelligheder mellem
Physiologien og Psykologien. Physiologien beskjæftiger
sig med Phænomener, der gaa for sig samtidig, medens
de Phænomener, Psykologien undersøger, danne en Række.
Overgangen fra Physiologien til Psykologien sker ved
den stigende Koncentration af Nervesystemet, hvorved de
spredte Indtryk samles og bringes i Forbindelse med hver%
% --- p. 160 ---
\setcounter{footnote}{0}%
\opage{160}andre. Kun naar der har dannet sig et Centralorgan, til
hvilket Indtrykkene kunne ledes, bliver Fornemmelse mulig.
Samtidigheden gaaer da over til Succession, og det er netop
en uafbrudt Række af Forandringer i det fælles Cemtrum, % PRINTED AS IS (Cemtrum: O-G28)
der konstituerer Bevidstheden. Dette betegner et høiere
Udviklingstrin; thi naar Indtrykkene voxe i Mængde og i
Sammensætning, kan Tillempelsen efter dem ikke ske, naar
de indtræffe paa engang; de ville da modarbeide og forvirre
hverandre. Kun naar de danne en Række, kan dette
undgaaes. Bevidstheden opstaaer altsaa i Kraft af den almindelige
Lov for Livets Udvikling: Tillempelsen af det Indre
efter det Ydre.

Den enkelte Fornemmelse er altsaa kun til ved at
sammenlignes med de andre Led i Fornemmelsernes og
Forestillingernes Række. Naar den bringes i nøiere Forbindelse
med et eller flere af disse Led, danner der sig
igjen en høiere Enhed, og saaledes i det Uendelige. Sammensætningsmaaden
er den samme gjennem Aandens hele
Bygning. Aanden bestaaer ikke i den blotte Samling af
Sandsninger og Fornemmelser; disse maa organiseres og
sondres i faste Grupper, til en af hvilke hver ny Fornemnemmelse kan henføres. Herpaa beroer, hvad vi kalde % PRINTED AS IS: „Fornemnemmelse“
Tanke og Erkjendelse. Under denne sin Udvikling frembyder
Aanden de samme Træk, som have vist sig at gjælde
for Udviklingen i det Hele. Fra den uklare Sandsefornemmelse,
beroende paa faa tilbagevendende Indtryk med
liden Forskjellighed og liden Kombination, som man finder
hos de lavere Dyreformer, der endnu ikhe have bestemt % PRINTED AS IS (ikhe: O-G29)
udformede Nerver og Nervecentre, komme vi Skridt for
Skridt, alt som Organisationen bliver fuldkomnere og Sandserne
flere, til større Kombinationer og videre gaaende
Artikulation, eller vi komme ogsaa her »fra en ubestemt
og usammenhængende Ensartethed til en bestemt og sammenhængende
Uensartethed.» De ydre Sandser (Syn, Hørelse
og Hudfølelse) staa i disse Henseender over dem, der angaa
Organismens Indre (emotions og visceral feelings); men
over al Sandsning staaer Forstanden i Klarhed og Bestemthed.
% --- p. 161 ---
\setcounter{footnote}{0}%
\opage{161}Tanken eller Forstanden er Et med Forestillingernes bestemt
udprægede Artikulation. Tankens Udvikling er altsaa
Udviklingen af Forestillingernes Forhold og Kombinationer.

Denne Udvikling maa finde Sted i de samme Henseender,
i hvilke Livets Udvikling i det Hele gaaer for sig.
Intet levende Væsen kan bestaa uden under Forudsætning
af Harmoni mellem Organismen og dens Omgivelser. Jo
høiere og rigere Livet er, des mangfoldigere og dybere er
denne Harmoni. Hvilken Forskjel mellem Indvoldsormen i
et organisk Væv og en Newton eller Shakespeare, hvis
Tanke omfatter Verden! Og dog er der en kontinuerlig
Overgang mellem disse Yderpunkter. Det laveste Trin i
Rækken er det, hvor det levende Væsen er aldeles ensartet
med det Medium, i hvilket det lever. Et høiere Trin
naaes allerede hos Zoophyterne, der trække sig sammen ved
Berøring og saaledes vise en Succession af Fornemmelse
og Bevægelse. Efterhaanden udvider Harmonien sig i
Rummet; der udfordres ikke længer umiddelbar Berøring,
for at en Bevægelse skal indtræde. Naar Lysindtrykket
koncentreres i Øiet, bliver videre gaaende Orientering mulig.
Brevduen finder sin Vei gjennem Luften; Styrmanden fører
Skibet over Verdenshavet; Astronomen udspeider Verdensrummet.
Lige saa med Hensyn til Tiden; det er saaledes,
for at tage et enkelt Exempel, et af Civilisitionens Karaktermærker, % PRINTED AS IS (Civilisitionens: O-H01)
at man tager et længere Tidsrum med i sine
Beregninger. Menneskene plante Træer, som først komme
deres Efterkommere til Gode, opføre Bygninger,  som skulle
staa i Aarhundreder og kaste i Astronomien og Geologien
deres Blik ud over umaadelige Tidsperioder. Paa lignende
Maade voxer Harmonien mellem Aanden og Naturen i
Specialitet, i Sammensætning og i systematisk Forbindelse.\footnote[1]{Spencer minder her om \emph{Leibnitz}. Forholdet mellem disse Tænkere har \emph{Ribot} angivet paa en træffende Maade. Han finder især to fælles Hovedpunkter i deres Lære. „For det Første Ideen om Kontinuiteten eller den universelle Kompenetration, der gjør, at Alt slutter sig sammen indbyrdes, at alle Ting ere indbyrdes Aarsager og Virkninger, og at den Maade, \opage{162}hvorpaa den menneskelige Aand adskiller dem, er vilkaarlig, skjøndt den er nødvendig. Denne Ide er egentlig Et med Fremskridtets Ide, er kun et andet Udtryk for denne; thi naar Alt omdannes og antager andre Skikkelser, saa hænger Alt sammen; der gives i Naturen ingen Hiatus, intet Brud paa Sammenhængen. Kun er Fremskridtets Ide snarest dynamisk, Kontinuitetens Ide derimod statisk. --- Man véd, at Leibnitz for at forklare Universet har udtænkt Monader, en Slags metaphysiske Atomer, der havde alle mulige Grader lige fra den simple Antitypi til den fuldkomne Apperception. Alt efter deres Natur konstituere de den livløse Materie eller det levende Væsen, Menneske, Dyr eller Engel. Og da Intet i Universet er isoleret, virker hele Universet paa enhver enkelt Monade, og denne udtrykker derfor Universet, bliver et Speil, som paa sin Vis reflekterer Universet. Afklæd nu denne store Ide den metaphysiske Form, og en ubestridelig positiv Sandhed vil træde frem. Stil forskjellige Væsener --- en Sten, et Træ, en Hund, en Vild, en Europæer, Newton eller Shakespeare --- i samme Medium, og ethvert af dem vil reflektere det paa sin egen Maade og efter sin Natur, det ene meget lidt, det andet i høi Grad. Der vil mellem Væsenet og dets Medium være, hvad Herbert Spencer kalder Korrespondents, og Livets Grad vil maales ved Korrespondentsens Grad, idet den ideale Livsform er den fuldkomne Korrespondents. Det Menneske, der kunde naa dette Høidepunkt, vilde fuldstændig gjenspeile Universets Virkelighed; han vilde være en med Makrokosmen adækvat Mikrokosme. Denne Ide om Korrespondents, som er et Hovedpunkt i Spencers Psykologi, staaer for mig som en Oversættelse paa den experimentale Psykologis Sprog af Leibnitz's Ord: „enhver Monade er et Speil, der reflekterer Universet.““ La psychologie Anglaise contemporaine. p. 149-150.}% the note runs over from p. 161 to p. 162 (joined by draft.py)


% --- p. 162 ---
\setcounter{footnote}{0}%
\opage{162}Da enhver aandelig Virksomhed bestaaer i en Lempelse
af det Indre efter det Ydre, og da de ydre Forhold
gjennem umærkelige Overgange voxe i Antal, Sammensætning
og Forskjellighed, saa maa Overgangene indenfor de
indre Forhold være ligesaa kontinuerlige og umærkelige.
Derfor kunne de sædvanlige psykologiske Klassifikationer kun
have overfladisk Gyldighed. Instinkt, Fornuft, Iagttagelse,
Begriben, Erindring, Indbildningskraft, Villie etc. --- disse
saakaldte Evner kunne ikke være Andet end forskjellige Grader og Former
af de indre Forholds Tillempelse efter de ydre. Den laveste Form
for sjæleligt Liv er \emph{Reflexbevægelsen}. Den staaer paa Grændsen
af det rent Physiske. Et simpelt indre Forhold reguleres her efter
et simpelt ydre Forhold; flere Reflexbevægelser kunne gaa for
% --- p. 163 ---
\setcounter{footnote}{0}%
\opage{163}sig samtidig; de foregaa derfor ikke med Bevidsthed. Naar
Reflexbevægelsen bliver mere sammensat, gaaer den over i
\emph{Instinktet.} Ved den hyppige Gjentagelse af flere
psykiske Tilstande (Indtryk og derpaa følgende Bevægelse)
danner der sig en arvelig Tendents til at foretage disse
sammensatte Funktioner. Først danne de simplere Forbindelser
sig, der ere hyppigere end de mere sammensatte.
Men jo mere sammensat Instinktet bliver, des mindre
sikkert bliver det; Funktionerne tabe da efterhaanden deres
automatiske Karakter. Flere forskjellige Funktioner ville
da under visse Forhold have en Tendents til at træde i
Kraft; og ved Kampen mellem disse Tendentser danner den
Række af Forestillinger sig, som udgjør \emph{Erindringens}
Indhold. Erindringen kan igjen gaa over til Instinkt,
nemlig naar den er tilstrækkelig fast og sikker; Exempel
herpaa haves i enhver tillært Færdighed. Lige saa gaaer
Instinktet over til, hvad vi kalde \emph{Fornuft,} uden at nogen
bestemt Grændse kan angives. Som Stuart Mill\footnote[1]{Og forresten allerede \emph{Leibnitz}. De anima brutorum § 14: Consecutiones empiricæ nobis sunt communes cum brutis, et in eo consistunt, ut sentiens ea qvæ aliqvoties conjuncta fuisse expertum est, rursus conjunctum iri exspectet. --- Men Leibnitz gjennemfører her ikke sin lex continui, idet denne empiriske Slutning skal være absolut (toto coelo) forskjellig fra Fornuftslutningen (consecutio rationalis), der beroer paa en Indsigt i Forholdets Nødvendighed.} har
vist, er den oprindelige Slutningsmaade aldeles instinktmæssig.
Naar et Phænomen fremstiller sig, slutter man
uvilkaarlig, at dets sædvanlige Egenskaber og andet Tilbehør
ogsaa maa være tilstede. Det er en Slutning fra et Enkelt
til et andet Enkelt, som allerede Dyr og smaa Børn foretage.
Herfra gaaes der successivt over til de mere indviklede
(induktive og deduktive) Slutningsformer. Den fornuftige
Tænkning gaaer lige som Erindringen ved Gjentagelse
og Øvelse ofte tilbage til det Instinktmæssige.
Saaledes tro vi umiddelbart at se Afstanden mellem Tingene,
skjøndt vi først gjennem lang Øvelse have lært at kom%
% --- p. 164 ---
\setcounter{footnote}{0}%
\opage{164}binere de forskjellige enkelte Forestillinger, hvoraf Forestillingen
om Afstand er sammensat.

Ogsaa \emph{Villien} opstaaer ved den voxende Sammensætning
af Reflexbevægelserne. Forskjellen mellem en frivillig
og en ufrivillig Bevægelse bestaaer i, at hin kun
indtræder efter at være bleven forestillet i Bevidstheden.
Paa Grund af det Sammensatte i de sjælelige Tilstande kan
Handlingen ikke fuldføres strax; en delvis Udførelse gaaer
forud for den fulde Udførelse, saa at Bevidstheden faaer
Tid til at beskjæftige sig med den og veie Mulighederne.
Den automatiske Handlens Ophør og Villiens Oprindelse
falde sammen. --- Fri er Villien kun i den Betydning, at
Handlingen udgaaer fra Individets egen Bevidsthedstilstand.
Derimod er den sædvanlige Lære om Villiens Frihed meningsløs,
fordi den forudsætter, at man har en Evne til
Handling, som kan beherske den forudgaaende Bevidsthedstilstand.
Men denne afhænger af de tidligere Erfaringer,
Individet enten selv har gjort eller har faaet nedarvet i sin
Konstitution. Dersom der var en saadan uberegnelig Faktor,
som den frie Villie efter den sædvanlige Opfattelse er, saa
vilde ikke blot Psykologien være umulig som Videnskab,
men hele den store Udvikling, vi kunne forfølge gjennem
Naturens og Aandens Verden, vilde være frugtesløs. Denne
Udvikling gaaer i Retning af Harmoni mellem det Indre
og det Ydre. Livet maa blive høiere og Lyksaligheden
større, fordi de indre Forhold bestemmes ved de ydre og
lempe sig efter dem. Men dersom de indre Forhold tildels
bestemtes ved en anden Faktor, vilde Harmonien hvert
Øieblik kunne forstyrres og Fremskridtet hindres. ---

Hvad der saaledes ad \emph{Synthesens} Vei er fremstillet,
--- Aandslivets Væxt ud af de mest uanselige Spirer
gjennem kontinuerlige Overgange, og saaledes at det helt
igjennem bestaaer af de samme Elementer, kun paa forskjellige
Udviklingstrin og i forskjellige Kombinationer, ---
det skal nu prøves ad \emph{Analysens} Vei. Efterat have set
Udviklingsgangen fra de simpleste Elementer til den høieste
Sammensætning, begynde vi nu omvendt med de høieste
% --- p. 165 ---
\setcounter{footnote}{0}%
\opage{165}Tankeformer og søge at reducere dem saaledes, at der
til Grund for dem alle viser sig at ligge et og samme
Element.

Som den høieste Form for Tænkning betragter Spencer
den \emph{kvantitative} eller mathematiske. Den behandler
de mest indviklede Forhold og dog med fuldkommen Nøiagtighed.
Den forudsætter Ensartethed hos sine Gjenstande
og undersøger kun Forholdet mellem deres Størrelser. Til
Grund for alle mathematiske Undersøgelser og Beviser ligger
det Axiom, at to Ting eller Forhold, der ere lige med et
og samme tredie, ere indbyrdes lige store. Dette Axiom
maa erkjendes ved Intuition, ved en Tankeakt, der ikke
kan reduceres til nogen simplere Form. Paa en saadan
Tankeakt bero hver for sig de successive Omdannelser af
en Ligning. Vi sætte her Ligestorheden af ingen anden
Grund end, fordi vi umiddelbart ere os den bevidste. ---
Naar de kvantitative Forhold ikke tilstede nogen exakt Bestemmelse,
staa vi ved Overgangen \emph{til den kvalitative}
\emph{Tænkning,} der angaaer Forbindelsen mellem Tingene og
Egenskaberne. Det er her ikke blotte Størrelser, der undersøges;
men der søges, hvilke almindelige Bestemmelser
der gjælde for visse Phænomener, eller hvorledes visse almindelige
Bestemmelser finde deres Anvendelse i de enkelte
Tilfælde. Formen for den kvalitative Tænkning er Slutningen,
som kun adskiller sig fra den mathematiske Proportion
ved Mangelen af de exakte Størrelsesbestemmelser
; det hidtil ubekjendte Faktum, man statuerer ved Hjælp af
en Slutning, svarer ganske til det ubekjendte fjerde Led i
en Proportion. Baade den induktive og den deduktive
Slutning bero tilsidst paa en Sammenligning mellem Forhold,
og i den primitive Slutningsform (aboriginal reasoning),
som ligger til Grund for begge, Slutningen fra et Enkelt
til et andet Enkelt, have vi en umiddelbar Henførelse af
et Phænomen til et andet Phænomen. Den Tankeakt, der
ligger til Grund, er altsaa ogsaa her en Intuition, umiddelbar
Opfattelse af Lighed.\footnote[1]{Spencer adskiller sig med Hensyn til Læren om Slutningen \opage{166}navnlig derved fra Mill, at han skarpt skjelner mellem Logiken og den psykologiske Lære om Slutningsprocessen. Hin er en Videnskab om universelle objektive Relationer, „de Betingelser, under hvilke vi kunne antage Forbindelser i Tid og Rum“, medens denne beskriver den subjektive Proces, hvorved vi erkjende hine nødvendige Forbindelser. Logiken er altsaa en Statik, medens den psykologiske Lære om Tænkningen er en Dynamik. Heraf følger det, at Spencer ikke som Mill gaaer i nominalistisk Retning, lige som han i det Hele indtager et høiere Standpunkt end Mills empirisk-subjektive. Se „Principles of psychology“ 2d ed. II. p. 87-98. The classification of sciences p. 12.}% the note runs over from p. 165 to p. 166 (joined by draft.py)


% --- p. 166 ---
\setcounter{footnote}{0}%
\opage{166}Al Tænkning, baade kvantitativ og kvalitativ, beroer
saaledes paa \emph{Klassifikation}, Henføren af beslægtede
Phænomener til samme Gruppe. Enhver Gjenkjendelse er
en Klassifikation, ja, allerede selve \emph{Opfattelsen} (\emph{Perceptionen})
af en Gjenstand er en saadan. Man kan ikke se
en Ting uden at henføre den til en vis Art. Der ligger i
selve den umiddelbare Anskuelse en Slutning: denne Ting
ser saaledes ud (f. Ex. den er rød og rund, hænger paa et
Træ af en vis Art o. s. v.), --- følgelig er det en saadan
Ting (et Æble). Slutningen er en indirekte, Perceptionen
en direkte Fastsættelse af Forhold mellem to Ting. ---
Selve Perceptionen eller Opfattelsen vindes ved Vexelvirkning
mellem Tingen og det erkjendende Subjekt. Alle
Tingens Egenskaber kunne føres tilbage til de forskjellige
Stillinger, Subjektet og Tingen staa i til hinanden. Men
disse Stillinger opfatte vi igjen kun ved Hjælp af vore egne
Lemmers Bevægelse og indbyrdes Forhold. Ved Sammenligningen
af disse Bevægelser opstaaer Rums- og \emph{Tidsanskuelserne}.
Alle Indtryk fornemmes som successive. Naar
nu to eller flere Indtryk kun kunne følge efter hinanden
i en vis Orden, have vi Tidsanskuelsen; er det derimod
muligt at vende denne Orden om, faa vi Forestillingen om
Koexistents, hvoraf Rumsanskuelsen udvikler sig. Denne
er altsaa lige saa lidt som Tidsanskuelsen en fra først af
færdig Anskuelsesform, men forudsætter Opfattelsen af Koexistents,
der igjen har Opfattelsen af Succession til sin
Forudsætning. --- Bestandig er det saaledes en Opfattelse
% --- p. 167 ---
\setcounter{footnote}{0}%
\opage{167}af Succession og Forandring, en Opfattelse af Forskjel og
Lighed mellem Bevidsthedtilstande, der ligger til Grund som % PRINTED AS IS (Bevidsthedtilstande: O-H07)
det for alle Bevidsthedsformer, fra de laveste til de høieste,
fælles Element.

Grundloven for den aandelige Udvikling er Et med
Loven for Forestillingernes Association. Den primitive og
den væsentligste Association er den mellem enhver enkelt
Fornemmelse og den Art eller Klasse af foregaaende Fornemmelser,
som den ligner og derfor henføres til. Denne
Association foregaaer aldeles automatisk. Den er ikke en
vilkaarlig Tankeakt, som kan indtræde eller ikke; men den
er Et med Gjenkjendelsen af en Fornemmelse, vi tidligere
have havt, og, som allerede vist, en Fornemmelse kan slet
ikke opstaa i Aanden uden for saa vidt den bringes
i Forbindelse med sine Forgængere af lignende Natur.
Kun naar Bestemmelsen gaaer i det Specielle, finder en
Overveielse Sted, og derved komme vi ud over det Automatiske.


Al Tænkning, lige fra den mest abstrakte og sammensatte
Slutning til den mest elementære Anskuelse bestaaer
i at statuere Ligheds- og Ulighedsforhold. Af disse er
igjen Opfattelsen af Ulighed (Forskjel, Forandring, Overgang)
den til Grund liggende. Opfattelse af Lighed forudsætter
nemlig Opfattelse af Ulighed. Naar Noget ses at
være ligt noget Andet, have vi tre Bevidsthedstilstande, af
hvilke den tredie vilde falde sammen med den første, hvis
ikke netop den anden (selve Overgangen) laa imellem.
Disse tre Tilstande kunne udtrykkes saaledes: a, x, a,
(Ulighedsforholdets tre Tilstande derimod ved a, x, b). Men
saa er der jo et Ulighedsforhold mellem a og x, et Ulighedsforhold,
der ei kan beskrives anderledes end som en
Forandring i Bevidstheden. Videre kan Analysen ikke
føres. ---

Efter at have  set, hvorledes Spencer ad Synthesens
og Analysens Vei udvikler Aandslivets Love og Elementer,
komme vi naturligvis til at opkaste det Spørgsmaal, hvorledes
han stiller sig til de to philosophiske Hovedret%
% --- p. 168 ---
\setcounter{footnote}{0}%
\opage{168}ninger og deres Forsøg paa et løse Problemet om Erkjendelsens % PRINTED AS IS (et løse: O-H04: for at)
Væsen.

Spencer stiller sig i et kritisk Forhold til begge Retninger.
Den ene betragter enten, med Leibnitz, de indre
Forhold som fra først af skabte i Harmoni med de ydre,
eller, med Kant, som subjektive Anskuelses- og Tankeformer,
hvori vi optage og hvorved vi opfatte alle Indtryk.
Men i første Tilfælde bliver Udviklingen betydningsløs, da
jo Harmonien er tilstede fra først af. De foregaaende Undersøgelser
have derimod vist os Harmonien som Resultatet
af en Udvikling, hvorved de indre Forhold lempes efter de
ydre og derved stedse antage høiere Former. Og naar
Kant antager oprindelige Anskuelses- og Tankeformer, begaaer
han den samme Feil, som den hvorpaa hele den ældre Psykologis
Ufuldkommenhed beroer, nemlig kun at betragte
det fuldt udviklede og civiliserede Menneskes Bevidsthed,
i Stedet for at forfølge Bevidsthedslivet tilbage til dets
primitive Former. Da vilde det, som vi have set, vise sig,
at Rum og Tid ikke ere oprindelige Former, men at de
lige som alle andre Anskuelser og Begreber udspringe fra
Bevidstheden om Forskjel og Lighed. Denne Bevidsthed
bliver altsaa tilbage som den eneste oprindelige Form, baade
for Anskuelse, Forstand og Fornuft. Naar Kant fremdeles
antager Erkjendelsesformerne for blot subjektive, saa at
Tingen i sig selv enten er formløs eller har Former, der
slet ikke paavirke os, saa bestemmes jo Bevidstheden
egentlig slet ikke ved Tingene, men udvikle sig rent indenfra,
og Kants Lære fører os saaledes tilbage til Leibnitz's
præstabilerede Harmoni. Begge synde mod Udviklingens
Lov, der ikke tillader nogen Sondring af det Indre fra det
Ydre, men ifølge hvilken Tanken voxer under sin Vexelvirkning
med Tingene.

For saa vidt har Empirismen Ret. Al Erkjendelse
skriver sig fra Erfaringen. De ydre Forhold frembringe de
indre. Det ses af de mest almindelige Kjendsgjerninger
: Erfaring gjør klog; Vanen er den anden Natur; Sammenhæng
og Overensstemmelse mellem Indtrykkene bevirke faste
% --- p. 169 ---
\setcounter{footnote}{0}%
\opage{169}Forbindelser mellem Forestillingerne. Men Erfaringshypothesen
i sin sædvanlige Form er ikke mindre uholdbar end
Leibnitz's og Kants Hypotheser. Dersom Aanden, som
Locke lærte, var en ubeskreven Tavle, hvorfra skulde da
Evnen komme til at organisere Erfaringerne, og hvorfra
skulde Gradsforskjellene i Henseende til denne Evne hos
forskjellige Arter af Organismer og hos forskjellige Individer
indenfor samme Art komme? Hvorfor kunde en Kat saa
ikke ved Opdragelse blive til en Hund? Og hvorfor var saa
i det Hele et Nervesystem nødvendigt? Leibnitz og Kant
have Ret i, at uden en oprindelig Organisation kan Evnen
til at ordne og bearbeide Indtrykkene og udøve tilsvarende
Handlinger ikke forklares. Erfaringerne danne kun det
Materiale, som Individet bearbeider. Endnu klarere viser
dette sig ved Følelserne. Følelser og Lidenskaber ytre sig,
førend der er gjort saa udstrakt en Erfaring, at deres Udvikling
i den kan finde sin Forklaring. Ikke at tale om de
mangfoldige Lidenskaber, der allerede ytre sig hos Barnet,
saa gaaer den mægtigste af alle Lidenskaber, Elskovslidenskaben,
absolut forud for enhver Erfaring.

Spencer bekæmper den rene Empirisme, saaledes som
især Stuart Mill konsekvent har udviklet den. Han hævder
navnlig, at Erfaringen ikke kan begrunde sig selv, men
forudsætter et Kriterium, et Princip, som staaer over Erfaringen,
et Princip, i Kraft af hvilket ethvert Skridt i
Erkjendelsens Udvikling gjøres, og uden hvilket intet kan
gjøres.\footnote[1]{Se ovenfor p. 14; 34.} Dette sidste, abstrakte Kriterium er: Umuligheden
af det Modsatte. Hvad vi ikke trods al mulig Anstrengelse
kunne forbinde i vor Tanke, maa vi forkaste. Enhver Tro
eller Antagelse er Resultatet af en Styrkeprøve mellem de
forskjellige Forbindelser af Forestillinger, en systematisk
Kamp, der tjener til at paavise, hvilke Forestillinger der
ere uopløselig forbundne. Hvad vi ikke kunne tænke adskilt,
\emph{maa} høre sammen. Og efter disse nødvendige
% --- p. 170 ---
\setcounter{footnote}{0}%
\opage{170}Tankeforbindelser bedømme vi saa alt Andet. --- De faste
og universelle Forhold mellem Tingene maa under Udviklingens
Gang danne faste og universelle Forbindelser i Organismen;
ved den uendelige Gjentagelse af absolute ydre
Ensformigheder avles der absolute indre Forbindelser, nødvendige
Erkjendelsesformer. Der kan naturligvis under Udviklingens
Gang opstaa urigtige Tankeforbindelser, der kun
ere Produkt af en begrændset og ensidig Vexelvirkning med
Tingene; men disse Forbindelser, som til en Tid kunne
synes uopløselige, opløses under den videre gaaende Udvikling.
Det rokker ikke Principet: der \emph{gives} uopløselige
Tankeforbindelser. Men et andet Spørgsmaal er, \emph{hvilke}
de ere. --- De uopløselige Tankeforbindelser udtrykke Netto
Resultatet af al Erfaring lige ned til den nærværende Tid,
maaske af Millioner Generationers Erfaring, og faa derved
en Værdi, som langt overgaaer Værdien af en Bevisførelse,
der kun støtter sig til den individuelle Erfaring.

Resultatet er, at al Erkjendelse og alle Bevidsthedsformer
stamme fra Erfaringen, men ikke fra den individuelle
Erfaring. Hvert enkelt Individ begynder ikke absolut forfra,
men modtager som Arv fra sine Forfædre en opsamlet og
fæstnet Erfaring, der har ledet Slægtens Organisation i en
vis Retning og til visse Former. Der gives et Apriorisk,
men kun for Individet; hvad der er apriorisk for Individet,
er aposteriorisk for Slægten. Under Slægtens Udviklingsgang
virke de ydre Forhold saaledes paa den, at de almindelige
Former, hvorunder Tingene fremstille sig, faa et
tilsvarende Udtryk i den indre Organisation, navnlig i Nervesystemet.
Der danner sig bestemte Nerveforbindelser, som
potentielt ere tilstede før Fødselen, som gaa forud for og
ere uafhængige af alle individuelle Erfaringer, og som automatisk
lægge sig for Dagen allerede ved de første Erkjendelsesakter.


Hvad selve Stridsspørgsmaalet mellem de to philosophiske
Skoler angaaer, da giver Spencer den idealistiske
Retning Medhold, idet han negter, at Bevidstheden forholder
sig passiv, og at Erkjendelsen kan forklares ud af Individets
% --- p. 171 ---
\setcounter{footnote}{0}%
\opage{171}Erfaring. For at Erfaringen kan blive Erkjendelse, maa den
formes og organiseres, og der forudsættes altsaa en organiserende
Kraft og Former, i hvilke Organisationen gaaer
for sig. Dette stemmer netop med Kants Lære. Spencer
afviger her kun fra Kant ved at reducere alle Anskuelses- og
Tankeformer til én oprindelig Funktion, Opfattelsen af
Forskjel og Lighed.

Men Spencer bliver ikke staaende ved den Form, hvorunder
Problemet tidligere stedse er bleven behandlet. Han
fører Undersøgelsen videre ved at spørge, hvorfra den aandelige
Organisation skriver sig. Dette Spørgsmaal var tidligere
ikke draget med ind i Debatten. Man afsluttede Diskussionen,
naar man havde paavist Nødvendigheden af en saadan
oprindelig Organisation som Forudsætning for al Erfaring
og Erkjendelse. Ved at paavise denne Forudsætning
havde man bestemt Menneskets væsentligste Karaktermærke.
Lige saa standsede man i Botaniken og Zoologien, naar
man havde paavist de væsentlige Karaktermærker for en
vis Art, og spurgte ikke om selve denne Arts og dens Karaktermærkers
Oprindelse. Men lige som Darwin har bragt
dette Spørgsmaal frem indenfor Naturvidenskaben, saaledes
har Spencer (uafhængig af og før Darwin) bragt det frem
indenfor Philosophien. Udviklingens Lov, som er begrundet
i vor Bevidstheds Væsen, og som Erfaringen paa ethvert
Omraade bekræfter, maa ogsaa finde sin Anvendelse her.
De høiere Bevidsthedsformer kave udviklet sig af de lavere; % PRINTED AS IS (kave: O-H08)
selve Bevidsthedsthedslivet har udviklet sig af de lavere % PRINTED AS IS (Bevidsthedsthedslivet: O-H09)
Livsformer; og i Løbet af denne Udvikling, der bestaaer i
det Indres Tillempelse efter det Ydre, have de ydre universelle
Forhold præget deres Former i de indre Forhold.
Spencer kan derfor erklære sig for en Tilhænger af Empirismen,
naar man ved Erfaring forstaaer hele Slægtens
Erfaring.

Hertil maa dog bemærkes, at Spencers egen Kritik af
Empirismen lige saa vel rammer, naar Spørgsmaalet er om
Slægtens Erfaring, som naar det gjælder Individets Erfaring.
Et virkeligt Nødvendigt og Universelt kan der intet af
% --- p. 172 ---
\setcounter{footnote}{0}%
\opage{172}Stederne naaes, naar Erfaring tages i Empirismens Forstand.
Dersom der ikke, ogsaa hvad Slægten i det Hele
angaaer, ligger en tilegnende og organiserende Aktivitet til
Grund, kan det ikke forklares, hvorledes de Former og
Funktioner, der danne sig under Udviklingens Gang, kunne
blive Elementer i vor Natur. Siger man: a priori for Individet,
maa man ogsaa sige: a priori for Slægten, thi
der er her Tale om et af Sagens Natur følgende Grundforhold, ikke om en Tilfældighed, som kunde forholde sig
paa én Maade for Individets, og paa en anden for Slægtens
Vedkommende. »Slægten» er jo desuden kun en relativ
Bestemmelse; den kan igjen opfattes som Individ i Forhold
til en mere omfattende Slægt, og saaledes i det
Uendelige.

At Spencer ikke klart har set, at det her dreier sig
om et Principspørgsmaal, har sin Grund i hans Standpunkt
paa Empirismens Grændse. Og dog laa det nær for ham
at faa Øie derpaa. Han staaer ganske vist endnu paa
»Association psychology»s Grund; men Ideassociationerne
staa hos ham ikke mere som mekaniske Forbindelser, men
optræde som Former og Resultater af den aandelige Grundvirksomhed.
Den mekaniske Psykologi er her lige som
smeltet. Aktivitet og Passivitet gaa i Et, idet enhver
enkelt Fornemmelse kun bliver mulig ved hin Grundvirksomhed.
Men hermed er da Principet givet: uden indre
Virksomhed ingen Tilegnelse af det Ydre. Dette fører os
lige til den Proces, hvorpaa ifølge Spencer Livets Udvikling
beroer, Tillempelsen af det Indre efter det Ydre. Det er
det Indre, der lemper sig; altsaa maa det opfattes som
aktivt, og det Ydre afgiver saa Maal og Maade for Aktivitetens
Udfoldelse. Men det Indre beholder stedse Prioriteten,
hvor langt vi end gaa tilbage.

Hvad er nu i sidste Instant dette Indre, som al Erkjendelse
og Erfaring forudsætter, og som altsaa er et
Princip, vi ikke kunne gaa ud over? Det kan, ifølge Spencers
egen Lære, ikke være Andet end det Absolute, den hemlighedsfulde
Magt, der aabenbarer sig i Alt. Kun ved den
% --- p. 173 ---
\setcounter{footnote}{0}%
\opage{173}ere Liv og Bevidsthed mulige. Den er derfor ogsaa den
sidste Kilde til Sandhed, den sidste Garanti for Erkjendelsesfunktionernes
Gyldighed. Den høieste Nødvendighed,
Tanken leder os til, er kun et Udtryk for vor høieste aandelige
Kraft, og enhver Kraft er atter kun Form for den
absolute Kraft. Det Absolute er altsaa, som af de efterhegelske
Philosopher især Weisse har hævdet,\footnote[1]{Philosophien i Tydskland. p. 187---89.} den sidste
Grund til Tankens Vished. Denne Overbevisning har Spencer
selv udtalt, naar han betragter den fremskridende Tænker
som Organ for den absolute Aarsag og udleder hans Tro
af denne Kilde (se ovenfor p. 131). Man ser, hvorledes
den tidligere paaviste Mangel hos Spencer i Opfattelsen af
det Ubetingede i dets Forhold til Phænomenerne, bliver af
indgribende Betydning ogsaa ved de specielle Problemers
Behandling.

Afset fra det Mangelfulde i Behandlingen maa det
erkjendes som en stor Fortjeneste, at Spencer har indledet
en Betragtning af Erkjendelsesproblemet fra denne nye Side.
Saa sandt som Mennesket er et relativt og endeligt Væsen,
maa der ogsaa nødvendigvis spørges om hans Væsensbestemmelsers
Oprindelse. Kun en mystisk og hyperidealistisk
Anskuelse kan afvise dette Spørgsmaal. Det Absolute og
i sig Gyldige rokkes ikke derved; thi den Opfattelse, vi
kunne have af det, kan dog kun være relativ, kan ikke
selv være absolut; og det er kun om denne eiendommelige
relative Opfattelsesmaades Oprindelse der spørges.

Spencers Lære om Erkjendelsesformernes Oprindelse
fra Slægtens nedarvede Erfaringsresultater er en Hypothese lige
som Darwins Lære om Arternes Oprindelse i det Hele. Dette
maa man nøie fastholde. Derimod er hans Kritik af Empirismen
ingen Hypothese, men Resultatet af en rationel
psykologisk Analyse. Han mener selv ved hin Hypothese
at høre Empirismen til; ved sin Kritik er han derimod
principielt gaaet ud over den. ---


% --- p. 174 ---
\setcounter{footnote}{0}%
\opage{174}Der er en stor Klasse af psykologiske Phænomener,
som Spencer kun har behandlet meget kort; det er Følelserne,
som allerede bemærket maaske det vanskeligste
Emne indenfor Psykologien. Spencer giver her kun en
Skizze (Princ. of psychol. II. p. 508 ff.), med det særlige
Øiemed at paavise det psykologiske Grundlag for Ethik
og Sociologi. Her som paa alle Punkter er Udviklingens
Ide den raadende.

Paa det laveste Trin ere Følelserne umiddelbart knyttede
til sandselige Fornemmelser; men efterhaanden fjerne de sig
fra dette Tilknytningspunkt og antage en mere ideal Karakter.
Denne Udvikling bevæger sig gjennem en Række
af Trin. Friheden og Bevægeligheden bliver større; de
oprindelige automatiske Fornemmelser, de umiddelbart sig
paatrængende Stemninger give Plads for en Modtagelighed
for nye og mere omfattende Synspunkter. Men navnlig
fører Udviklingen til, at de sympathetiske Følelser faa Overhaand
over de egoistiske. Kampen for Tilværelsen gjorde
Urmennesket (the aboriginal man) til en Røver. Han
maatte for at bestaa overfor sine mange Fjender være udrustet
med vilde Instinkter. En udviklet Sympathi vilde
ikke blot have været umulig, men ogsaa fordærvelig for
ham. Men indirekte baner allerede dette vilde Liv Veien
for Sympathien. Er Egoismen nødvendig til Individets
Opholdelse, saa er Sympathien nødvendig til Artens Opholdelse,
navnlig hos svagere Væsener; derfor se vi især
græsædende Dyr leve i Fællesskab. Den stadige Samværen
med beslægtede Væsener, fælles Frygt og fælles Glæde nærer
Sympathien og gjør den efterhaanden til et varigt Element
i Individernes Natur. Foruden Forholdet mellem Væsener
af samme Art virker nu ogsaa Forholdet mellem de to
Kjøn og mellem Forældre og Børn i samme Retning. De
her paapegede Aarsager føre fra de egoistiske Følelsers
Herredømme til et høiere Stadium, hvor de ego-altruistiske\footnote[1]{Udtrykket Altruisme har Spencer laant fra Comte til at betegne Modsætningen til Egoismen.}
Følelser herske.

% --- p. 175 ---
\setcounter{footnote}{0}%
\opage{175}Sympathien fortrænger nemlig ikke strax Egoismen;
tvertimod, den virker selv først til at udvikle og opretholde
nye Former for Egoismen. Sympathien fremkaldes ved
Samliv med Andre. Men denne Sammenslutning fremkalder
netop en stærkere og skarpere Modsætning til dem,
der staa udenfor dette Samliv. Dette viser sig i den ene Nations,
den ene Religions, den ene Stands Forhold til de
andre. --- Ogsaa indenfor Samfundet virker Egoismen i
nøie Forbindelse med Sympathien. Individet faaer nemlig
Motiverne til sin Handlen i den Myndighed, Samfundet
udøver over ham. Frygt for Høvdingens Magt og for
Stammegudernes Vrede bringer Individet til at udføre, hvad
der gjælder for Pligt, hvilket er, hvad der i Stammens
Historie har vist sig at være nyttigt og nødvendigt. Til
saadanne Handlinger er der knyttet Roes og Løn, til deres
Undladelse Straf. Herosdyrkelsen bliver saaledes et nødvendigt
Stadium i Udviklingen, et Udtryk for en nøie
Sammensmeltning af egoistiske og altruistiske Følelser. Selv
om den Løn eller Straf, der forventes, først antages at
skulle komme i en anden Verden, ere Motiverne dog verdslige,
skjøndt Verdsligheden er, hvad man træffende har
kaldet »other-worldliness.« Handlingerne betragtes ikke
efter deres indre Natur, afset fra nærmere eller fjernere
Hensyn til det handlende Selv. De ego-altruistiske Følelser
betegne altsaa en Overgangstid. De ere paa Udviklingens
nærværende Stadium de eneste, der kunne lede Menneskenes
Handlemaade. Det Indre lemper sig jo ifølge Livets almindelige
Lov stedse efter det Ydre og kan ikke bestaa
paa anden Maade. Enhver Tidsalder, enhver Samfundstilstand
uddanner derfor de Følelser og Forestillinger som
passe til Vilkaarene. Særlig ses dette i de religiøse Forestillinger.
Disse ere nødvendigvis anthropomorphistiske, og
de menneskelige Former, der overføres paa det Guddommelige,
ere naturligvis netop de, der paa given Tid og Sted herske
i Bevidstheden, fordi de fremkaldes af de indre og ydre
Forholds Vexelvirkning.

% --- p. 176 ---
\setcounter{footnote}{0}%
\opage{176}Vi befinde os endnu i Overgangstiden fra Vildskab og
Barbari til virkelig socialt Liv. Vel mene vi ofte, at
Barbariet ligger langt bag os, og at vor Periode er Civilisationens;
men det Samme have alle tidligere Tidsaldre
ment, og der gaves neppe noget Barbari i Fortiden, som
ikke under en eller anden Form endnu bestaaer. Røver- og
Krigerlivet, som tidligere herskede, nærede særlig de
egoistiske Følelser, medens det industrielle Liv, navnlig ved
Arbeidets Deling, vækker og nærer Følelsen af Sammenhøren
og Fællesskab. De store og vilde Forbrydelser
blive derfor sjeldnere. Sympathien hindrer den direkte
Tilførelse af Ondt; fremmed Smerte føles stærkere end
før, --- dog ikke saa stærkt, at den ikke ofte ad Omveie
skulde blive fremkaldt. Menneskenaturen har endnu
Elementer i sig, der skrive sig fra Urmenneskets Vildhed.
Det ser man bl. A. hos Børnene. Deres Egoisme, deres
Tilbøielighed til at tilegne sig Alt, deres Mishandlen
af lavere Væsener ere Træk, som egnede sig for den Vildes
Tilstand, men som maa forsvinde, efterhaanden som det
humane Samfundsliv udvikler sig. Det Umoralske beroer,
lige som det Onde i det Hele, paa en Uoverensstemmelse
mellem Livet og dets Vilkaar; de Onder, saa vel moralske
som physiske, vi lide under, ere Tegn paa, at vor Natur kun
ufuldkomment har lempet sig efter Betingelserne.

Menneskets oprindelige Vilkaar-medførte, at han maatte
ofre Andres Velfærd for sin egen; hans nuværende Vilkaar
fordre mere og mere, at han kun skal have saadanne ønsker,
hvis Tilfredsstillelse ikke hindrer Andres Velfærd. Derfor
trænges de egoistiske Følelser efterhaanden tilbage, og de
altruistiske udvikle sig. Frygt, Ærgjærrighed og Stræben
efter Bifald ere ikke længer fyldestgjørende Motiver; den
moralske Sands har nu udviklet sig saaledes, at Hensynet
til Andres Vel bliver det overveiende. Men Altruismen
forudsætter stedse Egoismen, idet de altruistiske Følelser
bestaa i sympathetiske Vækkelser af de egoistiske Følelser.
Sympathien i sig selv er derfor ikke tilstrækkelig til Moralprincip.
Adam Smith oversaa netop, at Sympathien altid
% --- p. 177 ---
\setcounter{footnote}{0}%
\opage{177}forudsætter, at man selv har erfaret det, hvorfor man har
Medfølelse med en Anden. Medlidenheden f. Ex. forudsætter,
at man selv har følt Smerte. En af de mest sammensatte
altruistiske Følelser er Retfærdighedsfølelsen. Den bestaaer
i en sympathetisk Vækkelse af Driften til at hævde sine
Rettigheder. Hvo som ikke selv kjender og tragter efter
personlig Frihed, kan og vil ikke anerkjende Andres; heraf
forklares slaviske Naturers Tyranni. Og paa den anden
Side vilde en absolut Altruisme, selv om den var mulig,
forudsætte, at der stadig var Nogle, der vare egoistiske nok
til at benytte sig af Andres Selvopgivelse og deraf følgende
Selvtilintetgjørelse. --- Udviklingen maa gaa henimod en
Tilstand, hvor Enhver arbeider for Andres Velfærd og Ret
som for sin egen. De ego-altruistiske Følelser svare til
de forbigaaende Tilstande i Samfundet, de altruistiske til
de blivende, fordi de angive Samfundslivets nødvendige
Betingelser.

Selve den psykologiske Udvikling fører os saaledes til
det ethiske Princip. Dette Princip kan ikke isoleres fra
fra Livets Princip i det Hele. Moralitet er transcendental % PRINTED AS IS (fra Livets: O-H13: fra doubled at the line turn)
Physiologi. Livets Gang, Tillempelsen af det Indre efter
det Ydre, fører Menneskeheden fra Barbari til sædeligt
Samfundsliv, og den moralske Lov er den Lov, som vilde
gjælde i det fuldkomne Samfund. Derfor er det ikke noget
Bevis mod et ethisk System, at det ikke stemmer med
Virkeligheden; det vilde tvertimod være et Bevis mod det,
dersom det stemmede med Virkeligheden, eftersom Loven
for den ideale Tilstand ikke kan fyldestgjøres under de
ufuldkomne Tilstande, under Overgangen fra Dyriskhed og
Barbari til Humanitet. Men selv om det ethiske Ideal ikke
kan virkeliggjøres under de givne Betingelser, er det dog
af stor Vigtighed at kjende og fastholde det, for at Udviklingen
kan ske i Retning af det og ikke bort fra det.
Det var derfor en stor Feil, naar Tænkerne tidligere søgte
Moralprincipets Begrundelse i en Naturtilstand, som gik
forud for Samfundslivet. En saadan Tilstand staaer netop
i fuldstændig Modsætning til det ethiske Princip. For at
% --- p. 178 ---
\setcounter{footnote}{0}%
\opage{178}finde dette maa man ikke se tilbage, men fremad. Den
absolute Ethik følger af Livets Grundlove, af Betingelserne
for den menneskelige Naturs høieste Udvikling. (Absolute
morality = conformity to the laws of complete life).

Men paa den anden Side maa vi dog bestandig tage
Hensyn til Menneskehedens faktiske Tilstand og Karakter,
til alle de Ufuldkommenheder i de nuværende Vilkaar, som
hindre den absolute Ethiks Virkeliggjørelse. Vi maa ikke
blot gaa apriorisk og deduktivt, ogsaa aposteriorisk og induktivt
til Værks.\footnote[1]{Essays II. p. 210-211.} Nyttens Princip supplerer da her det
ideale Moralprincip. Den fremadskridende Civilisation beroer
nødvendigvis paa en Række af Kompromiser mellem Gammelt
og Nyt, mellem det Ideale og det Praktiske. Under vor
Stræben efter at gjøre, hvad der netop passer for den nærværende
Tid, maa Tanken dog stedse være rettet imod,
hvad der abstrakt set er bedst. Nytteprincipet et et nødvendigt % PRINTED AS IS (et et: O-H14: for er et)
Supplement, skjøndt ikke det høieste Princip.

Det er ikke Erkjendelsen af det Nyttige, der fremkalder
de moralske Følelser, men omvendt de moralske Følelser,
der gjøre Erkjendelsen af det Nyttige mulig. Nytteprincipet
forudsætter altsaa stedse de dybere Motiver, der udspringe
af Slægtens Natur, som Individet har Del i. Utilitarianismen
stiller lige som Empirismen i det Hele Individet som isoleret
Forstandsvæsen uden for Slægten, den sig udviklende
Slægt, og kan derfor ikke gribe det Oprindelige i Individets
Natur. For at erfare at Sympathi gavner maa man først
handle sympathetisk. Kun ved at indse dette opdage vi det
sande Motiv, der lader Handlingen udføres for dens egen
Skyld, medens Nyttetheorien lader Formaalet falde udenfor
Handlingen.

Nyttetheorien lader det desuden uafgjort, hvad Nytte
eller Lykke egentlig er. Lykken maa ganske vist betragtes
som en guddommelig Ide, der skal realiseres. Men den
kan kun bestaa i vore Evners Øvelse og Udvikling. Det
% --- p. 179 ---
\setcounter{footnote}{0}%
\opage{179}er vor Pligt at udvikle vore Evner, og vi have Ret til
den Frihed, ved hvilken denne Udvikling ene bliver mulig.
Nyttetheorien viser altsaa ogsaa her tilbage til en mere
omfattende Anskuelse, til Udviklingens almindelige Ide.\footnote[1]{Stuart Mill fortæller (Moral p. 82), at Spencer har meddelt ham, at han ikke vil betragtes som Modstander af Nyttetheorien, skjøndt han har bekæmpet den baade i sin „Social Statics“ og i sin Psykologi. Nyttetheorien er jo ogsaa, som ovenstaaende Fremstilling viser, en \emph{Del} af Spencers ethiske Lære; men ifølge Bentham og Mill skulde den være \emph{hele} Moralen. Foruden i dette Principspørgsmaal er der ogsaa den Forskjel mellem Spencer og Utilitarianerne, at disse gaa rent empirisk til Værks, medens Spencer sætter Ethiken lige som Sociologien i Forbindelse med Livets almindelige Love og derved vinder Muligheden af en deduktiv Methode indenfor Ethiken. Dette hænger sammen med hans systematiske Tendents og omfattende Blik.}

Ikke blot i Utilitarianismen, ogsaa hos Comte finder
Spencer en Overvurdering af Erkjendelsen paa Følelsens
Bekostning. I Modsætning til Comtes Lære (i hans »Philosophie
positive»), at »Ideerne styre Verden,» at »hele den
sociale Mekanisme tilsidst hviler paa Meninger,» lærer
Spencer, at Verden styres af Følelserne, for hvilke Ideerne
kun tjene til Førere, og at den sociale Mekanisme saa godt
som udelukkende har Karaktererne til sit Grundlag. »Menneskenes
Drifter ere i Hovedsagen nedarvede, men deres
Overbevisninger ere i Hovedsagen selverhvervede og afhænge
af de omgivende Betingelser.» »I Virkeligheden beroer det
paa Folkets Karakter og Samfundets Tilstand, hvilke Ideer
der skal udbrede sig. Hovedaarsagen til det sociale Fremskridt
er derfor den Forandring af Menneskenes Natur, som
bevirkes ved det sociale Livs bestandige Tugt, der bringer den
mere og mere i Overensstemmelse med de sociale Forhold.»
Ifølge Comte er det »det intellektuelle Anarki,» ifølge Spencer
er det de moralske Modsætninger, der bevirke vor Tids
politiske Kriser.\footnote[2]{The classification of the sciences. p. 37-38.}

Denne Opfattelse udspringer af Spencers psykologiske
Grundanskuelse. Vi have set, at han særlig hævder
% --- p. 180 ---
\setcounter{footnote}{0}%
\opage{180}Følelsernes Oprindelighed, fordi de endnu mindre end Erkjendelsens
Former kunne være Frugten af det enkelte Individs
Erfaring og Arbeide. Det er derfor især gjennem
Følelserne, at Individet paavirkes af Slægtens tidligere Erfaringer,
og de Egenskaber og Karaktertræk, der herigjennem
udvikle sig, ere for stærke til at kunne ændres ved vilkaarlige Foranstaltninger og Paavirkninger. Spencer drager
heraf vigtige sociologiske Konsekventser. Man overvurderer
efter hans Mening de Virkninger, som ydre Reformer, særlig
i Opdragelse og Undervisning, kunne have. Fremfor Alt
kommer det an paa at forandre Grundegenskaberne, Karakterernes
Elementer --- men de forandres kun gjennem den
praktiske Kamp og Vexelvirkning med Livsvilkaarene. Samfundets
Udvikling er en Naturproces, hvis Resultatater man % PRINTED AS IS (Resultatater: O-H20)
ikke paa kunstig Maade kan fremskynde. Det er en Illusion,
som de Radikale dele med de Konservative, at Lovgivningen
er almægtig. Lovene maa jo fremgaa af Menneskenaturen
og kunne altsaa ikke være hævede over de intellektuelle
og moralske Ufuldkommenheder, hvoraf denne endnu lider.
Heller ikkke kan Opdragelsen reformeres, saa længe Forældrene % PRINTED AS IS (ikkke: O-H21)
endnu ere saa langt tilbage. »Al Overtro døer
langsomt, og vi frygte, at denne Tro paa Regeringens Almagt
ikke vil danne nogen Undtagelse.«

»Et videnskabeligt Studium af Sociologien fører til
klarere Bedømmelse af de forskjellige politiske og religiøse
Partier. Den Opfattelse som den sociale Videnskab begrunder
og udvikler, er paa éngang radikal og konservativ --- radikal
i en Grad, der overgaaer Alt, hvad den sædvanlige Radikalisme
fatter, og konservativ i en Grad, der overgaaer den
nuværende Konservatismes Forestillinger. Naar man har
tilegnet sig den Sandhed, at Samfundene ere Produkter af
Udviklingen, der til de forskjellige Tider og paa de forskjellige
Steder antage forskjellige Former i Henseende til Struktur
og Funktion, --- saa opstaaer den Overbevisning, at hvad
der for vore Tanker og Følelser staaer som yderlig slette
Indretninger, dog passede under visse Betingelser, der umuliggjorde
bedre Indretninger; og heraf udspringer da en
% --- p. 181 ---
\setcounter{footnote}{0}%
\opage{181}tolerant Fortolkning af Fortidens Tyrannier, over hvilke endog
den ivrigste Tory i vore Dage vilde forarges. Paa den
anden Side udspringer der af Iagttagelsen af, hvorledes de
Processer, der have bragt Tingene til deres nuværende Tilstand,
stedse gaa for sig: ikke med en aftagende Hurtighed, der
kunde tyde paa at de nærmede sig deres Afslutning, men med
en voxende Hurtighed, som tyder paa lang Vedvaren og
umaadelige Omdannelser, --- udspringer der af denne Iagttagelse
Overbevisningen om, at den fjerne Fremtid gjemmer i
sit Skjød Former for socialt Liv, som ere høiere end vi kunne
tænke os; og denne Tro gaaer langt ud over, hvad den
Radikale troer, der tragter efter en Forandring, som kan
sammenlignes med Organisationer, der allerede nu existere.
Og medens denne Opfattelse af Samfundet fører til at se,
at umaadelige Forandringer kunne hidføres under Tingenes
langsomme Gang, saa fører den ogsaa til at se, at indenfor
korte Perioder ere kun ringe Forandringer mulige.»\footnote[1]{The study of sociology. London 1874. p. 399-400.}

Man ser, at Spencer ikke som Mill træder ned paa
Arenaen for at deltage i Tidens Kampe, men holder sig paa
et rent theoretisk og videnskabeligt Standpunkt. Hans
Stræben gaaer netop ud paa at frigjøre Sociologien fra alle
subjektive og Partisynsmaader, Noget, der naturligvis, som
han selv erkjender, kun tilnærmelsesvis er muligt, da jo
den Forskende selv er Medlem af det Samfund, hvis Livsbetingelser
og Udviklingsgang, hvis Naturhistorie han undersøger.\footnote[2]{Som induktivt Grundlag for „Sociologiens Principer“ er Spencer i Forening med flere Medarbeidere begyndt at udgive en „deskriptiv Sociologi,“ en tabellarisk Oversigt over de forskjellige Kultursamfund, deres Former og deres Udvikling, som skal lette et komparativt Studium af de sociologiske Phænomener. --- Det ovenfor citerede Skrift, „The study of sociology,“ fremstiller Sociologiens almindelige Opgave og de subjektive og objektive Vanskeligheder ved dens Behandling.}
At denne hans Stræben efter videnskabelig
Objektivitet og hans Skepsis overfor Reformforslag ikke har
sin Grund i Mangel paa Indsigt i Mislighederne ved vore
nuværende Samfundsforhold, ses ikke blot af det nys Anførte,
% --- p. 182 ---
\setcounter{footnote}{0}%
\opage{182}men ogsaa af Udtalelser som disse: »Ethisk set har der
aldrig været anden Berettigelse for de Manges Underkastelse
under de Faa uden den, at den fremmede de Manges Velfærd.
--- Lige som den gamle Theori om politisk Regering
er bleven omdannet saaledes, at Styreren i Stedet for at
være Nationens Eiermand, betragtes som dens Tjener, saaledes
vil den gamle Theori om industriel og social Regering
undergaa en Forandring, der vil lade de styrende Klasser
føle, at deres Interesser kun ere underordnede i Forhold til
de Massers Interesser, hvis Arbeider de lede.»\footnote[1]{Ibid. p. 256.}

Ethiken har ifølge Spencer en absolut og ideal Karakter.
Den er derfor uvidende om det Onde og dets
Følger. Den ethiske Lov, som er Et med Udviklingens
Lov, fører netop selv ud over det Onde. Fremskridtet er
en Nødvendighed; Civilisationen gaaer for sig som en Naturproces.
Ufuldkommenhed vil ikke sige Andet end en
Uoverensstemmelse med Betingelserne, som enten bestaaer i
at man har Evner, der hverken passe til Vilkaarene eller kunne
finde Anvendelse under dem, eller ogsaa i at man mangler
de Evner, der udfordres for at kunne bestaa. Men af Livets
almindelige Lov, at det Indre lemper sig efter det Ydre og
derved udvikler sig selv, følger, at de overflødige Evner dø
bort af Mangel paa Øvelse, og at de nødvendige Kræfter
efterhaanden udvikle sig. Idealets Virkeliggjørelse er derfor
logisk vis. Det Onde visner bort, det Gode lever.\footnote[2]{Allerede Condorcet og I. G. Fichte have udtalt denne Tanke om „the evanescence of evil.“} Man
skal kun lade Virkeligheden raade, lade den øve sin Opdragelse
og bortskjære det Svage og Fordærvelige.

Overgangstiden fra Dyriskhed til Humanitet maa nødvendigvis
være ulykkelig, thi Ulykke er den nødvendige
Følge af Misforhold mellem den indre Organisation og de
ydre Betingelser. Disse Lidelser maa udholdes, skjøndt de
naturligvis kunne mildnes ved Menneskekjærlighedens Virksomhed.
Forsøg paa kunstig at unddrage Menneskene
% --- p. 183 ---
\setcounter{footnote}{0}%
\opage{183}Forholdenes Konsekventser tjene kun til at trække Sagen
i Langdrag og berøve dem tillige det Udbytte, Lidelserne
skulde bringe: Opdragelse og Udvikling.\footnote[1]{„Dersom vi, i Stedet for at tillade Menneskene at komme i Berøring med deres Stillings virkelige Betingelser, stille dem under kunstige Betingelser, under falske Vilkaar, saa ville de lempe sig efter disse, og ville saa tilsidst have at foretage en ny Tillempelse efter de virkelige Forhold.“ Social Statics. p. 387. Kfr. Spencers Skrift „Education, intellectual, moral and physical.“ p. 113 ff.} Det er kun
Kortsynethed og Blødagtighed, der over de øieblikkelige
Onder glemmer det store, universelle Maal. Selv i den
største Smerte og Nød øine vi en kæmpende Kjærlighed,
en Kraft, der virker til Sundhed, Godhed og Lykke. Enhver
maa arbeide med efter sine Evner i den Overbevisning,
at Intet gaaer tabt; selv en rent øieblikkelig Stemning af
Ædelmodighed og Velvillie føier et Atom til den ideale Bygning,
hvorpaa der arbeides. ---

Da Spencers speciellere Behandling af Ethiken og
Sociologien endnu ikke foreligger, ville vi ikke nærmere
undersøge eller kritisere de almindelige Træk af hans praktiske
Philosophi, som nu ere fremdragne. Det er klart, at
hans Stilling til Empirismen paa det praktiske Omraade
svarer til den, han indtager overfor den i theoretisk Henseende.
Den ethiske Lov hævdes i sin Idealitet, men ses
til samme Tid som den høieste Naturlov, den høieste Form
for Livets almindelige Lov. Herved faaer det Ideale den
faste og reale Begrundelse, som Idealismen saa ofte har
ladet det mangle paa. Det er Udviklingens Princip og
den Lov, som gjælder for dette Princips Virken, der ogsaa
forbinder det Ideale og det Reale. Og Opfattelsen af det
ethiske Ideal bliver kun mulig ved den uendeliggjørende
Kraft, som Bevidstheden eier, men som Empirismen ganske
underordner den passive Receptivitet. --- Ogsaa her viser
det sig da, at Spencers Anskuelse kun afslutter sig paa en
konsekvent Maade, dersom vi have Ret i, hvad der i den
nu givne Fremstilling og Kritik gjentagne Gange er frem%
% --- p. 184 ---
\setcounter{footnote}{0}%
\opage{184}hævet, at det ideelle Element i den er det overveiende,
er det Princip, som raader og virker i den empiriske
Realitet. ---

\closinghead{Slutning.}


Den Vei, vi have gjennemvandret, har ført os til et
bestemt Maal. Empirismen med sin Opfattelse af Bevidstheden
som passivt modtagende og med sin Opløsning af
Aandens Enhed i Atomer, der kun ved Ideassociationens
Love fik en vis, skjøndt tilfældig Sammenhæng, --- denne
Lære, som var en Arv fra det 18de Aarhundrede, have vi
set udvikle sig ved sin egen Konsekvents og ved Debat
med den kritiske og idealistiske Philosophi til et Punkt,
hvor ikke blot Bevidsthedens og Aandens aktive Enhed anerkjendes,
men hvor den indre Nødvendighed træder i Stedet
for den ydre Tilfældighed.

De Atomer, hvori Bevidsthedslivet opløste sig for den
empiriske »Association-psychology», vare de enkelte Fornemmelser
eller rettere Fornemmelseselementer. Men det
viste sig, at ingen enkelt Fornemmelse var mulig uden en
elementær aandelig Virken, der opfattede dens Forskjel fra
andre Fornemmelser og paa en eller anden Maade kombinerede
den med disse. Den sidste Faktor, hvortil Analysen
førte, blev altsaa selve denne ideelle Aktivitet, hvorpaa
Aandens Enhed og Sammenhæng beroer; den ligger til
Grund for al Fornemmelse og Erfaring, som kun ved den
ere mulige.

Den engelske Philosophi har saaledes paa sin eiendommelige
Maade arbeidet sig frem til den store Grundtanke,
som Descartes, Leibnitz, Kant og Hegel have udviklet,
og som den nyere Naturvidenskab saa slaaende og storartet
har bekræftet, at Subjektivitetens Ide er den sidste, hvortil
vi komme, og at i den Principet for det Objektive maa
% --- p. 185 ---
\setcounter{footnote}{0}%
\opage{185}søges. Men det har hidtil været den kontinentale
Philosophis Mangel og Ensidighed, at den har udviklet
Subjektivitetens Ide overveiende paa apriorisk og dogmatisk
Vis. Derved afskar man sig ikke blot den aposterioriske
Verifikation, uden hvilken ingen videnskabelig Sandhed kan
hævdes, men man kom ogsaa til at opfatte selve det ideelle
Princip abstrakt og isoleret fra den reale Virkelighed. Dette
traadte klart frem ved det Vendepunkt, som indtraadte
med Hegels System, det storartede Forsøg paa en absolut
Idealisme, en Fremstilling af Ideen i sin absolute Væren
i og for sig (theologisk udtrykt: af Gud som han var før
Verdens Skabelse).

Den engelske Philosophis Styrke ligger i dens skarpe
og sikre Blik for Betingelserne, som vel ofte har tilsløret
Blikket for det Ubetingede, men som til Gjengjæld har
gjort en mere real og konkret Opfattelse mulig. Den har
set, at naar Udviklingen ikke foregaaer i en ideel Ætherverden,
men i Virkelighedens og Relativitetens Verden, saa
maa der lægges en afgjørende Vægt, en principiel Accent
paa de enkelte Betingelser og Omstændigheder. Aanden
udvikler sig kun gjennem det Indhold, den optager. Først
gjennem en lang Kamp, en møisommelig Udviklingsproces,
bliver det optagne Indhold en Del af Aandens Natur.

Udviklingens Lov viser os saaledes det Ubetingede i
dets Enhed med Betingelserne, det Indre i dets nødvendige
Forhold til det Ydre. Sit bestemteste og prægnanteste
Udtryk finder Udviklingens Ide i Livets Lov, saaledes som
Spencer har formuleret den: det Indre lemper sig efter det
Ydre og udvikler sig selv gjennem denne Tillempelse. Det
ideelle Princips Gyldighed viser sig netop gjennem dets Afhængighed
eller tilsyneladende Forsvinden.

Kun ved Udviklingens Ide vinder den engelske Philosophi
Sammenhæng og Afslutning. Vi have søgt at paavise,
hvilke Modsigelser der, navnlig hos Stuart Mill, reiste
sig, naar denne Ide ikke lagdes til Grund. Det enkelte
reale Moment, eller paa Psykologiens Omraade den enkelte
Fornemmelse, bliver ved at opfattes isoleret dødt og ufrugt%
% --- p. 186 ---
\setcounter{footnote}{0}%
\opage{186}bart; det maa for at have virkelig Gyldighed og Realitet
ses som Led i den hele Sammenhæng; men Led i Sammenhængen
bliver det kun ved en Proces, og den for denne
Proces gjældende Lov bliver altsaa den sidste Forklaringsgrund,
hvortil vi kunne komme. Ved denne Lov forsones
de Bestemmelser, der, umiddelbart og isoleret tagne, stride
mod hverandre. Saaledes stride Ideassociationens Love,
naar de opfattes som mekaniske og tilfældige, mod Bevidsthedens
Enhed og Energi; men denne Strid ophæves, naar
man ser dem som Love for Bevidsthedslivets Udvikling, som
Udtryk for de Betingelser, ved hvilke Enheden udvikles og
beriges og Energien antager høiere Former.

Det følger af sig selv, at vi hos de fleste af de engelske
Philosopher, selv hos Spencer, der netop har gjort
Udviklingen til sin Grundidé og støttet til den bekæmper
den rene Empirisme, maa finde Spor og Eftervirkninger af
den empirisk-mekaniske Opfattelsesmaade, som danner Udgangspunktet
for den engelske Philosopheren i vort Aarhundrede.
Dette har den foregaaende Fremstilling gjentagne
Gange fremhævet, og vi skulle ikke dvæle ved det her, hvor
vi kun i Korthed ville fremhæve de almindelige Resultater,
vi kunne uddrage af den nu afsluttede Fremstilling og Kritik.

Den engelske Philosophi i vor Tid er i sin Karakter
beslægtet med den efterhegelske Philosophi i Tydskland og
har væsentlige Berøringspunkter med den. I England gaaer
Bevægelsen fra Empirisme og Realisme tilbage til Idealisme;
i Tydskland havde man den spekulative Idealisme
til historisk Forudsætning, men søgte nu at tage et nyt
Udgangspunkt i Erfaringen og det ved den givne Reale for
ad denne Vei at vende tilbage til den ideale Anskuelse.
I begge Retninger anerkjendes saaledes det experimentale
Udgangspunkt som væsentligt, skjøndt det anerkjendes, at
dette Udgangspunkt viser tilbage til dybere Forudsætninger
og selv er Led i en større, ved disse Forudsætninger bestemt
Sammenhæng.\footnote[1]{Herefter maa Udtalelsen i min Bog „Philosophien i Tydskland \opage{187}efter Hegel“ p. 323 berigtiges; den rammer Mill, men ikke den engelske Philosophi som Helhed.}% the note runs over from p. 186 to p. 187 (joined by draft.py)
 Man kan sige, at Idealismen har beseiret
% --- p. 187 ---
\setcounter{footnote}{0}%
\opage{187}Realismen, men at lige saa omvendt Realismen har beseiret
Idealismen. Fra begge Sider vises vi tilbage til en og
samme Grundidé.

Methoden er i den engelske Philosophi væsentlig induktiv.
Induktionens Logik skyldes jo ogsaa netop denne
Philosophi. Den samme Methode, som Naturvidenskaben
skylder sin Storhed, anerkjendes ogsaa mere og mere som
væsentlig for Philosophien. Idet denne som den almindelige
Videnskab skal sammenfatte de enkelte Videnskabers Resultater,
kan det kun ske ved den samme Methode, som anvendes
indenfor enhver af disse selv. Det sidste, strenge
Kriterium for, hvad der skal erkjendes som virkeligt, kan
kun søges i den exakte Erfaring. Det er Stuart Mills
store Fortjeneste, at han ikke blot har formuleret Induktionens
konkrete Methoder, men ogsaa, mere indtrængende
end nogen anden Logiker har hævdet Verifikationens Nødvendighed.
Han overser ikke Deduktionens Betydning, men
indskærper, at dens Resultater stedse maa prøves i Erfaringen.
Hvad han derimod ikke ser, det er, at Induktionen
selv forudsætter Principer og Grundsætninger, uden hvilke
den Intet vilde kunne opbygge. Hvad det Virkelige er, kan
kun Erfaringen lære os; men at vi gjennem Erfaringen
naa Erkjendelsen af det Virkelige, er begrundet i vor aandelige
og legemlige Organisation. Det Indre er ogsaa her
Principet for det Ydre; det maa vel i Erfaringen lempe sig
efter det Ydre, forholde sig modtagende til det; men dette
kan det kun i Kraft af sin ideelle Oprindelighed; den Maade,
hvorpaa det tilegner sig det Ydre, beroer paa dets egen
forud givne Natur; og gjennem Tillempelsen udvikler det
sig selv. I sin Methode som i sin Grundidé fører vor
Tids engelske Philosophi os altsaa ud over Empirismen; og
ogsaa ved Spørgsmaalet om Methoden er det ved Udviklingens
Lov, at Fremskridtet sker.

Staaer det fast, at vi kun have en sand Erkjendelse
% --- p. 188 ---
\setcounter{footnote}{0}%
\opage{188}af det Virkelige, naar vi kunne verificere den, ɔ: eftervise
dens Stadfæstelse ved i Erfaringen givne Kjendsgjerninger,
saa er det ogsaa vist, at vi ikke kunne vinde en virkelig
Erkjendelse af det absolute Princip. Vor Begrebsudvikling
kan paa dette Punkt ikke verificeres. Kants »Kritik eller
spekulativen Theologie» staaer for saa vidt ved Magt endnu
den Dag idag. Men naar Analysen af vor Bevidstheds
Væsen og af Erfaringens Resultater fører os til Forudsætningen
om et saadant absolut Princip (Spencer); naar Kilden
til Ideen om det Uendelige ligger i vor Aands oprindelige
Aktivitet (Bain); naar Aarsagsideen er et Udtryk for, at vi
med Nødvendighed se en indre Enhed og Sammenhæng
som Grundlag for den phænomenale Mangfoldighed og Bevægelse
(Hamilton); --- saa rokkes dette ikke ved Verifikationens
Umulighed. Vi kunne ikke deducere det Faktiske
og Betingede ud af det Ubetingedes Ide; men vi kunne
udvikle denne Ide saaledes, at den ses som den høieste
Forudsætning, mod hvilken al Erkjendelse peger tilbage, og
som det Ideal, mod hvilket al Udvikling i Naturens og
Aandens Verden peger fremad. ---

Men det er ikke blot ved sine Hovedideer og sin
Methode, at den engelske Philosophi er af stor Betydning,
dens Fortjeneste beroer ikke mindst paa den Række af
Enkeltundersøgelser, den har foretaget paa Logikens og Psykologiens
Omraade. Fremskridtet paa de almindelige Ideers
Omraade er stedse betinget ved et Fremskridt i Behandlingen
af de konkrete Problemer. Stuart Mills Formulering
af den experimentale Methode, Bains deskriptive Psykologi,
Hamiltons skarpsindige Theori om Kausalitetsprincipet,
Spencers psykologiske og metaphysiske Analyser ere Mønstre
paa videnskabelig Undersøgelse. Ogsaa denne Retning mod
Enkeltundersøgelser har den engelske Philosophi fælles med
den efterhegelske Philosophi i Tydskland; vi minde blot om
Fechners og Lotzes physiologiske og psykologiske Undersøgelser
og om Feuerbachs og Weisses religionsphilosophiske.
Jo mere man anerkjender Induktionens og Verifikationens
% --- p. 189 ---
\setcounter{footnote}{0}%
\opage{189}Betydning, des større Vægt maa der ogsaa ifølge Sagens
Natur lægges paa saadanne specielle Undersøgelser.

Vi staa her ved et Punkt, hvor Forskjellen mellem
Skolerne taber sin Betydning. Ved Enkeltundersøgelserne
lige som ved Tilegnelsen af det ved Erfaringsvidenskaben
vundne Indhold danner der sig efterhaanden en fælles Grund
for al philosophisk Forsken. For saa vidt er Philosophien
i Færd med at indtræde i, hvad Comte kalder det positive
Stadium. At denne Indtrædelse dog ikke, som Comte mente,
kan ske ved at opgive den dybere Begrundelse og den omfattende
Anskuelse, som kun naaes ad logisk og metaphysisk
Analyses Vei, det er en Sandhed, ved hvis Anerkjendelse
de engelske Philosopher adskille sig fra den franske Positivisme.
Men en skarpere Bevidsthed om hin fælles
Grund, en mere og mere konsekvent Gjennemførelse af de
philosophiske Principer vil, som Historien, og særlig Philosophiens
Historie i vor Tid, viser, ogsaa føre til større
og større Overensstemmelse i Resultaterne. Baade i Grundideerne
og i Methoden have vi saaledes fundet Konvergents
mellem de philosophiske Hovedretninger i vor Tid. Skolerne
tabe deres Betydning, jo mere paa den ene Side Betydningen
af den exakte Erfarings Resultater, paa den anden Side de
spekulative Deduktioners subjektive Karakter og Umuligheden
af deres Verifikation fremtræder for Bevidstheden.

Den efterhegelske Philosophi i Tydskland havde til
sin historiske Forudsætning den absolute Idealisme og kunde
kun ved et Brud med den vinde Betingelser for videre
Fremadskriden. Den engelske Philosophi derimod har udviklet
sig kontinuerlig fra det forrige Aarhundrede ind i
dette og har banet sig Veien til nye Resultater ved at
forfølge sin engang indslagne Bane samtidig med, at den
paavirkedes af den kontinentale Philosophis Resultater.

Begge Retninger vise os, hver fra sin Side, hen imod
det usynlige Midtpunkt, hvor Sandheden boer, som vel
aldrig vil kunne gribes af den endelige Tanke, men som
dog er det Maal, efter hvilket den bestandig stræber, og
som den i uendelig Tilnærmelse bevæger sig henimod. Det
% --- p. 190 ---
\setcounter{footnote}{0}%
\opage{190}har været Hensigten med den nu sluttede Skildring af
disse vigtige Retninger indenfor den philosophiske Forsken
i vor Tid, at vise, at der ogsaa nu arbeides og virkes henimod
dette Maal. Eftertrykket ligger ikke paa de til en
given Tid vundne Resultater; disse vise bestandig ud over
sig selv. Udviklingens Lov lærer os derimod, at Sandheden
ligger i den evige Stræben, og det er Videnskabens og
Philosophiens Kald at vække og nære denne Stræben, at
føre den fremad til nye Opgaver og nye Arbeider.

\end{document}
